% PROVENANCE-14-FERMION-FABLE-CANONICAL-QUANTIZATION.tex
% The LaTeX twin of provenance page 14 (the page of the same name at the repository root): the same
% sections, the same content, the same order.  Build it from this directory with
%     latexmk -pdf -interaction=nonstopmode -halt-on-error PROVENANCE-14-FERMION-FABLE-CANONICAL-QUANTIZATION.tex
% or all three pages with  bash build_all.sh .  The field equations, the anticommutator, J and the
% equations of state are \input from generated/, the TeX written by
% claude-fable/export_fermion_fable_tex.wls from the notebook's own Part VII objects.
% The assertion labels quoted below are copied from claude-fable/run_fermion_fable_part7.log (and,
% where marked, run_fable51_part6.log, run_fermion_fable_final.log and, for the three superseded
% Section 32 labels of section 10.3, run_fermion_fable_part8.log) by the script that assembled
% this file; none was retyped.  The log excerpts in the Verbatim blocks are spliced in unchanged.
% preamble.tex -- shared preamble of the LaTeX twins of the provenance pages
%   PROVENANCE-14-FERMION-FABLE-CANONICAL-QUANTIZATION
%   PROVENANCE-15-FERMION-FABLE-AND-THE-PRIMORDIAL-GRAVITATIONAL-FIELD
%   PROVENANCE-16-THE-COMPLETE-SOLUTION-AND-ITS-COMMANDS
% Each document starts with % preamble.tex -- shared preamble of the LaTeX twins of the provenance pages
%   PROVENANCE-14-FERMION-FABLE-CANONICAL-QUANTIZATION
%   PROVENANCE-15-FERMION-FABLE-AND-THE-PRIMORDIAL-GRAVITATIONAL-FIELD
%   PROVENANCE-16-THE-COMPLETE-SOLUTION-AND-ITS-COMMANDS
% Each document starts with % preamble.tex -- shared preamble of the LaTeX twins of the provenance pages
%   PROVENANCE-14-FERMION-FABLE-CANONICAL-QUANTIZATION
%   PROVENANCE-15-FERMION-FABLE-AND-THE-PRIMORDIAL-GRAVITATIONAL-FIELD
%   PROVENANCE-16-THE-COMPLETE-SOLUTION-AND-ITS-COMMANDS
% Each document starts with \input{preamble}.  Build all three from this directory with
%     bash build_all.sh            (Linux, macOS, Git Bash on Windows)
%     powershell -ExecutionPolicy Bypass -File build_all.ps1
% or one of them with
%     latexmk -pdf -interaction=nonstopmode -halt-on-error PROVENANCE-14-FERMION-FABLE-CANONICAL-QUANTIZATION.tex
% The style follows fable-cosmology/latex/fable_cosmology.tex (pdfLaTeX, Latin Modern, a4, 2.5 cm).
\documentclass[11pt,a4paper]{article}

\usepackage[T1]{fontenc}
\usepackage[utf8]{inputenc}
\usepackage{lmodern}
\usepackage[margin=2.5cm]{geometry}
\usepackage{microtype}
\usepackage{amsmath,amssymb,mathtools}
\usepackage{bm}
\usepackage{graphicx}
\usepackage{booktabs}
\usepackage{tabularx}
\usepackage{longtable}
\usepackage{array}
\usepackage{caption}
\usepackage{subcaption}
\usepackage{placeins}
\usepackage{enumitem}
\usepackage{fancyvrb}
\usepackage{upquote}
\usepackage{xcolor}
\usepackage[hyphens]{url}
\usepackage[hidelinks]{hyperref}

\graphicspath{{figures/}}
\captionsetup{font=small,labelfont=bf}
\setlist{itemsep=2pt,topsep=4pt}
\fvset{fontsize=\footnotesize,xleftmargin=1em}
\newcolumntype{Y}{>{\raggedright\arraybackslash}X}
\setcounter{tocdepth}{2}
\allowdisplaybreaks

% ---- the objects of the notebook, named as the notebook names them ----------------------------
\newcommand{\fable}{\textsf{fable}}
\newcommand{\fableScalar}{\textsf{fableScalar}}
\newcommand{\Lhat}{\hat{L}}
\newcommand{\sig}{\sigma_{16}}
\newcommand{\T}[1]{T_{16}[#1]}
\newcommand{\sgn}{\sqrt{|g|}}
\newcommand{\dd}{\mathrm{d}}
\newcommand{\diag}{\operatorname{diag}}
\newcommand{\sgnf}{\operatorname{sign}}
\newcommand{\Hobs}{H_{\mathrm{obs}}}
\newcommand{\Vhid}{V_{\mathrm{hid}}}
\newcommand{\Psibar}{\bar{\Psi}}
\newcommand{\cconj}{\ddagger}          % classical conjugation of a Grassmann field
\newcommand{\hadj}{\dagger}            % Hilbert-space adjoint of the J-quantization
\newcommand{\Jm}{\mathcal{J}}          % the Krein fundamental symmetry J = -i T0 T1 T2 T3 T4
\newcommand{\KS}{\mathrm{KS}}
\newcommand{\kF}{k_{\mathrm{F}}}
\newcommand{\wF}{\omega_{\mathrm{F}}}
\newcommand{\meff}{m_{\mathrm{eff}}}
\newcommand{\Tr}{\operatorname{Tr}}
\newcommand{\file}[1]{\path{#1}}
\newcommand{\assertion}[1]{\textsc{assert:}~\emph{#1}}
\pdfstringdefDisableCommands{\def\fableScalar{fableScalar}\def\fable{fable}}
.  Build all three from this directory with
%     bash build_all.sh            (Linux, macOS, Git Bash on Windows)
%     powershell -ExecutionPolicy Bypass -File build_all.ps1
% or one of them with
%     latexmk -pdf -interaction=nonstopmode -halt-on-error PROVENANCE-14-FERMION-FABLE-CANONICAL-QUANTIZATION.tex
% The style follows fable-cosmology/latex/fable_cosmology.tex (pdfLaTeX, Latin Modern, a4, 2.5 cm).
\documentclass[11pt,a4paper]{article}

\usepackage[T1]{fontenc}
\usepackage[utf8]{inputenc}
\usepackage{lmodern}
\usepackage[margin=2.5cm]{geometry}
\usepackage{microtype}
\usepackage{amsmath,amssymb,mathtools}
\usepackage{bm}
\usepackage{graphicx}
\usepackage{booktabs}
\usepackage{tabularx}
\usepackage{longtable}
\usepackage{array}
\usepackage{caption}
\usepackage{subcaption}
\usepackage{placeins}
\usepackage{enumitem}
\usepackage{fancyvrb}
\usepackage{upquote}
\usepackage{xcolor}
\usepackage[hyphens]{url}
\usepackage[hidelinks]{hyperref}

\graphicspath{{figures/}}
\captionsetup{font=small,labelfont=bf}
\setlist{itemsep=2pt,topsep=4pt}
\fvset{fontsize=\footnotesize,xleftmargin=1em}
\newcolumntype{Y}{>{\raggedright\arraybackslash}X}
\setcounter{tocdepth}{2}
\allowdisplaybreaks

% ---- the objects of the notebook, named as the notebook names them ----------------------------
\newcommand{\fable}{\textsf{fable}}
\newcommand{\fableScalar}{\textsf{fableScalar}}
\newcommand{\Lhat}{\hat{L}}
\newcommand{\sig}{\sigma_{16}}
\newcommand{\T}[1]{T_{16}[#1]}
\newcommand{\sgn}{\sqrt{|g|}}
\newcommand{\dd}{\mathrm{d}}
\newcommand{\diag}{\operatorname{diag}}
\newcommand{\sgnf}{\operatorname{sign}}
\newcommand{\Hobs}{H_{\mathrm{obs}}}
\newcommand{\Vhid}{V_{\mathrm{hid}}}
\newcommand{\Psibar}{\bar{\Psi}}
\newcommand{\cconj}{\ddagger}          % classical conjugation of a Grassmann field
\newcommand{\hadj}{\dagger}            % Hilbert-space adjoint of the J-quantization
\newcommand{\Jm}{\mathcal{J}}          % the Krein fundamental symmetry J = -i T0 T1 T2 T3 T4
\newcommand{\KS}{\mathrm{KS}}
\newcommand{\kF}{k_{\mathrm{F}}}
\newcommand{\wF}{\omega_{\mathrm{F}}}
\newcommand{\meff}{m_{\mathrm{eff}}}
\newcommand{\Tr}{\operatorname{Tr}}
\newcommand{\file}[1]{\path{#1}}
\newcommand{\assertion}[1]{\textsc{assert:}~\emph{#1}}
\pdfstringdefDisableCommands{\def\fableScalar{fableScalar}\def\fable{fable}}
.  Build all three from this directory with
%     bash build_all.sh            (Linux, macOS, Git Bash on Windows)
%     powershell -ExecutionPolicy Bypass -File build_all.ps1
% or one of them with
%     latexmk -pdf -interaction=nonstopmode -halt-on-error PROVENANCE-14-FERMION-FABLE-CANONICAL-QUANTIZATION.tex
% The style follows fable-cosmology/latex/fable_cosmology.tex (pdfLaTeX, Latin Modern, a4, 2.5 cm).
\documentclass[11pt,a4paper]{article}

\usepackage[T1]{fontenc}
\usepackage[utf8]{inputenc}
\usepackage{lmodern}
\usepackage[margin=2.5cm]{geometry}
\usepackage{microtype}
\usepackage{amsmath,amssymb,mathtools}
\usepackage{bm}
\usepackage{graphicx}
\usepackage{booktabs}
\usepackage{tabularx}
\usepackage{longtable}
\usepackage{array}
\usepackage{caption}
\usepackage{subcaption}
\usepackage{placeins}
\usepackage{enumitem}
\usepackage{fancyvrb}
\usepackage{upquote}
\usepackage{xcolor}
\usepackage[hyphens]{url}
\usepackage[hidelinks]{hyperref}

\graphicspath{{figures/}}
\captionsetup{font=small,labelfont=bf}
\setlist{itemsep=2pt,topsep=4pt}
\fvset{fontsize=\footnotesize,xleftmargin=1em}
\newcolumntype{Y}{>{\raggedright\arraybackslash}X}
\setcounter{tocdepth}{2}
\allowdisplaybreaks

% ---- the objects of the notebook, named as the notebook names them ----------------------------
\newcommand{\fable}{\textsf{fable}}
\newcommand{\fableScalar}{\textsf{fableScalar}}
\newcommand{\Lhat}{\hat{L}}
\newcommand{\sig}{\sigma_{16}}
\newcommand{\T}[1]{T_{16}[#1]}
\newcommand{\sgn}{\sqrt{|g|}}
\newcommand{\dd}{\mathrm{d}}
\newcommand{\diag}{\operatorname{diag}}
\newcommand{\sgnf}{\operatorname{sign}}
\newcommand{\Hobs}{H_{\mathrm{obs}}}
\newcommand{\Vhid}{V_{\mathrm{hid}}}
\newcommand{\Psibar}{\bar{\Psi}}
\newcommand{\cconj}{\ddagger}          % classical conjugation of a Grassmann field
\newcommand{\hadj}{\dagger}            % Hilbert-space adjoint of the J-quantization
\newcommand{\Jm}{\mathcal{J}}          % the Krein fundamental symmetry J = -i T0 T1 T2 T3 T4
\newcommand{\KS}{\mathrm{KS}}
\newcommand{\kF}{k_{\mathrm{F}}}
\newcommand{\wF}{\omega_{\mathrm{F}}}
\newcommand{\meff}{m_{\mathrm{eff}}}
\newcommand{\Tr}{\operatorname{Tr}}
\newcommand{\file}[1]{\path{#1}}
\newcommand{\assertion}[1]{\textsc{assert:}~\emph{#1}}
\pdfstringdefDisableCommands{\def\fableScalar{fableScalar}\def\fable{fable}}

\usepackage{fvextra}

% ---- this page's own macros -------------------------------------------------------------------
% a notebook assertion label or a piece of code, typeset verbatim (the argument is already escaped)
\newcommand{\cfL}[1]{{\ttfamily\spaceskip=0.5em plus 0.6em minus 0.1em\relax #1}}
% short code snippets, typeset verbatim (\detokenize: no % or # inside)
\newcommand{\cfcode}[1]{{\ttfamily\spaceskip=0.5em plus 0.6em minus 0.1em\relax\detokenize{#1}}}
\newcommand{\cflab}[1]{``{\ttfamily\spaceskip=0.5em plus 0.6em minus 0.1em\relax\detokenize{#1}}''}
\newcommand{\cfLabel}[1]{{\frenchspacing\texttt{\detokenize{#1}}}}
\newcommand{\cfproved}[1]{\textbf{[proved #1]}}
\newcommand{\cfdisplayed}[1]{\textbf{[displayed #1]}}
\newcommand{\cfprose}[1]{\textbf{[prose #1]}}
\newcommand{\cfderived}{\textbf{[derived here]}}
\newcommand{\cffile}{\textbf{[file]}}
\newcommand{\Tsix}[1]{T_{16}[#1]}
\newcommand{\IDs}{\mathrm{ID16}}
\newcommand{\Lsym}{L_{\mathrm{sym}}}
\newcommand{\Lsymd}{L_{\mathrm{sym},\partial}}
\newcommand{\Lun}{L_{\mathrm{un}}}
\newcommand{\Lund}{L_{\mathrm{un},\partial}}
\newcommand{\That}{\hat T}
\newcommand{\Tpar}{T_{\mathrm{par}}}
\newcommand{\Tcov}{T_{\mathrm{cov}}}
\newcommand{\Gt}{\tilde G}
\newcommand{\epsKS}{\varepsilon_{\KS}}
\newcommand{\sigKS}{\sigma_{\KS}}
\newcommand{\PKS}{P_{\KS}}
\newcommand{\Pobs}{P_{\mathrm{obs}}}
\newcommand{\Phid}{P_{\mathrm{hid}}}
\newcommand{\EL}{\mathrm{EL}}
\newcommand{\sgnm}{\operatorname{sgn}}
\newcommand{\asinh}{\operatorname{asinh}}
\newcommand{\binomc}[2]{\tbinom{#1}{#2}}
\newcommand{\secref}[1]{section~\ref{#1}}
\newcommand{\Secref}[1]{Section~\ref{#1}}
\fvset{fontsize=\footnotesize,xleftmargin=1em,breaklines=true,breakanywhere=true}
\setlength{\emergencystretch}{3em}
\tolerance=1500
\pdfstringdefDisableCommands{\def\cfcode#1{#1}\def\Tsix#1{T16[#1]}}

\title{Provenance 14 --- Fermion fable: the complex 16-spinor, its canonical quantization in
$4+4$ dimensions, its field equations in the primordial gravitational field, and its
energy--momentum tensor operator}
\author{Effort A of 2026-09-24}
\date{2026-09-24}

\begin{document}
\maketitle

\noindent\textbf{Effort A of 2026-09-24.} The classical 16-component spinor field \emph{fable} of
2026-09-16 is refined into a fermion: a complex 16-component Grassmann spinor. It is quantized
canonically with the observer's time $x_4$ as time, in the $4+4$-dimensional pre-universe. Its field
equations are written out component by component in the primordial gravitational field (the
author's canonical frame, with the free function $a_4$ never given a value). Its energy--momentum
tensor operator, energy density, pressures and equations of state are derived, and the canonical
spin connection is discussed completely.

This page is complete on its own. Everything a reader needs is restated here. Code, logs, the
notebook, and the TeX and PDF files are named only as the places where things are computed,
recorded or produced.

\paragraph{How to read the evidence.} Every mathematical statement on this page carries one of these
marks:

{\small
\begin{longtable}{@{}p{0.27\linewidth}p{0.67\linewidth}@{}}
\toprule
\raggedright mark & \raggedright meaning \tabularnewline
\midrule
\endhead
\raggedright \textbf{[A]}, \cfproved{P \S n} & \raggedright an assertion \cfcode{cfAssert[label, test]} of
the Mathematica notebook, Part P, Section n. The run log prints \cfcode{PASS  label}; the label is quoted
\textbf{verbatim}, in typewriter type, exactly as the log prints it \tabularnewline
\raggedright \textbf{[D]}, \cfdisplayed{P \S n} & \raggedright a displayed output of the notebook (a table,
a count, a closed form), not an assertion \tabularnewline
\raggedright \textbf{[P]}, \cfprose{P \S n} & \raggedright stated only in a \texttt{Text} cell (or a code
comment) of the notebook; not checked by an assertion \tabularnewline
\raggedright \cfderived & \raggedright a short derivation carried out on this page from marked
results, every step shown \tabularnewline
\raggedright \cffile & \raggedright read from a named file of the repository (code, CSV, log), which is
named \tabularnewline
\bottomrule
\end{longtable}
}

Numbers are quoted from the files and logs named next to them. Nothing is invented. $a_4$ is never
given a value anywhere.

\paragraph{Two kinds of ``section''.} ``Section n'' with a capital S is a section of the Mathematica
notebook (Sections 1--33; Part~VII is Sections 26--29). ``section n'' in lower case is a section of
this page. The mathematics is written in standard notation; where the notebook's Wolfram names
matter they are given in typewriter type. $\Tsix{a}$ (\cfcode{T16[a]}) are the $16\times16$ Dirac
matrices, $\sigma_{16}$ (\cfcode{sigma16}) the spinor metric, $\Psi^{\ddagger}$ the classical
conjugate, $\bar\Psi$ the Dirac conjugate, $\partial_\mu$ a partial derivative, $D_\mu$ the
spin-covariant derivative, and $\sqrt g$ (\cfcode{Sqrt[g]}) the volume factor $\sqrt{|\det g|}$.

\tableofcontents

\section*{Where things live}
\addcontentsline{toc}{section}{Where things live}
\label{sec:where}

{\small
\begin{longtable}{@{}p{0.42\linewidth}p{0.52\linewidth}@{}}
\toprule
\raggedright what & \raggedright where (relative to the repository root) \tabularnewline
\midrule
\endhead
\raggedright the Mathematica notebook (generated, never edited by hand) & \raggedright \file{claude-fable/claude-fable_Einstein-Rosen-2-Planes.nb} \tabularnewline
\raggedright its cell manifests; Part~VII (Sections 26--29) is this effort & \raggedright \file{claude-fable/cells_part1.wl} \dots\ \file{cells_part8.wl}; \textbf{\file{claude-fable/cells_part7.wl}} \tabularnewline
\raggedright the builder and the run harness & \raggedright \file{claude-fable/build_tools.py}, \file{claude-fable/runner_header.wl} \tabularnewline
\raggedright the runner, the checker, the section map & \raggedright \file{claude-fable/run_from_nb.wls}, \file{claude-fable/verify_nb.wls}, \file{claude-fable/nb_section_map.wls} \tabularnewline
\raggedright Part~VII's own acceptance run, the record of Part~VII: the notebook through Part~VII (198 Input cells, 528/528, 0 cells with messages) & \raggedright \file{claude-fable/run_fermion_fable_part7.log} \tabularnewline
\raggedright the acceptance run of Part~VIII, kept as a historical record (217 Input cells, 642/642: the Rust solver's CSVs did not exist yet, so Part~VIII's two Rust-vs-Mathematica comparisons were skipped, and three Section~32 labels still carry the superseded index notation; \secref{sec:10.3}) & \raggedright \file{claude-fable/run_fermion_fable_part8.log} \tabularnewline
\raggedright the final run of the whole notebook, Parts I--VIII, in its committed state (217 Input cells, 644/644, 0 cells with messages, 198.352\,s, 46 witnesses, 0 identities accepted on numerical evidence alone) & \raggedright \file{claude-fable/run_fermion_fable_final.log} \tabularnewline
\raggedright the checker's logs; the section map & \raggedright \file{claude-fable/verify_nb_part7.log}, \file{claude-fable/verify_nb_part8.log}; \file{claude-fable/nb_section_map.log} \tabularnewline
\raggedright the TeX/text export of the Part~VII results & \raggedright \file{claude-fable/export_fermion_fable_tex.wls} $\to$ \file{provenance-latex/generated/*.tex}, \file{*.txt} \tabularnewline
\raggedright the LaTeX twin of this page, its preamble and build scripts & \raggedright \file{provenance-latex/PROVENANCE-14-FERMION-FABLE-CANONICAL-QUANTIZATION.tex} $\to$ \file{.pdf}; \file{provenance-latex/preamble.tex}; \file{provenance-latex/build_all.sh}, \file{build_all.ps1} \tabularnewline
\raggedright the author's two helper packages, which must sit next to the notebook & \raggedright \file{claude-fable/ConvertMapleToMathematicaV2.wl}, \file{claude-fable/EtoExp.wl} \tabularnewline
\raggedright the numerically stable Kohn--Sham integrals and Dirac-sea energy used by the solver \file{fable_fermion} (final at \texttt{2bc936d}: 34/34 tests) & \raggedright \file{fable-cosmology/rust/fable_fermion/src/kohn_sham.rs} \tabularnewline
\bottomrule
\end{longtable}
}

\noindent\texttt{<repo>} below is wherever the repository is cloned. Every shell block finds it with
\cfcode{git rev-parse --show-toplevel}.

% ===============================================================================================
\setcounter{section}{-1}
\section{The task, in the author's words}
\label{sec:0}

On 2026-09-24 the author gave three instructions. Their operative parts, verbatim (the ellipses
[\dots] are cuts in the quotation, not in the instruction's meaning):

\begin{quote}
\textbf{A.} ``Refine the new 16-component fermi spinor field named fable to be a complex 16-spinor,
and quantize it using canonical quantization, extended to 4+4 dimensions. Write out the field
equations for the fermion fable field in the presence of a primordial Gravitational Field.
Create new provenance markdown file, latex, and pdf files, that contain the exact, correct,
fermion fable field equations, energy-momentum tensor operator, pressure, density and equations
of state for fermion fable field. Include a complete discussion of the canonical spin
connection.''

\textbf{B.} ``Include fable as a source in the Einstein Gravitational Field Equations for this
primordial Gravitational Field, and solve the coupled field equations over a time interval from
the beginning of the present universe until the present time. Try to use some of the relevant
ideas from Density Functional Theory to obtain the fable ground and first excited states. Create
another new provenance markdown file, latex, and pdf files [\dots] the interacting (fermion fable
field, primordial Gravitational Field) system [\dots] coupled equations, canonical spin connection
for the interacting system, energy-momentum tensor operator, pressure, density and equations of
state [\dots] complete discussion of the canonical spin connections. Answer: [1] Does this system
provide a physical mechanism for a time-varying dark energy equation of state (w)? [2] [\dots]
time-varying dark matter equation of state (w)?''

\textbf{C.} ``Thoroughly discuss the commands that you use to solve the coupled system, and to test,
verify, execute, and display your calculations of the coupled field equations. Create another
new provenance markdown file, latex, and pdf files that contain your complete solution. Each of
your major efforts should have its own section in the markdown file, latex, and pdf files. Do
not ever refer the USER to another markdown file for any reason. The full set of 100\% complete
ideas must be given its own section in each markdown file.'' And: push, check and verify the
repository. ``DO NOT TAKE SHORTCUTS.''
\end{quote}

\noindent\textbf{This page is Effort A.} Efforts B and C have their own pages. The complete set of
ideas of all three efforts is \secref{sec:1} below, identical on all three pages.

\paragraph{The author's choice.} Instruction A itself names the field: ``a complex 16-spinor''.
As \secref{sec:3} shows, that is forced by the author's own algebra: a real anticommuting
16-spinor with the author's spinor metric has no dynamics, a Majorana 16-spinor is massless with no
potential, and a one-chirality (Weyl or Majorana--Weyl) spinor cannot propagate. When these
alternatives were laid out on 2026-09-24, the author confirmed the choice: \textbf{a complex
16-spinor}, with Dirac conjugate $\bar\Psi=\Psi^{\ddagger}\sigma_{16}$. Every result below is for
that field.

\paragraph{The work, in the order it was done.} A design (revision~1) was written and then reviewed
adversarially by four reviewers --- quantization; field equations, energy--momentum tensor and spin
connection; Einstein equations and cosmology; density-functional theory and numerics --- each
followed by an independent skeptic. Every finding was adopted, some only in corrected form
(\secref{sec:10.4}). The physics was then implemented as Part~VII of the Mathematica notebook
(Sections 26--29), where every claim is an assertion. The notebook was run end to end (528/528
assertions, 0 messages). Part~VIII (the coupled system of Effort B, Sections 30--33) was added
afterwards; the final run of the whole notebook, Parts I--VIII, evaluates 217 Input cells with 644/644
assertions and 0 cells with messages, and its Part~VII lines are those of the Part~VII run, line for
line apart from the timings (\secref{sec:10.3}). The exact field equations, anticommutator, fundamental symmetry and
equations of state were exported from the notebook's own objects to TeX and plain text, and those
exports are quoted verbatim on this page.

% ===============================================================================================
\section{The complete set of ideas}
\label{sec:1}

\texttt{\{\{COMPLETE-SET-OF-IDEAS\}\}}

% ===============================================================================================
% Section 2: the shared foundations section (docs-shared/foundations.tex), with the Part VII
% placeholders filled from claude-fable/run_fermion_fable_part7.log.
\begingroup\sloppy

\section[Foundations: the pre-universe, its algebra, the canonical frame, the author's La and eLa, and the classical fable of Part VI]{Foundations: the pre-universe, its algebra, the canonical frame, the author's \texorpdfstring{$L_a$}{La} and \texttt{eLa}, and the classical fable of Part~VI}
\label{fnd:sec}\label{sec:2}

This section sets out everything that the new work on fermion fable stands on. It restates the
author's pre-universe (a curved $4+4$-dimensional spacetime), the real Clifford algebra and
split-octonion spinors built on it, the author's own wave function, Lagrangian and field
equations, the canonical frame field, and the classical field fable of 2026-09-16. It ends with
the reason why that classical field has to be refined into a complex 16-component spinor. The
reader needs nothing outside this document to follow it.

\subsection{How the evidence is cited}
\label{fnd:evidence}\label{sec:2.1}

All the mathematics lives in one Mathematica notebook,
\texttt{claude-fable/claude-fable\_Einstein-Rosen-2-Planes.nb}. That notebook is never edited by
hand. Its Parts I--VI, which this section restates, are generated from the cell manifests
\texttt{claude-fable/cells\_part1.wl} \dots\ \texttt{cells\_part6.wl}. In a manifest,
\texttt{Text} cells explain and \texttt{Input} cells compute. Every claim the notebook checks is
an assertion \texttt{cfAssert[label, test]}, which prints \texttt{PASS~~label} or
\texttt{FAIL~~label}. The complete run of Parts I--VI is recorded in
\texttt{claude-fable/run\_fable51\_part6.log}. It evaluated 172 Input cells in 162.5\,s with no
cell raising a message, and its final tally reads (the witness line is shortened here; in the
log it continues with a one-sentence explanation of what a witness is):
\begin{footnotesize}
\begin{verbatim}
identities accepted on numerical evidence alone : 0   (stage 3 returned True)
non-vanishing witnesses                         : 27
  PASS  no identity in this notebook was accepted on numerical evidence alone
Assertions run: 358   passed: 358   failed: 0
\end{verbatim}
\end{footnotesize}

Each statement below is tagged with one of four kinds of evidence:
\begin{itemize}
  \item \textbf{[A]} an assertion. Its label is quoted exactly as the log prints it after
        \texttt{PASS}.
  \item \textbf{[D]} a displayed result: an Output the notebook prints but does not assert.
        Where a displayed result was re-read for this document from the saved kernel state,
        that is said.
  \item \textbf{[P]} prose only: a statement that appears in a \texttt{Text} cell and is not
        checked by any assertion.
  \item \textbf{[VII]} proved in Part VII, the new part of the notebook written for fermion
        fable (Sections 26--29). Its labels are quoted verbatim from
        \file{claude-fable/run_fermion_fable_part7.log}, which reports
        \cfcode{Assertions run: 528   passed: 528   failed: 0} for Parts I--VII together.
\end{itemize}

How an identity is certified (restated from Section~15 of the notebook). \texttt{cfZeroQ} and
\texttt{cfZeroArrayQ} try three stages in turn:
\begin{enumerate}
  \item exact structural zero;
  \item \texttt{Simplify} under the geometric assumptions, with a 5-second budget;
  \item a 30-digit numerical evaluation at three rational probe points, with test functions
        substituted for the undetermined $a_4$ (and for the undetermined rapidity of the
        notebook's boosted-frame bridge); the result is accepted as zero only if its magnitude
        is below $10^{-22}$.
\end{enumerate}
An identity accepted at stage~3 counts as a \emph{numerical certificate}, and the run reports
zero of those: every identity in Parts I--VI is closed symbolically. A claim that something is
\emph{not} zero is made with \texttt{cfNonZeroWitnessQ}. It passes only if, at some probe point,
every entry evaluates to a number and at least one of them has magnitude above $10^{-10}$. That
is a complete proof of non-vanishing; the run has 27 such witnesses. The probe functions are
test inputs only, never definitions.

The session policy is the original notebook's, unchanged: \texttt{Simplify} has a 1-second budget
and \texttt{FullSimplify} a 3-second budget. \texttt{DIM8 = 8}, and the model's three constants
$M$ (a mass), $K$ and $H$ (a Hubble-like constant) are \texttt{Protect}ed, so no cell can assign
them (Section~1: a \texttt{Text} cell and the \texttt{Input} cells that set this up; no
assertion) [P].

\subsection{Coordinates, the flat 4+4 metric, and which directions are timelike}
\label{fnd:coords}\label{sec:2.2}

The coordinates are $X=\{x_0,x_1,\dots,x_7\}$. The flat tangent-space metric is
\begin{equation}
  \texttt{eta4488} = \eta = \operatorname{diag}(+1,+1,+1,+1,-1,-1,-1,-1),\qquad
  \text{signature }(4,4).
\end{equation}

\begin{center}
\small
\begin{tabular}{@{}llp{0.30\textwidth}p{0.30\textwidth}@{}}
\toprule
coordinate & $\eta$ & name in the original notebook (Section 2) & reading used in this work\\
\midrule
$x_0$ & $+1$ & hidden space & the hidden \textbf{spacelike} coordinate\\
$x_1,x_2,x_3$ & $+1$ & ordinary 3-space & the \textbf{observed sheet}\\
$x_4$ & $-1$ & time & the \textbf{observer's time}\\
$x_5,x_6,x_7$ & $-1$ & superluminal deflating time directions & the hidden \textbf{timelike sheet}\\
\bottomrule
\end{tabular}
\end{center}

[A] \cfLabel{eta4488 . eta4488 == ID8}; \cfLabel{eta4488 is symmetric};
\cfLabel{signature of eta4488 is (4,4)}.

Index conventions (restated from Section~2). Latin indices $a,b,A,B$ are flat and run $0..7$.
Greek indices $\mu,\nu$ are curved and run $0..7$. Wolfram arrays are 1-based, so index $A$ sits
at array position $A+1$. For example, $x_4$ is entry 5 of $X$, so Part VI's energy density
$\rho=T_{44}$ is the array element \texttt{T[[5,5]]}.

The notebook uses two sets of assumptions:
\begin{itemize}
  \item \texttt{ssX = H > 0 \&\& sX0 \&\& 6 H x0 > 0 \&\& 2 la[x4] > 0 \&\& Cot[6 H x0] > 0
        \&\& Sin[6 H x0] > 0}, with \texttt{sX0 = x0 > 0 \&\& x1 > 0 \&\& ... \&\& x7 > 0} (all
        eight coordinates positive), for the $(x_0,x_4)$ chart;
  \item \texttt{constraintVars = x0 > 0 \&\& x4 > 0 \&\& z > 0 \&\& t > 0}, for the chart
        $z = 6Hx_0$, $t = Hx_4$.
\end{itemize}
\textbf{\texttt{la} is a free function inherited from the original notebook.} It occurs only
inside \texttt{ssX}. It is never given a value, and nothing here invents one [P, Section~2].

\subsection{SO(4): the self-dual and anti-self-dual blocks}
\label{fnd:so4}\label{sec:2.3}

The six antisymmetric real $4\times4$ matrices split under the Hodge star into a self-dual
3-space and an anti-self-dual 3-space, the two commuting $\mathfrak{su}(2)$ factors of
$\mathfrak{so}(4)$. The original notebook builds both from two elementary tensors (Section~3):
\begin{gather}
  Q_a[h,p,q] = \operatorname{Signature}[\{h,p,q,4\}],\qquad
  Q_b[h,p,q] = \delta_{p4}\delta_{qh} - \delta_{ph}\delta_{q4},\\
  \texttt{s4by4}[h] = Q_a - Q_b \ \text{(self-dual)},\qquad
  \texttt{t4by4}[h] = Q_a + Q_b \ \text{(anti-self-dual)},\\
  h=1,2,3.
\end{gather}
Explicitly [D]:
\begin{gather*}
s_1=\left(\begin{smallmatrix}0&0&0&1\\0&0&1&0\\0&-1&0&0\\-1&0&0&0\end{smallmatrix}\right),\quad
s_2=\left(\begin{smallmatrix}0&0&-1&0\\0&0&0&1\\1&0&0&0\\0&-1&0&0\end{smallmatrix}\right),\quad
s_3=\left(\begin{smallmatrix}0&1&0&0\\-1&0&0&0\\0&0&0&1\\0&0&-1&0\end{smallmatrix}\right),\\
t_1=\left(\begin{smallmatrix}0&0&0&-1\\0&0&1&0\\0&-1&0&0\\1&0&0&0\end{smallmatrix}\right),\quad
t_2=\left(\begin{smallmatrix}0&0&-1&0\\0&0&0&-1\\1&0&0&0\\0&1&0&0\end{smallmatrix}\right),\quad
t_3=\left(\begin{smallmatrix}0&1&0&0\\-1&0&0&0\\0&0&0&-1\\0&0&1&0\end{smallmatrix}\right),
\end{gather*}
with $s_h = \texttt{s4by4}[h]$ and $t_h = \texttt{t4by4}[h]$.

[A] \cfLabel{s4by4 is self-dual}; \cfLabel{t4by4 is anti-self-dual};
\cfLabel{s4by4 and t4by4 are antisymmetric}; \cfLabel{the two su(2) factors commute}; and, for the
nine mixed products $\texttt{st}[J,K]=s_J\,t_K$, \cfLabel{st[J,K] is symmetric}.

\subsection{The real \texorpdfstring{$8\times8$}{8x8} Clifford generators \texorpdfstring{$\tau$, $\bar\tau$}{tau, taubar}, and the spinor metric \texorpdfstring{$\sigma$}{sigma}}
\label{fnd:tau}\label{sec:2.4}

The O(4,4) spinor metric is
\begin{equation}
  \sigma = \begin{pmatrix} 0 & I_4\\ I_4 & 0\end{pmatrix}.
\end{equation}
[A] \cfLabel{sigma . sigma == ID8}; \cfLabel{sigma is symmetric}; \cfLabel{Tr[sigma] == 0}. The
type-1 and type-2 split-octonion spinors both carry this same $\sigma$, because
$\sigma=\sigma^{-1}$ [P].

The generators (Section~4) are
\begin{align}
  \tau[0] &= I_8, \\
  \tau[h] &= \begin{pmatrix} 0 & s_h\\ s_h & 0\end{pmatrix}, \qquad h=1,2,3,\\
  \tau[3+k] &= \begin{pmatrix} 0 & t_{4-k}\\ -t_{4-k} & 0\end{pmatrix}, \qquad k=1,2,3\quad(t_3,t_2,t_1),\\
  \tau[7] &= \tau[1]\tau[2]\tau[3]\tau[4]\tau[5]\tau[6] = \operatorname{diag}(-I_4,\,+I_4)\quad\text{[D]},\\
  \bar\tau[A] &= \sigma\,(\sigma\,\tau[A])^{T},\qquad A=0,\dots,7\quad(\text{the $\sigma$-adjoint}).
\end{align}
The order $3,2,1$ is the original's. It is what makes $\tau[7]$ the product of the first six and
$\sigma=\tau[1]\tau[2]\tau[3]$ [P]. The original calls the list of the six blocks
\texttt{sixAntiSymmetric8by8}, but that name fits only half of it: [A]
\cfLabel{blocks 1-3 (built from s4by4) are antisymmetric};
\cfLabel{blocks 4-6 (built from t4by4) are symmetric}.

The defining relation is the split Clifford algebra of the $4+4$ quadratic form in its real
$8\times8$ ``half-spinor'' form:
\begin{equation}
  \tfrac12\bigl(\tau[A]\,\bar\tau[B] + \tau[B]\,\bar\tau[A]\bigr) = \eta_{AB}\,I_8 .
\end{equation}
[A] \cfLabel{(1/2)(tau[A].taubar[B] + tau[B].taubar[A]) == eta4488[[A+1,B+1]] ID8}.

The ordering identities [A]:
\begin{itemize}
  \item \cfLabel{sigma == tau[1].tau[2].tau[3]}
  \item \cfLabel{sigma == tau[4].tau[5].tau[6].tau[7]}
  \item \cfLabel{tau[1]...tau[7] == tau[0] == ID8}
  \item \cfLabel{sigma . taubar[A] == Transpose[sigma . tau[A]]}
  \item \cfLabel{-eta4488[[A+1,A+1]] tau[A] == Transpose[tau[A]] for A = 1..7}
\end{itemize}
The last one says that $\tau[1..3]$ are antisymmetric and $\tau[4..7]$ are symmetric. The second
invariant of the original is $\Omega=\sigma\,\tau[7]$: [A] \cfLabel{Omega == tau[4].tau[5].tau[6]}.

A complete basis of the $8\times8$ algebra (Section~7) is the 63 ordered products of 1, 2 or 3 of
$\tau[1..7]$, together with $I_8$. [A] \cfLabel{the 8x8 basis has 64 elements};
\cfLabel{the 64 elements are distinct}; \cfLabel{(1/8) Tr[bas64[A].bas64[B]] == eta64[[A,B]]};
\cfLabel{symmetric 8x8 basis elements number 36};
\cfLabel{antisymmetric 8x8 basis elements number 28}.

\subsection{The \texorpdfstring{$16\times16$}{16x16} Dirac matrices, chirality, \texorpdfstring{$\sigma_{16}$}{sigma16}, and \texorpdfstring{$\mathfrak{so}(4,4)$}{so(4,4)}}
\label{fnd:T16}\label{sec:2.5}

Following E.~A.~Lord's reduced Brauer--Weyl construction (Math.\ Proc.\ Camb.\ Phil.\ Soc.\ 64
(1968) 765--778), the $8\times8$ generators are doubled off-diagonally (Section~5):
\begin{gather}
  T_{16}[A] = \begin{pmatrix} 0 & \bar\tau[A]\\ \tau[A] & 0\end{pmatrix}\quad (A=0,\dots,7),\\
  T_{16}[8] = T_{16}[0]T_{16}[1]\cdots T_{16}[7],\qquad
  \sigma_{16} = T_{16}[0]T_{16}[1]T_{16}[2]T_{16}[3].
\end{gather}
The $T_{16}[A]$ are real $16\times16$ matrices, and ``$\gamma[a]$'' means $T_{16}[a]$. $T_{16}[8]$
is the chirality operator and $\sigma_{16}$ is the $16\times16$ spinor metric.

Assertions [A]:
\begin{itemize}
  \item \cfLabel{{T16[A], T16[B]}/2 == eta4488[[A+1,B+1]] ID16}
  \item \cfLabel{sigma16 == ArrayFlatten[{{-sigma,0},{0,sigma}}]}
  \item \cfLabel{sigma16 . T16[A] is antisymmetric for A = 0..7}
  \item \cfLabel{T16[8] == sigma16 . T16[4].T16[5].T16[6].T16[7]}
  \item \cfLabel{sigma16 is symmetric and sigma16 . sigma16 == ID16}
  \item \cfLabel{sigma16 . covariantDiffMatrix is symmetric}, with
        $\texttt{covariantDiffMatrix}=T_{16}[5]T_{16}[6]T_{16}[7]$
\end{itemize}

\paragraph{The two chiralities are the two split-octonion spinors.} The notebook defines
$P_L=\tfrac12(I_{16}-T_{16}[8])$ and $P_R=\tfrac12(I_{16}+T_{16}[8])$. [A]
\cfLabel{PL + PR == ID16}; \cfLabel{PL is idempotent}; \cfLabel{PR is idempotent};
\cfLabel{PL . PR == PR . PL == 0};
\cfLabel{PL projects onto the upper (type-1) components, PR onto the lower (type-2) components}.

The last assertion is equivalent to
\begin{equation}
  T_{16}[8] = \operatorname{diag}(-I_8,\,+I_8).
\end{equation}
A 16-component spinor is therefore the direct sum $\Psi=(\psi_1,\psi_2)$:
\begin{itemize}
  \item the upper eight components are the \textbf{type-1} split-octonion spinor, of
        chirality $-1$;
  \item the lower eight components are the \textbf{type-2} spinor, of chirality $+1$.
\end{itemize}
Every $T_{16}[A]$ exchanges the two halves, since it is block off-diagonal [P].

\paragraph{The so(4,4) generators} are
\begin{equation}
  S^{AB} = \texttt{SAB[[A+1,B+1]]} = \tfrac14\,[T_{16}[A],T_{16}[B]].
\end{equation}
Assertions [A]:
\begin{itemize}
  \item \cfLabel{SAB is antisymmetric in its two flat indices}
  \item \cfLabel{sigma16 . SAB is antisymmetric}
  \item \cfLabel{T16[8] commutes with every SAB: the chiral halves are each Spin(4,4)-invariant}
  \item \cfLabel{[SAB,SAB] closes on the so(4,4) algebra}, which is the relation
        $[S_{AB},S_{CD}] = -(\eta_{AC}S_{BD}-\eta_{AD}S_{BC}-\eta_{BC}S_{AD}+\eta_{BD}S_{AC})$
  \item \cfLabel{[SAB, T16] reproduces the vector representation}, which is
        $[S_{AB},\gamma_C] = -\eta_{CA}\gamma_B+\eta_{CB}\gamma_A$
\end{itemize}
The spin connection of the frame field enters as
$\tfrac18\,\omega_{\mu ab}[\gamma^a,\gamma^b]=\tfrac12\,\omega_{\mu ab}S^{ab}$.

\paragraph{The 256-element basis} (Section~6) is the ordered products of $k$ distinct
generators, $k=1..8$, followed by $I_{16}$, taken in the original's order. [A]
\cfLabel{base16[[93]] is sigma16} (the word $\{0,1,2,3\}$);
\cfLabel{base16[[255]] is the chirality operator T16[8]}; \cfLabel{base16[[256]] is the identity};
\cfLabel{Tr[M.M]/16 is +1 or -1 for every basis element};
\cfLabel{136 basis elements have Tr[M.M] > 0}; \cfLabel{120 basis elements have Tr[M.M] < 0};
\cfLabel{symmetric basis elements number 16*17/2 = 136};
\cfLabel{antisymmetric basis elements number 16*15/2 = 120};
\cfLabel{Tr[M.M] > 0 exactly when M is symmetric};
\cfLabel{Tr[M.M] < 0 exactly when M is antisymmetric}.

Section~8 also builds the ordinary $4\times4$ Dirac algebra of the $(x_1,x_2,x_3,x_4)$ block,
whose metric is $g_{44}=\operatorname{diag}(1,1,1,-1)$. For example [A]:
\cfLabel{{gamma4[h], gamma4[k]} == 2 g44[[h,k]] ID4  for h,k = 1..4};
\cfLabel{[S44,S44] closes on the so(3,1) algebra}. It plays no role below.

\subsection{Cartan triality and the split octonions (Section 9): what is proved}
\label{fnd:triality}\label{sec:2.6}

The vector, the type-1 spinor and the type-2 spinor are three 8-dimensional representations of
Spin(4,4) that triality relates. The original makes the vector-to-spinor bridge concrete with
one distinguished spinor, \texttt{unit}. It is row 8 of the eigenvectors of $\sigma$, rescaled by
$1/\sqrt2$. Assertions [A]:
\begin{itemize}
  \item \cfLabel{unit is pinned to the author's spinor {1/Sqrt[2],0,0,0,1/Sqrt[2],0,0,0} -- the row order of Eigensystem is not relied on silently}
  \item \cfLabel{each uEig row is sigma-normalized to +/-1}
  \item \cfLabel{unit is sigma-normalized to +1}
\end{itemize}
The bridge and its inverse are
\begin{gather}
  \texttt{triVecToSpin}[[A+1,a+1]] = (\texttt{unit}\cdot\sigma\cdot\bar\tau[A])[[a+1]],\\
  \texttt{triSpinToVec}=\texttt{triVecToSpin}^{-1},
\end{gather}
where the row index $A$ of \texttt{triVecToSpin} is a flat vector index and the column index $a$
is a type-1 spinor index.
Assertions [A]:
\begin{itemize}
  \item \cfLabel{triVecToSpin . triSpinToVec == ID8}
  \item \cfLabel{triSpinToVec . triVecToSpin == ID8}
  \item \cfLabel{triSpinToVec == Transpose[ eta_AA tau[A].unit ]  (the original's closed form)}
  \item \cfLabel{etaTri is symmetric}; \cfLabel{etaTri has signature (4,4)}
  \item \cfLabel{triVecToSpin transports etaTri back to eta4488}
  \item \cfLabel{THE TRIALITY METRIC IS EXACTLY THE SPINOR METRIC:  etaTri == sigma}
\end{itemize}
The split-octonion structure constants are read off from the generators through the bridge. In
the type-1 spinor basis they are \texttt{mabc}; in the vector basis they are \texttt{mABC}, the
usual multiplication table. Assertions [A]:
\begin{itemize}
  \item \cfLabel{eA[1] is a two-sided identity for the vector-basis product}
  \item \cfLabel{<eA, eB> composition is compatible with eta (basis-unit norm form)}
  \item \cfLabel{m_{ABC} restricted to imaginary directions is totally antisymmetric}
\end{itemize}
The scope of what is proved: the norm is multiplicative \emph{on the basis units}, and the
imaginary structure constants are totally antisymmetric. The general composition law
$N(xy)=N(x)N(y)$ for arbitrary elements, and alternativity, are \textbf{not} asserted.

\subsection{The wave function of the un-universe and the author's Lagrangian \texorpdfstring{$L_a$}{La} (Section 11)}
\label{fnd:La}\label{sec:2.7}

The original notebook calls \texttt{Psi16} the wave function of the ``un-universe''. It is
sixteen scalar functions of the hidden coordinate $x_0$ and the time $x_4$:
\begin{equation}
  \Psi_{16} = (f_0,\dots,f_{15}),\qquad f_k \equiv \texttt{f16[k][x0, x4]},
\end{equation}
where the upper eight components are the type-1 spinor and the lower eight the type-2 spinor.
[A] \cfLabel{Psi16 is the direct sum of the type-1 and type-2 spinors}. Two re-parametrizations
are recorded:
\begin{itemize}
  \item \texttt{f16[k] -> Z[k][6 H \#1, H \#2] \&}, which is the one used downstream;
  \item \texttt{f16[k] -> nZ[k][6 H \#1, H \#2]/Sqrt[Sin[6 H \#1]] \&}, the author's alternative
        normalization.
\end{itemize}

\paragraph{The author's Lagrangian, exactly as defined in Section 11:}
\begin{footnotesize}
\begin{verbatim}
La[] := ((1/H) Transpose[Psi16] . sigma16 .
          Sum[T16[alpha1 - 1] . D[Psi16, X[[alpha1]]], {alpha1, 1, Length[X]}]
        + 2 (M/H) Transpose[Psi16] . sigma16 . Psi16
       ) // Simplify[#, constraintVars] &
\end{verbatim}
\end{footnotesize}
that is,
\begin{align}
  L_a &= \frac1H\,\Psi_{16}^{T}\sigma_{16}\,T_{16}[\alpha]\,\partial_\alpha\Psi_{16}
        + \frac{2M}{H}\,\Psi_{16}^{T}\sigma_{16}\Psi_{16} \\
      &= \frac1H\,\Psi_{16}^{T}\sigma_{16}\bigl(T_{16}[0]\,\partial_{x_0}+T_{16}[4]\,\partial_{x_4}\bigr)\Psi_{16}
        + \frac{2M}{H}\,s,\qquad s=\Psi_{16}^{T}\sigma_{16}\Psi_{16}.
\end{align}
The second line holds because $\Psi_{16}$ depends only on $x_0$ and $x_4$. It was checked for
this document against the evaluated \texttt{La[]} (a \texttt{Simplify} difference of exactly 0);
it is not a notebook assertion. The mass term is $(2M/H)\,s$: it is $+2M/H$ times the scalar
bilinear $s$.

The evaluated \texttt{La[]} [D] is shown below. Write $f_k=\texttt{f16[k][x0,x4]}$,
$\partial_0 f_k=\texttt{Derivative[1,0][f16[k]][x0,x4]}$ and
$\partial_4 f_k=\texttt{Derivative[0,1][f16[k]][x0,x4]}$. Then $H\,L_a$ is exactly
\begin{align*}
H L_a ={}& -4M\,(f_0f_4+f_1f_5+f_2f_6+f_3f_7-f_8f_{12}-f_9f_{13}-f_{10}f_{14}-f_{11}f_{15})\\
&+f_{12}(\partial_4f_5+\partial_0f_0)+f_{13}(-\partial_4f_4+\partial_0f_1)+f_{14}(-\partial_4f_7+\partial_0f_2)\\
&+f_{15}(\partial_4f_6+\partial_0f_3)+f_{8}(-\partial_4f_1+\partial_0f_4)+f_{9}(\partial_4f_0+\partial_0f_5)\\
&+f_{10}(\partial_4f_3+\partial_0f_6)+f_{11}(-\partial_4f_2+\partial_0f_7)+f_{4}(\partial_4f_{13}-\partial_0f_8)\\
&-f_{5}(\partial_4f_{12}+\partial_0f_9)-f_{6}(\partial_4f_{15}+\partial_0f_{10})+f_{7}(\partial_4f_{14}-\partial_0f_{11})\\
&-f_{0}(\partial_4f_9+\partial_0f_{12})+f_{1}(\partial_4f_8-\partial_0f_{13})+f_{2}(\partial_4f_{11}-\partial_0f_{14})\\
&-f_{3}(\partial_4f_{10}+\partial_0f_{15}).
\end{align*}

\paragraph{The structural identity behind it.} $\sigma_{16}T_{16}[A]$ is antisymmetric
(Section~5). An antisymmetric form vanishes on the diagonal, so for commuting components
$\Psi_{16}^{T}\sigma_{16}T_{16}[\alpha]\Psi_{16}=0$ identically. The kinetic term therefore equals
minus its transpose, with no boundary term. [A]
\cfLabel{Psi16 . sigma16 . T16[alpha] . Psi16 == 0 identically (antisymmetric form)};
\cfLabel{the kinetic term equals minus its transpose, exactly};
\cfLabel{La[] == its conjugate form when the coefficients match}. The last one is
$L_a = -\tfrac1H(\partial_\alpha\Psi_{16})^{T}\sigma_{16}T_{16}[\alpha]\Psi_{16}
+\tfrac{2M}{H}\Psi_{16}^{T}\sigma_{16}\Psi_{16}$. The original's scratch cell compared $L_a$ with a
form carrying mismatched factors of $H$ and so found a non-zero difference [P].

\subsection{The Euler--Lagrange equations \texttt{eLa} (Section 12)}
\label{fnd:eLa}\label{sec:2.8}

The original's Euler--Lagrange operator, transcribed exactly:
\begin{footnotesize}
\begin{verbatim}
eL[Lagrangian_Symbol, detsqrt_] := Module[{L, t},
  L = Lagrangian[];
  t = Table[
    FullSimplify[
      (1/detsqrt) ( D[L, f16[k][x0, x4]]
                    - D[D[L, Derivative[1, 0][f16[k]][x0, x4]], x0]
                    - D[D[L, Derivative[0, 1][f16[k]][x0, x4]], x4] ),
      constraintVars],
    {k, 0, 15}];
  Return[t /. subsDefects]];                (* subsDefects = {} *)
eLa = eL[La, 1]
\end{verbatim}
\end{footnotesize}
[A] \cfLabel{there are sixteen field equations}. $\texttt{eLa[[k]]}$ is the Euler--Lagrange
expression for the component $f_{k-1}$, and the field equations are $\texttt{eLa[[k]]}=0$.
\textbf{The sixteen expressions, exactly as the notebook computes them} [D], in the notation above:
\begin{align*}
\texttt{eLa[[1]]} &= -\tfrac{2}{H}\bigl(2Mf_{4}+\partial_4f_{9}+\partial_0f_{12}\bigr) &&(\text{varied: }f_{0})\\
\texttt{eLa[[2]]} &= -\tfrac{2}{H}\bigl(2Mf_{5}-\partial_4f_{8}+\partial_0f_{13}\bigr) &&(\text{varied: }f_{1})\\
\texttt{eLa[[3]]} &= -\tfrac{2}{H}\bigl(2Mf_{6}-\partial_4f_{11}+\partial_0f_{14}\bigr) &&(\text{varied: }f_{2})\\
\texttt{eLa[[4]]} &= -\tfrac{2}{H}\bigl(2Mf_{7}+\partial_4f_{10}+\partial_0f_{15}\bigr) &&(\text{varied: }f_{3})\\
\texttt{eLa[[5]]} &= -\tfrac{2}{H}\bigl(2Mf_{0}-\partial_4f_{13}+\partial_0f_{8}\bigr) &&(\text{varied: }f_{4})\\
\texttt{eLa[[6]]} &= -\tfrac{2}{H}\bigl(2Mf_{1}+\partial_4f_{12}+\partial_0f_{9}\bigr) &&(\text{varied: }f_{5})\\
\texttt{eLa[[7]]} &= -\tfrac{2}{H}\bigl(2Mf_{2}+\partial_4f_{15}+\partial_0f_{10}\bigr) &&(\text{varied: }f_{6})\\
\texttt{eLa[[8]]} &= -\tfrac{2}{H}\bigl(2Mf_{3}-\partial_4f_{14}+\partial_0f_{11}\bigr) &&(\text{varied: }f_{7})\\
\texttt{eLa[[9]]} &= \phantom{-}\tfrac{2}{H}\bigl(2Mf_{12}-\partial_4f_{1}+\partial_0f_{4}\bigr) &&(\text{varied: }f_{8})\\
\texttt{eLa[[10]]} &= \phantom{-}\tfrac{2}{H}\bigl(2Mf_{13}+\partial_4f_{0}+\partial_0f_{5}\bigr) &&(\text{varied: }f_{9})\\
\texttt{eLa[[11]]} &= \phantom{-}\tfrac{2}{H}\bigl(2Mf_{14}+\partial_4f_{3}+\partial_0f_{6}\bigr) &&(\text{varied: }f_{10})\\
\texttt{eLa[[12]]} &= \phantom{-}\tfrac{2}{H}\bigl(2Mf_{15}-\partial_4f_{2}+\partial_0f_{7}\bigr) &&(\text{varied: }f_{11})\\
\texttt{eLa[[13]]} &= \phantom{-}\tfrac{2}{H}\bigl(2Mf_{8}+\partial_4f_{5}+\partial_0f_{0}\bigr) &&(\text{varied: }f_{12})\\
\texttt{eLa[[14]]} &= \phantom{-}\tfrac{2}{H}\bigl(2Mf_{9}-\partial_4f_{4}+\partial_0f_{1}\bigr) &&(\text{varied: }f_{13})\\
\texttt{eLa[[15]]} &= \phantom{-}\tfrac{2}{H}\bigl(2Mf_{10}-\partial_4f_{7}+\partial_0f_{2}\bigr) &&(\text{varied: }f_{14})\\
\texttt{eLa[[16]]} &= \phantom{-}\tfrac{2}{H}\bigl(2Mf_{11}+\partial_4f_{6}+\partial_0f_{3}\bigr) &&(\text{varied: }f_{15})
\end{align*}
In Wolfram InputForm the first reads, character for character (broken over two lines here),
\begin{footnotesize}
\begin{verbatim}
eLa[[1]] = (-2*(2*M*f16[4][x0, x4] + Derivative[0, 1][f16[9]][x0, x4]
             + Derivative[1, 0][f16[12]][x0, x4]))/H
\end{verbatim}
\end{footnotesize}
and the other fifteen follow the same pattern with the indices and signs above.

These were printed for this document by loading the notebook's own \texttt{DumpSave} file
\texttt{claude-fable/claude-fable\_Einstein-Rosen-2-Planes-eLa.mx}. That file is \texttt{SameQ} to
the author's original
\texttt{Pre-gravityPre-Big\_Bang\_M6=3-\allowbreak Generations\_of\_\allowbreak Einstein-Rosen-2-Planes-eLa.mx} in the
repository root; this was re-checked for this document with \texttt{SameQ}, which returned
\texttt{True}. So they are the author's equations, unchanged. In matrix form they read
\begin{gather}
  \texttt{eLa} = \frac2H\,\sigma_{16}\bigl(T_{16}[0]\,\partial_{x_0}+T_{16}[4]\,\partial_{x_4}+2M\bigr)\Psi_{16},
  \qquad\text{i.e.}\\
  \bigl(T_{16}[0]\,\partial_{x_0}+T_{16}[4]\,\partial_{x_4}\bigr)\Psi_{16} = -2M\,\Psi_{16}.
\end{gather}
This matrix form was checked for this document (a \texttt{Simplify} difference of exactly 0); it
is not a notebook assertion. In Part VI's language (below) it is the flat-frame Dirac equation
$\gamma^\mu\partial_\mu\Psi = H\,V'(s)\,\Psi$ with $\gamma^\mu=T_{16}[\mu]$ and $V'(s)=-2M/H$. What
Part VI asserts is the fidelity anchor below: the Euler--Lagrange expressions of its fable
Lagrangian with $V=-(2M/H)s$, on the flat frame and with no volume factor, are exactly
\texttt{eLa}.

\paragraph{The coupling pattern.} Iterating \texttt{MergeSetsStep} to a fixed point merges any
two sets of component indices that share an index. The result partitions the sixteen components
into [D]
\begin{equation}
  \texttt{eLaCouplings} = \{\{0,5,8,13\},\{1,4,9,12\},\{2,7,10,15\},\{3,6,11,14\}\}.
\end{equation}
[A] \cfLabel{the coupled blocks partition all sixteen components};
\cfLabel{there are four blocks of four}. The notebook writes \texttt{eLa} and \texttt{eLazt} to
\texttt{.mx} files beside itself, as the original does.

\subsection{The \texorpdfstring{$(z,t)$}{(z,t)} chart, four blocks of four, and the closed-form solutions (Section 13)}
\label{fnd:zt}\label{sec:2.9}

\paragraph{The chart.} The chart is $z=6Hx_0$, $t=Hx_4$. Substituting
$\texttt{f16[k]}\to\texttt{Z[k][6 H x0, H x4]}$ and then $x_0\to z/(6H)$, $x_4\to t/H$, with
the original's normalization $1/(2H)$, gives $\texttt{eLazt}=\tfrac{1}{2H}\,\texttt{eLa}$ in the
new chart. The first equation, for example, reads [D]
\begin{equation}
  \texttt{eLazt[[1]]} = -\frac{2M\,Z_4 + H\,\partial_t Z_9 + 6H\,\partial_z Z_{12}}{H^2},
  \qquad Z_k\equiv\texttt{Z[k][z,t]}.
\end{equation}

\paragraph{Decoupling.} Solving $0=\texttt{eLazt}$ for the sixteen $t$-derivatives gives [A]
\cfLabel{all sixteen t-derivatives are determined}. Relabelling each coupling set as a
consecutive block, \texttt{sZtOyZ} maps $Z_k$ to $yZ_j$ as follows [D]:
\begin{itemize}
  \item $Z_0,Z_5,Z_8,Z_{13}\to yZ_0..yZ_3$;
  \item $Z_1,Z_4,Z_9,Z_{12}\to yZ_4..yZ_7$;
  \item $Z_2,Z_7,Z_{10},Z_{15}\to yZ_8..yZ_{11}$;
  \item $Z_3,Z_6,Z_{11},Z_{14}\to yZ_{12}..yZ_{15}$.
\end{itemize}
Assertions [A]:
\begin{itemize}
  \item \cfLabel{the relabelling is a bijection of the sixteen components}
  \item \cfLabel{the system splits into four blocks of four equations}
  \item \cfLabel{each 4x4 block is closed in its own four components}
\end{itemize}
The relabelling is a linear map. [A] \cfLabel{the yZ coefficient matrix is the identity};
\cfLabel{ORTHOGONAL: Transpose[S] . S == ID16}; \cfLabel{ORTHOGONAL: S . Transpose[S] == ID16};
\cfLabel{but NOT a O(8,8) transformation: it does not preserve sigma16};
\cfLabel{and NOT a direct sum: it mixes type-1 with type-2}.

\paragraph{Closed forms.} \texttt{DSolve}, wrapped in a 60-second \texttt{TimeConstrained},
returns every block unevaluated. The log prints \texttt{block 1: returned unevaluated} through
\texttt{block 4: returned unevaluated}. The author obtained closed forms with Maple and pasted
them into the original as strings. The notebook reproduces them verbatim, parses them with the
author's \texttt{ConvertMapleToMathematicaV2}, and then substitutes them back into the four
blocks. That substitution, not the parser, is what certifies them. [A]
\cfLabel{all sixteen components have a closed form}; \cfLabel{there are sixteen pure-function rules};
\cfLabel{THE MAPLE CLOSED FORMS SOLVE ALL SIXTEEN EQUATIONS}.

Block $j$ has one constant $C_j$ and four constants $c_{j1}..c_{j4}$. Every component of block $j$
carries the same exponential
\begin{align}
  E_j&=\exp\!\left(\frac{6(z/6+t)H^2C_j^2-6(-z/6+t)M^2}{6C_jH^2}\right)\notag\\
     &=\exp\!\Bigl(C_j\bigl(t+\tfrac z6\bigr)-\frac{M^2}{C_jH^2}\bigl(t-\tfrac z6\bigr)\Bigr).
\end{align}
Block 1, verbatim from the Maple string:
\begin{align*}
yZ_0 &= \phantom{-}E_1\,\frac{H^2(c_{11}c_{12}-c_{13}c_{14})C_1^2 - M^2(c_{11}c_{12}+c_{13}c_{14})}{2HC_1M}, &
yZ_2 &= c_{13}c_{14}\,E_1,\\
yZ_1 &= -E_1\,\frac{H^2(c_{11}c_{12}-c_{13}c_{14})C_1^2 + M^2(c_{11}c_{12}+c_{13}c_{14})}{2HC_1M}, &
yZ_3 &= c_{11}c_{12}\,E_1.
\end{align*}
Blocks 2, 3 and 4 have the same shape with their own constants and signs (Section~13 prints all
four strings).

\subsection{Bilinears, the solution branches, and \texorpdfstring{$M_6=3$}{M6 = 3} generations of Einstein--Rosen 2-planes (Section 14)}
\label{fnd:M6}\label{sec:2.10}

Five bilinears are formed from the solution written in the relabelled basis.
\texttt{psi1sol} and \texttt{psi2sol} are its first and second halves of eight components
\textbf{in the relabelled $yZ$ basis}, that is, coupling blocks 1--2 and coupling blocks 3--4.
They are not the type-1 and type-2 spinors. [A]
\cfLabel{the yZ split at 8 is NOT the type-1/type-2 split};
\cfLabel{the first relabelled block holds four type-1 and four type-2 components}.

For the full sixteen-component norm, [A]:
\begin{itemize}
  \item \cfLabel{psi . sigma16 . psi == psi2.sigma.psi2 - psi1.sigma.psi1}
  \item \cfLabel{sigma16 alone is SYMMETRIC, so psi . sigma16 . psi is not forced to vanish}
  \item \cfLabel{and for this solution it does NOT vanish: it is not the zero expression}
\end{itemize}

With the light-cone coordinates $\texttt{ztadv}=6t+z$ and $\texttt{ztret}=6t-z$, the notebook
requires the two block norms to be light-cone constants. [A]
\cfLabel{the light-cone-constant condition has solutions}. \texttt{Solve} returns four branches
[D], among them $\{C_2\to-C_1,\ C_4\to-C_3\}$, the original's branch. [A]
\cfLabel{on the branch C2 = -C1, C4 = -C3 the two norms are independent of ztadv and ztret}.

The displayed values on that branch [D, re-read for this document from the saved kernel state]
are sharper than the assertion. There, $\psi\cdot\sigma_{16}\cdot\psi$,
$\psi_1\cdot\sigma\cdot\psi_1$ and $\psi_2\cdot\sigma\cdot\psi_2$ are all identically $0$, while
$\psi_2\cdot\sigma\cdot\psi_1$ is not.

\paragraph{The two mass branches} [D] are, both taken on the branch
$C_2\to-C_1$, $C_4\to-C_3$:
\begin{itemize}
  \item \texttt{Psi16ab}: $M\to H\,M_{ab}\sqrt{C_1C_3}$, which keeps a free dimensionless ratio
        $M_{ab}$;
  \item \texttt{Psi16aa}: $M\to H\sqrt{C_1C_3}$, the case $M_{ab}=1$.
\end{itemize}
In the $(x_0,x_4)$ chart the advanced and retarded coordinates are $\xi_{\rm adv}=x_0-x_4$ and
$\xi_{\rm ret}=x_0+x_4$. Every component of \texttt{Psi16aa} is a constant times one of
$\exp(\pm H(C_1\xi_{\rm ret}+C_3\xi_{\rm adv}))$ or $\exp(\pm H(C_3\xi_{\rm ret}+C_1\xi_{\rm adv}))$
[D, re-read for this document from the saved kernel state; the notebook prints \texttt{Short}
forms]. The special case $C_3\to C_1$, \texttt{Psi16x0ONLY}, depends on $x_0$ alone. It is
exactly [D, re-read in full from the saved kernel state], with $e_\pm \equiv e^{\pm 2C_1Hx_0}$:
\begin{align*}
\texttt{Psi16x0ONLY}=\bigl(&{-c_{13}c_{14}e_+},\ c_{23}c_{24}e_-,\ -c_{33}c_{34}e_+,\ c_{43}c_{44}e_-,\\
&\phantom{-}c_{21}c_{22}e_-,\ -c_{11}c_{12}e_+,\ c_{41}c_{42}e_-,\ -c_{31}c_{32}e_+,\\
&\phantom{-}c_{13}c_{14}e_+,\ c_{23}c_{24}e_-,\ c_{33}c_{34}e_+,\ c_{43}c_{44}e_-,\\
&\phantom{-}c_{21}c_{22}e_-,\ c_{11}c_{12}e_+,\ c_{41}c_{42}e_-,\ c_{31}c_{32}e_+\bigr).
\end{align*}

\paragraph{Triality turns each spinor of the solution into a vector.} The genuine type-1 and
type-2 spinors are $\psi_1=\texttt{Psi16ab[[1;;8]]}$ and $\psi_2=\texttt{Psi16ab[[9;;16]]}$. From
them:
\begin{gather}
  \texttt{vectorFormOfType2}=\texttt{triVecToSpin}\cdot\psi_2,\\
  \texttt{vectorFormOfType1}=\eta\cdot(\psi_1^{T}\cdot\sigma\cdot\texttt{triSpinToVec}).
\end{gather}
Assertions [A]:
\begin{itemize}
  \item \cfLabel{both triality vectors have eight components}
  \item \cfLabel{vectorFormOfType2 == psi2 contracted on the SPINOR index of triVecToSpin}
  \item \cfLabel{triVecToSpin and triSpinToVec really are mutually inverse}
\end{itemize}
Each vector is normal to a 7-plane through the origin of the split-octonion algebra, the set of
position vectors \texttt{Zoct} with $\langle\text{vector},\texttt{Zoct}\rangle=0$. [A]
\cfLabel{sevenPlane2 is annihilated by vectorFormOfType2};
\cfLabel{sevenPlane1 is annihilated by vectorFormOfType1}.

Two non-parallel 7-planes in an 8-space meet in a 6-plane, \textbf{M6}. [A]
\cfLabel{M6 is a non-empty solution set}; \cfLabel{M6 has six free parameters}. The intersection
determines \texttt{zoct[7]} and \texttt{zoct[8]} as linear combinations of the free
\texttt{zoct[1]}..\texttt{zoct[6]}, with coefficients that depend on $x_0$, $x_4$ and the
constants $C_1$, $C_3$, $M_{ab}$, $c_{jk}$ [D].

What is proved is that M6 exists and is a 6-parameter set. The reading of M6 as the home of
``the three generations of Einstein--Rosen 2-plane bridges'', which gives the original notebook
its title, is the author's interpretation and is stated in prose only [P].

Finally [A]: \cfLabel{cfResolution[unit]: the closed-form bridge equals the matrix inverse};
\cfLabel{cfResolution[unit]: the two bridges are mutually inverse};
\cfLabel{cfResolution[unit] reproduces triVecToSpin}.

\subsection{The canonical frame field (Section 15), restated}
\label{fnd:frame}\label{sec:2.11}

At every point of the curved $4+4$ spacetime there is a local flat $4+4$ Minkowski coordinate
system. The frame field installs it; in eight dimensions a vielbein is simply an 8-dimensional
vierbein. The frame is stored with the curved index $\mu$ as the row and the flat index $a$ as
the column. The canonical frame is the one the original notebook records as the value of its
symbol \texttt{gtrye}:
\begin{gather}
  e = \texttt{frameCanonical} = \operatorname{diag}\bigl(\tan 6Hx_0,\ q,\ q,\ q,\ 1,\ p,\ p,\ p\bigr),\\
  q = \frac{e^{-a_4(Hx_4)}}{(\sin 6Hx_0)^{1/6}},\qquad p = \frac{e^{+a_4(Hx_4)}}{(\sin 6Hx_0)^{1/6}},\qquad
  g = e\,\eta\,e^{T},\\
  ds^2 = \tan^2(6Hx_0)\,dx_0^2 + q^2(dx_1^2+dx_2^2+dx_3^2) - dx_4^2 - p^2(dx_5^2+dx_6^2+dx_7^2).
\end{gather}
Assertions [A]:
\begin{itemize}
  \item \cfLabel{CANONICAL [structural]: frame . coframe == ID8  (certifies only that the frame is invertible)}
  \item \cfLabel{CANONICAL [definition]: g == frame . eta4488 . Transpose[frame]  (bridge [d]; this IS how g was defined, so it cannot fail)}
  \item \cfLabel{CANONICAL [has content]: g == diag(Tan[6 H x0]^2, q^2, q^2, q^2, -1, -p^2, -p^2, -p^2), the stated line element}
  \item \cfLabel{CANONICAL [has content]: Diagonal[g] == Diagonal[eta4488] * Diagonal[frame]^2 exactly}
  \item \cfLabel{CANONICAL [has content]: g has signature (4,4), GIVEN a4 real and 0 < 6 H x0 < Pi/2}
  \item \cfLabel{Inverse[g] == Transpose[coframe] . eta4488 . coframe}
  \item \cfLabel{g is symmetric}
  \item \cfLabel{the metric is non-degenerate}
  \item \cfLabel{FABLE-5.1: det g == Sec[6 H x0]^2  (positive, as signature 4+4 requires)}
\end{itemize}
So $\sqrt{|\det g|}=\sec 6Hx_0$ on $0<6Hx_0<\pi/2$. This is the notebook's volume factor
\texttt{sqrtDetgTele = Sec[6 H x0]}, defined in Section~21 and used throughout Part VI
(Section~15 itself displays \texttt{Abs[Sec[6 H x0]]}).

Two facts follow from the definitions of $q$ and $p$. First, $qp=(\sin 6Hx_0)^{-1/3}$: the two
exponentials cancel [P, Section~15]. Hence $q^3p^3=1/\sin 6Hx_0$, and the 8-volume element
factorizes. [A]:
\begin{itemize}
  \item \cfLabel{FABLESCALAR [has content]: Sqrt[det g_8] == q^3 * (Tan[6 H x0] p^3): observed 3-volume density times hidden 4-volume density}
  \item \cfLabel{FABLESCALAR [has content]: the 8-volume element is independent of x4: d(q^3 p^3)/dx4 == 0}
  \item \cfLabel{FABLESCALAR [has content]: the hidden 4-volume density scales as Exp[+3 a4[H x4]] and the observed 3-volume as Exp[-3 a4[H x4]]: Vhid Exp[-3 a4] and q^3 Exp[+3 a4] are x4-independent}
\end{itemize}
One Text cell of Section~23 writes $q^3p^3=1/\sqrt{\sin 6Hx_0}$. That is a typo in the prose:
from $qp=(\sin 6Hx_0)^{-1/3}$ the cube is $1/\sin 6Hx_0$, which is what makes
$\sqrt{\det g}=q^3\tan(6Hx_0)\,p^3=\sec 6Hx_0$. The assertions above are unaffected.

\paragraph{The curved Dirac matrices} are
$\gamma^\mu=\sum_a \texttt{coframe[[a,}\mu\texttt{]]}\,T_{16}[a]$, with $\texttt{coframe}=e^{-1}$.
On the canonical frame they are as follows (read for this document from the saved kernel
state; they follow directly from the diagonal coframe):
\begin{align}
  \gamma^0&=\cot(6Hx_0)\,T_{16}[0], &
  \gamma^i&=\tfrac1q\,T_{16}[i]=e^{a_4}(\sin 6Hx_0)^{1/6}\,T_{16}[i],\notag\\
  \gamma^4&=T_{16}[4], &
  \gamma^h&=\tfrac1p\,T_{16}[h]=e^{-a_4}(\sin 6Hx_0)^{1/6}\,T_{16}[h],
\end{align}
for $i=1,2,3$ and $h=5,6,7$. [A]
\cfLabel{{gammaCurved[mu], gammaCurved[nu]} == 2 gInv[mu,nu] ID16}; and, from Part VI,
\cfLabel{FABLE [has content]: (gamma^4)^2 == -ID16, so gamma^4 d_4 Psi' == H V'(s) Psi' is solved by d_4 Psi' == -H V'(s) gamma^4 Psi'}.

\paragraph{The free function $a_4$.} $a_4$ is a free, differentiable scalar function of $Hx_4$.
The original notebook never defines it; it occurs only inside the recorded frame. It is carried
symbolically everywhere and is never given a value, and every result holds for arbitrary $a_4$.
The notebook says so on screen: the log prints
\texttt{NOTE~~a4 is an undefined scalar function inherited from the original notebook} together
with the explanation, which also names \texttt{la}. Near its end Part VI checks [A]
\cfLabel{PART VI [control]: a4 is STILL UNDEFINED -- nothing in Part VI gave it a value; and the FLRW scale factor aF never entered the canonical frame}.

\paragraph{The physical reading.} The observer is the geodesic observer $u=\partial/\partial x_4$:
$g_{44}=-1$, and $x_0$ is constant along the worldline [P, Section~23]. The cosmological time
is $t=x_4$, the proper time of that observer [P, Section~25]. The observed sheet is $x_1,x_2,x_3$; the hidden spacelike coordinate is $x_0$;
the hidden timelike sheet is $x_5,x_6,x_7$. Part VI makes the kinematic identification
$a(t)\sim q\sim e^{-a_4(Ht)}$, $H_{\rm obs}=-H\,a_4'(Ht)$. Assertions [A]:
\begin{itemize}
  \item \cfLabel{KINEMATIC [definition]: ln q == -a4[H x4] + const(x0):  D[Log[q], x4] == -H a4'[H x4] == H_obs}
  \item \cfLabel{KINEMATIC [has content]: the second sheet contracts at exactly the opposite rate:  D[Log[p], x4] == +H a4'[H x4] == -H_obs}
  \item \cfLabel{KINEMATIC [has content]: H_obs > 0 exactly when a4' < 0 (and H_obs < 0 when a4' > 0) -- the sign HYPOTHESIS of this Part, stated, not derived}
\end{itemize}
So which sheet grows is set by the sign of $a_4'$, which is never fixed. Section~15's prose
describes the 3-space contracting while the deflating directions expand. Part VI reads the same
frame the other way round (the observed sheet expands when $a_4'<0$) and states that sign as a
hypothesis [P].

\paragraph{What the canonical spin connection contributes.} The spin connection of this frame
follows from the zero-torsion vielbein postulate. Only the facts that the classical fable uses
are restated here:
\begin{equation}
  \Gamma^{\rm spin}_\mu=\tfrac18\,\omega_{\mu ab}\,[T_{16}[a],T_{16}[b]],\qquad
  D_\mu\Psi=\partial_\mu\Psi+\Gamma^{\rm spin}_\mu\Psi ,
\end{equation}
where each $\Gamma^{\rm spin}_\mu$ is block diagonal: it never mixes type-1 with type-2.
Its 24 non-zero components [D, Section~16] are listed below, with $c=\cos 6Hx_0$,
$s=\sin 6Hx_0$, $A=a_4(Hx_4)$, $A'=a_4'(Hx_4)$, $j=1,2,3$ and $k=5,6,7$:
\begin{center}
\begin{tabular}{@{}ll@{}}
\toprule
component & value\\
\midrule
$\omega_{j;0j}=-\omega_{j;j0}$ & $H c^2/(e^{A}s^{13/6})$\\
$\omega_{j;4j}=-\omega_{j;j4}$ & $H A'/(e^{A}s^{1/6})$\\
$\omega_{k;0k}=-\omega_{k;k0}$ & $-e^{A}H c^2/s^{13/6}$\\
$\omega_{k;4k}=-\omega_{k;k4}$ & $e^{A}H A'/s^{1/6}$\\
\bottomrule
\end{tabular}
\end{center}
Assertions [A]:
\begin{itemize}
  \item \cfLabel{CANONICAL [has content]: omega[mu,a,b] == -omega[mu,b,a]  (metric compatibility)}
  \item \cfLabel{INDEPENDENT CHECK: the frame-only derivation reproduces omegaCanonical exactly}
  \item \cfLabel{[has content] the spinor connection is compatible with omega: [Gamma^spin_mu, gamma^a] + omega_mu^a_b gamma^b == 0}
  \item \cfLabel{Gamma^spin[mu] is block diagonal: it never mixes type-1 with type-2}
  \item \cfLabel{[has content] and therefore the curved Dirac matrices are covariantly constant: D_mu gammaCurved^nu == 0}
\end{itemize}
From Part V (the fable-5.1 bridge), three results are used below [A]:
\begin{itemize}
  \item \cfLabel{FABLE-5.1 [THE RESULT]: gamma^mu Gamma^spin[mu] == -(1/2) T[mu] gamma^mu, one vector term},
        with the Weitzenb\"ock torsion vector $T_\mu=\partial_\mu\ln\sin 6Hx_0$, i.e.
        \cfLabel{FABLE-5.1 torsion vector has exactly one non-zero component, 6 H Cot[6 H x0] along x0}.
        Hence $\gamma^\mu\Gamma^{\rm spin}_\mu=-3H\cot^2(6Hx_0)\,T_{16}[0]$.
  \item \cfLabel{FABLE-5.1 [THE RESULT]: DiracCanonical[ Sqrt[Sin[6Hx0]] Psi' ] == Sqrt[Sin[6Hx0]] DiracFable51[Psi'], for a generic Psi'}
  \item \cfLabel{FABLE-5.1 [THE RESULT]: Psi^T sigma16 gamma^mu Gamma^spin[mu] Psi == 0 for EVERY Psi -- the spin connection drops out of the Lagrangian},
        because \cfLabel{FABLE-5.1: the reason is Section 11's identity, Psi^T sigma16 T16[a] Psi == 0 (sigma16 T16[a] antisymmetric)}.
        Note that this reason is a statement about \textbf{commuting} components.
\end{itemize}
The Ricci scalar of the canonical metric is [D, Section~16]
\begin{equation}
  R=-3H^2\Bigl((55+7\cos 12Hx_0)\cot^2(6Hx_0)\csc^2(6Hx_0)-2\,a_4'(Hx_4)^2\Bigr).
\end{equation}

\subsection{The classical field fable of 2026-09-16 (Part VI, Section 24)}
\label{fnd:fable}\label{sec:2.12}

Part VI (Sections 23--25) places two matter fields on the pre-universe. It asks of each the
cosmologist's question: what are the energy density $\rho$, the pressure $P$, and $w=P/\rho$? The
first field, \texttt{fableScalar}, is a real scalar with potential $V(\phi)$. The second,
\textbf{fable}, is the subject of this document.

\paragraph{fableScalar, in brief.} It has
$\hat L_\phi=-\tfrac12 g^{\mu\nu}\partial_\mu\phi\,\partial_\nu\phi-V(\phi)$ and the Hilbert tensor
$T_{\mu\nu}=\partial_\mu\phi\,\partial_\nu\phi+g_{\mu\nu}\hat L_\phi$. Its central theorem is [A]
\cfLabel{FABLESCALAR [THE RESULT]: the null energy condition rho + P_1 == 2 KE holds IDENTICALLY, for every phi(x0, x4, x5) and every V},
so it cannot cross $w=-1$ with positive energy density.

\subsubsection{The field and its Lagrangian}
\label{sec:2.12.1}
fable is a \textbf{real, commuting} 16-component spinor $\Psi(x)$ that transforms under
Spin(4,4) exactly as $\Psi_{16}$. It uses $T_{16}[a]$, $\sigma_{16}$, the curved $\gamma^\mu$,
the scalar bilinear $s=\Psi^{T}\sigma_{16}\Psi$, and a general potential $V(s)$:
\begin{gather}
  L_\Psi=\sqrt{\det g}\,\hat L_\Psi,\qquad
  \hat L_\Psi=\frac1H\,\Psi^{T}\sigma_{16}\gamma^\mu D_\mu\Psi-V(s),\\
  D_\mu=\partial_\mu+\Gamma^{\rm spin}_\mu,\qquad s=\Psi^{T}\sigma_{16}\Psi .
\end{gather}
In Part VI's generic computations $\Psi=\Psi(x_0,x_4)$ has sixteen arbitrary component
functions. Assertions [A]:
\begin{itemize}
  \item \cfLabel{FABLE [definition]: s = Psi^T sigma16 Psi is a genuine scalar with content: on the constant spinor e1 + e5, s == -2}
  \item \cfLabel{FABLE [has content]: Lhat with the covariant derivative == Lhat with the partial derivative, for the generic Psi(x0, x4) (Part V's theorem, re-asserted here)}
\end{itemize}

\paragraph{The fidelity anchor.} The author's Lagrangian is the special case $V(s)=-(2M/H)s$,
taken on the flat frame ($e\to I_8$) and with no volume factor ($\sqrt g\to1$). With that
potential $-V=+(2M/H)s$, which is exactly $L_a$'s mass term. [A]:
\begin{itemize}
  \item \cfLabel{FABLE [fidelity]: with V = -(2M/H) s, frame -> ID8 and Sqrt[g] -> 1, on the model's Psi16, the Lagrangian helper IS the author's La[] of Section 11}
  \item \cfLabel{FABLE [fidelity]: and its Euler-Lagrange equations, formed with the author's own eL operator, ARE the sixteen equations eLa of Section 12}
  \item \cfLabel{FABLE [solver regression]: the new helper cfEulerLagrange reproduces eLa too}
  \item \cfLabel{FABLE-5.1: the covariant Lagrangian with frame -> ID8 and volume factor -> 1 IS the author's La of Section 11, term for term}
\end{itemize}

\subsubsection{The field equation}
\label{sec:2.12.2}
The equation is obtained by explicit variation of $\sqrt g\,\hat L_\Psi$. The kinetic matrix is
$A^\mu=\sigma_{16}\gamma^\mu$, which is antisymmetric; the mass matrix $\sigma_{16}$ is
symmetric. [A]:
\begin{itemize}
  \item \cfLabel{FABLE [THE RESULT]: the Euler-Lagrange equations by explicit variation == Sqrt[g] [ (2/H)(A^mu d_mu Psi + (1/2) D Psi) - 2 V'(s) sigma16 Psi ]},
        where $\mathcal D=\frac{1}{\sqrt g}\partial_\mu(\sqrt g\,A^\mu)$
  \item \cfLabel{FABLE [has content]: A^mu = sigma16 gamma^mu is antisymmetric and sigma16 is symmetric -- the structure that makes the variation close}
  \item \cfLabel{FABLE [has content]: in general (1/Sqrt[g]) d_mu(Sqrt[g] gamma^mu) == [gamma^mu, Gamma^spin_mu]  (covariant constancy of gamma^mu, Section 17)}
  \item \cfLabel{FABLE [has content]: the HYPOTHESIS that makes it a product -- the anticommutator {gamma^mu, Gamma^spin_mu} == 0 on the canonical frame (no totally antisymmetric part of the connection)}
  \item \cfLabel{FABLE [has content]: hence (1/2)(1/Sqrt[g]) d_mu(Sqrt[g] gamma^mu) == gamma^mu Gamma^spin_mu, and on this geometry == -3 H Cot[6 H x0]^2 T16[0]}
  \item \cfLabel{FABLE [THE RESULT]: EL == Sqrt[g] (2/H) sigma16 . ( gamma^mu D_mu Psi - H V'(s) Psi ): the field equation of fable is EXACTLY the Levi-Civita covariant Dirac equation  gamma^mu D_mu Psi == H V'(s) Psi}
  \item \cfLabel{FABLE [has content]: the connection term does not vanish on an x0-independent spinor -- witnessed on the constant spinor e1: a spinor homogeneous in everything but x4 is NOT a solution}
\end{itemize}
So the field equation of the classical fable is
\begin{gather}
  \gamma^\mu D_\mu\Psi = H\,V'(s)\,\Psi,\qquad\text{i.e. on the canonical frame}\\
  \cot(6Hx_0)\,T_{16}[0]\,\partial_0\Psi+\tfrac1q\,T_{16}[i]\,\partial_i\Psi+T_{16}[4]\,\partial_4\Psi
  +\tfrac1p\,T_{16}[h]\,\partial_h\Psi\notag\\
  -3H\cot^2(6Hx_0)\,T_{16}[0]\,\Psi = H\,V'(s)\,\Psi
\end{gather}
(summed over $i=1,2,3$ and $h=5,6,7$; for $\Psi(x_0,x_4)$ only the $\partial_0$ and $\partial_4$
terms survive).

\paragraph{The rescaling.} [A]:
\begin{itemize}
  \item \cfLabel{FABLE [definition]: under Psi = Sqrt[Sin[6 H x0]] Psi', s == Sin[6 H x0] Psi'^T sigma16 Psi'}
  \item \cfLabel{FABLE [THE RESULT]: (gamma^mu D_mu - H V'(s)) [Sqrt[Sin[6 H x0]] Psi'] == Sqrt[Sin[6 H x0]] (gamma^mu d_mu Psi' - H V'(s) Psi') for a GENERIC Psi'(x0, x4): the rescaling removes the connection term exactly}
  \item \cfLabel{FABLE [has content]: for Psi' = Psi'(x4) the field equation reduces to Sqrt[Sin[6 H x0]] (gamma^4 d_4 Psi' - H V'(s) Psi'), s = Sin[6 H x0] Psi'^T sigma16 Psi'}
  \item \cfLabel{FABLE [has content]: the reduced equation still depends on x0 through V'(Sin[6 H x0] s'): for a NONLINEAR V a Psi'(x4) alone is not a solution -- witnessed with V = s^2 on Psi' = e1 + e5}
  \item \cfLabel{FABLE [has content]: for the mass term V = -(2M/H) s, V' is constant and Psi'(x4) IS an exact solution}
\end{itemize}

\subsubsection{The energy--momentum tensor}
\label{sec:2.12.3}
The tensor is built with the covariant derivative and symmetrized:
\begin{equation}
  T^{\rm cov}_{\mu\nu}=-\frac{1}{2H}\bigl(\Psi^{T}\sigma_{16}\gamma_\mu D_\nu\Psi+\Psi^{T}\sigma_{16}\gamma_\nu D_\mu\Psi\bigr)
  +g_{\mu\nu}\hat L_\Psi,
\end{equation}
with $\gamma_\mu=g_{\mu\nu}\gamma^\nu$.
Assertions [A]:
\begin{itemize}
  \item \cfLabel{FABLE [definition]: T_cov is symmetric by construction}
  \item \cfLabel{FABLE [has content]: the diagonal of T_cov equals the diagonal of the partial-derivative tensor, for the generic Psi(x0, x4): rho and every pressure are the same in both}
  \item \cfLabel{FABLE [has content]: but the two tensors DIFFER off the diagonal -- witnessed on the (x3, x4) component T[[4,5]], a momentum density along the observed sheet, for the constant spinor e1 + e13 and V = s^2}
  \item \cfLabel{FABLE [THE RESULT]: nabla^mu T_{mu nu} == 0 ON SHELL in all eight components, for a generic Psi(x0, x4) -- the covariant tensor is the conserved one}
\end{itemize}
``On shell'' means the following. The field equation is solved for
$\partial_4\Psi=F=-\gamma^4\bigl(HV'(s)\Psi-\gamma^0\partial_0\Psi-\tfrac12\,\mathrm{divG}\,\Psi\bigr)$,
where $\mathrm{divG}=\frac{1}{\sqrt g}\partial_\mu(\sqrt g\,\gamma^\mu)$ (on this geometry
$\tfrac12\,\mathrm{divG}=-3H\cot^2(6Hx_0)\,T_{16}[0]$, by the assertion quoted above), and
\texttt{cfOnShell} replaces every $x_4$-derivative of the field by the corresponding
derivative of $F$, repeatedly, until none is left. [A]
\cfLabel{FABLE [definition]: F solves the field equation for d_4 Psi: gamma^4 F + gamma^0 d_0 Psi + (1/2) divG Psi - H V'(s) Psi == 0};
\cfLabel{FABLE [definition]: cfOnShell leaves no x4-derivative of the field in rho}. The
partial-derivative tensor was checked symbolically while Part VI was designed: its $x_0$ and
$x_4$ components are conserved and the other six are not. That check takes minutes and is not
repeated in the notebook [P].

\subsubsection{\texorpdfstring{$\rho$, $P$, and $w$}{rho, P, and w} on shell}
\label{sec:2.12.4}
The observer is $u=\partial/\partial x_4$, so $\rho=T_{44}$ (\texttt{T[[5,5]]}) and
$P_i=T^i{}_i$ (no sum, $i=1,2,3$). Two kinetic bilinears organize the answer:
\begin{align}
  K_\Psi&=\frac1H\,\Psi^{T}\sigma_{16}\gamma^4\partial_4\Psi && (\text{along the observer's time}),\\
  K_h&=-\frac1H\,\Psi^{T}\sigma_{16}\gamma^0\partial_0\Psi && (\text{along the hidden coordinate } x_0).
\end{align}
Assertions [A]:
\begin{itemize}
  \item \cfLabel{FABLE [THE RESULT]: on shell, Lhat == s V'(s) - V(s)}
  \item \cfLabel{FABLE [THE RESULT]: on shell, rho == V(s) + K_h,   K_h = -(1/H) Psi^T sigma16 gamma^0 d_0 Psi}
  \item \cfLabel{FABLE [THE RESULT]: on shell, P_1 == s V'(s) - V(s),  and P_2 == P_3 == P_1}
  \item \cfLabel{FABLE [has content]: on shell, K_Psi == s V'(s) + K_h -- the x4 kinetic bilinear is fixed by the potential and the hidden one}
  \item \cfLabel{FABLE [has content]: K_h == 0 for Psi = Sqrt[Sin[6 H x0]] Psi'(x4)  (d_0 Psi is proportional to Psi, and Psi^T sigma16 gamma^0 Psi == 0)}
  \item \cfLabel{FABLE [THE RESULT]: so for those solutions rho == V(s), hence w == s V'(s)/V(s) - 1, and rho > 0 requires V(s) > 0}
\end{itemize}
In summary:
\begin{gather}
  \rho=V(s)+K_h,\qquad P=sV'(s)-V(s),\\
  \text{and for }\Psi=\sqrt{\sin 6Hx_0}\,\Psi'(x_4):\qquad \rho=V(s),\qquad w=\frac{sV'(s)}{V(s)}-1 .
\end{gather}

\subsubsection{No dilution on the pre-universe}
\label{sec:2.12.5}
Assertions [A]:
\begin{itemize}
  \item \cfLabel{FABLE [THE RESULT]: on shell, ds/dx4 == 2 Psi^T sigma16 gamma^4 gamma^0 d_0 Psi for a generic Psi(x0, x4)}
  \item \cfLabel{FABLE [has content]: the two other terms of d_4 Psi drop out because sigma16 gamma^4 and sigma16 gamma^4 gamma^0 are ANTISYMMETRIC}
  \item \cfLabel{FABLE [control]: that bilinear is NOT identically zero for an x0-dependent Psi -- witnessed on Psi = e1 + x0 e2: 'ds/dx4 == 0 for generic Psi(x0, x4)' would be FALSE}
  \item \cfLabel{FABLE [THE RESULT]: for Psi = Sqrt[Sin[6 H x0]] Psi'(x4) it vanishes: s is CONSTANT IN x4 -- no dilution on the pre-universe}
  \item \cfLabel{FABLE [THE RESULT]: the same from the reduced equation: ds'/dx4 == 0 on shell for Psi'(x4), because Psi'^T sigma16 gamma^4 Psi' == 0}
\end{itemize}
So on those solutions $s$, $V$, $\rho$ and $w$ are all constant along $x_4$. This is the spinor
analogue of the scalar's ``no Hubble friction'' (for the scalar, [A]
\cfLabel{FABLESCALAR [THE RESULT]: for phi = phi(x4) the field equation is phi'' == -V'(phi): NO HUBBLE FRICTION on the pre-universe}).
For the scalar, Section~23's prose gives the reason as geometric: the factor
$\sqrt g\,g^{44}=-\sec 6Hx_0$ does not depend on $x_4$, because the 8-volume element is
$x_4$-independent [P].

\subsubsection{Dust for the author's mass term}
\label{sec:2.12.6}
[A]:
\begin{itemize}
  \item \cfLabel{FABLE [THE RESULT]: the author's mass term V = -(2M/H) s is DUST: on shell P_1 == 0 identically and rho == -(2M/H) s  (w == 0, exactly)}
  \item \cfLabel{FABLE [has content]: that density is not zero -- witnessed on Psi' = e1 + e5 at M = 1 (rho = (2M/H) * 2 Sin[6 H x0] there)}
\end{itemize}
The density $\rho=-(2M/H)s$ is positive only on the branch $Ms<0$ [P, Section~24].

\subsubsection{\texorpdfstring{$w=sV'/V-1$}{w = sV'/V - 1} for the six potentials, and the classical phantom crossing}
\label{sec:2.12.7}
With $u\equiv(s/s_1)^2$:
\begin{align*}
&\text{mass:} & V&=ms, & w&=0,\\
&\text{lambda-mass:} & V&=V_0+ms, & w&=-\frac{V_0}{V_0+ms},\\
&\text{power:} & V&=ms+\lambda s^{n}, & w&=\frac{ms+n\lambda s^{n}}{ms+\lambda s^{n}}-1,\\
&\text{hilltop:} & V&=V_0-\mu(s-s_*)^2, & w&=\frac{-2\mu s(s-s_*)}{V_0-\mu(s-s_*)^2}-1,\\
&\text{lorentz:} & V&=V_0+\frac{ms}{1+u}, & w&=\frac{ms(1-u)}{(1+u)^2\bigl(V_0+ms/(1+u)\bigr)}-1,\\
&\text{expdamp:} & V&=V_0+ms\,e^{-s/s_1}, & w&=\frac{ms\,e^{-s/s_1}(1-s/s_1)}{V_0+ms\,e^{-s/s_1}}-1 .
\end{align*}
Assertions [A]:
\begin{itemize}
  \item \cfLabel{FABLE [has content]: the closed forms of w_Psi(s) = s V'(s)/V(s) - 1 for the six potentials of the numerical work (mass, lambda-mass, power, hilltop, lorentz, expdamp)}
  \item \cfLabel{FABLE [THE RESULT]: mass -> w == 0 (dust);  V = lam s^n alone -> w == n - 1 (accelerating for n < 2/3, phantom for n < 0);  n = 0.236 -> w == -0.764 exactly}
  \item \cfLabel{FABLE [THE RESULT]: w + 1 == s V'/V vanishes where V' does, with V > 0 there: lorentz and expdamp cross w = -1 at s = s1, hilltop at s = s* -- the phantom divide is crossed with no wrong-sign kinetic term}
  \item \cfLabel{FABLE [has content]: lorentz and expdamp are bounded below by V0 > 0 for s >= 0 (rho > 0 on the whole history), and w < -1 for s > s1 (the past), w > -1 for 0 < s < s1 (the present)}
  \item \cfLabel{FABLE [has content]: the limits -- lorentz and expdamp: w -> -1 both as s -> infinity and as s -> 0;  lambda-mass: w -> 0 (dust) as s -> infinity and -> -1 (Lambda) as s -> 0}
  \item \cfLabel{FABLE [has content]: expdamp with (V0, m, s1) = (1.566, 0.839, 2.21) at s = 1 (today) gives w == -0.8608, Unite's w0 = -0.861 to three decimals}
\end{itemize}
The assumptions for the sign statements are $V_0>0$, $m>0$, $s_1>0$. This is the
\textbf{classical phantom crossing}: $w+1=sV'/V$ changes sign where $V'$ does, with $V>0$ there,
and no kinetic term has the wrong sign. It is the spinor-quintom mechanism of Cai and Wang,
Class.\ Quantum Grav.\ 25 (2008) 165014. The stability of perturbations is not examined in
Part VI [P]. The $V_0$ terms of lambda-mass, lorentz and expdamp are a bare cosmological
constant and are reported as such [D, the summary table printed in Section~25].

\subsubsection{The FLRW reference model (Model B)}
\label{sec:2.12.8}
Because $ds/dx_4=0$ on the pre-universe, a history $w(a)$ needs a different background. Part VI
builds it as a \textbf{separate} frame,
$\texttt{frameFLRW}=\operatorname{diag}(1,a_F(x_4),a_F(x_4),a_F(x_4),1,1,1,1)$, with its own
scale factor $a_F$, and feeds it to the same spin-connection solver. The canonical frame and
$a_4$ are not touched. Assertions [A]:
\begin{itemize}
  \item \cfLabel{FLRW [definition]: the reference frame does not mention a4, and the canonical frame is untouched}
  \item \cfLabel{FLRW [solver regression]: vielbein postulate residual is zero}
  \item \cfLabel{FLRW [has content]: omega is antisymmetric (metric compatible)}
  \item \cfLabel{FLRW [THE RESULT]: gamma^mu Gamma^spin_mu == (3/2) (a'/a) gamma^4 -- the Hubble term of the FLRW Dirac equation}
  \item \cfLabel{FLRW [definition]: the rule solves gamma^mu D_mu Psi == H V'(s) Psi for d_4 Psi}
  \item \cfLabel{FLRW [THE RESULT]: d(a^3 s)/dx4 == 0 on shell, for EVERY V: s = s0 a^-3, the dust dilution law, and w_Psi(a) = s V'(s)/V(s) - 1 on it}
  \item \cfLabel{FLRW [control]: without the Hubble term, d(a^3 s)/dx4 == 3 (a'/a) a^3 s instead (i.e. ds/dx4 == 0, the pre-universe law): the dilution IS the Hubble term}
\end{itemize}

\paragraph{Model B, solved independently with NDSolve.} The units are $8\pi G=1$, $H_0=1$ and
$\rho_{\rm crit}=3$. The independent variable is $N=\ln a$, and $s=e^{-3N}$ is exact, so only the
cosmic time needs integrating:
\begin{gather}
  \frac{dt}{dN}=\frac1H,\qquad H^2=\Omega_m e^{-3N}+\Omega_r e^{-4N}+\frac{V(s)}{3},\qquad
  \Omega_m=\tfrac{3}{10},\quad \Omega_r=\tfrac{84}{10^6},\\
  V=\text{expdamp with }(V_0,m,s_1)=\bigl(\tfrac{1566}{1000},\tfrac{839}{1000},\tfrac{221}{100}\bigr),\\
  N\in[-7,0],\quad 701\text{ rows},\quad t(N_0)=\tfrac{1}{2H(N_0)}\text{ added}.
\end{gather}
All inputs are exact rationals, and \texttt{WorkingPrecision -> 24}. Assertions [A]:
\begin{itemize}
  \item \cfLabel{MODEL B [definition]: 701 rows from N = -7 to 0, s(0) == 1 and rho_Psi(0) == V(1) == 1.566 + 0.839 Exp[-1/2.21]}
  \item \cfLabel{MODEL B [THE RESULT]: w <= -1 at every grid point with s > s1 and w > -1 at every grid point with s < s1: the crossing is at s = s1, a = s1^(-1/3) = 0.7677}
  \item \cfLabel{MODEL B [has content]: a genuine phantom phase: min w = -1.302 (to 0.001) at a = 0.543, and w(a = 1) = -0.8608}
  \item \cfLabel{MODEL B [has content]: V(s) > 0 at every row -- rho_Psi > 0 on the whole history}
  \item \cfLabel{MODEL B [fidelity]: the reference CSV has the expected header and the same 701 values of N}
  \item % one label, split only so that the line can break after a slash
        \cfLabel{MODEL B [fidelity]: w agrees with fable-cosmology/}\allowbreak\cfLabel{reference/}\allowbreak\cfLabel{mathematica_spinor_expdamp.csv at EVERY row to 1e-9}
  \item \cfLabel{MODEL B [fidelity]: so do t, H and rho_Psi, to 1e-9 (s differs only by the file's 16-digit rounding of numbers of order 1e9)}
\end{itemize}
The run prints [D] (one line in the log, broken here):
\begin{footnotesize}
\begin{verbatim}
Model B (expdamp):
  w(a=1) = -0.86084564037764806560795392619028548152`20.
  min w  = -1.30209881490918676547015655743122703056`20.  at N = -0.61
           (a = 0.54335086907449978711266408781657974806`20.)
  t(a=1) = 0.96433642136311524576147498023448633382`20.
  H(a=1) = 0.99998205003431888956547950340514648567`20.
\end{verbatim}
\end{footnotesize}
Model B applies no normalization: $V(1)=1.566+0.839\,e^{-1/2.21}\approx2.0996=3\,\Omega_{\Psi0}$
with $\Omega_{\Psi0}=0.6999$ [P, Section~24]. That is why the printed $H(a{=}1)$ is 0.99998 and
not exactly 1: $H(a{=}1)^2=3/10+84/10^6+V(1)/3\approx0.99996$ (this arithmetic was added for this
document; it is not a notebook statement).

\subsubsection{What is proved on the pre-universe, and what belongs to the reference model}
\label{sec:2.12.9}
Proved on the 8-manifold, for arbitrary $a_4$, with every identity closed symbolically:
\begin{itemize}
  \item the Lagrangian and its fidelity to $L_a$ and \texttt{eLa};
  \item the covariant Dirac equation and its rescaling;
  \item the covariant energy--momentum tensor and its conservation;
  \item $\rho=V+K_h$ and $P=sV'-V$;
  \item dust for the author's mass term;
  \item the constancy of $s$ in $x_4$;
  \item the volume factorization.
\end{itemize}
The history $s=s_0a^{-3}$, the phantom crossing at $a=0.768$, and all contact with the Unite
fits $(w_0,w_a)=(-0.861,-0.60)$ belong to 4-dimensional physics on the separate FLRW frame
[P, Section~25].

The volume bookkeeping is proved [A]:
\begin{itemize}
  \item \cfLabel{VOLUME [has content]: the reduced density rho_4 = Vhid rho_8 is NOT separately conserved:  d rho_4/dx4 + 3 H_obs (rho_4 + P_4) == 3 H_obs P_4, for every w}
  \item \cfLabel{VOLUME [has content]: equivalently d rho_4/dx4 == -3 H_obs rho_4: rho_4 falls as a^-3 for EVERY w -- a volume effect, the same for a cosmological constant as for dust}
\end{itemize}
Which of $\rho_4$ and $\rho_8$ gravitates is left undetermined in Part VI, because Part VI has no
gravitational sector: no Einstein equations, and $a_4$ free [P, Section~25].

\subsection{Why the classical fable must be refined}
\label{fnd:refine}\label{sec:2.13}

The fable of Part VI is a field of \textbf{real, commuting} numbers. A fermion needs
\textbf{anticommuting} (Grassmann) components, and the change is not cosmetic. It rests on two
algebraic facts about the author's Clifford algebra.

\paragraph{F1, the symmetry of $C\Gamma$.} Write $\Gamma_{(r)}$ for an antisymmetrized product of
$r$ of the $T_{16}[a]$. With $C_-=\sigma_{16}$:
\begin{itemize}
  \item $\sigma_{16}$ is symmetric;
  \item $\sigma_{16}T_{16}[a]$ is \textbf{antisymmetric} (this much is already Part I's
        \cfLabel{sigma16 . T16[A] is antisymmetric for A = 0..7});
  \item $\sigma_{16}\Gamma_{ab}$ is antisymmetric;
  \item $\sigma_{16}\Gamma_{abc}$ and $\sigma_{16}\Gamma_{abcd}$ are symmetric;
  \item $\sigma_{16}T_{16}[8]$ is symmetric;
  \item $\sigma_{16}T_{16}[a]\sigma_{16}^{-1}=-T_{16}[a]^{T}$.
\end{itemize}
The second conjugation is $C_+=\sigma_{16}T_{16}[8]$:
\begin{itemize}
  \item it is symmetric, and $C_+T_{16}[a]$ is \textbf{symmetric};
  \item $C_+\Gamma_{ab}$ and $C_+\Gamma_{abc}$ are antisymmetric;
  \item $C_+T_{16}[a]C_+^{-1}=+T_{16}[a]^{T}$.
\end{itemize}
The Spin(4,4)-invariant bilinear forms on the 16-spinor form a two-dimensional space, spanned by
$\sigma_{16}$ and $\sigma_{16}T_{16}[8]$ ($C_+T_{16}[8]=\sigma_{16}$). Both are symmetric, and
both are chirality-diagonal. [VII] \cfL{REFINEMENT [definition]: T16[8]\textasciicircum{}2 == sigma16\textasciicircum{}2 == C+\textasciicircum{}2 == ID16, C+ is symmetric, and C+ . T16[8] == sigma16 -- the four candidate mass matrices are only two},
\cfL{REFINEMENT [has content]: C- T16[a] C-\textasciicircum{}-1 == -T16[a]\textasciicircum{}T and C+ T16[a] C+\textasciicircum{}-1 == +T16[a]\textasciicircum{}T for a = 0..7}, \cfL{REFINEMENT [THE RESULT]: the symmetry of C Gamma depends only on the rank r of Gamma:  C-: S A A S S A A S S,  C+: S S A A S S A A S  (r = 0..8)},
\cfL{REFINEMENT [THE RESULT]: the Spin(4,4)-invariant bilinear forms on the 16-spinor form a 2-dimensional space, spanned by C- = sigma16 and C+ = sigma16 T16[8]},
\cfL{REFINEMENT [has content]: every invariant form is SYMMETRIC and chirality-diagonal, and the two chiral pieces PL sigma16 PL and PR sigma16 PR are invariant separately}.

\paragraph{F2, what anticommuting components do to a bilinear.} For Grassmann $\theta$:
\begin{itemize}
  \item $\theta^{T}N\theta$ sees only the \textbf{antisymmetric} part of $N$;
  \item $\theta^{T}N\,d\theta=\tfrac12\,d(\theta^{T}N\theta)+\theta^{T}N_{\rm sym}\,d\theta$.
\end{itemize}
For commuting components the roles are exactly reversed: $\theta^{T}N\theta$ sees only the
symmetric part, and the Euler--Lagrange expression of $\theta^{T}N\phi$ is
$2N_{\rm antisym}\phi$. That reversal is why Part VI's real commuting fable has both a mass term
($\sigma_{16}$ symmetric) and a kinetic term ($\sigma_{16}T_{16}[a]$ antisymmetric). [VII]
\cfL{GRASSMANN [THE RESULT]: theta\textasciicircum{}T N theta == theta\textasciicircum{}T N\_A theta and theta\textasciicircum{}T N\_S theta == 0 for a general 16x16 N -- an anticommuting bilinear sees only the ANTISYMMETRIC part}, \cfL{GRASSMANN [THE RESULT]: theta\textasciicircum{}T N phi - (1/2) d(theta\textasciicircum{}T N theta) == theta\textasciicircum{}T N\_S phi for a general N},
\cfL{GRASSMANN [THE RESULT]: the Euler-Lagrange expression of L = theta\textasciicircum{}T N phi (left derivatives) is (N + N\textasciicircum{}T) phi: EMPTY field equations exactly when N is antisymmetric}, \cfL{GRASSMANN [control]: for COMMUTING components the roles are reversed -- theta\textasciicircum{}T N theta sees only N\_S, and the Euler-Lagrange expression of theta\textasciicircum{}T N phi is 2 N\_A phi}.

The consequences:
\begin{itemize}
  \item[(a)] \textbf{A real Grassmann 16-spinor with the author's $\sigma_{16}$ has no dynamics.}
        Its mass term $\theta^{T}\sigma_{16}\theta$ vanishes identically, because $\sigma_{16}$ is
        symmetric. Its kinetic term $\theta^{T}\sigma_{16}T_{16}[a]\,d\theta$ is a total
        derivative, because $\sigma_{16}T_{16}[a]$ is antisymmetric. So its Euler--Lagrange
        expression is identically empty, even though the kinetic term itself is not zero. The
        statement holds on \textbf{every} frame. $\sigma_{16}\Gamma_{abc}$ is symmetric, so the
        connection term $\theta^{T}\sigma_{16}\{\gamma^\mu,\Gamma^{\rm spin}_\mu\}\theta$ vanishes
        too. [VII] \cfL{REFINEMENT (a) [THE RESULT]: real Grassmann + sigma16: both invariant mass terms vanish, and theta\textasciicircum{}T sigma16 T16[a] phi == (1/2) d(theta\textasciicircum{}T sigma16 T16[a] theta) for every a (a total derivative)},
        \cfL{REFINEMENT (a) [has content]: its Euler-Lagrange expression is identically empty in every direction -- while the kinetic term itself is NOT zero},
        \cfL{REFINEMENT (a) [THE RESULT]: FRAME-INDEPENDENT -- sigma16 \{T16[a], S\textasciicircum{}bc\} is SYMMETRIC for all a, b, c (so theta\textasciicircum{}T sigma16 \{gamma\textasciicircum{}mu, Gamma\_mu\} theta == 0 on every frame), and d\_mu(Sqrt[g] sigma16 gamma\textasciicircum{}mu) == Sqrt[g] sigma16 [gamma\textasciicircum{}mu, Gamma\_mu] on the canonical frame}.
  \item[(b)] \textbf{A Majorana 16-spinor built with $C_+$ propagates but is massless.} Its
        kinetic term is genuine, but every invariant bilinear vanishes on it, so there is no $s$
        and $V(s)$ is undefined. It has exactly one quartic Lorentz scalar: the quartic scalars
        $v\cdot v$, $B\cdot B$ and $C\cdot C$ built from it are all proportional to one quartic.
        Its commuting shadow has no kinetic term at all. [VII]
        \cfL{REFINEMENT (b) [THE RESULT]: Majorana with C+: the kinetic term is genuine -- Euler-Lagrange expression 2 C+ T16[a] phi with C+ T16[a] invertible -- but every invariant bilinear vanishes (massless, no V(s))}, \cfL{REFINEMENT (b) [has content]: the scope -- the quartic Lorentz scalars v.v, B.B, C.C of the Majorana spinor are NOT zero and are all proportional to ONE quartic (same monomials, one constant ratio each)},
        \cfL{REFINEMENT (b) [has content]: and the commuting shadow of the C+ Majorana spinor has NO kinetic term: for commuting components theta\textasciicircum{}T C+ T16[a] phi is the total derivative (1/2) d(theta\textasciicircum{}T C+ T16[a] theta)}.
  \item[(c)] \textbf{Weyl and Majorana--Weyl 8-spinors (one chirality) cannot propagate.} Both
        $C$s are chirality-diagonal and $\gamma^\mu$ is chirality-odd, so
        $P_L\,C\,T_{16}[a]\,P_L=P_R\,C\,T_{16}[a]\,P_R=0$. [VII]
        \cfL{REFINEMENT (c) [THE RESULT]: Weyl and Majorana-Weyl: PL C T16[a] PL == PR C T16[a] PR == 0 for C = C-, C+ and every a (C chirality-diagonal, T16[a] chirality-odd) -- no kinetic term}.
  \item[(d)] \textbf{Fermion fable is the complex 16-component spinor.} It is
        $\Psi=(\theta_1+i\theta_2)/\sqrt2$, built from two real Grassmann 16-spinors, with the
        conjugate $\bar\Psi=\Psi^{\ddagger}\sigma_{16}$. Here ${}^{\ddagger}$ is the conjugation
        of the complex Grassmann algebra; the Hilbert-space adjoint is fixed in the quantization,
        \secref{sec:6}.
        \begin{itemize}
          \item It has two scalar bilinears. $s=\bar\Psi\Psi=i\,\theta_1^{T}\sigma_{16}\theta_2$ is
                non-zero and Hermitian under ${}^{\ddagger}$, and so is
                $p=\bar\Psi\,T_{16}[8]\,\Psi$. (With the Hilbert-space adjoint that the
                quantization fixes, $p$ is odd under the fundamental symmetry $J$ and becomes an
                anti-Hermitian operator, while $s$ stays Hermitian.)
          \item The symmetrized kinetic term is Hermitian and genuine.
          \item It reduces \textbf{exactly} to Part VI's real commuting fable: its Lagrangian,
                its field equation, and its covariant energy--momentum tensor. Through that
                reduction it reduces to the author's $L_a$ and \texttt{eLa}.
        \end{itemize}
        [VII] \cfL{REFINEMENT (d) [THE RESULT]: s = Psibar Psi == i theta1\textasciicircum{}T sigma16 theta2: NON-ZERO (16 monomials) and HERMITIAN; p = Psibar T16[8] Psi is non-zero and Hermitian too},
        \cfL{REFINEMENT (d) [THE RESULT]: its symmetric part is A\_S == (i/2) Omega, Omega = [[0, K], [-K, 0]] real SYMMETRIC and invertible: a genuine kinetic term}, \cfL{REFINEMENT (d) [has content]: the symmetrized kinetic term is HERMITIAN and equals Psi\textasciicircum{}ddag K phi - (1/2) d(Psi\textasciicircum{}ddag K Psi); the unsymmetrized one is not Hermitian}, \cfL{FIDELITY [1] [fidelity]: Psi -> psi, Psibar -> psi\textasciicircum{}T sigma16 turns the symmetrized covariant Lhat into Part VI's cfLhatFable EXACTLY (canonical frame, generic V)},
        \cfL{FIDELITY [2] [fidelity]: on the flat frame with V = -(2M/H) s and no volume factor, the same limit on Psi16 IS the author's La[] of Section 11}, \cfL{FIDELITY [3] [fidelity]: sigma16 . EL\_Psibar + EL\_Psi in the real limit == Part VI's explicit variation cfELFable, component for component},
        \cfL{FIDELITY [3] [fidelity]: and on the flat frame with V = -(2M/H) s the same combination IS the author's eLa of Section 12}, \cfL{FIDELITY [4] [fidelity]: the energy-momentum tensor That in the real limit IS Part VI's covariant tensor cfTCov, all 64 components}.
\end{itemize}

The claim must be stated with its scope. The complex 16-spinor is the \textbf{minimal} field
that:
\begin{enumerate}
  \item uses the author's $\sigma_{16}$ conjugation;
  \item propagates;
  \item admits a potential $V(s)$;
  \item reduces to Part VI's classical fable as its real commuting shadow.
\end{enumerate}
It is \textbf{not} the unique field with dynamics: a Majorana spinor with $C_+$ also propagates.
The author chose the complex 16-spinor on 2026-09-24.

One more consequence matters for everything that follows. Part VI's mechanism for the spin
connection dropping out of the Lagrangian was $\Psi^{T}\sigma_{16}T_{16}[a]\Psi=0$, which is a
property of commuting components. It does not carry over to complex Grassmann fields. There, the
drop of the connection is a property of the matrix $\{\gamma^\mu,\Gamma^{\rm spin}_\mu\}$ itself,
and it has to be proved again. Whether the classical results of this section, among them $w$ and
the phantom crossing, survive quantization is decided in sections~\ref{sec:7} and~\ref{sec:8}. The
complete refinement, with every assertion of Section~26, is \secref{sec:3}.

\endgroup

% ===============================================================================================
\section{The refinement: why fable must be a complex 16-spinor}
\label{sec:3}

This is Section~26 of the notebook (manifest \file{claude-fable/cells_part7.wl}, Input cells
173--179, 38 assertions). It turns the classical, real, commuting fable of \secref{sec:2.12} into a
fermion.

\paragraph{Two conjugations, two symbols.} $\Psi^{\ddagger}$ is the \textbf{classical} conjugate of
the field, the conjugation of the Grassmann algebra, the one that appears in the Lagrangian; the
Dirac conjugate is $\bar\Psi=\Psi^{\ddagger}\sigma_{16}$. After quantization (\secref{sec:6}) the
Hilbert-space adjoint is a \textbf{different} operation, $\Psi^{\dagger}=\Psi^{\ddagger}J$. The two
are never written with the same symbol \cfprose{VII, introduction of Part~VII}.

\paragraph{Two kinds of algebra, and when each is used} \cfprose{VII}. Statements whose content
\emph{is} the anticommutation of the field (that a mass term or a kinetic term vanishes, is
Hermitian, or is a total derivative) are proved in a genuine Grassmann (exterior) algebra, built and
tested in Section~26 (subsection~\ref{sec:3.3}). Statements about \textbf{bilinear} expressions
(Lagrangians, Euler--Lagrange equations, currents, energy--momentum tensors, conservation laws) are
proved with $\Psi$ and $\bar\Psi$ represented by two \textbf{independent} vectors of ordinary
functions of the coordinates, with $\bar\Psi$ always written to the \textbf{left} of $\Psi$. That
commuting proxy is exact for such identities: in every term $\bar\Psi$ is the leftmost and $\Psi$
the rightmost odd factor, so the left derivative with respect to $\bar\Psi$ and the right
derivative with respect to $\Psi$ are the ordinary derivatives, and no reordering of odd factors
ever occurs. Where reality matters (Hermiticity under ${}^{\ddagger}$), $\Psi=u+iv$ with real $u,v$
and $\bar\Psi=(u-iv)^{T}\sigma_{16}$.

\subsection{The Clifford symmetry table}
\label{sec:3.1}

A spinor bilinear $\psi^{T}C\,\Gamma\,\psi$ is built from a ``charge conjugation'' matrix $C$ and
one of the 256 ordered products $\Gamma$ of the Dirac matrices (Section~6's basis \cfcode{base16},
grouped by the rank $r$, the number of factors). There are two candidates for $C$:
\begin{align*}
  C_- &= \sigma_{16}, & C_-\,\Tsix{a}\,C_-^{-1} &= -\Tsix{a}^{T} &&\text{(the author's spinor metric)},\\
  C_+ &= \sigma_{16}\,\Tsix{8}, & C_+\,\Tsix{a}\,C_+^{-1} &= +\Tsix{a}^{T} .
\end{align*}
Every product $C\Gamma$ is either symmetric (S) or antisymmetric (A), and the answer depends only
on the rank:

\begin{center}
\begin{tabular}{@{}lccccccccc@{}}
\toprule
rank $r$ & 0 & 1 & 2 & 3 & 4 & 5 & 6 & 7 & 8 \\
\midrule
number of products, $\binom{8}{r}$ & 1 & 8 & 28 & 56 & 70 & 56 & 28 & 8 & 1 \\
$C_-\Gamma=\sigma_{16}\Gamma$ & S & A & A & S & S & A & A & S & S \\
$C_+\Gamma=\sigma_{16}\Tsix{8}\Gamma$ & S & S & A & A & S & S & A & A & S \\
\bottomrule
\end{tabular}
\end{center}

For each $C$, 136 products are symmetric and 120 antisymmetric, as it must be for $16\times16$
matrices ($16\cdot17/2=136$, $16\cdot15/2=120$). The ``four candidate'' mass matrices
$\sigma_{16}$, $C_+$, $\sigma_{16}\Tsix{8}$, $C_+\Tsix{8}$ are only two, because
$\Tsix{8}^2=\IDs$: $\sigma_{16}\Tsix{8}=C_+$ and $C_+\Tsix{8}=\sigma_{16}$. Why this table decides
everything: for \textbf{commuting} components $\psi^{T}N\psi$ sees only the symmetric part of $N$;
for \textbf{anticommuting} (Grassmann) components only the antisymmetric part.

\cfproved{VII \S26}
\begin{itemize}
  \item \cfL{REFINEMENT [definition]: T16[8]\textasciicircum{}2 == sigma16\textasciicircum{}2 == C+\textasciicircum{}2 == ID16, C+ is symmetric, and C+ . T16[8] == sigma16 -- the four candidate mass matrices are only two}
  \item \cfL{REFINEMENT [has content]: C- T16[a] C-\textasciicircum{}-1 == -T16[a]\textasciicircum{}T and C+ T16[a] C+\textasciicircum{}-1 == +T16[a]\textasciicircum{}T for a = 0..7}
  \item \cfL{REFINEMENT [THE RESULT]: the symmetry of C Gamma depends only on the rank r of Gamma:  C-: S A A S S A A S S,  C+: S S A A S S A A S  (r = 0..8)} (the assertion also checks that the
        counts per rank are $\binom{8}{r}$)
  \item \cfL{REFINEMENT [has content]: for each C, 136 of the 256 products are symmetric and 120 antisymmetric}
\end{itemize}

\subsection{The Lorentz-invariant bilinears: exactly two, both symmetric}
\label{sec:3.2}

A bilinear $\psi^{T}N\psi$ (real field) or $\Psi^{\ddagger}N\Psi$ (complex field) is a Spin(4,4)
scalar exactly when $S^{T}N+NS=0$ for all 28 generators $S=$ \cfcode{SAB} of Section~5 (the
generators are real, so the condition is the same in both cases). Solving these 7168 linear
equations ($28\times256$) for the 256 entries of $N$ gives a \textbf{two}-dimensional space, spanned
by $C_-$ and $C_+$, equivalently by the two chiral pieces $P_L\sigma_{16}P_L$ and
$P_R\sigma_{16}P_R$, one for each split-octonion spinor type. Every invariant $N$ is
\textbf{symmetric} and chirality-\textbf{diagonal} (it commutes with $\Tsix{8}$). That is the
complete list of candidates for a mass term; nothing is left out \cfprose{VII \S26}.

\cfproved{VII \S26}
\begin{itemize}
  \item \cfL{REFINEMENT [THE RESULT]: the Spin(4,4)-invariant bilinear forms on the 16-spinor form a 2-dimensional space, spanned by C- = sigma16 and C+ = sigma16 T16[8]}
  \item \cfL{REFINEMENT [has content]: every invariant form is SYMMETRIC and chirality-diagonal, and the two chiral pieces PL sigma16 PL and PR sigma16 PR are invariant separately}
\end{itemize}

\subsection{A Grassmann algebra, built and tested}
\label{sec:3.3}

To prove statements about anticommuting fields the notebook uses an honest exterior algebra, not a
sign convention applied by hand. An element is an \cfcode{Association} from \textbf{monomials} to
coefficients; a monomial is the sorted list of the generator indices it contains, \cfcode{{}} being
the unit. The product of two monomials is zero if they share a generator; otherwise it is the
sorted union with the sign of the permutation that sorts the concatenation (\cfcode{Signature}).
The operations are \cfprose{VII \S26}:

{\small
\begin{longtable}{@{}p{0.33\linewidth}p{0.61\linewidth}@{}}
\toprule
\raggedright operation & \raggedright what it computes \tabularnewline
\midrule
\endhead
\raggedright \cfcode{grAdd}, \cfcode{grScale}, \cfcode{grMul} & \raggedright sums, scalar multiples and products \tabularnewline
\raggedright \cfcode{grSame[x, y]} & \raggedright equality of two elements ($x-y$ has no monomial left) \tabularnewline
\raggedright \cfcode{grBil[xs, N, ys]} & \raggedright the bilinear $\sum_{ab}x_a N_{ab}\,y_b$ of two lists of elements \tabularnewline
\raggedright \cfcode{grD[x, dmap]} & \raggedright the \textbf{even} derivation that maps generator $i$ to generator \cfcode{dmap[i]} (with $\theta_a\mapsto\phi_a$ it is $d/dt$, with $\phi_a=d\theta_a/dt$) \tabularnewline
\raggedright \cfcode{grDL[x, i]} & \raggedright the \textbf{left} derivative with respect to generator $i$ \tabularnewline
\raggedright \cfcode{grConj[x]} & \raggedright the classical conjugation ${}^{\ddagger}$: $(xy)^{\ddagger}=y^{\ddagger}x^{\ddagger}$, generators self-conjugate (real), coefficients complex-conjugated \tabularnewline
\raggedright \cfcode{grEL[L, thIdx, phIdx, dmap]} & \raggedright the Euler--Lagrange expression $\partial L/\partial\theta_c-d\,(\partial L/\partial\phi_c)$, with left derivatives \tabularnewline
\bottomrule
\end{longtable}
}

Before any physics the algebra is tested on facts known independently (random elements drawn with
\cfcode{SeedRandom[20260924]}). \cfproved{VII \S26}
\begin{itemize}
  \item \cfL{GRASSMANN [solver regression]: generators anticommute and square to zero (all pairs of 1..6)}
  \item \cfL{GRASSMANN [solver regression]: the product is associative, (x y) z == x (y z), on elements of degree 1, 2, 3 (and the products tested are not zero)}
  \item \cfL{GRASSMANN [solver regression]: graded commutativity  x y == (-1)\textasciicircum{}(|x| |y|) y x  (degrees 1, 2, 3), and x\textasciicircum{}2 == 0 for x odd}
  \item \cfL{GRASSMANN [solver regression]: the conjugation is an involutive anti-automorphism,  (x y)\textasciicircum{}ddag == y\textasciicircum{}ddag x\textasciicircum{}ddag,  (x\textasciicircum{}ddag)\textasciicircum{}ddag == x}
  \item \cfL{GRASSMANN [solver regression]: grD is an even derivation, d(x y) == d(x) y + x d(y);  grDL obeys d\_i(x y) == d\_i(x) y + (-1)\textasciicircum{}|x| x d\_i(y)}
\end{itemize}

Then the two lemmas that drive the section are proved for a completely general, symbolic
$16\times16$ matrix $N$ ($N_A$, $N_S$ its antisymmetric and symmetric parts, $\theta$ sixteen real
Grassmann generators, $\phi=d\theta/dt$ sixteen more):
\begin{align*}
  \theta^{T}N\theta &= \theta^{T}N_A\,\theta \qquad(\text{and } \theta^{T}N_S\,\theta=0),\\
  \theta^{T}N\phi-\tfrac12\,d\bigl(\theta^{T}N\theta\bigr) &= \theta^{T}N_S\,\phi,\\
  \text{Euler--Lagrange expression of } L=\theta^{T}N\phi &= (N+N^{T})\,\phi=2N_S\,\phi .
\end{align*}
So a first-order Grassmann kinetic term is a \textbf{total derivative}, with empty field equations,
precisely when its matrix is antisymmetric. For \textbf{commuting} variables the roles are exactly
reversed (the control). \cfproved{VII \S26}
\begin{itemize}
  \item \cfL{GRASSMANN [THE RESULT]: theta\textasciicircum{}T N theta == theta\textasciicircum{}T N\_A theta and theta\textasciicircum{}T N\_S theta == 0 for a general 16x16 N -- an anticommuting bilinear sees only the ANTISYMMETRIC part}
  \item \cfL{GRASSMANN [THE RESULT]: theta\textasciicircum{}T N phi - (1/2) d(theta\textasciicircum{}T N theta) == theta\textasciicircum{}T N\_S phi for a general N}
  \item \cfL{GRASSMANN [THE RESULT]: the Euler-Lagrange expression of L = theta\textasciicircum{}T N phi (left derivatives) is (N + N\textasciicircum{}T) phi: EMPTY field equations exactly when N is antisymmetric}
  \item \cfL{GRASSMANN [control]: for COMMUTING components the roles are reversed -- theta\textasciicircum{}T N theta sees only N\_S, and the Euler-Lagrange expression of theta\textasciicircum{}T N phi is 2 N\_A phi}
\end{itemize}

\subsection{The real options, and why each fails}
\label{sec:3.4}

With the table and the lemmas the verdicts are immediate, and each is proved in the Grassmann
algebra.

\paragraph{(a) A real Grassmann 16-spinor with the author's $\sigma_{16}$ has no dynamics, on any
frame.} Its only invariant mass terms vanish, $\theta^{T}\sigma_{16}\theta=\theta^{T}C_+\theta=0$
(both matrices are symmetric). In every direction its flat kinetic term
$\theta^{T}\sigma_{16}\Tsix{a}\,d\theta$ is the total derivative
$\tfrac12\,d\bigl(\theta^{T}\sigma_{16}\Tsix{a}\,\theta\bigr)$ ($\sigma_{16}\Tsix{a}$ is
antisymmetric), with an identically empty Euler--Lagrange expression; the total derivative itself is
not zero. On a curved frame the covariant kinetic term is
\[
  \sqrt g\,\theta^{T}\sigma_{16}\gamma^\mu D_\mu\theta=(\text{total derivative})
  +\tfrac12\sqrt g\,\theta^{T}\sigma_{16}\{\gamma^\mu,\Gamma_\mu\}\,\theta ,
\]
because $\partial_\mu(\sqrt g\,\sigma_{16}\gamma^\mu)=\sqrt g\,\sigma_{16}[\gamma^\mu,\Gamma_\mu]$;
and the second term vanishes on \textbf{every} frame, because $\{\Tsix{a},S^{bc}\}$ is a rank-3
product (or zero) and $\sigma_{16}$ times a rank-3 product is symmetric. \cfproved{VII \S26}
\begin{itemize}
  \item \cfL{REFINEMENT (a) [THE RESULT]: real Grassmann + sigma16: both invariant mass terms vanish, and theta\textasciicircum{}T sigma16 T16[a] phi == (1/2) d(theta\textasciicircum{}T sigma16 T16[a] theta) for every a (a total derivative)}
  \item \cfL{REFINEMENT (a) [has content]: its Euler-Lagrange expression is identically empty in every direction -- while the kinetic term itself is NOT zero}
  \item \cfL{REFINEMENT (a) [THE RESULT]: FRAME-INDEPENDENT -- sigma16 \{T16[a], S\textasciicircum{}bc\} is SYMMETRIC for all a, b, c (so theta\textasciicircum{}T sigma16 \{gamma\textasciicircum{}mu, Gamma\_mu\} theta == 0 on every frame), and d\_mu(Sqrt[g] sigma16 gamma\textasciicircum{}mu) == Sqrt[g] sigma16 [gamma\textasciicircum{}mu, Gamma\_mu] on the canonical frame}
\end{itemize}

\paragraph{(b) A Majorana 16-spinor built with $C_+$ propagates, but is massless and has no $V(s)$.}
Now $C_+\Tsix{a}$ is symmetric, so the kinetic term is genuine: its Euler--Lagrange expression is
$2C_+\Tsix{a}\,\phi$, and $C_+\Tsix{a}$ is invertible. But both invariant forms are symmetric, so
there is \textbf{no} Lorentz-scalar bilinear: no mass, no $s$, no $V(s)$. Its quartic Lorentz
scalars are not zero, and they are all \textbf{one} scalar: with the vector
$v_a=\theta^{T}\sigma_{16}\Tsix{a}\,\theta$, the 2-form $B_{ab}=\theta^{T}C_+\Tsix{a}\Tsix{b}\,\theta$
and the 3-form $C_{abc}=\theta^{T}C_+\Tsix{a}\Tsix{b}\Tsix{c}\,\theta$, the contractions $v\cdot v$,
$B\cdot B$ and $C\cdot C$ are proportional to each other (same monomials, one constant ratio each).
Its commuting shadow has no kinetic term ($C_+\Tsix{a}$ symmetric is a total derivative for
commuting components). \cfproved{VII \S26}
\begin{itemize}
  \item \cfL{REFINEMENT (b) [THE RESULT]: Majorana with C+: the kinetic term is genuine -- Euler-Lagrange expression 2 C+ T16[a] phi with C+ T16[a] invertible -- but every invariant bilinear vanishes (massless, no V(s))}
  \item \cfL{REFINEMENT (b) [has content]: the scope -- the quartic Lorentz scalars v.v, B.B, C.C of the Majorana spinor are NOT zero and are all proportional to ONE quartic (same monomials, one constant ratio each)}
  \item \cfL{REFINEMENT (b) [has content]: and the commuting shadow of the C+ Majorana spinor has NO kinetic term: for commuting components theta\textasciicircum{}T C+ T16[a] phi is the total derivative (1/2) d(theta\textasciicircum{}T C+ T16[a] theta)}
\end{itemize}

\paragraph{(c) A Weyl or Majorana--Weyl 8-spinor (one chirality, $P_L$ or $P_R$) cannot propagate.}
Both invariant forms are chirality-\textbf{diagonal} and every $\Tsix{a}$ is
chirality-\textbf{off}-diagonal, so $P_L\,C\,\Tsix{a}\,P_L=P_R\,C\,\Tsix{a}\,P_R=0$: no kinetic term
at all. \cfproved{VII \S26}
\begin{itemize}
  \item \cfL{REFINEMENT (c) [THE RESULT]: Weyl and Majorana-Weyl: PL C T16[a] PL == PR C T16[a] PR == 0 for C = C-, C+ and every a (C chirality-diagonal, T16[a] chirality-odd) -- no kinetic term}
\end{itemize}

\paragraph{(d) Why Part~VI's real commuting fable escapes (a), and why it cannot be quantized as a
fermion.} For commuting components the symmetric $\sigma_{16}$ gives a non-zero mass term and the
antisymmetric $\sigma_{16}\Tsix{a}$ gives a genuine kinetic term (the control of
subsection~\ref{sec:3.3}). That is why Part~VI's classical theory is consistent, and why it cannot be
quantized as a fermion: a commuting spinor violates the spin--statistics connection
\cfprose{VII \S26}.

\subsection{The complex 16-spinor}
\label{sec:3.5}

Take two real Grassmann 16-spinors $\theta_1$, $\theta_2$ (32 real generators) and
\[
  \Psi=\frac{\theta_1+i\,\theta_2}{\sqrt2},\qquad
  \Psi^{\ddagger}=\frac{(\theta_1-i\,\theta_2)^{T}}{\sqrt2},\qquad
  \bar\Psi=\Psi^{\ddagger}\sigma_{16} .
\]
Now everything the real options lacked is present:
\begin{itemize}
  \item $s=\bar\Psi\Psi=i\,\theta_1^{T}\sigma_{16}\theta_2$ is \textbf{non-zero} (16 monomials) and
        \textbf{Hermitian} (real under ${}^{\ddagger}$). The second invariant
        $p=\bar\Psi\,\Tsix{8}\,\Psi=\Psi^{\ddagger}C_+\Psi$ is non-zero and Hermitian as well: the
        complex 16-spinor has \textbf{two} scalar bilinears.
  \item $s$ is a sum of 16 \textbf{commuting nilpotent} even elements $x_a$ ($x_a^2=0$), so every
        polynomial $V(s)$ is a finite polynomial: $s^{17}=0$.
  \item The kinetic term is genuine: $\Psi^{\ddagger}K\phi=\Theta^{T}A\,\Phi$ with
        $\Theta=(\theta_1,\theta_2)$, $\Phi=d\Theta$, $K=\sigma_{16}\Tsix{a}$, and the symmetric part
        of the $32\times32$ matrix $A$ is
        $A_S=\tfrac{i}{2}\bigl(\begin{smallmatrix}0&K\\-K&0\end{smallmatrix}\bigr)=:\tfrac{i}{2}\,\Omega$,
        which is \textbf{not} zero. $\Omega$ is real symmetric and invertible (it is the symplectic
        matrix of \secref{sec:6.3}).
  \item The symmetrized kinetic term $\tfrac12\bigl(\Psi^{\ddagger}K\phi-\phi^{\ddagger}K\Psi\bigr)$
        is \textbf{Hermitian}, and it differs from $\Psi^{\ddagger}K\phi$ by the total derivative
        $-\tfrac12\,d\bigl(\Psi^{\ddagger}K\Psi\bigr)$; the unsymmetrized one is \textbf{not}
        Hermitian.
\end{itemize}
Invariance of $s$ and $p$ under Spin(4,4) is the statement $S^{T}N+NS=0$ of
subsection~\ref{sec:3.2} (the Lorentz generators are real). The self-interaction is taken to depend
on $s$ only; \secref{sec:6.7} shows why $p$ cannot enter. \cfproved{VII \S26}
\begin{itemize}
  \item \cfL{REFINEMENT (d) [definition]: Psi\textasciicircum{}ddag == (theta1 - i theta2)\textasciicircum{}T/Sqrt[2] in the algebra}
  \item \cfL{REFINEMENT (d) [THE RESULT]: s = Psibar Psi == i theta1\textasciicircum{}T sigma16 theta2: NON-ZERO (16 monomials) and HERMITIAN; p = Psibar T16[8] Psi is non-zero and Hermitian too}
  \item \cfL{REFINEMENT (d) [has content]: s is a sum of 16 commuting nilpotent even elements x\_a (x\_a x\_b == x\_b x\_a, x\_a\textasciicircum{}2 == 0), so s\textasciicircum{}17 == 0 and V(s) is a finite polynomial}
  \item \cfL{REFINEMENT (d) [definition]: the 32x32 matrix A read off the algebra reproduces Psi\textasciicircum{}ddag K phi exactly}
  \item \cfL{REFINEMENT (d) [THE RESULT]: its symmetric part is A\_S == (i/2) Omega, Omega = [[0, K], [-K, 0]] real SYMMETRIC and invertible: a genuine kinetic term}
  \item \cfL{REFINEMENT (d) [has content]: the symmetrized kinetic term is HERMITIAN and equals Psi\textasciicircum{}ddag K phi - (1/2) d(Psi\textasciicircum{}ddag K Psi); the unsymmetrized one is not Hermitian}
  \item \cfL{REFINEMENT (d) [has content]: the same holds in every direction a = 0..7: the symmetrized Psi\textasciicircum{}ddag sigma16 T16[a] phi is Hermitian, with symmetric part (i/2)[[0, sigma16 T16[a]], [-sigma16 T16[a], 0]] != 0}
\end{itemize}

\subsection{Four four-dimensional Dirac fermions}
\label{sec:3.6}

The observer's Clifford algebra is generated by $\Tsix{1}$, $\Tsix{2}$, $\Tsix{3}$ (the observed
space) and $\Tsix{4}$ (the observer's time). Its 16 ordered products are linearly independent, so
over the complex numbers it is the full matrix algebra $M_4(\mathbb C)$, the algebra of $4\times4$
Dirac matrices. Its \textbf{commutant} in $M_{16}(\mathbb C)$, computed by solving
$X\Tsix{i}=\Tsix{i}X$ ($i=1..4$), is again 16-dimensional, generated by the four ``hidden''
(flavour) gammas
\begin{gather*}
  h_A=\Omega_4\,\Tsix{A},\quad A=0,5,6,7,\qquad
  \Omega_4=\Tsix{1}\Tsix{2}\Tsix{3}\Tsix{4},\quad \Omega_4^2=-\IDs,\\
  \{h_A,h_B\}=-2\,\eta_{AB}\,\IDs ,
\end{gather*}
which commute with $\Tsix{1..4}$: a second Clifford algebra, of the hidden directions. The two
algebras meet only in the multiples of the identity (the rank of their union is $16+16-1=31$). By
the double-commutant theorem
\[
  \mathbb C^{16}=\mathbb C^{4}\,(\text{the observer's Dirac index})\;\otimes\;\mathbb C^{4}\,(\text{the hidden flavour index}) .
\]
Under the observer's Lorentz group Spin(3,1), generated by $\Tsix{i}\Tsix{j}$ with $i,j$ in
$\{1,2,3,4\}$, fermion fable is \textbf{four 4-dimensional Dirac fermions}. Per momentum that is
$4\times2=8$ particle states (and 8 antiparticle states): the degeneracy $g=8$ of sections
\ref{sec:6} and~\ref{sec:8}. \cfproved{VII \S26}
\begin{itemize}
  \item \cfL{4D [THE RESULT]: the 16 ordered products of T16[1..4] are linearly independent (the observer's algebra is M4(C)), and its commutant in M16(C) is 16-dimensional}
  \item \cfL{4D [has content]: Omega4\textasciicircum{}2 == -ID16, the hidden gammas h\_A = Omega4 T16[A] (A = 0,5,6,7) commute with T16[1..4] and obey \{h\_A, h\_B\} == -2 eta\_AB}
  \item \cfL{4D [THE RESULT]: the 16 products of the hidden gammas span the commutant, and the two algebras meet only in the multiples of ID16 (rank of the union 16 + 16 - 1 = 31): C\textasciicircum{}16 = C\textasciicircum{}4 (x) C\textasciicircum{}4, four 4D Dirac fermions}
  \item \cfL{4D [has content]: the observer's Lorentz generators T16[i] T16[j] (i < j in 1..4) commute with every hidden gamma}
\end{itemize}

\subsection{The refined Lagrangian, and the fidelity chain to Part VI, \texorpdfstring{$L_a$}{La} and \texttt{eLa}}
\label{sec:3.7}

The Lagrangian of fermion fable is the \textbf{symmetrized covariant} one (\secref{sec:4} shows why
no other form is admissible for a complex field):
\begin{gather*}
  L=\sqrt{\det g}\;\hat L,\qquad
  \hat L=\frac{1}{2H}\Bigl[\bar\Psi\gamma^\mu D_\mu\Psi-(D_\mu\bar\Psi)\gamma^\mu\Psi\Bigr]-V(s),\qquad s=\bar\Psi\Psi,\\
  D_\mu\Psi=\partial_\mu\Psi+\Gamma^{\mathrm{spin}}_\mu\Psi,\qquad
  D_\mu\bar\Psi=\partial_\mu\bar\Psi-\bar\Psi\,\Gamma^{\mathrm{spin}}_\mu ,
\end{gather*}
and its energy--momentum tensor (\secref{sec:7}) is
\[
  \hat T_{\mu\nu}=-\frac{1}{4H}\Bigl[\bar\Psi\gamma_\mu D_\nu\Psi-(D_\nu\bar\Psi)\gamma_\mu\Psi+(\mu\leftrightarrow\nu)\Bigr]+g_{\mu\nu}\hat L .
\]
Four helpers implement them for an arbitrary frame (curved gammas and spin connection passed in):
\cfcode{cfLhatFermion}, \cfcode{cfEquationFermion} ($\gamma^\mu D_\mu\Psi-HV'(s)\Psi$),
\cfcode{cfAdjointEquationFermion} ($(D_\mu\bar\Psi)\gamma^\mu+HV'(s)\bar\Psi$) and
\cfcode{cfTFermion}. Before they are used for anything new they are anchored to what the notebook
already knows. Set $\Psi\to\psi$ (a real commuting spinor) and $\bar\Psi\to\psi^{T}\sigma_{16}$ (the
``bosonic shadow''). Then:
\begin{enumerate}
  \item $\hat L$ becomes Part~VI's \cfcode{cfLhatFable} \textbf{exactly}, on the canonical frame,
        with generic $V$;
  \item on the flat frame with $V=-(2M/H)\,s$ and no volume factor, it becomes the author's
        \cfcode{La[]} of Section~11 (\secref{sec:2.7});
  \item the Euler--Lagrange expressions of the complex field, combined as
        $\sigma_{16}\cdot\EL_{\bar\Psi}+\EL_{\Psi}$ (the chain rule for
        $\bar\Psi=\psi^{T}\sigma_{16}$), become Part~VI's explicit variation \cfcode{cfELFable},
        and on the flat frame the author's sixteen equations \cfcode{eLa} of Section~12
        (\secref{sec:2.8});
  \item $\hat T$ becomes Part~VI's covariant tensor \cfcode{cfTCov} (\secref{sec:2.12.3}), all 64
        components.
\end{enumerate}
\cfproved{VII \S26}
\begin{itemize}
  \item \cfL{FIDELITY [1] [fidelity]: Psi -> psi, Psibar -> psi\textasciicircum{}T sigma16 turns the symmetrized covariant Lhat into Part VI's cfLhatFable EXACTLY (canonical frame, generic V)}
  \item \cfL{FIDELITY [2] [fidelity]: on the flat frame with V = -(2M/H) s and no volume factor, the same limit on Psi16 IS the author's La[] of Section 11}
  \item \cfL{FIDELITY [3] [fidelity]: sigma16 . EL\_Psibar + EL\_Psi in the real limit == Part VI's explicit variation cfELFable, component for component}
  \item \cfL{FIDELITY [3] [fidelity]: and on the flat frame with V = -(2M/H) s the same combination IS the author's eLa of Section 12}
  \item \cfL{FIDELITY [4] [fidelity]: the energy-momentum tensor That in the real limit IS Part VI's covariant tensor cfTCov, all 64 components}
\end{itemize}
Every new result of this effort therefore contains Part~VI, and through it the author's $L_a$ and
\cfcode{eLa}, as its real commuting limit.

\subsection{The conclusion, with its scope}
\label{sec:3.8}

The complex 16-spinor is the \textbf{minimal} field that
\begin{enumerate}
  \item uses the author's $\sigma_{16}$ conjugation,
  \item propagates,
  \item admits a potential $V(s)$ with $s=\bar\Psi\Psi$, and
  \item reduces to Part~VI's classical fable as its real commuting shadow.
\end{enumerate}
It is \textbf{not} the unique field with dynamics: the $C_+$ Majorana spinor propagates too, but it
is massless, has no bilinear $V(s)$, and has no commuting shadow with a kinetic term. The statement
of minimality is the conclusion of Section~26 as a whole; no single assertion states it. It rests
on the assertions of subsections \ref{sec:3.1}--\ref{sec:3.7} \cfprose{VII \S26, assembled from the
assertions quoted}.

% ===============================================================================================
\section{The Lagrangian and its Hermiticity}
\label{sec:4}

This is the first cell of Section~27 (Input cell 180, 10 assertions). From here on $\Psi$ and
$\bar\Psi$ are two independent 16-component fields of \textbf{all eight} coordinates (the commuting
proxy of \secref{sec:3}, $\bar\Psi$ always on the left).

\subsection{The four candidates}
\label{sec:4.1}

Four candidate Lagrangians differ by where the derivative acts and whether it is covariant:
\begin{align*}
  \Lsym &= \frac{1}{2H}\Bigl[\bar\Psi\gamma^\mu D_\mu\Psi-(D_\mu\bar\Psi)\gamma^\mu\Psi\Bigr]-V(s)
          &&\text{(the Lagrangian of fermion fable)},\\
  \Lsymd &= \text{the same with } \partial_\mu \text{ in place of } D_\mu,\\
  \Lun &= \frac1H\,\bar\Psi\gamma^\mu D_\mu\Psi-V(s),\\
  \Lund &= \frac1H\,\bar\Psi\gamma^\mu\partial_\mu\Psi-V(s),
\end{align*}
with $D_\mu\Psi=\partial_\mu\Psi+\Gamma_\mu\Psi$, $D_\mu\bar\Psi=\partial_\mu\bar\Psi-\bar\Psi\Gamma_\mu$,
and $\Gamma_\mu=$ \cfcode{GammaSpinCanonical[[mu+1]]}, the canonical spin connection of
\secref{sec:9}. (The notebook's labels call them \cfcode{Lsym}, \cfcode{Lsym_d}, \cfcode{Lun} and
\cfcode{Lun_d}.) The action is $S=\int d^8x\,\sqrt{\det g}\,\Lsym$, with
$\sqrt{\det g}=\sec 6Hx_0$ on the canonical frame. There is \textbf{no factor $i$} in front of the
kinetic term, as in the author's $L_a$: \secref{sec:4.2}(b) shows that none is needed.

\subsection{The facts, each asserted on the canonical frame for generic \texorpdfstring{$\Psi(x_0..x_7)$, $\bar\Psi(x_0..x_7)$}{Psi(x0..x7), Psibar(x0..x7)}}
\label{sec:4.2}

\paragraph{(a) The connection is real and $\sigma_{16}\Gamma_\mu$ is antisymmetric,} i.e.
$\Gamma_\mu^{T}\sigma_{16}=-\sigma_{16}\Gamma_\mu$. So $D_\mu\bar\Psi=\partial_\mu\bar\Psi-\bar\Psi\Gamma_\mu$
is exactly the Dirac conjugate of $D_\mu\Psi$. \cfproved{VII \S27}
\cfL{LAGRANGIAN (a) [definition]: every Gamma\textasciicircum{}spin\_mu is real and sigma16 Gamma\_mu is antisymmetric (Gamma\_mu\textasciicircum{}T sigma16 == -sigma16 Gamma\_mu): D\_mu Psibar is the Dirac conjugate of D\_mu Psi}

\paragraph{(b) $\Lsym$ is Hermitian; $\Lun$ is not.} With $\Psi=u+iv$ and
$\bar\Psi=(u-iv)^{T}\sigma_{16}$, $\Lsym$ is real. No factor $i$ is needed, because
$\sigma_{16}\gamma^\mu$ is real and antisymmetric:
$\bigl(\Psi^{\ddagger}K\,d\Psi\bigr)^{\ddagger}=-(d\Psi^{\ddagger})K\Psi$ for real antisymmetric $K$, so
the symmetrized combination is real. $\Lun$ is not Hermitian: its imaginary part is not zero (a
witness). \cfproved{VII \S27}
\begin{itemize}
  \item \cfL{LAGRANGIAN (b) [THE RESULT]: Lsym is HERMITIAN -- real for Psi = u + i v, Psibar = (u - i v)\textasciicircum{}T sigma16, generic u, v of all eight coordinates}
  \item \cfL{LAGRANGIAN (b) [control]: Lun is NOT Hermitian -- its imaginary part is not zero, witnessed on the constant spinors u = e1, v = e13 (with V = s\textasciicircum{}2)}
\end{itemize}

\paragraph{(c) $\Lsym$ contains the connection only through $\{\gamma^\mu,\Gamma_\mu\}$, and that
matrix vanishes for each $\mu$.} Expanding $D$,
\[
  \Lsym-\Lsymd=\frac{1}{2H}\sum_\mu\bar\Psi\,\{\gamma^\mu,\Gamma_\mu\}\,\Psi\qquad\text{(identically, off shell)}.
\]
On the canonical frame $\{\gamma^\mu,\Gamma_\mu\}=0$ for \textbf{each} $\mu$ separately: every
$\Gamma_\mu$ is a combination of $\Tsix{\mu}\Tsix{0}$ and $\Tsix{\mu}\Tsix{4}$, with no totally
antisymmetric part (\secref{sec:9.8}). So $\Lsym=\Lsymd$ there: the connection drops out of the
symmetrized Lagrangian on every frame on which the \textbf{matrix} $\{\gamma^\mu,\Gamma_\mu\}$
vanishes (diagonal frames), and \textbf{not} in general. \cfproved{VII \S27}
\cfL{LAGRANGIAN (c) [THE RESULT]: Lsym - Lsym\_d == (1/(2H)) Sum\_mu Psibar \{gamma\textasciicircum{}mu, Gamma\_mu\} Psi identically, and \{gamma\textasciicircum{}mu, Gamma\_mu\} == 0 for EACH mu on the canonical frame: Lsym == Lsym\_d}

\paragraph{(d) $\Lun$ and $\Lsym$ differ by a total divergence.} \cfproved{VII \S27}
\cfL{LAGRANGIAN (d) [has content]: Lun - Lsym == (1/(2H)) (1/Sqrt[g]) d\_mu( Sqrt[g] Psibar gamma\textasciicircum{}mu Psi ), a total divergence}

\paragraph{(e) But $\Lun$ and $\Lund$ differ by a term that is not zero for a complex field.}
\[
  \Lun-\Lund=\frac1H\,\bar\Psi\gamma^\mu\Gamma_\mu\Psi=-3\cot^2(6Hx_0)\;\bar\Psi\,\Tsix{0}\,\Psi .
\]
This vanishes only in the real commuting limit, which is the scope of the statement ``$\hat L$ with
$D$ $=$ $\hat L$ with $\partial$'' of Parts V and VI (\secref{sec:2.12.1}). \cfproved{VII \S27}
\begin{itemize}
  \item \cfL{LAGRANGIAN (e) [has content]: Lun - Lun\_d == (1/H) Psibar gamma\textasciicircum{}mu Gamma\_mu Psi == -3 Cot[6 H x0]\textasciicircum{}2 Psibar T16[0] Psi}
  \item \cfL{LAGRANGIAN (e) [control]: that term is NOT zero for a complex field (witness u = e1, v = e13), and IS zero in the real commuting limit (the scope of Parts V-VI)}
\end{itemize}

\paragraph{(f) So $\Lund$ is inadmissible.} Varying $\bar\Psi$ gives
$\gamma^\mu\partial_\mu\Psi=HV'\Psi$ with the connection term \textbf{missing}. Varying $\Psi$
gives an adjoint equation that \textbf{contains} it,
$(\partial_\mu\bar\Psi)\gamma^\mu+\bar\Psi[\gamma^\mu,\Gamma_\mu]=-HV'\bar\Psi$. The two are
\textbf{not} Dirac conjugates of each other: they differ by
\[
  \bar\Psi[\gamma^\mu,\Gamma_\mu]=2\bar\Psi\gamma^\mu\Gamma_\mu=-6H\cot^2(6Hx_0)\;\bar\Psi\,\Tsix{0}
\]
(a witness). Only a Lagrangian whose two variations are conjugate defines a consistent theory.
\cfproved{VII \S27}
\begin{itemize}
  \item \cfL{LAGRANGIAN (f) [THE RESULT]: Lun\_d varied with respect to Psibar gives (Sqrt[g]/H)(gamma\textasciicircum{}mu d\_mu Psi - H V' Psi): the connection term is MISSING}
  \item \cfL{LAGRANGIAN (f) [THE RESULT]: varied with respect to Psi it gives -(Sqrt[g]/H)((d\_mu Psibar) gamma\textasciicircum{}mu + Psibar [gamma\textasciicircum{}mu, Gamma\_mu] + H V' Psibar), with Psibar [gamma\textasciicircum{}mu, Gamma\_mu] == 2 Psibar gamma\textasciicircum{}mu Gamma\_mu: the connection term is PRESENT}
  \item \cfL{LAGRANGIAN (f) [control]: so Lun\_d is INADMISSIBLE -- its two field equations are not Dirac conjugates: they differ by 2 Psibar gamma\textasciicircum{}mu Gamma\_mu = -6 H Cot\textasciicircum{}2 Psibar T16[0], non-zero (witness u = e1, v = 0, V = s\textasciicircum{}2)}
\end{itemize}

\subsection{Why the connection drops out, and why this is not Part VI's mechanism}
\label{sec:4.3}

For the real commuting field of Part~VI the connection dropped out because
$\Psi^{T}\sigma_{16}\Tsix{a}\Psi=0$ for commuting components ($\sigma_{16}\Tsix{a}$ antisymmetric;
\secref{sec:2.11}). For a complex field that bilinear identity is false:
$\bar\Psi\,\Tsix{0}\,\Psi=\Psi^{\ddagger}(\sigma_{16}\Tsix{0})\Psi$ with $\sigma_{16}\Tsix{0}$ real
antisymmetric is $i$ times the Hermitian form of $-i\sigma_{16}\Tsix{0}$, which is not zero in
general \cfderived. The drop-out of the connection from $\Lsym$ is therefore a property of the
\textbf{matrix} $\{\gamma^\mu,\Gamma_\mu\}$ (fact (c)), not of a bilinear identity. The general
identity behind it is $\sum_\mu\{\gamma^\mu,\Gamma_\mu\}=\tfrac12\,\omega_{[cab]}\,\gamma^{cab}$:
only the totally antisymmetric part of the connection can enter the symmetrized Lagrangian
(\secref{sec:9.8}).

Hermiticity rests on two facts: $\sigma_{16}\gamma^a$ and $\sigma_{16}\Gamma_\mu$ are \textbf{real
antisymmetric}. The first is Part~I's \cfL{sigma16 . T16[A] is antisymmetric for A = 0..7}; the second is
fact (a).

% ===============================================================================================
\section{The field equations in the primordial gravitational field, written out}
\label{sec:5}

This is Section~27 of the notebook (Input cells 181, 182 and 184) and item (2) of the pre-universe
cell of Section~29 (Input cell 197). The \textbf{primordial gravitational field} is the author's
canonical frame of \secref{sec:2.11}:
\begin{gather*}
  \texttt{frameCanonical}=\diag\bigl(\tan 6Hx_0,\ q,\ q,\ q,\ 1,\ p,\ p,\ p\bigr),\\
  q=\frac{e^{-a_4(Hx_4)}}{\sin^{1/6}(6Hx_0)},\qquad p=\frac{e^{+a_4(Hx_4)}}{\sin^{1/6}(6Hx_0)}\qquad(a_4\text{ free, never given a value}),\\
  ds^2=\tan^2(6Hx_0)\,dx_0^2+q^2\bigl(dx_1^2+dx_2^2+dx_3^2\bigr)-dx_4^2-p^2\bigl(dx_5^2+dx_6^2+dx_7^2\bigr),\\
  \sqrt{|\det g|}=\sec 6Hx_0,\\
  \gamma^0=\cot(6Hx_0)\,\Tsix{0},\qquad \gamma^i=\tfrac1q\,\Tsix{i}\ \ (i=1,2,3),\\
  \gamma^4=\Tsix{4},\qquad \gamma^h=\tfrac1p\,\Tsix{h}\ (h=5,6,7),
\end{gather*}
on $0<6Hx_0<\pi/2$. The spin connection enters through
$\gamma^\mu\Gamma_\mu=-3H\cot^2(6Hx_0)\,\Tsix{0}$ (\secref{sec:9.8}).

\subsection{The derivation}
\label{sec:5.1}

Vary $\bar\Psi$ and $\Psi$ independently in $L=\sqrt g\,\Lsym$ (the left derivative with respect to
$\bar\Psi$ and the right derivative with respect to $\Psi$ are the ordinary derivatives of the
commuting proxy). For generic fields of all eight coordinates:
\begin{align*}
  \EL_{\bar\Psi} &= \frac{\sqrt g}{H}\,\bigl(\gamma^\mu D_\mu\Psi-H\,V'(s)\,\Psi\bigr),\\
  \EL_{\Psi} &= -\frac{\sqrt g}{H}\,\bigl((D_\mu\bar\Psi)\gamma^\mu+H\,V'(s)\,\bar\Psi\bigr),
\end{align*}
so \textbf{the field equation and the adjoint field equation of fermion fable} are
\[
  \boxed{\;\gamma^\mu D_\mu\Psi=H\,V'(s)\,\Psi,\qquad (D_\mu\bar\Psi)\gamma^\mu=-H\,V'(s)\,\bar\Psi\;}
\]
The \textbf{sign} of the adjoint equation is derived, not assumed, and it is exactly the one
consistency requires: with $\Psi=u+iv$ and $\bar\Psi=(u-iv)^{T}\sigma_{16}$, the adjoint equation is
the Dirac conjugate of the field equation, (adjoint residual) $=-$(field residual)${}^{*}\cdot\sigma_{16}$,
because $\Tsix{a}^{T}\sigma_{16}=-\sigma_{16}\Tsix{a}$. The opposite sign fails that test (control).
The derivation uses only the covariant constancy of $\gamma^\mu$, in the form
$\frac{1}{\sqrt g}\,\partial_\mu\bigl(\sqrt g\,\gamma^\mu\bigr)=[\gamma^\mu,\Gamma_\mu]$ (Section~17);
the anticommutator is not needed for the symmetrized Lagrangian. \cfproved{VII \S27}
\begin{itemize}
  \item \cfL{FIELD EQUATION [has content]: the identity the derivation uses -- (1/Sqrt[g]) d\_mu(Sqrt[g] gamma\textasciicircum{}mu) == [gamma\textasciicircum{}mu, Gamma\_mu] (covariant constancy, Section 17)}
  \item \cfL{FIELD EQUATION [THE RESULT]: EL\_Psibar == (Sqrt[g]/H)(gamma\textasciicircum{}mu D\_mu Psi - H V'(s) Psi) for generic Psi, Psibar of all eight coordinates: the field equation is gamma\textasciicircum{}mu D\_mu Psi == H V'(s) Psi}
  \item \cfL{FIELD EQUATION [THE RESULT]: EL\_Psi == -(Sqrt[g]/H)((D\_mu Psibar) gamma\textasciicircum{}mu + H V'(s) Psibar): the adjoint equation is (D\_mu Psibar) gamma\textasciicircum{}mu == -H V'(s) Psibar}
  \item \cfL{FIELD EQUATION [has content]: the adjoint equation IS the Dirac conjugate of the field equation: with Psi = u + i v, Psibar = (u - i v)\textasciicircum{}T sigma16, adjoint residual == -(field residual)\textasciicircum{}* . sigma16}
  \item \cfL{FIELD EQUATION [control]: the opposite sign, (D Psibar) gamma == +H V' Psibar, is NOT the conjugate -- witnessed on u = e1 + e5, v = 0 with V = s\textasciicircum{}2 (s = -2 there)}
\end{itemize}

\subsection{The \texorpdfstring{$16\times16$}{16x16} matrix form on the canonical frame}
\label{sec:5.2}

For fields of all eight coordinates, with $\partial_\mu=\partial/\partial x_\mu$, sums over $i=1,2,3$
and $h=5,6,7$:
\begin{align*}
  &\cot(6Hx_0)\,\Tsix{0}\,\partial_0\Psi+\frac1q\sum_i\Tsix{i}\,\partial_i\Psi+\Tsix{4}\,\partial_4\Psi
  +\frac1p\sum_h\Tsix{h}\,\partial_h\Psi\\
  &\qquad-3H\cot^2(6Hx_0)\,\Tsix{0}\,\Psi=H\,V'(s)\,\Psi,\\[6pt]
  &\cot(6Hx_0)\,\partial_0\bar\Psi\,\Tsix{0}+\frac1q\sum_i\partial_i\bar\Psi\,\Tsix{i}+\partial_4\bar\Psi\,\Tsix{4}
  +\frac1p\sum_h\partial_h\bar\Psi\,\Tsix{h}\\
  &\qquad-3H\cot^2(6Hx_0)\,\bar\Psi\,\Tsix{0}=-H\,V'(s)\,\bar\Psi .
\end{align*}
In the adjoint equation the connection enters as
$-\bar\Psi\Gamma_\mu\gamma^\mu=+\bar\Psi\gamma^\mu\Gamma_\mu$ (per-$\mu$ anticommutator), so it carries
the \textbf{same} coefficient $-3H\cot^2$ as in the field equation. \cfproved{VII \S27}
\begin{itemize}
  \item \cfL{WRITTEN OUT (i) [THE RESULT]: the field equation on the canonical frame is Cot T16[0] d\_0 Psi + (1/q) T16[i] d\_i Psi + T16[4] d\_4 Psi + (1/p) T16[h] d\_h Psi - 3 H Cot\textasciicircum{}2 T16[0] Psi == H V'(s) Psi}
  \item \cfL{WRITTEN OUT (i) [THE RESULT]: the adjoint equation is Cot d\_0 Psibar T16[0] + ... - 3 H Cot\textasciicircum{}2 Psibar T16[0] == -H V'(s) Psibar (the connection enters with the sign fixed by \{gamma\textasciicircum{}mu, Gamma\_mu\} = 0)}
\end{itemize}

\subsection{The split-octonion 8+8 form}
\label{sec:5.3}

Write $\Psi=(\psi_1,\psi_2)$: $\psi_1$ = the upper eight components (type-1 split-octonion spinor,
chirality $-1$), $\psi_2$ = the lower eight (type-2, chirality $+1$). With
$\Tsix{a}=\bigl(\begin{smallmatrix}0&\bar\tau_a\\ \tau_a&0\end{smallmatrix}\bigr)$ and
$\Gamma^{\mathrm{spin}}_\mu=\diag\bigl(\Gamma^{(1)}_\mu,\Gamma^{(2)}_\mu\bigr)$ (Section~17; the spin
connection never mixes the two types), \textbf{covariantly, on any frame}:
\begin{align*}
  e_a{}^\mu\,\bar\tau_a\bigl(\partial_\mu+\Gamma^{(2)}_\mu\bigr)\psi_2&=H\,V'(s)\,\psi_1 &&\text{(rows 1..8)},\\
  e_a{}^\mu\,\tau_a\bigl(\partial_\mu+\Gamma^{(1)}_\mu\bigr)\psi_1&=H\,V'(s)\,\psi_2 &&\text{(rows 9..16)},\\
  s=-\psi_1^{\ddagger}\sigma\psi_1+\psi_2^{\ddagger}\sigma\psi_2,&\qquad \bar\psi_1=-\psi_1^{\ddagger}\sigma,\quad \bar\psi_2=\psi_2^{\ddagger}\sigma,\\
  \bigl(\partial_\mu\bar\psi_2-\bar\psi_2\Gamma^{(2)}_\mu\bigr)\tau_a\,e_a{}^\mu&=-H\,V'(s)\,\bar\psi_1,\\
  \bigl(\partial_\mu\bar\psi_1-\bar\psi_1\Gamma^{(1)}_\mu\bigr)\bar\tau_a\,e_a{}^\mu&=-H\,V'(s)\,\bar\psi_2
\end{align*}
($\bar\Psi=(\bar\psi_1,\bar\psi_2)=\Psi^{\ddagger}\sigma_{16}$, $\sigma_{16}=\diag(-\sigma,\sigma)$;
$\tau_a=$ \cfcode{tau[a]}, $\bar\tau_a=$ \cfcode{taubar[a]}.) On the canonical frame
($\tau_0=\bar\tau_0=\mathrm{ID8}$), the connection term is $-3H\cot^2\,\psi_2$ in the upper rows and
$-3H\cot^2\,\psi_1$ in the lower rows:
\begin{align*}
  \cot\,\partial_0\psi_2+\frac1q\sum_i\bar\tau_i\,\partial_i\psi_2+\bar\tau_4\,\partial_4\psi_2+\frac1p\sum_h\bar\tau_h\,\partial_h\psi_2-3H\cot^2\psi_2&=H\,V'(s)\,\psi_1,\\
  \cot\,\partial_0\psi_1+\frac1q\sum_i\tau_i\,\partial_i\psi_1+\tau_4\,\partial_4\psi_1+\frac1p\sum_h\tau_h\,\partial_h\psi_1-3H\cot^2\psi_1&=H\,V'(s)\,\psi_2 ,
\end{align*}
with $\cot=\cot(6Hx_0)$. \cfproved{VII \S27}
\begin{itemize}
  \item \cfL{WRITTEN OUT (ii) [THE RESULT]: the split-octonion 8+8 form -- the type-1 rows are the taubar equations for psi2, the type-2 rows the tau equations for psi1 (tau[0] == taubar[0] == ID8)}
  \item \cfL{WRITTEN OUT (ii) [THE RESULT]: the covariant 8+8 form -- rows 1..8 are e\_a\textasciicircum{}mu taubar[a] (d\_mu + Gamma\textasciicircum{}(2)\_mu) psi2 - H V' psi1, rows 9..16 are e\_a\textasciicircum{}mu tau[a] (d\_mu + Gamma\textasciicircum{}(1)\_mu) psi1 - H V' psi2; and the conjugate rows}
  \item \cfL{WRITTEN OUT (ii) [has content]: with Psibar = Psi\textasciicircum{}ddag sigma16 and sigma16 = diag(-sigma, sigma): s == -psi1\textasciicircum{}ddag sigma psi1 + psi2\textasciicircum{}ddag sigma psi2 (psibar1 = -psi1\textasciicircum{}ddag sigma, psibar2 = psi2\textasciicircum{}ddag sigma)}
\end{itemize}

\subsection{All sixteen component equations}
\label{sec:5.4}

Each $\Tsix{a}$ is a signed permutation matrix, so component equation $r$ contains exactly
\textbf{one} derivative per coordinate, of the component that $\Tsix{a}$ maps to $r$; the connection
term multiplies the same component as the $x_0$-derivative; and $-HV'$ multiplies $\psi_r$ itself.
For fields of all eight coordinates the sixteen equations form \textbf{one} coupled block. (The four
blocks of four of Section~13, \secref{sec:2.9}, belong to the reduction to fields of $(x_0,x_4)$
only.) The term lists \cfcode{cfFieldEqTerms} and \cfcode{cfAdjEqTerms} from which the equations
below are rendered reproduce all sixteen field equations and all sixteen adjoint equations exactly.
\cfproved{VII \S27}
\begin{itemize}
  \item \cfL{WRITTEN OUT (iii) [definition]: the term lists reproduce all sixteen field equations and all sixteen adjoint equations exactly}
  \item \cfL{WRITTEN OUT (iii) [has content]: every component equation has exactly one derivative per coordinate x0..x7, the connection term multiplies the component carrying the x0-derivative, and -H V' multiplies psi\_r itself}
  \item \cfL{WRITTEN OUT (iii) [has content]: for fields of all eight coordinates the sixteen equations form ONE coupled block (the four blocks of four of Section 13 belong to the (x0, x4) reduction)}
\end{itemize}

\paragraph{How to read the equations below.} They are the TeX export written by
\file{claude-fable/export_fermion_fable_tex.wls} from the notebook's own asserted objects
(\file{provenance-latex/generated/fermion_fable_field_equations.tex}, included here unchanged; the
same equations in plain text are the \file{.txt} file of the same name).
\begin{itemize}
  \item $\psi_k$ ($k=0..15$) is component $k$ of $\Psi$; $\bar\psi_k$ is component $k$ of $\bar\Psi$.
        The index is the notebook's 0-based component index (array position $k+1$).
  \item $\partial_\mu\psi_k$ is $\partial\psi_k/\partial x_\mu$, $\mu=0..7$.
  \item $a_4$ is the free function. Hence $e^{a_4(Hx_4)}\sqrt[6]{\sin(6Hx_0)}=1/q$ multiplies the
        observed derivatives $\partial_1,\partial_2,\partial_3$, and
        $e^{-a_4(Hx_4)}\sqrt[6]{\sin(6Hx_0)}=1/p$ multiplies the hidden derivatives
        $\partial_5,\partial_6,\partial_7$.
  \item $V'(s)$ is evaluated at $s=\bar\Psi\Psi$; $H$ is the author's Protected constant.
  \item Equation $E_r$ is row $r$ of $\gamma^\mu D_\mu\Psi-HV'(s)\Psi=0$, with the $V'$ term moved to
        the right.
\end{itemize}

{\small
\input{generated/fermion_fable_field_equations}
}

\subsection{All sixteen adjoint equations}
\label{sec:5.5}

Equation $\bar E_r$ is column $r$ of $(D_\mu\bar\Psi)\gamma^\mu+HV'(s)\bar\Psi=0$, with the $V'$ term
moved to the right (\file{provenance-latex/generated/fermion_fable_adjoint_equations.tex}, unchanged):

{\small
% fermion fable: the sixteen component adjoint equations (D_mu Psibar) gamma^mu = -H V'(s) Psibar
% generated by claude-fable/export_fermion_fable_tex.wls from the notebook's Part VII objects, 2026-09-24
% notation: x_0..x_7 are the coordinates (x_4 = the observer's time), a_4 = the free function a4 (never given a value),
% q = e^{-a_4(H x_4)}/sin^{1/6}(6 H x_0), p = e^{+a_4(H x_4)}/sin^{1/6}(6 H x_0), s = bar(psi) psi.
\begin{align*}
(\bar E_{0})\quad & \cot \left(6 H x_0\right)\,\partial_{0}\bar\psi_{8} - e^{-a_4\left(H x_4\right)} \sqrt[6]{\sin \left(6 H x_0\right)}\,\partial_{7}\bar\psi_{8} - e^{a_4\left(H x_4\right)} \sqrt[6]{\sin \left(6 H x_0\right)}\,\partial_{3}\bar\psi_{13} \\
 &  + \partial_{4}\bar\psi_{13} + e^{a_4\left(H x_4\right)} \sqrt[6]{\sin \left(6 H x_0\right)}\,\partial_{2}\bar\psi_{14} - e^{-a_4\left(H x_4\right)} \sqrt[6]{\sin \left(6 H x_0\right)}\,\partial_{5}\bar\psi_{14} \\
 &  - e^{a_4\left(H x_4\right)} \sqrt[6]{\sin \left(6 H x_0\right)}\,\partial_{1}\bar\psi_{15} - e^{-a_4\left(H x_4\right)} \sqrt[6]{\sin \left(6 H x_0\right)}\,\partial_{6}\bar\psi_{15} - 3 H \cot ^2\left(6 H x_0\right)\,\bar\psi_{8} \\
 & = -H V'(s)\,\bar\psi_{0}
\end{align*}
\begin{align*}
(\bar E_{1})\quad & \cot \left(6 H x_0\right)\,\partial_{0}\bar\psi_{9} - e^{-a_4\left(H x_4\right)} \sqrt[6]{\sin \left(6 H x_0\right)}\,\partial_{7}\bar\psi_{9} + e^{a_4\left(H x_4\right)} \sqrt[6]{\sin \left(6 H x_0\right)}\,\partial_{3}\bar\psi_{12} \\
 &  - \partial_{4}\bar\psi_{12} - e^{a_4\left(H x_4\right)} \sqrt[6]{\sin \left(6 H x_0\right)}\,\partial_{1}\bar\psi_{14} + e^{-a_4\left(H x_4\right)} \sqrt[6]{\sin \left(6 H x_0\right)}\,\partial_{6}\bar\psi_{14} \\
 &  - e^{a_4\left(H x_4\right)} \sqrt[6]{\sin \left(6 H x_0\right)}\,\partial_{2}\bar\psi_{15} - e^{-a_4\left(H x_4\right)} \sqrt[6]{\sin \left(6 H x_0\right)}\,\partial_{5}\bar\psi_{15} - 3 H \cot ^2\left(6 H x_0\right)\,\bar\psi_{9} \\
 & = -H V'(s)\,\bar\psi_{1}
\end{align*}
\begin{align*}
(\bar E_{2})\quad & \cot \left(6 H x_0\right)\,\partial_{0}\bar\psi_{10} - e^{-a_4\left(H x_4\right)} \sqrt[6]{\sin \left(6 H x_0\right)}\,\partial_{7}\bar\psi_{10} - e^{a_4\left(H x_4\right)} \sqrt[6]{\sin \left(6 H x_0\right)}\,\partial_{2}\bar\psi_{12} \\
 &  + e^{-a_4\left(H x_4\right)} \sqrt[6]{\sin \left(6 H x_0\right)}\,\partial_{5}\bar\psi_{12} + e^{a_4\left(H x_4\right)} \sqrt[6]{\sin \left(6 H x_0\right)}\,\partial_{1}\bar\psi_{13} - e^{-a_4\left(H x_4\right)} \sqrt[6]{\sin \left(6 H x_0\right)}\,\partial_{6}\bar\psi_{13} \\
 &  - e^{a_4\left(H x_4\right)} \sqrt[6]{\sin \left(6 H x_0\right)}\,\partial_{3}\bar\psi_{15} - \partial_{4}\bar\psi_{15} - 3 H \cot ^2\left(6 H x_0\right)\,\bar\psi_{10} \\
 & = -H V'(s)\,\bar\psi_{2}
\end{align*}
\begin{align*}
(\bar E_{3})\quad & \cot \left(6 H x_0\right)\,\partial_{0}\bar\psi_{11} - e^{-a_4\left(H x_4\right)} \sqrt[6]{\sin \left(6 H x_0\right)}\,\partial_{7}\bar\psi_{11} + e^{a_4\left(H x_4\right)} \sqrt[6]{\sin \left(6 H x_0\right)}\,\partial_{1}\bar\psi_{12} \\
 &  + e^{-a_4\left(H x_4\right)} \sqrt[6]{\sin \left(6 H x_0\right)}\,\partial_{6}\bar\psi_{12} + e^{a_4\left(H x_4\right)} \sqrt[6]{\sin \left(6 H x_0\right)}\,\partial_{2}\bar\psi_{13} + e^{-a_4\left(H x_4\right)} \sqrt[6]{\sin \left(6 H x_0\right)}\,\partial_{5}\bar\psi_{13} \\
 &  + e^{a_4\left(H x_4\right)} \sqrt[6]{\sin \left(6 H x_0\right)}\,\partial_{3}\bar\psi_{14} + \partial_{4}\bar\psi_{14} - 3 H \cot ^2\left(6 H x_0\right)\,\bar\psi_{11} \\
 & = -H V'(s)\,\bar\psi_{3}
\end{align*}
\begin{align*}
(\bar E_{4})\quad & -e^{a_4\left(H x_4\right)} \sqrt[6]{\sin \left(6 H x_0\right)}\,\partial_{3}\bar\psi_{9} - \partial_{4}\bar\psi_{9} + e^{a_4\left(H x_4\right)} \sqrt[6]{\sin \left(6 H x_0\right)}\,\partial_{2}\bar\psi_{10} \\
 &  + e^{-a_4\left(H x_4\right)} \sqrt[6]{\sin \left(6 H x_0\right)}\,\partial_{5}\bar\psi_{10} - e^{a_4\left(H x_4\right)} \sqrt[6]{\sin \left(6 H x_0\right)}\,\partial_{1}\bar\psi_{11} + e^{-a_4\left(H x_4\right)} \sqrt[6]{\sin \left(6 H x_0\right)}\,\partial_{6}\bar\psi_{11} \\
 &  + \cot \left(6 H x_0\right)\,\partial_{0}\bar\psi_{12} + e^{-a_4\left(H x_4\right)} \sqrt[6]{\sin \left(6 H x_0\right)}\,\partial_{7}\bar\psi_{12} - 3 H \cot ^2\left(6 H x_0\right)\,\bar\psi_{12} \\
 & = -H V'(s)\,\bar\psi_{4}
\end{align*}
\begin{align*}
(\bar E_{5})\quad & e^{a_4\left(H x_4\right)} \sqrt[6]{\sin \left(6 H x_0\right)}\,\partial_{3}\bar\psi_{8} + \partial_{4}\bar\psi_{8} - e^{a_4\left(H x_4\right)} \sqrt[6]{\sin \left(6 H x_0\right)}\,\partial_{1}\bar\psi_{10} \\
 &  - e^{-a_4\left(H x_4\right)} \sqrt[6]{\sin \left(6 H x_0\right)}\,\partial_{6}\bar\psi_{10} - e^{a_4\left(H x_4\right)} \sqrt[6]{\sin \left(6 H x_0\right)}\,\partial_{2}\bar\psi_{11} + e^{-a_4\left(H x_4\right)} \sqrt[6]{\sin \left(6 H x_0\right)}\,\partial_{5}\bar\psi_{11} \\
 &  + \cot \left(6 H x_0\right)\,\partial_{0}\bar\psi_{13} + e^{-a_4\left(H x_4\right)} \sqrt[6]{\sin \left(6 H x_0\right)}\,\partial_{7}\bar\psi_{13} - 3 H \cot ^2\left(6 H x_0\right)\,\bar\psi_{13} \\
 & = -H V'(s)\,\bar\psi_{5}
\end{align*}
\begin{align*}
(\bar E_{6})\quad & -e^{a_4\left(H x_4\right)} \sqrt[6]{\sin \left(6 H x_0\right)}\,\partial_{2}\bar\psi_{8} - e^{-a_4\left(H x_4\right)} \sqrt[6]{\sin \left(6 H x_0\right)}\,\partial_{5}\bar\psi_{8} + e^{a_4\left(H x_4\right)} \sqrt[6]{\sin \left(6 H x_0\right)}\,\partial_{1}\bar\psi_{9} \\
 &  + e^{-a_4\left(H x_4\right)} \sqrt[6]{\sin \left(6 H x_0\right)}\,\partial_{6}\bar\psi_{9} - e^{a_4\left(H x_4\right)} \sqrt[6]{\sin \left(6 H x_0\right)}\,\partial_{3}\bar\psi_{11} + \partial_{4}\bar\psi_{11} \\
 &  + \cot \left(6 H x_0\right)\,\partial_{0}\bar\psi_{14} + e^{-a_4\left(H x_4\right)} \sqrt[6]{\sin \left(6 H x_0\right)}\,\partial_{7}\bar\psi_{14} - 3 H \cot ^2\left(6 H x_0\right)\,\bar\psi_{14} \\
 & = -H V'(s)\,\bar\psi_{6}
\end{align*}
\begin{align*}
(\bar E_{7})\quad & e^{a_4\left(H x_4\right)} \sqrt[6]{\sin \left(6 H x_0\right)}\,\partial_{1}\bar\psi_{8} - e^{-a_4\left(H x_4\right)} \sqrt[6]{\sin \left(6 H x_0\right)}\,\partial_{6}\bar\psi_{8} + e^{a_4\left(H x_4\right)} \sqrt[6]{\sin \left(6 H x_0\right)}\,\partial_{2}\bar\psi_{9} \\
 &  - e^{-a_4\left(H x_4\right)} \sqrt[6]{\sin \left(6 H x_0\right)}\,\partial_{5}\bar\psi_{9} + e^{a_4\left(H x_4\right)} \sqrt[6]{\sin \left(6 H x_0\right)}\,\partial_{3}\bar\psi_{10} - \partial_{4}\bar\psi_{10} \\
 &  + \cot \left(6 H x_0\right)\,\partial_{0}\bar\psi_{15} + e^{-a_4\left(H x_4\right)} \sqrt[6]{\sin \left(6 H x_0\right)}\,\partial_{7}\bar\psi_{15} - 3 H \cot ^2\left(6 H x_0\right)\,\bar\psi_{15} \\
 & = -H V'(s)\,\bar\psi_{7}
\end{align*}
\begin{align*}
(\bar E_{8})\quad & \cot \left(6 H x_0\right)\,\partial_{0}\bar\psi_{0} + e^{-a_4\left(H x_4\right)} \sqrt[6]{\sin \left(6 H x_0\right)}\,\partial_{7}\bar\psi_{0} + e^{a_4\left(H x_4\right)} \sqrt[6]{\sin \left(6 H x_0\right)}\,\partial_{3}\bar\psi_{5} \\
 &  - \partial_{4}\bar\psi_{5} - e^{a_4\left(H x_4\right)} \sqrt[6]{\sin \left(6 H x_0\right)}\,\partial_{2}\bar\psi_{6} + e^{-a_4\left(H x_4\right)} \sqrt[6]{\sin \left(6 H x_0\right)}\,\partial_{5}\bar\psi_{6} \\
 &  + e^{a_4\left(H x_4\right)} \sqrt[6]{\sin \left(6 H x_0\right)}\,\partial_{1}\bar\psi_{7} + e^{-a_4\left(H x_4\right)} \sqrt[6]{\sin \left(6 H x_0\right)}\,\partial_{6}\bar\psi_{7} - 3 H \cot ^2\left(6 H x_0\right)\,\bar\psi_{0} \\
 & = -H V'(s)\,\bar\psi_{8}
\end{align*}
\begin{align*}
(\bar E_{9})\quad & \cot \left(6 H x_0\right)\,\partial_{0}\bar\psi_{1} + e^{-a_4\left(H x_4\right)} \sqrt[6]{\sin \left(6 H x_0\right)}\,\partial_{7}\bar\psi_{1} - e^{a_4\left(H x_4\right)} \sqrt[6]{\sin \left(6 H x_0\right)}\,\partial_{3}\bar\psi_{4} \\
 &  + \partial_{4}\bar\psi_{4} + e^{a_4\left(H x_4\right)} \sqrt[6]{\sin \left(6 H x_0\right)}\,\partial_{1}\bar\psi_{6} - e^{-a_4\left(H x_4\right)} \sqrt[6]{\sin \left(6 H x_0\right)}\,\partial_{6}\bar\psi_{6} \\
 &  + e^{a_4\left(H x_4\right)} \sqrt[6]{\sin \left(6 H x_0\right)}\,\partial_{2}\bar\psi_{7} + e^{-a_4\left(H x_4\right)} \sqrt[6]{\sin \left(6 H x_0\right)}\,\partial_{5}\bar\psi_{7} - 3 H \cot ^2\left(6 H x_0\right)\,\bar\psi_{1} \\
 & = -H V'(s)\,\bar\psi_{9}
\end{align*}
\begin{align*}
(\bar E_{10})\quad & \cot \left(6 H x_0\right)\,\partial_{0}\bar\psi_{2} + e^{-a_4\left(H x_4\right)} \sqrt[6]{\sin \left(6 H x_0\right)}\,\partial_{7}\bar\psi_{2} + e^{a_4\left(H x_4\right)} \sqrt[6]{\sin \left(6 H x_0\right)}\,\partial_{2}\bar\psi_{4} \\
 &  - e^{-a_4\left(H x_4\right)} \sqrt[6]{\sin \left(6 H x_0\right)}\,\partial_{5}\bar\psi_{4} - e^{a_4\left(H x_4\right)} \sqrt[6]{\sin \left(6 H x_0\right)}\,\partial_{1}\bar\psi_{5} + e^{-a_4\left(H x_4\right)} \sqrt[6]{\sin \left(6 H x_0\right)}\,\partial_{6}\bar\psi_{5} \\
 &  + e^{a_4\left(H x_4\right)} \sqrt[6]{\sin \left(6 H x_0\right)}\,\partial_{3}\bar\psi_{7} + \partial_{4}\bar\psi_{7} - 3 H \cot ^2\left(6 H x_0\right)\,\bar\psi_{2} \\
 & = -H V'(s)\,\bar\psi_{10}
\end{align*}
\begin{align*}
(\bar E_{11})\quad & \cot \left(6 H x_0\right)\,\partial_{0}\bar\psi_{3} + e^{-a_4\left(H x_4\right)} \sqrt[6]{\sin \left(6 H x_0\right)}\,\partial_{7}\bar\psi_{3} - e^{a_4\left(H x_4\right)} \sqrt[6]{\sin \left(6 H x_0\right)}\,\partial_{1}\bar\psi_{4} \\
 &  - e^{-a_4\left(H x_4\right)} \sqrt[6]{\sin \left(6 H x_0\right)}\,\partial_{6}\bar\psi_{4} - e^{a_4\left(H x_4\right)} \sqrt[6]{\sin \left(6 H x_0\right)}\,\partial_{2}\bar\psi_{5} - e^{-a_4\left(H x_4\right)} \sqrt[6]{\sin \left(6 H x_0\right)}\,\partial_{5}\bar\psi_{5} \\
 &  - e^{a_4\left(H x_4\right)} \sqrt[6]{\sin \left(6 H x_0\right)}\,\partial_{3}\bar\psi_{6} - \partial_{4}\bar\psi_{6} - 3 H \cot ^2\left(6 H x_0\right)\,\bar\psi_{3} \\
 & = -H V'(s)\,\bar\psi_{11}
\end{align*}
\begin{align*}
(\bar E_{12})\quad & e^{a_4\left(H x_4\right)} \sqrt[6]{\sin \left(6 H x_0\right)}\,\partial_{3}\bar\psi_{1} + \partial_{4}\bar\psi_{1} - e^{a_4\left(H x_4\right)} \sqrt[6]{\sin \left(6 H x_0\right)}\,\partial_{2}\bar\psi_{2} \\
 &  - e^{-a_4\left(H x_4\right)} \sqrt[6]{\sin \left(6 H x_0\right)}\,\partial_{5}\bar\psi_{2} + e^{a_4\left(H x_4\right)} \sqrt[6]{\sin \left(6 H x_0\right)}\,\partial_{1}\bar\psi_{3} - e^{-a_4\left(H x_4\right)} \sqrt[6]{\sin \left(6 H x_0\right)}\,\partial_{6}\bar\psi_{3} \\
 &  + \cot \left(6 H x_0\right)\,\partial_{0}\bar\psi_{4} - e^{-a_4\left(H x_4\right)} \sqrt[6]{\sin \left(6 H x_0\right)}\,\partial_{7}\bar\psi_{4} - 3 H \cot ^2\left(6 H x_0\right)\,\bar\psi_{4} \\
 & = -H V'(s)\,\bar\psi_{12}
\end{align*}
\begin{align*}
(\bar E_{13})\quad & -e^{a_4\left(H x_4\right)} \sqrt[6]{\sin \left(6 H x_0\right)}\,\partial_{3}\bar\psi_{0} - \partial_{4}\bar\psi_{0} + e^{a_4\left(H x_4\right)} \sqrt[6]{\sin \left(6 H x_0\right)}\,\partial_{1}\bar\psi_{2} \\
 &  + e^{-a_4\left(H x_4\right)} \sqrt[6]{\sin \left(6 H x_0\right)}\,\partial_{6}\bar\psi_{2} + e^{a_4\left(H x_4\right)} \sqrt[6]{\sin \left(6 H x_0\right)}\,\partial_{2}\bar\psi_{3} - e^{-a_4\left(H x_4\right)} \sqrt[6]{\sin \left(6 H x_0\right)}\,\partial_{5}\bar\psi_{3} \\
 &  + \cot \left(6 H x_0\right)\,\partial_{0}\bar\psi_{5} - e^{-a_4\left(H x_4\right)} \sqrt[6]{\sin \left(6 H x_0\right)}\,\partial_{7}\bar\psi_{5} - 3 H \cot ^2\left(6 H x_0\right)\,\bar\psi_{5} \\
 & = -H V'(s)\,\bar\psi_{13}
\end{align*}
\begin{align*}
(\bar E_{14})\quad & e^{a_4\left(H x_4\right)} \sqrt[6]{\sin \left(6 H x_0\right)}\,\partial_{2}\bar\psi_{0} + e^{-a_4\left(H x_4\right)} \sqrt[6]{\sin \left(6 H x_0\right)}\,\partial_{5}\bar\psi_{0} - e^{a_4\left(H x_4\right)} \sqrt[6]{\sin \left(6 H x_0\right)}\,\partial_{1}\bar\psi_{1} \\
 &  - e^{-a_4\left(H x_4\right)} \sqrt[6]{\sin \left(6 H x_0\right)}\,\partial_{6}\bar\psi_{1} + e^{a_4\left(H x_4\right)} \sqrt[6]{\sin \left(6 H x_0\right)}\,\partial_{3}\bar\psi_{3} - \partial_{4}\bar\psi_{3} \\
 &  + \cot \left(6 H x_0\right)\,\partial_{0}\bar\psi_{6} - e^{-a_4\left(H x_4\right)} \sqrt[6]{\sin \left(6 H x_0\right)}\,\partial_{7}\bar\psi_{6} - 3 H \cot ^2\left(6 H x_0\right)\,\bar\psi_{6} \\
 & = -H V'(s)\,\bar\psi_{14}
\end{align*}
\begin{align*}
(\bar E_{15})\quad & -e^{a_4\left(H x_4\right)} \sqrt[6]{\sin \left(6 H x_0\right)}\,\partial_{1}\bar\psi_{0} + e^{-a_4\left(H x_4\right)} \sqrt[6]{\sin \left(6 H x_0\right)}\,\partial_{6}\bar\psi_{0} - e^{a_4\left(H x_4\right)} \sqrt[6]{\sin \left(6 H x_0\right)}\,\partial_{2}\bar\psi_{1} \\
 &  + e^{-a_4\left(H x_4\right)} \sqrt[6]{\sin \left(6 H x_0\right)}\,\partial_{5}\bar\psi_{1} - e^{a_4\left(H x_4\right)} \sqrt[6]{\sin \left(6 H x_0\right)}\,\partial_{3}\bar\psi_{2} + \partial_{4}\bar\psi_{2} \\
 &  + \cot \left(6 H x_0\right)\,\partial_{0}\bar\psi_{7} - e^{-a_4\left(H x_4\right)} \sqrt[6]{\sin \left(6 H x_0\right)}\,\partial_{7}\bar\psi_{7} - 3 H \cot ^2\left(6 H x_0\right)\,\bar\psi_{7} \\
 & = -H V'(s)\,\bar\psi_{15}
\end{align*}

}

For every $r$, $E_r$ and $\bar E_r$ contain the same eight (derivative direction, component index)
pairs, $\partial_\mu\psi_k$ against $\partial_\mu\bar\psi_k$, with the same coefficients. The terms
with derivatives along the spacelike directions $x_0..x_3$ have the \textbf{same} sign in both; those
along the timelike directions $x_4..x_7$ have the \textbf{opposite} sign; the connection term
$-3H\cot^2$ multiplies the same component with the same sign; and the $V'$ term changes sign. This is
what the transposition requires: $\Tsix{a}$ is a real signed permutation matrix with
$\Tsix{a}^2=\eta_{aa}$, so $\Tsix{a}^{T}=\eta_{aa}\Tsix{a}$ [derived here; the rule was checked for
this page on all 128 derivative terms and 16 connection terms of the two exported files by a short
script; it is not a notebook assertion]. The adjoint equations are the Dirac conjugates of the field
equations (\cfL{FIELD EQUATION [has content]: the adjoint equation IS the Dirac conjugate of the field equation: with Psi = u + i v, Psibar = (u - i v)\textasciicircum{}T sigma16, adjoint residual == -(field residual)\textasciicircum{}* . sigma16}, \secref{sec:5.1}).

\subsection{The 8+8 form, component by component}
\label{sec:5.6}

Rows 0--7 are the $\bar\tau$ equations for $\psi_2$, rows 8--15 the $\tau$ equations for $\psi_1$;
$\psi^{(1)}_k=\psi_k$ and $\psi^{(2)}_k=\psi_{k+8}$ ($k=0..7$)
(\file{provenance-latex/generated/fermion_fable_field_equations_8plus8.tex}, unchanged):

{\small
% fermion fable: the split-octonion 8+8 form, Psi = (psi^(1), psi^(2))
% generated by claude-fable/export_fermion_fable_tex.wls from the notebook's Part VII objects, 2026-09-24
% notation: x_0..x_7 are the coordinates (x_4 = the observer's time), a_4 = the free function a4 (never given a value),
% q = e^{-a_4(H x_4)}/sin^{1/6}(6 H x_0), p = e^{+a_4(H x_4)}/sin^{1/6}(6 H x_0), s = bar(psi) psi.
\begin{align*}
& \cot \left(6 H x_0\right)\,\partial_0\psi^{(2)} + e^{a_4\left(H x_4\right)} \sqrt[6]{\sin \left(6 H x_0\right)}\sum_{i=1}^{3}\bar\tau_i\partial_i\psi^{(2)} + \bar\tau_4\partial_4\psi^{(2)} \\
& \quad + e^{-a_4\left(H x_4\right)} \sqrt[6]{\sin \left(6 H x_0\right)}\sum_{h=5}^{7}\bar\tau_h\partial_h\psi^{(2)} - 3H\cot ^2\left(6 H x_0\right)\,\psi^{(2)} = H V'(s)\,\psi^{(1)}, \\[4pt]
& \cot \left(6 H x_0\right)\,\partial_0\psi^{(1)} + e^{a_4\left(H x_4\right)} \sqrt[6]{\sin \left(6 H x_0\right)}\sum_{i=1}^{3}\tau_i\partial_i\psi^{(1)} + \tau_4\partial_4\psi^{(1)} \\
& \quad + e^{-a_4\left(H x_4\right)} \sqrt[6]{\sin \left(6 H x_0\right)}\sum_{h=5}^{7}\tau_h\partial_h\psi^{(1)} - 3H\cot ^2\left(6 H x_0\right)\,\psi^{(1)} = H V'(s)\,\psi^{(2)}, \\[4pt]
& s = -\psi^{(1)\ddagger}\sigma\,\psi^{(1)} + \psi^{(2)\ddagger}\sigma\,\psi^{(2)} .
\end{align*}
Component form (rows 0--7: the $\bar\tau$ equations for $\psi^{(2)}$; rows 8--15: the $\tau$ equations for $\psi^{(1)}$):
\begin{align*}
(E_{0})\quad & \cot \left(6 H x_0\right)\,\partial_{0}\psi^{(2)}_{0} + e^{-a_4\left(H x_4\right)} \sqrt[6]{\sin \left(6 H x_0\right)}\,\partial_{7}\psi^{(2)}_{0} - e^{a_4\left(H x_4\right)} \sqrt[6]{\sin \left(6 H x_0\right)}\,\partial_{3}\psi^{(2)}_{5} \\
 &  - \partial_{4}\psi^{(2)}_{5} + e^{a_4\left(H x_4\right)} \sqrt[6]{\sin \left(6 H x_0\right)}\,\partial_{2}\psi^{(2)}_{6} + e^{-a_4\left(H x_4\right)} \sqrt[6]{\sin \left(6 H x_0\right)}\,\partial_{5}\psi^{(2)}_{6} \\
 &  - e^{a_4\left(H x_4\right)} \sqrt[6]{\sin \left(6 H x_0\right)}\,\partial_{1}\psi^{(2)}_{7} + e^{-a_4\left(H x_4\right)} \sqrt[6]{\sin \left(6 H x_0\right)}\,\partial_{6}\psi^{(2)}_{7} - 3 H \cot ^2\left(6 H x_0\right)\,\psi^{(2)}_{0} \\
 & = H V'(s)\,\psi^{(1)}_{0}
\end{align*}
\begin{align*}
(E_{1})\quad & \cot \left(6 H x_0\right)\,\partial_{0}\psi^{(2)}_{1} + e^{-a_4\left(H x_4\right)} \sqrt[6]{\sin \left(6 H x_0\right)}\,\partial_{7}\psi^{(2)}_{1} + e^{a_4\left(H x_4\right)} \sqrt[6]{\sin \left(6 H x_0\right)}\,\partial_{3}\psi^{(2)}_{4} \\
 &  + \partial_{4}\psi^{(2)}_{4} - e^{a_4\left(H x_4\right)} \sqrt[6]{\sin \left(6 H x_0\right)}\,\partial_{1}\psi^{(2)}_{6} - e^{-a_4\left(H x_4\right)} \sqrt[6]{\sin \left(6 H x_0\right)}\,\partial_{6}\psi^{(2)}_{6} \\
 &  - e^{a_4\left(H x_4\right)} \sqrt[6]{\sin \left(6 H x_0\right)}\,\partial_{2}\psi^{(2)}_{7} + e^{-a_4\left(H x_4\right)} \sqrt[6]{\sin \left(6 H x_0\right)}\,\partial_{5}\psi^{(2)}_{7} - 3 H \cot ^2\left(6 H x_0\right)\,\psi^{(2)}_{1} \\
 & = H V'(s)\,\psi^{(1)}_{1}
\end{align*}
\begin{align*}
(E_{2})\quad & \cot \left(6 H x_0\right)\,\partial_{0}\psi^{(2)}_{2} + e^{-a_4\left(H x_4\right)} \sqrt[6]{\sin \left(6 H x_0\right)}\,\partial_{7}\psi^{(2)}_{2} - e^{a_4\left(H x_4\right)} \sqrt[6]{\sin \left(6 H x_0\right)}\,\partial_{2}\psi^{(2)}_{4} \\
 &  - e^{-a_4\left(H x_4\right)} \sqrt[6]{\sin \left(6 H x_0\right)}\,\partial_{5}\psi^{(2)}_{4} + e^{a_4\left(H x_4\right)} \sqrt[6]{\sin \left(6 H x_0\right)}\,\partial_{1}\psi^{(2)}_{5} + e^{-a_4\left(H x_4\right)} \sqrt[6]{\sin \left(6 H x_0\right)}\,\partial_{6}\psi^{(2)}_{5} \\
 &  - e^{a_4\left(H x_4\right)} \sqrt[6]{\sin \left(6 H x_0\right)}\,\partial_{3}\psi^{(2)}_{7} + \partial_{4}\psi^{(2)}_{7} - 3 H \cot ^2\left(6 H x_0\right)\,\psi^{(2)}_{2} \\
 & = H V'(s)\,\psi^{(1)}_{2}
\end{align*}
\begin{align*}
(E_{3})\quad & \cot \left(6 H x_0\right)\,\partial_{0}\psi^{(2)}_{3} + e^{-a_4\left(H x_4\right)} \sqrt[6]{\sin \left(6 H x_0\right)}\,\partial_{7}\psi^{(2)}_{3} + e^{a_4\left(H x_4\right)} \sqrt[6]{\sin \left(6 H x_0\right)}\,\partial_{1}\psi^{(2)}_{4} \\
 &  - e^{-a_4\left(H x_4\right)} \sqrt[6]{\sin \left(6 H x_0\right)}\,\partial_{6}\psi^{(2)}_{4} + e^{a_4\left(H x_4\right)} \sqrt[6]{\sin \left(6 H x_0\right)}\,\partial_{2}\psi^{(2)}_{5} - e^{-a_4\left(H x_4\right)} \sqrt[6]{\sin \left(6 H x_0\right)}\,\partial_{5}\psi^{(2)}_{5} \\
 &  + e^{a_4\left(H x_4\right)} \sqrt[6]{\sin \left(6 H x_0\right)}\,\partial_{3}\psi^{(2)}_{6} - \partial_{4}\psi^{(2)}_{6} - 3 H \cot ^2\left(6 H x_0\right)\,\psi^{(2)}_{3} \\
 & = H V'(s)\,\psi^{(1)}_{3}
\end{align*}
\begin{align*}
(E_{4})\quad & -e^{a_4\left(H x_4\right)} \sqrt[6]{\sin \left(6 H x_0\right)}\,\partial_{3}\psi^{(2)}_{1} + \partial_{4}\psi^{(2)}_{1} + e^{a_4\left(H x_4\right)} \sqrt[6]{\sin \left(6 H x_0\right)}\,\partial_{2}\psi^{(2)}_{2} \\
 &  - e^{-a_4\left(H x_4\right)} \sqrt[6]{\sin \left(6 H x_0\right)}\,\partial_{5}\psi^{(2)}_{2} - e^{a_4\left(H x_4\right)} \sqrt[6]{\sin \left(6 H x_0\right)}\,\partial_{1}\psi^{(2)}_{3} - e^{-a_4\left(H x_4\right)} \sqrt[6]{\sin \left(6 H x_0\right)}\,\partial_{6}\psi^{(2)}_{3} \\
 &  + \cot \left(6 H x_0\right)\,\partial_{0}\psi^{(2)}_{4} - e^{-a_4\left(H x_4\right)} \sqrt[6]{\sin \left(6 H x_0\right)}\,\partial_{7}\psi^{(2)}_{4} - 3 H \cot ^2\left(6 H x_0\right)\,\psi^{(2)}_{4} \\
 & = H V'(s)\,\psi^{(1)}_{4}
\end{align*}
\begin{align*}
(E_{5})\quad & e^{a_4\left(H x_4\right)} \sqrt[6]{\sin \left(6 H x_0\right)}\,\partial_{3}\psi^{(2)}_{0} - \partial_{4}\psi^{(2)}_{0} - e^{a_4\left(H x_4\right)} \sqrt[6]{\sin \left(6 H x_0\right)}\,\partial_{1}\psi^{(2)}_{2} \\
 &  + e^{-a_4\left(H x_4\right)} \sqrt[6]{\sin \left(6 H x_0\right)}\,\partial_{6}\psi^{(2)}_{2} - e^{a_4\left(H x_4\right)} \sqrt[6]{\sin \left(6 H x_0\right)}\,\partial_{2}\psi^{(2)}_{3} - e^{-a_4\left(H x_4\right)} \sqrt[6]{\sin \left(6 H x_0\right)}\,\partial_{5}\psi^{(2)}_{3} \\
 &  + \cot \left(6 H x_0\right)\,\partial_{0}\psi^{(2)}_{5} - e^{-a_4\left(H x_4\right)} \sqrt[6]{\sin \left(6 H x_0\right)}\,\partial_{7}\psi^{(2)}_{5} - 3 H \cot ^2\left(6 H x_0\right)\,\psi^{(2)}_{5} \\
 & = H V'(s)\,\psi^{(1)}_{5}
\end{align*}
\begin{align*}
(E_{6})\quad & -e^{a_4\left(H x_4\right)} \sqrt[6]{\sin \left(6 H x_0\right)}\,\partial_{2}\psi^{(2)}_{0} + e^{-a_4\left(H x_4\right)} \sqrt[6]{\sin \left(6 H x_0\right)}\,\partial_{5}\psi^{(2)}_{0} + e^{a_4\left(H x_4\right)} \sqrt[6]{\sin \left(6 H x_0\right)}\,\partial_{1}\psi^{(2)}_{1} \\
 &  - e^{-a_4\left(H x_4\right)} \sqrt[6]{\sin \left(6 H x_0\right)}\,\partial_{6}\psi^{(2)}_{1} - e^{a_4\left(H x_4\right)} \sqrt[6]{\sin \left(6 H x_0\right)}\,\partial_{3}\psi^{(2)}_{3} - \partial_{4}\psi^{(2)}_{3} \\
 &  + \cot \left(6 H x_0\right)\,\partial_{0}\psi^{(2)}_{6} - e^{-a_4\left(H x_4\right)} \sqrt[6]{\sin \left(6 H x_0\right)}\,\partial_{7}\psi^{(2)}_{6} - 3 H \cot ^2\left(6 H x_0\right)\,\psi^{(2)}_{6} \\
 & = H V'(s)\,\psi^{(1)}_{6}
\end{align*}
\begin{align*}
(E_{7})\quad & e^{a_4\left(H x_4\right)} \sqrt[6]{\sin \left(6 H x_0\right)}\,\partial_{1}\psi^{(2)}_{0} + e^{-a_4\left(H x_4\right)} \sqrt[6]{\sin \left(6 H x_0\right)}\,\partial_{6}\psi^{(2)}_{0} + e^{a_4\left(H x_4\right)} \sqrt[6]{\sin \left(6 H x_0\right)}\,\partial_{2}\psi^{(2)}_{1} \\
 &  + e^{-a_4\left(H x_4\right)} \sqrt[6]{\sin \left(6 H x_0\right)}\,\partial_{5}\psi^{(2)}_{1} + e^{a_4\left(H x_4\right)} \sqrt[6]{\sin \left(6 H x_0\right)}\,\partial_{3}\psi^{(2)}_{2} + \partial_{4}\psi^{(2)}_{2} \\
 &  + \cot \left(6 H x_0\right)\,\partial_{0}\psi^{(2)}_{7} - e^{-a_4\left(H x_4\right)} \sqrt[6]{\sin \left(6 H x_0\right)}\,\partial_{7}\psi^{(2)}_{7} - 3 H \cot ^2\left(6 H x_0\right)\,\psi^{(2)}_{7} \\
 & = H V'(s)\,\psi^{(1)}_{7}
\end{align*}
\begin{align*}
(E_{8})\quad & \cot \left(6 H x_0\right)\,\partial_{0}\psi^{(1)}_{0} - e^{-a_4\left(H x_4\right)} \sqrt[6]{\sin \left(6 H x_0\right)}\,\partial_{7}\psi^{(1)}_{0} + e^{a_4\left(H x_4\right)} \sqrt[6]{\sin \left(6 H x_0\right)}\,\partial_{3}\psi^{(1)}_{5} \\
 &  + \partial_{4}\psi^{(1)}_{5} - e^{a_4\left(H x_4\right)} \sqrt[6]{\sin \left(6 H x_0\right)}\,\partial_{2}\psi^{(1)}_{6} - e^{-a_4\left(H x_4\right)} \sqrt[6]{\sin \left(6 H x_0\right)}\,\partial_{5}\psi^{(1)}_{6} \\
 &  + e^{a_4\left(H x_4\right)} \sqrt[6]{\sin \left(6 H x_0\right)}\,\partial_{1}\psi^{(1)}_{7} - e^{-a_4\left(H x_4\right)} \sqrt[6]{\sin \left(6 H x_0\right)}\,\partial_{6}\psi^{(1)}_{7} - 3 H \cot ^2\left(6 H x_0\right)\,\psi^{(1)}_{0} \\
 & = H V'(s)\,\psi^{(2)}_{0}
\end{align*}
\begin{align*}
(E_{9})\quad & \cot \left(6 H x_0\right)\,\partial_{0}\psi^{(1)}_{1} - e^{-a_4\left(H x_4\right)} \sqrt[6]{\sin \left(6 H x_0\right)}\,\partial_{7}\psi^{(1)}_{1} - e^{a_4\left(H x_4\right)} \sqrt[6]{\sin \left(6 H x_0\right)}\,\partial_{3}\psi^{(1)}_{4} \\
 &  - \partial_{4}\psi^{(1)}_{4} + e^{a_4\left(H x_4\right)} \sqrt[6]{\sin \left(6 H x_0\right)}\,\partial_{1}\psi^{(1)}_{6} + e^{-a_4\left(H x_4\right)} \sqrt[6]{\sin \left(6 H x_0\right)}\,\partial_{6}\psi^{(1)}_{6} \\
 &  + e^{a_4\left(H x_4\right)} \sqrt[6]{\sin \left(6 H x_0\right)}\,\partial_{2}\psi^{(1)}_{7} - e^{-a_4\left(H x_4\right)} \sqrt[6]{\sin \left(6 H x_0\right)}\,\partial_{5}\psi^{(1)}_{7} - 3 H \cot ^2\left(6 H x_0\right)\,\psi^{(1)}_{1} \\
 & = H V'(s)\,\psi^{(2)}_{1}
\end{align*}
\begin{align*}
(E_{10})\quad & \cot \left(6 H x_0\right)\,\partial_{0}\psi^{(1)}_{2} - e^{-a_4\left(H x_4\right)} \sqrt[6]{\sin \left(6 H x_0\right)}\,\partial_{7}\psi^{(1)}_{2} + e^{a_4\left(H x_4\right)} \sqrt[6]{\sin \left(6 H x_0\right)}\,\partial_{2}\psi^{(1)}_{4} \\
 &  + e^{-a_4\left(H x_4\right)} \sqrt[6]{\sin \left(6 H x_0\right)}\,\partial_{5}\psi^{(1)}_{4} - e^{a_4\left(H x_4\right)} \sqrt[6]{\sin \left(6 H x_0\right)}\,\partial_{1}\psi^{(1)}_{5} - e^{-a_4\left(H x_4\right)} \sqrt[6]{\sin \left(6 H x_0\right)}\,\partial_{6}\psi^{(1)}_{5} \\
 &  + e^{a_4\left(H x_4\right)} \sqrt[6]{\sin \left(6 H x_0\right)}\,\partial_{3}\psi^{(1)}_{7} - \partial_{4}\psi^{(1)}_{7} - 3 H \cot ^2\left(6 H x_0\right)\,\psi^{(1)}_{2} \\
 & = H V'(s)\,\psi^{(2)}_{2}
\end{align*}
\begin{align*}
(E_{11})\quad & \cot \left(6 H x_0\right)\,\partial_{0}\psi^{(1)}_{3} - e^{-a_4\left(H x_4\right)} \sqrt[6]{\sin \left(6 H x_0\right)}\,\partial_{7}\psi^{(1)}_{3} - e^{a_4\left(H x_4\right)} \sqrt[6]{\sin \left(6 H x_0\right)}\,\partial_{1}\psi^{(1)}_{4} \\
 &  + e^{-a_4\left(H x_4\right)} \sqrt[6]{\sin \left(6 H x_0\right)}\,\partial_{6}\psi^{(1)}_{4} - e^{a_4\left(H x_4\right)} \sqrt[6]{\sin \left(6 H x_0\right)}\,\partial_{2}\psi^{(1)}_{5} + e^{-a_4\left(H x_4\right)} \sqrt[6]{\sin \left(6 H x_0\right)}\,\partial_{5}\psi^{(1)}_{5} \\
 &  - e^{a_4\left(H x_4\right)} \sqrt[6]{\sin \left(6 H x_0\right)}\,\partial_{3}\psi^{(1)}_{6} + \partial_{4}\psi^{(1)}_{6} - 3 H \cot ^2\left(6 H x_0\right)\,\psi^{(1)}_{3} \\
 & = H V'(s)\,\psi^{(2)}_{3}
\end{align*}
\begin{align*}
(E_{12})\quad & e^{a_4\left(H x_4\right)} \sqrt[6]{\sin \left(6 H x_0\right)}\,\partial_{3}\psi^{(1)}_{1} - \partial_{4}\psi^{(1)}_{1} - e^{a_4\left(H x_4\right)} \sqrt[6]{\sin \left(6 H x_0\right)}\,\partial_{2}\psi^{(1)}_{2} \\
 &  + e^{-a_4\left(H x_4\right)} \sqrt[6]{\sin \left(6 H x_0\right)}\,\partial_{5}\psi^{(1)}_{2} + e^{a_4\left(H x_4\right)} \sqrt[6]{\sin \left(6 H x_0\right)}\,\partial_{1}\psi^{(1)}_{3} + e^{-a_4\left(H x_4\right)} \sqrt[6]{\sin \left(6 H x_0\right)}\,\partial_{6}\psi^{(1)}_{3} \\
 &  + \cot \left(6 H x_0\right)\,\partial_{0}\psi^{(1)}_{4} + e^{-a_4\left(H x_4\right)} \sqrt[6]{\sin \left(6 H x_0\right)}\,\partial_{7}\psi^{(1)}_{4} - 3 H \cot ^2\left(6 H x_0\right)\,\psi^{(1)}_{4} \\
 & = H V'(s)\,\psi^{(2)}_{4}
\end{align*}
\begin{align*}
(E_{13})\quad & -e^{a_4\left(H x_4\right)} \sqrt[6]{\sin \left(6 H x_0\right)}\,\partial_{3}\psi^{(1)}_{0} + \partial_{4}\psi^{(1)}_{0} + e^{a_4\left(H x_4\right)} \sqrt[6]{\sin \left(6 H x_0\right)}\,\partial_{1}\psi^{(1)}_{2} \\
 &  - e^{-a_4\left(H x_4\right)} \sqrt[6]{\sin \left(6 H x_0\right)}\,\partial_{6}\psi^{(1)}_{2} + e^{a_4\left(H x_4\right)} \sqrt[6]{\sin \left(6 H x_0\right)}\,\partial_{2}\psi^{(1)}_{3} + e^{-a_4\left(H x_4\right)} \sqrt[6]{\sin \left(6 H x_0\right)}\,\partial_{5}\psi^{(1)}_{3} \\
 &  + \cot \left(6 H x_0\right)\,\partial_{0}\psi^{(1)}_{5} + e^{-a_4\left(H x_4\right)} \sqrt[6]{\sin \left(6 H x_0\right)}\,\partial_{7}\psi^{(1)}_{5} - 3 H \cot ^2\left(6 H x_0\right)\,\psi^{(1)}_{5} \\
 & = H V'(s)\,\psi^{(2)}_{5}
\end{align*}
\begin{align*}
(E_{14})\quad & e^{a_4\left(H x_4\right)} \sqrt[6]{\sin \left(6 H x_0\right)}\,\partial_{2}\psi^{(1)}_{0} - e^{-a_4\left(H x_4\right)} \sqrt[6]{\sin \left(6 H x_0\right)}\,\partial_{5}\psi^{(1)}_{0} - e^{a_4\left(H x_4\right)} \sqrt[6]{\sin \left(6 H x_0\right)}\,\partial_{1}\psi^{(1)}_{1} \\
 &  + e^{-a_4\left(H x_4\right)} \sqrt[6]{\sin \left(6 H x_0\right)}\,\partial_{6}\psi^{(1)}_{1} + e^{a_4\left(H x_4\right)} \sqrt[6]{\sin \left(6 H x_0\right)}\,\partial_{3}\psi^{(1)}_{3} + \partial_{4}\psi^{(1)}_{3} \\
 &  + \cot \left(6 H x_0\right)\,\partial_{0}\psi^{(1)}_{6} + e^{-a_4\left(H x_4\right)} \sqrt[6]{\sin \left(6 H x_0\right)}\,\partial_{7}\psi^{(1)}_{6} - 3 H \cot ^2\left(6 H x_0\right)\,\psi^{(1)}_{6} \\
 & = H V'(s)\,\psi^{(2)}_{6}
\end{align*}
\begin{align*}
(E_{15})\quad & -e^{a_4\left(H x_4\right)} \sqrt[6]{\sin \left(6 H x_0\right)}\,\partial_{1}\psi^{(1)}_{0} - e^{-a_4\left(H x_4\right)} \sqrt[6]{\sin \left(6 H x_0\right)}\,\partial_{6}\psi^{(1)}_{0} - e^{a_4\left(H x_4\right)} \sqrt[6]{\sin \left(6 H x_0\right)}\,\partial_{2}\psi^{(1)}_{1} \\
 &  - e^{-a_4\left(H x_4\right)} \sqrt[6]{\sin \left(6 H x_0\right)}\,\partial_{5}\psi^{(1)}_{1} - e^{a_4\left(H x_4\right)} \sqrt[6]{\sin \left(6 H x_0\right)}\,\partial_{3}\psi^{(1)}_{2} - \partial_{4}\psi^{(1)}_{2} \\
 &  + \cot \left(6 H x_0\right)\,\partial_{0}\psi^{(1)}_{7} + e^{-a_4\left(H x_4\right)} \sqrt[6]{\sin \left(6 H x_0\right)}\,\partial_{7}\psi^{(1)}_{7} - 3 H \cot ^2\left(6 H x_0\right)\,\psi^{(1)}_{7} \\
 & = H V'(s)\,\psi^{(2)}_{7}
\end{align*}

}

\subsection{The rescaled form: the spin connection disappears}
\label{sec:5.7}

Put $\Psi=\sqrt{\sin 6Hx_0}\;\Psi'$ and $\bar\Psi=\sqrt{\sin 6Hx_0}\;\bar\Psi'$. Then
$s=\sin(6Hx_0)\,s'$ with $s'=\bar\Psi'\Psi'$, and
\begin{multline*}
  \gamma^\mu D_\mu\Psi-H\,V'(s)\,\Psi=\sqrt{\sin 6Hx_0}\;\Bigl(\cot(6Hx_0)\,\Tsix{0}\,\partial_0\Psi'
  +\frac1q\sum_i\Tsix{i}\,\partial_i\Psi'+\Tsix{4}\,\partial_4\Psi'\\
  +\frac1p\sum_h\Tsix{h}\,\partial_h\Psi'-H\,V'\bigl(\sin(6Hx_0)\,s'\bigr)\,\Psi'\Bigr)
\end{multline*}
with \textbf{no connection term at all}, and the same for the adjoint equation. This works because
the connection term is a gradient: $\gamma^\mu\Gamma_\mu=-\tfrac12\,T_\mu\gamma^\mu$ with the
Weitzenb\"ock torsion vector $T_\mu=\partial_\mu\ln\sin(6Hx_0)$ (\secref{sec:9.10}), and multiplying by
$\exp\bigl(\tfrac12\ln\sin\bigr)=\sqrt{\sin}$ removes it. \cfproved{VII \S27}
\begin{itemize}
  \item \cfL{WRITTEN OUT (iv) [THE RESULT]: Psi = Sqrt[Sin[6Hx0]] Psi' turns the field equation into Sqrt[Sin] (Cot T16[0] d\_0 Psi' + (1/q) T16[i] d\_i Psi' + T16[4] d\_4 Psi' + (1/p) T16[h] d\_h Psi' - H V'(Sin s') Psi'): NO connection term}
  \item \cfL{WRITTEN OUT (iv) [THE RESULT]: and the adjoint equation into Sqrt[Sin] (Cot d\_0 Psi'bar T16[0] + ... + H V'(Sin s') Psi'bar), with no connection term either}
  \item \cfL{WRITTEN OUT (iv) [definition]: under the rescaling s == Sin[6 H x0] s',  s' = Psi'bar Psi'}
\end{itemize}

The sixteen rescaled equations ($\psi'_k$ = component $k$ of $\Psi'$; $V'$ evaluated at
$s=\sin(6Hx_0)\,\bar\Psi'\Psi'$;
\file{provenance-latex/generated/fermion_fable_field_equations_rescaled.tex}, unchanged):

{\small
% fermion fable: the rescaled field equations, Psi = Sqrt[Sin[6 H x0]] Psi' (the spin connection disappears); V' is evaluated at s = Sin[6 H x0] s'
% generated by claude-fable/export_fermion_fable_tex.wls from the notebook's Part VII objects, 2026-09-24
% notation: x_0..x_7 are the coordinates (x_4 = the observer's time), a_4 = the free function a4 (never given a value),
% q = e^{-a_4(H x_4)}/sin^{1/6}(6 H x_0), p = e^{+a_4(H x_4)}/sin^{1/6}(6 H x_0), s = bar(psi) psi.
\begin{align*}
(E'_{0})\quad & \cot \left(6 H x_0\right)\,\partial_{0}\psi'_{8} + e^{-a_4\left(H x_4\right)} \sqrt[6]{\sin \left(6 H x_0\right)}\,\partial_{7}\psi'_{8} - e^{a_4\left(H x_4\right)} \sqrt[6]{\sin \left(6 H x_0\right)}\,\partial_{3}\psi'_{13} \\
 &  - \partial_{4}\psi'_{13} + e^{a_4\left(H x_4\right)} \sqrt[6]{\sin \left(6 H x_0\right)}\,\partial_{2}\psi'_{14} + e^{-a_4\left(H x_4\right)} \sqrt[6]{\sin \left(6 H x_0\right)}\,\partial_{5}\psi'_{14} \\
 &  - e^{a_4\left(H x_4\right)} \sqrt[6]{\sin \left(6 H x_0\right)}\,\partial_{1}\psi'_{15} + e^{-a_4\left(H x_4\right)} \sqrt[6]{\sin \left(6 H x_0\right)}\,\partial_{6}\psi'_{15} \\
 & = H V'(s)\,\psi'_{0}
\end{align*}
\begin{align*}
(E'_{1})\quad & \cot \left(6 H x_0\right)\,\partial_{0}\psi'_{9} + e^{-a_4\left(H x_4\right)} \sqrt[6]{\sin \left(6 H x_0\right)}\,\partial_{7}\psi'_{9} + e^{a_4\left(H x_4\right)} \sqrt[6]{\sin \left(6 H x_0\right)}\,\partial_{3}\psi'_{12} \\
 &  + \partial_{4}\psi'_{12} - e^{a_4\left(H x_4\right)} \sqrt[6]{\sin \left(6 H x_0\right)}\,\partial_{1}\psi'_{14} - e^{-a_4\left(H x_4\right)} \sqrt[6]{\sin \left(6 H x_0\right)}\,\partial_{6}\psi'_{14} \\
 &  - e^{a_4\left(H x_4\right)} \sqrt[6]{\sin \left(6 H x_0\right)}\,\partial_{2}\psi'_{15} + e^{-a_4\left(H x_4\right)} \sqrt[6]{\sin \left(6 H x_0\right)}\,\partial_{5}\psi'_{15} \\
 & = H V'(s)\,\psi'_{1}
\end{align*}
\begin{align*}
(E'_{2})\quad & \cot \left(6 H x_0\right)\,\partial_{0}\psi'_{10} + e^{-a_4\left(H x_4\right)} \sqrt[6]{\sin \left(6 H x_0\right)}\,\partial_{7}\psi'_{10} - e^{a_4\left(H x_4\right)} \sqrt[6]{\sin \left(6 H x_0\right)}\,\partial_{2}\psi'_{12} \\
 &  - e^{-a_4\left(H x_4\right)} \sqrt[6]{\sin \left(6 H x_0\right)}\,\partial_{5}\psi'_{12} + e^{a_4\left(H x_4\right)} \sqrt[6]{\sin \left(6 H x_0\right)}\,\partial_{1}\psi'_{13} + e^{-a_4\left(H x_4\right)} \sqrt[6]{\sin \left(6 H x_0\right)}\,\partial_{6}\psi'_{13} \\
 &  - e^{a_4\left(H x_4\right)} \sqrt[6]{\sin \left(6 H x_0\right)}\,\partial_{3}\psi'_{15} + \partial_{4}\psi'_{15} \\
 & = H V'(s)\,\psi'_{2}
\end{align*}
\begin{align*}
(E'_{3})\quad & \cot \left(6 H x_0\right)\,\partial_{0}\psi'_{11} + e^{-a_4\left(H x_4\right)} \sqrt[6]{\sin \left(6 H x_0\right)}\,\partial_{7}\psi'_{11} + e^{a_4\left(H x_4\right)} \sqrt[6]{\sin \left(6 H x_0\right)}\,\partial_{1}\psi'_{12} \\
 &  - e^{-a_4\left(H x_4\right)} \sqrt[6]{\sin \left(6 H x_0\right)}\,\partial_{6}\psi'_{12} + e^{a_4\left(H x_4\right)} \sqrt[6]{\sin \left(6 H x_0\right)}\,\partial_{2}\psi'_{13} - e^{-a_4\left(H x_4\right)} \sqrt[6]{\sin \left(6 H x_0\right)}\,\partial_{5}\psi'_{13} \\
 &  + e^{a_4\left(H x_4\right)} \sqrt[6]{\sin \left(6 H x_0\right)}\,\partial_{3}\psi'_{14} - \partial_{4}\psi'_{14} \\
 & = H V'(s)\,\psi'_{3}
\end{align*}
\begin{align*}
(E'_{4})\quad & -e^{a_4\left(H x_4\right)} \sqrt[6]{\sin \left(6 H x_0\right)}\,\partial_{3}\psi'_{9} + \partial_{4}\psi'_{9} + e^{a_4\left(H x_4\right)} \sqrt[6]{\sin \left(6 H x_0\right)}\,\partial_{2}\psi'_{10} \\
 &  - e^{-a_4\left(H x_4\right)} \sqrt[6]{\sin \left(6 H x_0\right)}\,\partial_{5}\psi'_{10} - e^{a_4\left(H x_4\right)} \sqrt[6]{\sin \left(6 H x_0\right)}\,\partial_{1}\psi'_{11} - e^{-a_4\left(H x_4\right)} \sqrt[6]{\sin \left(6 H x_0\right)}\,\partial_{6}\psi'_{11} \\
 &  + \cot \left(6 H x_0\right)\,\partial_{0}\psi'_{12} - e^{-a_4\left(H x_4\right)} \sqrt[6]{\sin \left(6 H x_0\right)}\,\partial_{7}\psi'_{12} \\
 & = H V'(s)\,\psi'_{4}
\end{align*}
\begin{align*}
(E'_{5})\quad & e^{a_4\left(H x_4\right)} \sqrt[6]{\sin \left(6 H x_0\right)}\,\partial_{3}\psi'_{8} - \partial_{4}\psi'_{8} - e^{a_4\left(H x_4\right)} \sqrt[6]{\sin \left(6 H x_0\right)}\,\partial_{1}\psi'_{10} \\
 &  + e^{-a_4\left(H x_4\right)} \sqrt[6]{\sin \left(6 H x_0\right)}\,\partial_{6}\psi'_{10} - e^{a_4\left(H x_4\right)} \sqrt[6]{\sin \left(6 H x_0\right)}\,\partial_{2}\psi'_{11} - e^{-a_4\left(H x_4\right)} \sqrt[6]{\sin \left(6 H x_0\right)}\,\partial_{5}\psi'_{11} \\
 &  + \cot \left(6 H x_0\right)\,\partial_{0}\psi'_{13} - e^{-a_4\left(H x_4\right)} \sqrt[6]{\sin \left(6 H x_0\right)}\,\partial_{7}\psi'_{13} \\
 & = H V'(s)\,\psi'_{5}
\end{align*}
\begin{align*}
(E'_{6})\quad & -e^{a_4\left(H x_4\right)} \sqrt[6]{\sin \left(6 H x_0\right)}\,\partial_{2}\psi'_{8} + e^{-a_4\left(H x_4\right)} \sqrt[6]{\sin \left(6 H x_0\right)}\,\partial_{5}\psi'_{8} + e^{a_4\left(H x_4\right)} \sqrt[6]{\sin \left(6 H x_0\right)}\,\partial_{1}\psi'_{9} \\
 &  - e^{-a_4\left(H x_4\right)} \sqrt[6]{\sin \left(6 H x_0\right)}\,\partial_{6}\psi'_{9} - e^{a_4\left(H x_4\right)} \sqrt[6]{\sin \left(6 H x_0\right)}\,\partial_{3}\psi'_{11} - \partial_{4}\psi'_{11} \\
 &  + \cot \left(6 H x_0\right)\,\partial_{0}\psi'_{14} - e^{-a_4\left(H x_4\right)} \sqrt[6]{\sin \left(6 H x_0\right)}\,\partial_{7}\psi'_{14} \\
 & = H V'(s)\,\psi'_{6}
\end{align*}
\begin{align*}
(E'_{7})\quad & e^{a_4\left(H x_4\right)} \sqrt[6]{\sin \left(6 H x_0\right)}\,\partial_{1}\psi'_{8} + e^{-a_4\left(H x_4\right)} \sqrt[6]{\sin \left(6 H x_0\right)}\,\partial_{6}\psi'_{8} + e^{a_4\left(H x_4\right)} \sqrt[6]{\sin \left(6 H x_0\right)}\,\partial_{2}\psi'_{9} \\
 &  + e^{-a_4\left(H x_4\right)} \sqrt[6]{\sin \left(6 H x_0\right)}\,\partial_{5}\psi'_{9} + e^{a_4\left(H x_4\right)} \sqrt[6]{\sin \left(6 H x_0\right)}\,\partial_{3}\psi'_{10} + \partial_{4}\psi'_{10} \\
 &  + \cot \left(6 H x_0\right)\,\partial_{0}\psi'_{15} - e^{-a_4\left(H x_4\right)} \sqrt[6]{\sin \left(6 H x_0\right)}\,\partial_{7}\psi'_{15} \\
 & = H V'(s)\,\psi'_{7}
\end{align*}
\begin{align*}
(E'_{8})\quad & \cot \left(6 H x_0\right)\,\partial_{0}\psi'_{0} - e^{-a_4\left(H x_4\right)} \sqrt[6]{\sin \left(6 H x_0\right)}\,\partial_{7}\psi'_{0} + e^{a_4\left(H x_4\right)} \sqrt[6]{\sin \left(6 H x_0\right)}\,\partial_{3}\psi'_{5} \\
 &  + \partial_{4}\psi'_{5} - e^{a_4\left(H x_4\right)} \sqrt[6]{\sin \left(6 H x_0\right)}\,\partial_{2}\psi'_{6} - e^{-a_4\left(H x_4\right)} \sqrt[6]{\sin \left(6 H x_0\right)}\,\partial_{5}\psi'_{6} \\
 &  + e^{a_4\left(H x_4\right)} \sqrt[6]{\sin \left(6 H x_0\right)}\,\partial_{1}\psi'_{7} - e^{-a_4\left(H x_4\right)} \sqrt[6]{\sin \left(6 H x_0\right)}\,\partial_{6}\psi'_{7} \\
 & = H V'(s)\,\psi'_{8}
\end{align*}
\begin{align*}
(E'_{9})\quad & \cot \left(6 H x_0\right)\,\partial_{0}\psi'_{1} - e^{-a_4\left(H x_4\right)} \sqrt[6]{\sin \left(6 H x_0\right)}\,\partial_{7}\psi'_{1} - e^{a_4\left(H x_4\right)} \sqrt[6]{\sin \left(6 H x_0\right)}\,\partial_{3}\psi'_{4} \\
 &  - \partial_{4}\psi'_{4} + e^{a_4\left(H x_4\right)} \sqrt[6]{\sin \left(6 H x_0\right)}\,\partial_{1}\psi'_{6} + e^{-a_4\left(H x_4\right)} \sqrt[6]{\sin \left(6 H x_0\right)}\,\partial_{6}\psi'_{6} \\
 &  + e^{a_4\left(H x_4\right)} \sqrt[6]{\sin \left(6 H x_0\right)}\,\partial_{2}\psi'_{7} - e^{-a_4\left(H x_4\right)} \sqrt[6]{\sin \left(6 H x_0\right)}\,\partial_{5}\psi'_{7} \\
 & = H V'(s)\,\psi'_{9}
\end{align*}
\begin{align*}
(E'_{10})\quad & \cot \left(6 H x_0\right)\,\partial_{0}\psi'_{2} - e^{-a_4\left(H x_4\right)} \sqrt[6]{\sin \left(6 H x_0\right)}\,\partial_{7}\psi'_{2} + e^{a_4\left(H x_4\right)} \sqrt[6]{\sin \left(6 H x_0\right)}\,\partial_{2}\psi'_{4} \\
 &  + e^{-a_4\left(H x_4\right)} \sqrt[6]{\sin \left(6 H x_0\right)}\,\partial_{5}\psi'_{4} - e^{a_4\left(H x_4\right)} \sqrt[6]{\sin \left(6 H x_0\right)}\,\partial_{1}\psi'_{5} - e^{-a_4\left(H x_4\right)} \sqrt[6]{\sin \left(6 H x_0\right)}\,\partial_{6}\psi'_{5} \\
 &  + e^{a_4\left(H x_4\right)} \sqrt[6]{\sin \left(6 H x_0\right)}\,\partial_{3}\psi'_{7} - \partial_{4}\psi'_{7} \\
 & = H V'(s)\,\psi'_{10}
\end{align*}
\begin{align*}
(E'_{11})\quad & \cot \left(6 H x_0\right)\,\partial_{0}\psi'_{3} - e^{-a_4\left(H x_4\right)} \sqrt[6]{\sin \left(6 H x_0\right)}\,\partial_{7}\psi'_{3} - e^{a_4\left(H x_4\right)} \sqrt[6]{\sin \left(6 H x_0\right)}\,\partial_{1}\psi'_{4} \\
 &  + e^{-a_4\left(H x_4\right)} \sqrt[6]{\sin \left(6 H x_0\right)}\,\partial_{6}\psi'_{4} - e^{a_4\left(H x_4\right)} \sqrt[6]{\sin \left(6 H x_0\right)}\,\partial_{2}\psi'_{5} + e^{-a_4\left(H x_4\right)} \sqrt[6]{\sin \left(6 H x_0\right)}\,\partial_{5}\psi'_{5} \\
 &  - e^{a_4\left(H x_4\right)} \sqrt[6]{\sin \left(6 H x_0\right)}\,\partial_{3}\psi'_{6} + \partial_{4}\psi'_{6} \\
 & = H V'(s)\,\psi'_{11}
\end{align*}
\begin{align*}
(E'_{12})\quad & e^{a_4\left(H x_4\right)} \sqrt[6]{\sin \left(6 H x_0\right)}\,\partial_{3}\psi'_{1} - \partial_{4}\psi'_{1} - e^{a_4\left(H x_4\right)} \sqrt[6]{\sin \left(6 H x_0\right)}\,\partial_{2}\psi'_{2} \\
 &  + e^{-a_4\left(H x_4\right)} \sqrt[6]{\sin \left(6 H x_0\right)}\,\partial_{5}\psi'_{2} + e^{a_4\left(H x_4\right)} \sqrt[6]{\sin \left(6 H x_0\right)}\,\partial_{1}\psi'_{3} + e^{-a_4\left(H x_4\right)} \sqrt[6]{\sin \left(6 H x_0\right)}\,\partial_{6}\psi'_{3} \\
 &  + \cot \left(6 H x_0\right)\,\partial_{0}\psi'_{4} + e^{-a_4\left(H x_4\right)} \sqrt[6]{\sin \left(6 H x_0\right)}\,\partial_{7}\psi'_{4} \\
 & = H V'(s)\,\psi'_{12}
\end{align*}
\begin{align*}
(E'_{13})\quad & -e^{a_4\left(H x_4\right)} \sqrt[6]{\sin \left(6 H x_0\right)}\,\partial_{3}\psi'_{0} + \partial_{4}\psi'_{0} + e^{a_4\left(H x_4\right)} \sqrt[6]{\sin \left(6 H x_0\right)}\,\partial_{1}\psi'_{2} \\
 &  - e^{-a_4\left(H x_4\right)} \sqrt[6]{\sin \left(6 H x_0\right)}\,\partial_{6}\psi'_{2} + e^{a_4\left(H x_4\right)} \sqrt[6]{\sin \left(6 H x_0\right)}\,\partial_{2}\psi'_{3} + e^{-a_4\left(H x_4\right)} \sqrt[6]{\sin \left(6 H x_0\right)}\,\partial_{5}\psi'_{3} \\
 &  + \cot \left(6 H x_0\right)\,\partial_{0}\psi'_{5} + e^{-a_4\left(H x_4\right)} \sqrt[6]{\sin \left(6 H x_0\right)}\,\partial_{7}\psi'_{5} \\
 & = H V'(s)\,\psi'_{13}
\end{align*}
\begin{align*}
(E'_{14})\quad & e^{a_4\left(H x_4\right)} \sqrt[6]{\sin \left(6 H x_0\right)}\,\partial_{2}\psi'_{0} - e^{-a_4\left(H x_4\right)} \sqrt[6]{\sin \left(6 H x_0\right)}\,\partial_{5}\psi'_{0} - e^{a_4\left(H x_4\right)} \sqrt[6]{\sin \left(6 H x_0\right)}\,\partial_{1}\psi'_{1} \\
 &  + e^{-a_4\left(H x_4\right)} \sqrt[6]{\sin \left(6 H x_0\right)}\,\partial_{6}\psi'_{1} + e^{a_4\left(H x_4\right)} \sqrt[6]{\sin \left(6 H x_0\right)}\,\partial_{3}\psi'_{3} + \partial_{4}\psi'_{3} \\
 &  + \cot \left(6 H x_0\right)\,\partial_{0}\psi'_{6} + e^{-a_4\left(H x_4\right)} \sqrt[6]{\sin \left(6 H x_0\right)}\,\partial_{7}\psi'_{6} \\
 & = H V'(s)\,\psi'_{14}
\end{align*}
\begin{align*}
(E'_{15})\quad & -e^{a_4\left(H x_4\right)} \sqrt[6]{\sin \left(6 H x_0\right)}\,\partial_{1}\psi'_{0} - e^{-a_4\left(H x_4\right)} \sqrt[6]{\sin \left(6 H x_0\right)}\,\partial_{6}\psi'_{0} - e^{a_4\left(H x_4\right)} \sqrt[6]{\sin \left(6 H x_0\right)}\,\partial_{2}\psi'_{1} \\
 &  - e^{-a_4\left(H x_4\right)} \sqrt[6]{\sin \left(6 H x_0\right)}\,\partial_{5}\psi'_{1} - e^{a_4\left(H x_4\right)} \sqrt[6]{\sin \left(6 H x_0\right)}\,\partial_{3}\psi'_{2} - \partial_{4}\psi'_{2} \\
 &  + \cot \left(6 H x_0\right)\,\partial_{0}\psi'_{7} + e^{-a_4\left(H x_4\right)} \sqrt[6]{\sin \left(6 H x_0\right)}\,\partial_{7}\psi'_{7} \\
 & = H V'(s)\,\psi'_{15}
\end{align*}

}

\subsection{The \texorpdfstring{$x_4$}{x4}-only family and the condition \texorpdfstring{$V''s'=0$}{V'' s' = 0}}
\label{sec:5.8}

The rescaled field with no $x_0$-dependence, $\Psi=\sqrt{\sin 6Hx_0}\;\Psi'(x_4)$, satisfies the field
equation at every $x_0$ \textbf{if and only if} $\partial_{x_0}V'\bigl(\sin(6Hx_0)\,s'\bigr)=0$, and
$\partial_{x_0}V'(\sin\,s')=6H\cos(6Hx_0)\,s'\,V''(\sin\,s')$. So it is a solution exactly when
$V''s'=0$:
\begin{itemize}
  \item either $V$ is \textbf{linear} (the author's mass term $V=-(2M/H)\,s$, $V'=-2M/H$ constant), or
  \item the data are \textbf{null}, $s'=\bar\Psi'\Psi'=0$.
\end{itemize}
For the complex field $s'$ is \textbf{conserved} by the $x_4$-only equations
$\partial_4\Psi'=-HV'\,\Tsix{4}\Psi'$, $\partial_4\bar\Psi'=HV'\,\bar\Psi'\Tsix{4}$, for \textbf{any}
$V$. So null data give $x_4$-only solutions for any $V$, with $\rho=V(0)$. For a nonlinear $V$ and
$s'\neq0$ the family fails (witness: $V=s^2$). \cfproved{VII \S29}
\begin{itemize}
  \item \cfL{PRE-UNIVERSE (2) [THE RESULT]: Psi = Sqrt[Sin] Psi'(x4) satisfies the field equation iff d\_x0 V'(Sin[6Hx0] s') == 0, and d\_x0 V'(Sin s') == 6 H Cos[6 H x0] s' V''(Sin s'): V linear OR s' == 0}
  \item \cfL{PRE-UNIVERSE (2) [has content]: s' = Psi'bar Psi' is conserved by the x4-only equations (d\_4 Psi' = -H V' T16[4] Psi', d\_4 Psi'bar = H V' Psi'bar T16[4]) for ANY V; so null data s' = 0 are solutions for any V, with rho == V(0); and for V = s\textasciicircum{}2 with s' != 0 the x0-dependence is real (witness)}
\end{itemize}

\subsection{The operator form of the field equation, and its ordering}
\label{sec:5.9}

These are definitions, not identities \cfprose{VII \S27}. The Hamiltonian (\secref{sec:6.11})
contains $V(:\!s\!:)$, defined by the spectral calculus of the Hermitian operator
$s(x)=\Psi^{\dagger}\beta\Psi$ of \secref{sec:6.5} (regularized and normal-ordered). No ordering
prescription is needed there, and a non-polynomial $V$ is covered. The operator field equation is
\[
  \gamma^\mu D_\mu\Psi=H\,{:}V'(s)\,\Psi{:}
\]
where $V'(s)\Psi$ is the \textbf{Weyl-symmetrized} product, the average over the positions at which
$\Psi$ can be inserted in $V'(s)$.
\begin{itemize}
  \item It is \textbf{exact} for linear $V$ (the author's mass term).
  \item For nonlinear $V$ it differs from the naive product by Hartree and Fock contractions with the
        coincident $J$-vacuum propagator. That propagator is a c-number \textbf{matrix} with two
        parts: a \textbf{scalar} part, which is divergent and renormalizes $V'$; and a
        $\gamma^4$ part, the charge density of the Dirac sea, which is a constant
        chemical-potential shift. The $\gamma^4$ part is removed by charge-symmetric ordering or by
        a phase $\Psi\to e^{i\mu x_4}\Psi$, \textbf{not} by renormalizing $V$.
  \item In the Kohn--Sham mean field of \secref{sec:8}, $V'(s)$ is replaced by $V'(\langle s\rangle)$
        and the equation is linear.
\end{itemize}

\subsection{The U(1) current and the particle-number current}
\label{sec:5.10}

$\Lsym$ is invariant under $\Psi\to e^{i\alpha}\Psi$, $\bar\Psi\to e^{-i\alpha}\bar\Psi$. The Noether
current is
\[
  J^\mu=\frac{\partial L}{\partial(\partial_\mu\Psi)}\,(i\Psi)+(-i\bar\Psi)\,\frac{\partial L}{\partial(\partial_\mu\bar\Psi)}
  =\frac{\sqrt g}{H}\;i\,\bar\Psi\gamma^\mu\Psi ,
\]
so $j^\mu=i\,\bar\Psi\gamma^\mu\Psi$ is the U(1) current. It is real, because $i\sigma_{16}\gamma^\mu$ is
Hermitian. \textbf{Off shell}, for generic fields of all eight coordinates,
\[
  \frac{1}{\sqrt g}\,\partial_\mu\bigl(\sqrt g\,j^\mu\bigr)=i\,\bigl(\bar E\,\Psi+\bar\Psi\,E\bigr),
  \qquad E,\ \bar E=\text{the two field-equation residuals},
\]
so the current is conserved on shell. \cfproved{VII \S27}
\begin{itemize}
  \item \cfL{CURRENT [THE RESULT]: the Noether current of the U(1) phase is (Sqrt[g]/H) i Psibar gamma\textasciicircum{}mu Psi}
  \item \cfL{CURRENT [has content]: i sigma16 gamma\textasciicircum{}mu is Hermitian for every mu, so j\textasciicircum{}mu is real}
  \item \cfL{CURRENT [THE RESULT]: OFF SHELL (1/Sqrt[g]) d\_mu(Sqrt[g] j\textasciicircum{}mu) == i (Ebar Psi + Psibar E) for generic fields of all eight coordinates -- conserved on shell}
\end{itemize}

The \textbf{particle-number current} is $n^\mu:=-\frac{i}{H}\,\bar\Psi\gamma^\mu\Psi=-j^\mu/H$.
As \secref{sec:6.15} shows, $N=\int\sqrt g\,n^4\,d^7x$ is the normal-ordered number of particles
minus antiparticles, so the charge of $j$ carries a conventional minus sign:
$\int\sqrt g\,j^4\,d^7x=-HN$ \cfprose{VII \S27}. So one particle (one occupied positive-frequency
$J$-mode, $N=+1$) carries $j$-charge $-H$, and one antiparticle $+H$ [derived here from the two
definitions; the design review, QK-10, states the same sign as ``a positive-frequency $J$-mode
carries $j$-charge $-1$'', that is, in units of $H$].

\subsection{The only internal symmetry is the vector phase}
\label{sec:5.11}

The commutant of all eight $\Tsix{a}$ in $M_{16}(\mathbb C)$ is one-dimensional (multiples of
$\IDs$), so the only internal U(1) is the vector phase. In particular the \textbf{axial} phase
$\Psi\to\exp(i\alpha\,\Tsix{8})\,\Psi$ leaves $s=\bar\Psi\Psi$ invariant but \textbf{not} the kinetic
term ($\Tsix{8}$ anticommutes with every gamma):
$\exp(-i\alpha T_8)\,\Tsix{a}\,\exp(i\alpha T_8)=\Tsix{a}\exp(2i\alpha T_8)$, so
$\bar\Psi\gamma^\mu\partial_\mu\Psi$ changes at first order by
$2i\alpha\,\bar\Psi\gamma^\mu\Tsix{8}\,\partial_\mu\Psi$, which is not zero. This is the reverse of
the four-dimensional intuition, where the axial phase preserves the kinetic term and breaks the mass
term. \cfproved{VII \S27}
\begin{itemize}
  \item \cfL{SYMMETRY [THE RESULT]: the only internal U(1) is the vector phase -- the commutant of all eight T16[a] in M16(C) is one-dimensional (multiples of ID16)}
  \item \cfL{SYMMETRY [has content]: the axial phase exp(i alpha T16[8]) leaves s invariant but not the kinetic term: exp(-i alpha T8) T16[a] exp(i alpha T8) == T16[a] exp(2 i alpha T8), so Psibar gamma\textasciicircum{}mu d\_mu Psi changes at first order by 2 i alpha Psibar gamma\textasciicircum{}mu T16[8] d\_mu Psi, which is not zero (witness)}
\end{itemize}

% ===============================================================================================
\section{Canonical quantization in 4+4 dimensions}
\label{sec:6}

This is Section~28 of the notebook (Input cells 185--191, 56 assertions). ``Extended to 4+4
dimensions'' means: the time is one of the four timelike directions, $x_4$, and the slice of
constant time is seven-dimensional and contains the other three timelike directions. Everything
below follows from the symmetrized Lagrangian of \secref{sec:4}; the result is summarized first.

\paragraph{The result in one paragraph.} The Dirac bracket gives the anticommutator
$\{\Psi_a(x),\Psi^{\ddagger}_b(y)\}=\frac{H}{\sqrt{|g|}}\,(-i\sigma_{16}\gamma^4)_{ab}\,\delta^7(x-y)$.
Its matrix is $J=-i\,\Tsix{0}\Tsix{1}\Tsix{2}\Tsix{3}\Tsix{4}$, which is Hermitian with square 1 and
has signature $(8,8)$: the naive Fock space would have negative-norm states. $J$ acts only on a
hidden flavour index, $\sigma_{16}=J\beta$ with $\beta=-i\,\Tsix{4}$, and the Hilbert adjoint
$\Psi^{\dagger}:=\Psi^{\ddagger}J$ makes the anticommutator positive,
$\{\Psi,\Psi^{\dagger}\}=\frac{H}{\sqrt{|g|}}\,\delta^7$. This works exactly when the field carries no
momentum along the three hidden timelike directions $x_5,x_6,x_7$ (the admissibility theorem), and
there $J$ is unique. The admissible theory is a dimensional reduction with a finite hidden coordinate
volume $\Vhid$. Its Fock space is built on $J$-eigenmodes: per momentum, 8 positive-frequency and 8
negative-frequency modes; the ground state is the filled Dirac sea; the first excitations are 8
particles and 8 antiparticles. A finite Jordan--Wigner Fock space checks all of this exactly.

\subsection{Time, slices, and why the slice is not a Cauchy surface}
\label{sec:6.1}

Time is $x_4$, the observer's proper time. On the canonical frame $g_{44}=-1$ (lapse 1) and
$g_{4\mu}=0$ for $\mu\neq4$ (shift 0). A slice $x_4=\text{const}$ is \textbf{seven}-dimensional with
signature $(4,3)$: $x_0,x_1,x_2,x_3$ spacelike and $x_5,x_6,x_7$ \textbf{timelike}. It is therefore
not a Cauchy surface, and the quantization is carried out where it is well posed (the admissible
sector, \secref{sec:6.10}). \cfproved{VII \S28} \cfL{QUANTIZATION [definition]: lapse 1 and shift 0 (g\_44 == -1, g\_4mu == 0), and the slice x4 = const has signature (4,3) -- x0..x3 spacelike, x5..x7 timelike (given a4 real and 0 < 6 H x0 < Pi/2)}

\subsection{The momenta, the kinetic form \texorpdfstring{$G$}{G}, and the signature \texorpdfstring{$(8,8)$}{(8,8)}}
\label{sec:6.2}

From $\Lsym$ the canonical momenta are
\[
  \Pi_\Psi=\frac{\partial L}{\partial(\partial_4\Psi)}=\frac{\sqrt g}{2H}\,\bar\Psi\gamma^4,\qquad
  \Pi_{\bar\Psi}=\frac{\partial L}{\partial(\partial_4\bar\Psi)}=-\frac{\sqrt g}{2H}\,\gamma^4\Psi .
\]
The symplectic potential $\Pi_\Psi\,d\Psi+d\bar\Psi\,\Pi_{\bar\Psi}$ differs from that of the
unsymmetrized Lagrangian, $\frac{\sqrt g}{H}\,\bar\Psi\gamma^4\,d\Psi$, by the exact variation
$\delta\bigl[\frac{\sqrt g}{2H}\,\bar\Psi\gamma^4\Psi\bigr]$: the two define the \textbf{same}
symplectic form, hence the same brackets. With $\bar\Psi=\Psi^{\ddagger}\sigma_{16}$ the
time-derivative part of the Lagrangian is
\[
  \Psi^{\ddagger}K\,\partial_4\Psi,\qquad K=\frac{\sqrt g}{H}\,\sigma_{16}\gamma^4=i\,\Gt,\qquad
  \Gt=\frac{\sqrt g}{H}\,G,\qquad G=-i\,\sigma_{16}\gamma^4 .
\]
On the canonical frame $\gamma^4=\Tsix{4}$, so
\[
  G=-i\,\sigma_{16}\Tsix{4}=-i\,\Tsix{0}\Tsix{1}\Tsix{2}\Tsix{3}\Tsix{4}=:J .
\]
$G$ is Hermitian, $G^2=\IDs$, with eigenvalues $+1$ (8 times) and $-1$ (8 times): the kinetic form is
\textbf{indefinite}, signature $(8,8)$. The canonical frame has $\Gamma_4=0$, so the time-derivative
part, which alone determines the anticommutator, contains no spin connection. \cfproved{VII \S28}
\begin{itemize}
  \item \cfL{QUANTIZATION [THE RESULT]: the momenta of Lsym are Pi\_Psi == (Sqrt[g]/(2H)) Psibar gamma\textasciicircum{}4 and Pi\_Psibar == -(Sqrt[g]/(2H)) gamma\textasciicircum{}4 Psi}
  \item \cfL{QUANTIZATION [has content]: theta\_un - theta\_sym == delta[(Sqrt[g]/(2H)) Psibar gamma\textasciicircum{}4 Psi] exactly: the same symplectic form, the same brackets}
  \item \cfL{QUANTIZATION [THE RESULT]: G = -i sigma16 gamma\textasciicircum{}4 == J = -i T16[0].T16[1].T16[2].T16[3].T16[4] on the canonical frame; J is Hermitian, J\textasciicircum{}2 == ID16, eigenvalues +1 (8) and -1 (8): signature (8,8)}
  \item \cfL{QUANTIZATION [has content]: Gamma\_4 == 0 -- the time-derivative part of the Lagrangian, which fixes the anticommutator, contains no spin connection}
\end{itemize}

\subsection{The Dirac bracket and the canonical anticommutator}
\label{sec:6.3}

Write $\Psi=(\theta_1+i\theta_2)/\sqrt2$ with real Grassmann $\theta_1,\theta_2$ and
$\Theta=(\theta_1,\theta_2)$. In the Grassmann algebra of subsection~\ref{sec:3.3}, for
$K=c\,\sigma_{16}\Tsix{4}$ with any scalar $c$,
\begin{gather*}
  \Psi^{\ddagger}K\,d\Psi=\frac{i}{2}\,\Theta^{T}\Omega_c\,d\Theta+\frac14\,d\bigl(\theta_1^{T}K\theta_1+\theta_2^{T}K\theta_2\bigr),\\
  \Omega_c=\begin{pmatrix}0&K\\-K&0\end{pmatrix}\quad\text{(real, symmetric, invertible;}\quad
  \Omega_c^{-1}=\begin{pmatrix}0&-K^{-1}\\K^{-1}&0\end{pmatrix}\text{)} .
\end{gather*}
So up to a total time derivative the Lagrangian is the first-order system
$\frac{i}{2}\,\Theta^{T}\Omega\,\Theta'-h$. Its momenta are proportional to $\Theta$; the constraints
$\pi-(\text{momentum as a function of }\Theta)$ are second class with constraint matrix proportional
to $\Omega$; and the Dirac bracket of the $\Theta$'s is proportional to $\Omega^{-1}$. The constant is
fixed by the elementary case $\Omega=1$, $\frac{i}{2}\,\xi\xi'$, whose quantization is $\{\xi,\xi\}=1$.
Hence
\begin{align*}
  \{\Theta_i,\Theta_j\}&=(\Omega^{-1})_{ij},\\
  \{\Psi_a,\Psi^{\ddagger}_b\}&=\tfrac12\bigl[(\Omega^{-1})_{11}+(\Omega^{-1})_{22}-i\,(\Omega^{-1})_{12}+i\,(\Omega^{-1})_{21}\bigr]_{ab}
    =i\,(K^{-1})_{ab}=(\Gt^{-1})_{ab},\\
  \{\Psi_a,\Psi_b\}&=0 .
\end{align*}
The same sign and normalization are confirmed independently, with no convention for graded Poisson
brackets, by the Fock space of \secref{sec:6.14}: with this anticommutator the Heisenberg equation
generated by the canonical Hamiltonian reproduces the field equation, and with the opposite sign it
gives the time-reversed one. On the canonical frame $J^{-1}=J$ and $\sqrt g=\sec 6Hx_0$, so
\[
  \boxed{\;\{\Psi_a(x),\Psi^{\ddagger}_b(y)\}_{x_4=y_4}=\frac{H}{\sqrt g}\,J_{ab}\,\delta^7(x-y)=H\cos(6Hx_0)\,J_{ab}\,\delta^7(x-y),
  \qquad\{\Psi_a,\Psi_b\}=0\;}
\]
which is $\frac{H}{\sqrt{|g|}}\,(-i\sigma_{16}\gamma^4)_{ab}$, because
$(-i\sigma_{16}\Tsix{4})^{-1}=-i\sigma_{16}\Tsix{4}$. The spin connection does not appear: the
anticommutator comes from the time-derivative term alone, and $\Gamma_4=0$. \cfproved{VII \S28}
\begin{itemize}
  \item \cfL{DIRAC BRACKET [THE RESULT]: in the Grassmann algebra Psi\textasciicircum{}ddag K phi == (i/2) Theta\textasciicircum{}T Omega\_c Phi + (1/4) d(theta1\textasciicircum{}T K theta1 + theta2\textasciicircum{}T K theta2), for K = c sigma16 T16[4] with a free scalar c}
  \item \cfL{DIRAC BRACKET [has content]: Omega\_c is real symmetric and invertible, Omega\_c\textasciicircum{}-1 == [[0, -K\textasciicircum{}-1], [K\textasciicircum{}-1, 0]]}
  \item \cfL{DIRAC BRACKET [THE RESULT]: \{Theta\_i, Theta\_j\} = (Omega\textasciicircum{}-1)\_ij gives \{Psi\_a, Psi\textasciicircum{}ddag\_b\} == i (K\textasciicircum{}-1)\_ab == (Gtilde\textasciicircum{}-1)\_ab and \{Psi\_a, Psi\_b\} == 0}
  \item \cfL{ANTICOMMUTATOR [THE RESULT]: on the canonical frame (c = Sqrt[g]/H) \{Psi, Psi\textasciicircum{}ddag\} == (H/Sqrt[g]) J == H Cos[6 H x0] J  (J\textasciicircum{}-1 == J), which IS the design's (H/Sqrt[g])(-i sigma16 gamma\textasciicircum{}4)}
\end{itemize}

The exported anticommutator, with its matrix
(\file{provenance-latex/generated/fermion_fable_anticommutator.tex}, unchanged):

{\small
% fermion fable: the canonical anticommutator on the canonical frame
% generated by claude-fable/export_fermion_fable_tex.wls from the notebook's Part VII objects, 2026-09-24
% notation: x_0..x_7 are the coordinates (x_4 = the observer's time), a_4 = the free function a4 (never given a value),
% q = e^{-a_4(H x_4)}/sin^{1/6}(6 H x_0), p = e^{+a_4(H x_4)}/sin^{1/6}(6 H x_0), s = bar(psi) psi.
\begin{align*}
\{\Psi_a(x),\Psi^{\ddagger}_b(y)\}_{x_4=y_4} &= \frac{H}{\sqrt{|g|}}\,J_{ab}\,\delta^7(x-y) = H \cos \left(6 H x_0\right)\,J_{ab}\,\delta^7(x-y), \qquad \{\Psi_a,\Psi_b\}=0, \\
\{\Psi_a(x),\Psi^{\dagger}_b(y)\}_{x_4=y_4} &= \frac{H}{\sqrt{|g|}}\,\delta_{ab}\,\delta^7(x-y)
\end{align*}
with $\Psi^\dagger = \Psi^\ddagger J$; the anticommutator matrix is $\frac{H}{\sqrt{|g|}}J = H \cos \left(6 H x_0\right)\,J$ with
\[ J = \setcounter{MaxMatrixCols}{16}
\left(\begin{smallmatrix}
0 & 0 & 0 & 0 & 0 & 0 & 0 & 0 & 0 & i & 0 & 0 & 0 & 0 & 0 & 0 \\
0 & 0 & 0 & 0 & 0 & 0 & 0 & 0 & -i & 0 & 0 & 0 & 0 & 0 & 0 & 0 \\
0 & 0 & 0 & 0 & 0 & 0 & 0 & 0 & 0 & 0 & 0 & -i & 0 & 0 & 0 & 0 \\
0 & 0 & 0 & 0 & 0 & 0 & 0 & 0 & 0 & 0 & i & 0 & 0 & 0 & 0 & 0 \\
0 & 0 & 0 & 0 & 0 & 0 & 0 & 0 & 0 & 0 & 0 & 0 & 0 & -i & 0 & 0 \\
0 & 0 & 0 & 0 & 0 & 0 & 0 & 0 & 0 & 0 & 0 & 0 & i & 0 & 0 & 0 \\
0 & 0 & 0 & 0 & 0 & 0 & 0 & 0 & 0 & 0 & 0 & 0 & 0 & 0 & 0 & i \\
0 & 0 & 0 & 0 & 0 & 0 & 0 & 0 & 0 & 0 & 0 & 0 & 0 & 0 & -i & 0 \\
0 & i & 0 & 0 & 0 & 0 & 0 & 0 & 0 & 0 & 0 & 0 & 0 & 0 & 0 & 0 \\
-i & 0 & 0 & 0 & 0 & 0 & 0 & 0 & 0 & 0 & 0 & 0 & 0 & 0 & 0 & 0 \\
0 & 0 & 0 & -i & 0 & 0 & 0 & 0 & 0 & 0 & 0 & 0 & 0 & 0 & 0 & 0 \\
0 & 0 & i & 0 & 0 & 0 & 0 & 0 & 0 & 0 & 0 & 0 & 0 & 0 & 0 & 0 \\
0 & 0 & 0 & 0 & 0 & -i & 0 & 0 & 0 & 0 & 0 & 0 & 0 & 0 & 0 & 0 \\
0 & 0 & 0 & 0 & i & 0 & 0 & 0 & 0 & 0 & 0 & 0 & 0 & 0 & 0 & 0 \\
0 & 0 & 0 & 0 & 0 & 0 & 0 & i & 0 & 0 & 0 & 0 & 0 & 0 & 0 & 0 \\
0 & 0 & 0 & 0 & 0 & 0 & -i & 0 & 0 & 0 & 0 & 0 & 0 & 0 & 0 & 0
\end{smallmatrix}\right) \]

}

\subsection{The admissible truncation, \texorpdfstring{$V_{\mathrm{hid}}$}{V\_hid}, and the rescaled field}
\label{sec:6.4}

The admissible theory (\secref{sec:6.10}) is a \textbf{truncation}: a dimensional reduction to fields
$\Psi=\Psi(x_0,x_1,x_2,x_3,x_4)$, independent of $x_5,x_6,x_7$. It needs a \textbf{finite} hidden
coordinate volume $\Vhid=\int dx_5\,dx_6\,dx_7$. That is an assumption, and it has a consequence that
must be stated: compact timelike directions contain closed timelike curves. In the truncation the
delta function becomes $\delta^4(x_{0..3}-y_{0..3})$ and the factor $H/\sqrt g$ becomes
$H/(\Vhid\sqrt g)$:
\[
  \{\Psi_a(x),\Psi^{\ddagger}_b(y)\}_{x_4=y_4}=\frac{H}{\Vhid\sqrt{|g|}}\;G_{ab}\;\delta^4(x_{0..3}-y_{0..3}),
  \qquad G=-i\,\sigma_{16}\gamma^4 .
\]
For the rescaled field $\Psi'=\Psi/\sqrt{\sin 6Hx_0}$ of \secref{sec:5.7} the factor is
$H\cos(6Hx_0)/\sin(6Hx_0)=H\cot(6Hx_0)$. On a frame whose $\sqrt g$ depends on $x_4$ one quantizes
$\chi=(\sqrt g/H)^{1/2}\,\Psi$, with $\{\chi,\chi^{\ddagger}\}=G\,\delta^7$; on the canonical frame
$\sqrt g$ does not depend on $x_4$ and the two coincide \cfprose{VII \S28}. \cfproved{VII \S28}
\cfL{ANTICOMMUTATOR [has content]: the rescaled field Psi' = Psi/Sqrt[Sin[6 H x0]] has \{Psi', Psi'\textasciicircum{}ddag\} == H Cot[6 H x0] J delta\textasciicircum{}7, and Sqrt[g] is x4-independent (so quantizing chi = (Sqrt[g]/H)\textasciicircum{}(1/2) Psi changes nothing on the canonical frame)}

(For the time-dependent warped frames of the interacting system, Part~VIII of the notebook,
Section~32, proves that the rescaled field has an $x_4$-independent anticommutator:
\cfL{RESCALING [THE RESULT]: \{Psi, Psi\textasciicircum{}ddag\} == (H/Sqrt[|det g|]) G on the warped frame (Section 28's cfAntiPsiPsiDd at c = Sqrt[|det g|]/H), and for chi = Psi/f, f = Sqrt[Sin] (A\textasciicircum{}3 B\textasciicircum{}3 C)\textasciicircum{}(-1/2): \{chi, chi\textasciicircum{}ddag\} == H Cot[6 H x0] G -- the Sqrt[det g] cancels, independent of x4 and of A, B, C} and \cfL{RESCALING [has content]: on the warped frame G = -i sigma16 gamma\textasciicircum{}4 is Section 28's cfGK (= J) -- the same matrix on every member of the family (gamma\textasciicircum{}4 == T16[4], lapse 1)}.)

\subsection{What \texorpdfstring{$J$}{J} is, and the Hilbert adjoint}
\label{sec:6.5}

$J=-i\,\Tsix{0}\cdots\Tsix{4}=-i\,\Omega_4\Tsix{0}=-i\,h_0$, with
$\Omega_4=\Tsix{1}\Tsix{2}\Tsix{3}\Tsix{4}$ of \secref{sec:3.6}. So $J$ is one of the hidden
\textbf{flavour} gammas: it lies in the commutant of the observer's Clifford algebra, and on
$\mathbb C^{16}=\mathbb C^4(\text{observer})\otimes\mathbb C^4(\text{flavour})$ it acts on the
\textbf{flavour factor only}, where it has signature $(2,2)$ (verified on the image of a rank-4
primitive idempotent of the observer's algebra). With $\beta:=-i\,\Tsix{4}$ (Hermitian,
$\beta^2=1$):
\[
  \sigma_{16}=J\beta,\qquad \beta\,\Tsix{j}\,\beta=-\Tsix{j}\ \text{ and }\ \Tsix{j}\ \text{Hermitian}\quad(j=0..3),
\]
so $\beta$ is the observer's Dirac-adjoint matrix. Define the \textbf{Hilbert adjoint} by the
$J$-involution, $\Psi^{\dagger}:=\Psi^{\ddagger}J$. Then
\begin{align*}
  \bar\Psi&=\Psi^{\ddagger}\sigma_{16}=\Psi^{\dagger}\beta &&\text{(the standard 4D Dirac adjoint)},\\
  s&=\Psi^{\dagger}\beta\,\Psi,\\
  \Psi^{\ddagger}G\,\partial_4\Psi&=\Psi^{\dagger}\partial_4\Psi &&\text{(the time-kinetic term becomes } \tfrac{\sqrt g}{H}\,i\,\Psi^{\dagger}\partial_4\Psi\text{)},\\
  \{\Psi_a,\Psi^{\dagger}_b\}&=\frac{H}{\sqrt g}\,\delta_{ab}\,\delta^7 &&\text{(POSITIVE)}.
\end{align*}
So fermion fable is four flavours times a 4D Dirac field, and $J$-quantization replaces the classical
flavour metric of signature $(2,2)$ by the identity. The adjoint is \textbf{mode-independent}, and it
must be defined this way: a mode-by-mode prescription is basis-dependent unless every mode is a
$J$-eigenvector (\secref{sec:6.13}). \cfproved{VII \S28}
\begin{itemize}
  \item \cfL{J [THE RESULT]: J == -i Omega4 T16[0] == -i h\_0 is a hidden flavour gamma: it lies in the commutant of T16[1..4] (in the span of the hidden-gamma products of Section 26)}
  \item \cfL{J [has content]: on the flavour factor J has signature (2,2): cfIdemObs is a rank-4 projector commuting with J, and J restricted to its image has eigenvalues +1, +1, -1, -1}
  \item \cfL{J [THE RESULT]: sigma16 == J beta with beta = -i T16[4] Hermitian, beta\textasciicircum{}2 == 1; T16[0..3] Hermitian with beta T16[j] beta == -T16[j]: beta is the observer's Dirac adjoint}
  \item \cfL{J [THE RESULT]: with Psi\textasciicircum{}dagger := Psi\textasciicircum{}ddag J:  Psibar = Psi\textasciicircum{}dagger beta (J sigma16 == beta),  Psi\textasciicircum{}ddag G = Psi\textasciicircum{}dagger (J G == ID16),  \{Psi, Psi\textasciicircum{}dagger\} == (H/Sqrt[g]) ID16 (G\textasciicircum{}-1 J == ID16): POSITIVE}
\end{itemize}

$J$ and $\beta$, as exported (\file{provenance-latex/generated/fermion_fable_krein_J.tex}, unchanged):

{\small
% fermion fable: the Krein fundamental symmetry J, the Dirac-adjoint matrix beta, sigma16 = J beta
% generated by claude-fable/export_fermion_fable_tex.wls from the notebook's Part VII objects, 2026-09-24
% notation: x_0..x_7 are the coordinates (x_4 = the observer's time), a_4 = the free function a4 (never given a value),
% q = e^{-a_4(H x_4)}/sin^{1/6}(6 H x_0), p = e^{+a_4(H x_4)}/sin^{1/6}(6 H x_0), s = bar(psi) psi.
\[ J = -i\,T_{16}[0]T_{16}[1]T_{16}[2]T_{16}[3]T_{16}[4] = -i\,\sigma_{16}T_{16}[4] = \setcounter{MaxMatrixCols}{16}
\left(\begin{smallmatrix}
0 & 0 & 0 & 0 & 0 & 0 & 0 & 0 & 0 & i & 0 & 0 & 0 & 0 & 0 & 0 \\
0 & 0 & 0 & 0 & 0 & 0 & 0 & 0 & -i & 0 & 0 & 0 & 0 & 0 & 0 & 0 \\
0 & 0 & 0 & 0 & 0 & 0 & 0 & 0 & 0 & 0 & 0 & -i & 0 & 0 & 0 & 0 \\
0 & 0 & 0 & 0 & 0 & 0 & 0 & 0 & 0 & 0 & i & 0 & 0 & 0 & 0 & 0 \\
0 & 0 & 0 & 0 & 0 & 0 & 0 & 0 & 0 & 0 & 0 & 0 & 0 & -i & 0 & 0 \\
0 & 0 & 0 & 0 & 0 & 0 & 0 & 0 & 0 & 0 & 0 & 0 & i & 0 & 0 & 0 \\
0 & 0 & 0 & 0 & 0 & 0 & 0 & 0 & 0 & 0 & 0 & 0 & 0 & 0 & 0 & i \\
0 & 0 & 0 & 0 & 0 & 0 & 0 & 0 & 0 & 0 & 0 & 0 & 0 & 0 & -i & 0 \\
0 & i & 0 & 0 & 0 & 0 & 0 & 0 & 0 & 0 & 0 & 0 & 0 & 0 & 0 & 0 \\
-i & 0 & 0 & 0 & 0 & 0 & 0 & 0 & 0 & 0 & 0 & 0 & 0 & 0 & 0 & 0 \\
0 & 0 & 0 & -i & 0 & 0 & 0 & 0 & 0 & 0 & 0 & 0 & 0 & 0 & 0 & 0 \\
0 & 0 & i & 0 & 0 & 0 & 0 & 0 & 0 & 0 & 0 & 0 & 0 & 0 & 0 & 0 \\
0 & 0 & 0 & 0 & 0 & -i & 0 & 0 & 0 & 0 & 0 & 0 & 0 & 0 & 0 & 0 \\
0 & 0 & 0 & 0 & i & 0 & 0 & 0 & 0 & 0 & 0 & 0 & 0 & 0 & 0 & 0 \\
0 & 0 & 0 & 0 & 0 & 0 & 0 & i & 0 & 0 & 0 & 0 & 0 & 0 & 0 & 0 \\
0 & 0 & 0 & 0 & 0 & 0 & -i & 0 & 0 & 0 & 0 & 0 & 0 & 0 & 0 & 0
\end{smallmatrix}\right) \]
\[ \beta = -i\,T_{16}[4] = \setcounter{MaxMatrixCols}{16}
\left(\begin{smallmatrix}
0 & 0 & 0 & 0 & 0 & 0 & 0 & 0 & 0 & 0 & 0 & 0 & 0 & i & 0 & 0 \\
0 & 0 & 0 & 0 & 0 & 0 & 0 & 0 & 0 & 0 & 0 & 0 & -i & 0 & 0 & 0 \\
0 & 0 & 0 & 0 & 0 & 0 & 0 & 0 & 0 & 0 & 0 & 0 & 0 & 0 & 0 & -i \\
0 & 0 & 0 & 0 & 0 & 0 & 0 & 0 & 0 & 0 & 0 & 0 & 0 & 0 & i & 0 \\
0 & 0 & 0 & 0 & 0 & 0 & 0 & 0 & 0 & -i & 0 & 0 & 0 & 0 & 0 & 0 \\
0 & 0 & 0 & 0 & 0 & 0 & 0 & 0 & i & 0 & 0 & 0 & 0 & 0 & 0 & 0 \\
0 & 0 & 0 & 0 & 0 & 0 & 0 & 0 & 0 & 0 & 0 & i & 0 & 0 & 0 & 0 \\
0 & 0 & 0 & 0 & 0 & 0 & 0 & 0 & 0 & 0 & -i & 0 & 0 & 0 & 0 & 0 \\
0 & 0 & 0 & 0 & 0 & -i & 0 & 0 & 0 & 0 & 0 & 0 & 0 & 0 & 0 & 0 \\
0 & 0 & 0 & 0 & i & 0 & 0 & 0 & 0 & 0 & 0 & 0 & 0 & 0 & 0 & 0 \\
0 & 0 & 0 & 0 & 0 & 0 & 0 & i & 0 & 0 & 0 & 0 & 0 & 0 & 0 & 0 \\
0 & 0 & 0 & 0 & 0 & 0 & -i & 0 & 0 & 0 & 0 & 0 & 0 & 0 & 0 & 0 \\
0 & i & 0 & 0 & 0 & 0 & 0 & 0 & 0 & 0 & 0 & 0 & 0 & 0 & 0 & 0 \\
-i & 0 & 0 & 0 & 0 & 0 & 0 & 0 & 0 & 0 & 0 & 0 & 0 & 0 & 0 & 0 \\
0 & 0 & 0 & -i & 0 & 0 & 0 & 0 & 0 & 0 & 0 & 0 & 0 & 0 & 0 & 0 \\
0 & 0 & i & 0 & 0 & 0 & 0 & 0 & 0 & 0 & 0 & 0 & 0 & 0 & 0 & 0
\end{smallmatrix}\right), \qquad \sigma_{16} = J\beta, \quad J^2=\beta^2=1, \quad \bar\Psi = \Psi^\dagger\beta . \]

}

\subsection{Why this is legitimate for fermions, and why it would fail for bosons}
\label{sec:6.6}

For Grassmann fields the ${}^{*}$-structure is a free choice: the Berezin integral treats $\Psi$ and
$\bar\Psi$ as independent, and the $J$-involution is that choice. The filled Dirac sea makes the
Hamiltonian bounded below. For a \textbf{commuting} 16-component field the same involution gives
$[b,b^{\dagger}]=1$, but the Hamiltonian $\sum_u\omega_u\,b_u^{\dagger}b_u$ then has 8
negative-frequency modes per momentum and no exclusion principle to fill them: it is unbounded
below. The only local involutions compatible with a Hermitian $s$ and a Hermitian Hamiltonian are
$+J$ and $-J$ (the uniqueness theorem, \secref{sec:6.10}), and $-J$ gives every mode negative norm;
making the bosonic Hamiltonian bounded below would turn 4 modes into negative-norm modes. For
fermions the negative-frequency modes are filled and the obstruction disappears \cfprose{VII \S28;
design review QK-5}.

The $J$-odd bilinears ($\hat T_{\mu h}$, $j^h$, $\bar\Psi\,\Tsix{8}\,\Psi$) become
\textbf{anti-Hermitian} operators (\secref{sec:6.8}). So the classical real theory of Part~VI is the
classical limit of the $J$-\textbf{even} sector only. (A reflection-positivity claim that appeared in
the design was dropped by the review: it was not verified.)

\subsection{Chirality, and why \texorpdfstring{$V$}{V} depends on \texorpdfstring{$s$}{s} only}
\label{sec:6.7}

$J$ anticommutes with $\Tsix{8}$: $J$ exchanges the two split-octonion types. Each chirality subspace
is $G$-\textbf{null} ($P_L G P_L=P_R G P_R=0$: the Krein form pairs type-1 with type-2 only), so the
$J$-vacuum is not chirality-graded; but the two chirality subspaces \textbf{are} orthogonal in the
Hilbert product ($P_LP_R=0$, with $P_L$, $P_R$ Hermitian). Neither $p=\bar\Psi\,\Tsix{8}\,\Psi$
(anti-Hermitian under $J$) nor $ip$ (anti-Hermitian under the classical conjugation ${}^{\ddagger}$,
so it would make the Lagrangian complex) can enter $V$ consistently: $\boldsymbol{V=V(s)}$.
\cfproved{VII \S28} \cfL{J [has content]: chirality -- \{J, T16[8]\} == 0, each chirality subspace is G-NULL (PL G PL == PR G PR == 0), and the two chiralities are orthogonal in the Hilbert product (PL, PR Hermitian, PL PR == 0)}

\subsection{The Hermiticity rule and the \texorpdfstring{$J$}{J} table}
\label{sec:6.8}

Under the $J$-adjoint, a bilinear $\Psi^{\ddagger}M\Psi$ with $M$ classically Hermitian has
$(\Psi^{\ddagger}M\Psi)^{\dagger}=\Psi^{\ddagger}(JMJ)\Psi$: it is \textbf{Hermitian} iff $[J,M]=0$ and
\textbf{anti-Hermitian} iff $\{J,M\}=0$. The table records $J$ against every matrix that occurs
\cfproved{VII \S28}:

{\small
\begin{longtable}{@{}p{0.36\linewidth}p{0.17\linewidth}p{0.41\linewidth}@{}}
\toprule
\raggedright matrix & \raggedright relation to $J$ & \raggedright consequence \tabularnewline
\midrule
\endhead
\raggedright $\sigma_{16}$ & \raggedright commutes & \raggedright $s$ is Hermitian \tabularnewline
\raggedright $\Tsix{0},\dots,\Tsix{4}$ & \raggedright commute & \tabularnewline
\raggedright $\Tsix{5}$, $\Tsix{6}$, $\Tsix{7}$ & \raggedright anticommute & \tabularnewline
\raggedright $\sigma_{16}\gamma^\mu$, $\mu=0..4$ & \raggedright commute & \raggedright $j^\mu$, $\mu=0..4$, Hermitian \tabularnewline
\raggedright $\sigma_{16}\gamma^h$, $h=5,6,7$ & \raggedright anticommute & \raggedright $j^h$ anti-Hermitian \tabularnewline
\raggedright each $\gamma^\mu\Gamma_\mu$ (no sum), $\mu=0..7$ & \raggedright commute & \raggedright the connection term of the Dirac operator is $J$-even \tabularnewline
\raggedright $\Gamma_0$, $\Gamma_4$ (both zero), $\Gamma_1$, $\Gamma_2$, $\Gamma_3$ & \raggedright commute & \tabularnewline
\raggedright $\Gamma_5$, $\Gamma_6$, $\Gamma_7$ & \raggedright anticommute & \raggedright boosts mixing $x_0$ or $x_4$ with the hidden timelike sheet \tabularnewline
\raggedright $\Tsix{8}$ and $C_+=\sigma_{16}\Tsix{8}$ & \raggedright anticommute & \raggedright $p=\bar\Psi\,\Tsix{8}\,\Psi$ is anti-Hermitian \tabularnewline
\bottomrule
\end{longtable}
}
\begin{itemize}
  \item \cfL{J TABLE [THE RESULT]: J commutes with sigma16, T16[0..4], sigma16 gamma\textasciicircum{}mu (mu = 0..4), every gamma\textasciicircum{}mu Gamma\_mu and Gamma\_1..Gamma\_3; it ANTIcommutes with T16[5..7], sigma16 gamma\textasciicircum{}h (h = 5..7), Gamma\_5..Gamma\_7, T16[8] and C+}
  \item \cfL{J [has content]: the Hermiticity rule -- for classically Hermitian M, (J M)\textasciicircum{}dagger == J M iff [J, M] == 0 and == -J M iff \{J, M\} == 0: so s = Psi\textasciicircum{}ddag sigma16 Psi is Hermitian and p = Psi\textasciicircum{}ddag C+ Psi is ANTI-Hermitian}
\end{itemize}
Consequences: $s$ is Hermitian and $p$ anti-Hermitian (so $V$ may depend on $s$ but not on $p$); the
hidden current components $j^h$ and the energy--momentum components $\hat T_{\mu h}$ ($\mu=0..4$) are
built from anticommuting matrices and are anti-Hermitian: their expectation values are purely
imaginary in every state and vanish in every state invariant under the hidden rotations (the
$J$-vacuum and the Kohn--Sham states of \secref{sec:8}), consistent with $G_{\mu h}=0$ on the
canonical frame. $\hat T_{hh'}$ ($h\neq h'$) vanishes identically for fields independent of
$x_5..x_7$, and $\hat T_{hh}=g_{hh}\hat L$ is Hermitian (\secref{sec:7.5}).

(The same $J$ table holds on the dynamical frames of the interacting system; Part~VIII, Section~32,
proves it: \cfL{J [THE RESULT]: [J, gamma\textasciicircum{}mu Gamma\_mu] == 0 for EACH mu separately (h = 5,6,7 included), on the warped and on the Bianchi-I frame} and \cfL{J [THE RESULT]: [J, Gamma\_mu] == 0 for mu = 0..4 and \{J, Gamma\_h\} == 0 for h = 5, 6, 7, on the warped and on the Bianchi-I frame}.)

\subsection{The symmetry that survives, and what is lost}
\label{sec:6.9}

Of the 28 Lorentz generators $S_{ab}$, \textbf{13 commute} with $J$: Spin(4,1) of $x_0..x_4$ (10
generators) and Spin(3) of $x_5..x_7$ (3). The other \textbf{15 anticommute}: the 12 boosts between
$x_0..x_3$ and $x_5..x_7$, and the 3 \textbf{rotations} between $x_4$ and $x_5..x_7$ ($x_4..x_7$ are
all timelike, so those planes are rotations). The Krein form $G$ itself is invariant
($S^{T}G+GS=0$) under exactly the 21 generators with $a,b\neq4$: Spin(4,3), the group of the slice.
\cfproved{VII \S28} (The $J$ structure is defined in the diagonal gauge of the canonical frame and
transforms as $J\to S^{-1}JS$ under a change of frame gauge $S$; design review QK-6.)
\begin{itemize}
  \item \cfL{J [THE RESULT]: 13 of the 28 Lorentz generators commute with J (a, b both in \{0..4\}: Spin(4,1), or both in \{5,6,7\}: Spin(3)); the other 15 anticommute -- 12 boosts x0..x3 <-> x5..x7 and 3 rotations x4 <-> x5..x7}
  \item \cfL{J [has content]: the Krein form G is invariant (S\textasciicircum{}T G + G S == 0) under exactly the 21 generators with a, b != 4 -- Spin(4,3), the group of the slice}
\end{itemize}

\subsection{The admissible sector: a theorem}
\label{sec:6.10}

Take a plane wave at momentum $k$ (flat local analysis, the frame factors frozen). The one-particle
Hamiltonian is $h_k=\sigma_{16}\,(m-i\,\gamma\cdot k)$, with $\gamma\cdot k=\sum_{a\neq4}k_a\Tsix{a}$,
and the evolution generator is $E_k=G^{-1}h_k=G\,h_k$, $i\,\partial_4\Psi=E_k\Psi$. Write
$k_s^2=k_0^2+k_1^2+k_2^2+k_3^2$ and $k_h^2=k_5^2+k_6^2+k_7^2$.
\begin{enumerate}
  \item For \textbf{every} $k$, $E_k^2=(k_s^2-k_h^2+m^2)\,\IDs$: the dispersion relation is
        $\omega^2=k_s^2-k_h^2+m^2$.
  \item \textbf{Admissible} $k$ ($k_h=0$): $[J,E_k]=0$ and $J$ itself gives the positive quantization
        ($GJ=\IDs$).
  \item Hidden momentum \textbf{below} threshold ($k_h^2<k_s^2+m^2$): a momentum-dependent, boosted
        $J'=S^{-1}JS$ ($S$ the spin boost in the plane of $k_s$ and $k_h$ that removes $k_h$) commutes
        with $E_k$ and with $\beta$, $J'^2=1$, and $GJ'=S^2$ is positive definite: each such momentum
        can be quantized positively on its own (checked at $k=5e_1+3e_5$, $m=3$,
        $\omega^2=25-9+9=25$, boost with $\tanh\theta=3/5$).
  \item \textbf{At} threshold $E_k$ is not zero but $E_k^2=0$: $E_k$ is nilpotent and cannot be
        diagonalized (Jordan blocks of size 2, rank 8); there is no mode expansion (checked at
        $k=4e_1+5e_5$, $m=3$).
  \item \textbf{Above} threshold $\omega^2<0$: the frequencies are imaginary (growth along $x_4$, the
        ill-posed ultrahyperbolic Cauchy problem) and the modes are $G$-\textbf{neutral}
        ($u^{\dagger}Gu=0$), so they carry no norm at all. Example: $k=(0,\tfrac34,1,0\mid k_5=2)$,
        $m=1$: $\omega^2=-23/16$, $\omega=\pm i\sqrt{23}/4=\pm1.19896\,i$.
  \item \textbf{No} single momentum-independent positive $J'$ exists once a hidden momentum is
        present. If $V(s)$ and the local energy--momentum tensor are to be Hermitian, $J'$ must commute
        with $\beta$ (because $s=\Psi^{\ddagger}G\beta\Psi$) and with $E_k$ for all $k$. Certificate:
        for the four admissible unit momenta and one hidden momentum, \textbf{every} $X$ commuting
        with $\beta$ and all five $E_k$ has $\Tr(GX)=0$, so $GX$ is never positive definite. Analytic
        proof for a purely hidden momentum: $[E_k,\beta]=2i\,k_5\,\Tsix{5}$, so $J'$ commutes with
        $\Tsix{5}$; $\Tsix{5}$ anticommutes with $G$, hence
        $\Tsix{5}^{\dagger}(GJ')\,\Tsix{5}=-(GJ')$: a positive matrix would be congruent to a negative
        one.
  \item $J$ is \textbf{unique}: the joint commutant of $\beta$ and all admissible $E_k$ is
        8-dimensional and every one of its elements commutes with $J$; if $X$ is in it with $X^2=1$
        and $GX=JX$ definite, then $(JX)^2=1$, and a definite Hermitian matrix with square 1 is
        $\pm\IDs$: $X=J$ (positive metric) or $X=-J$ (every mode of negative norm).
\end{enumerate}
The general identity behind (6) is $[E_k,\beta]=2i\,T_k$, $T_k=\sum_a k_a\Tsix{a}$, for every $k$.

\paragraph{THEOREM (admissibility).} Let $W$ be the linear span of the momenta carried by the modes,
a subspace of the slice $\mathbb R^{(4,3)}$. A positive (Hilbert-space) Fock quantization in which
the local $s(x)$, the local energy--momentum tensor and the free Hamiltonian are all Hermitian exists
\textbf{if and only if $W$ is spacelike-definite}. Then $J'=S^{-1}JS$ with $S$ a slice boost carrying
$W$ into $\operatorname{span}(x_0..x_3)$, and $J'$ is unique (with $J'^2=1$) when
$W=\operatorname{span}(x_0..x_3)$. On the canonical frame the fields depend on all of $x_0..x_3$, so
the condition is: \textbf{no momentum along $x_5,x_6,x_7$}, and
$J=-i\,\Tsix{0}\Tsix{1}\Tsix{2}\Tsix{3}\Tsix{4}$ is unique.

The admissible theory is a \textbf{truncation} (a dimensional reduction $\Psi=\Psi(x_0,\dots,x_4)$),
not a superselection sector: $V(s)$ contains terms like $b^{\dagger}_{k_h}b^{\dagger}_{-k_h}b_0b_0$
that would connect it to hidden momenta. What survives of the multi-time (Tomonaga--Schwinger)
picture is the isometry fact: the metric, the frame and the connection do not depend on
$x_5,x_6,x_7$, so the hidden translations commute with the $x_4$-evolution, $[P_4,P_h]=0$
\cfprose{VII \S28}. \cfproved{VII \S28}
\begin{itemize}
  \item \cfL{ADMISSIBILITY (1) [THE RESULT]: E\_k\textasciicircum{}2 == (k\_s\textasciicircum{}2 - k\_h\textasciicircum{}2 + m\textasciicircum{}2) ID16 for every k: omega\textasciicircum{}2 = k\_s\textasciicircum{}2 - k\_h\textasciicircum{}2 + m\textasciicircum{}2}
  \item \cfL{ADMISSIBILITY (2) [THE RESULT]: for admissible k (k\_h = 0), [J, E\_k] == 0, E\_k is J-Hermitian, and G J == ID16 is positive: J quantizes every admissible momentum positively}
  \item \cfL{ADMISSIBILITY (3) [THE RESULT]: below threshold (k = 5 e1 + 3 e5, m = 3) the boosted J' = S\textasciicircum{}-1 J S commutes with E\_k and beta, J'\textasciicircum{}2 == 1, and G J' == S\textasciicircum{}2 is Hermitian positive definite}
  \item \cfL{ADMISSIBILITY (4) [THE RESULT]: AT threshold (k = 4 e1 + 5 e5, m = 3) E\_k != 0 but E\_k\textasciicircum{}2 == 0: nilpotent, not diagonalizable (rank 8: Jordan blocks of size 2)}
  \item \cfL{ADMISSIBILITY (5) [THE RESULT]: ABOVE threshold (k = (0, 3/4, 1, 0 | k5 = 2), m = 1) omega\textasciicircum{}2 == -23/16: every eigenvalue of E\_k is +-i Sqrt[23]/4 = +-1.19896 i and every eigenvector is G-NEUTRAL (u\textasciicircum{}dagger G u == 0)}
  \item \cfL{ADMISSIBILITY [has content]: the general identity [E\_k, beta] == 2 i T\_k, T\_k = Sum\_a k\_a T16[a], for every k (admissible and hidden components alike)}
  \item \cfL{ADMISSIBILITY (6) [THE RESULT]: traceless certificate -- every X commuting with beta, the four admissible E\_k and one hidden-momentum E\_k has Tr(G X) == 0, so no positive G J' exists; J itself is not in that commutant}
  \item \cfL{ADMISSIBILITY (6) [THE RESULT]: the analytic proof for a purely hidden momentum -- [E\_k, beta] == 2 i k5 T16[5]; T16[5] anticommutes with G and T16[5]\textasciicircum{}dagger G T16[5] == -G (so a J' commuting with T16[5] has T5\textasciicircum{}dagger (G J') T5 == -(G J'))}
  \item \cfL{ADMISSIBILITY (7) [THE RESULT]: J is UNIQUE -- the joint commutant of beta and the admissible E\_k is 8-dimensional, contains J, and every element of it commutes with J (so X\textasciicircum{}2 == 1 with J X > 0 forces X == J)}
  \item \cfL{ADMISSIBILITY [has content]: [P\_4, P\_h] == 0 -- the metric, the frame and the spin connection do not depend on x5, x6, x7}
\end{itemize}
The theorem statement combines items (1)--(7); the items are asserted individually, the ``if and
only if'' is their combination \cfprose{VII \S28}.

\subsection{The Hamiltonian on the canonical frame, and the spin connection}
\label{sec:6.11}

The Hamiltonian \textbf{density} of the symmetrized Lagrangian,
$\Pi_\Psi\,\partial_4\Psi+\partial_4\bar\Psi\,\Pi_{\bar\Psi}-L$, is
\[
  \sqrt g\,\Bigl[-\frac{1}{2H}\sum_{k\neq4}\bigl(\bar\Psi\gamma^k\partial_k\Psi-\partial_k\bar\Psi\,\gamma^k\Psi\bigr)+V(s)\Bigr]
\]
with \textbf{no spin connection} in it, for fields of all eight coordinates. (This corrects design
revision~1, which said that the connection ``enters the Hamiltonian''.) Varying it gives the
one-particle operator. In the admissible sector ($\Psi=\Psi(x_0,\dots,x_4)$) the field equation is
$i\,\Gt\,\partial_4\Psi=h\Psi$ with
\[
  h\Psi=\sqrt g\,\sigma_{16}\Bigl[V'(s)\,\Psi-\frac1H\Bigl(\sum_{j=0}^{3}\gamma^j\partial_j\Psi+\sum_\mu\gamma^\mu\Gamma_\mu\Psi\Bigr)\Bigr]
\]
($V'(s)$ is the mean-field $V'(\langle s\rangle)$ here). The connection appears here, through the
equation of motion, and it has a job: since
$\tfrac12\sum_{\mu\neq4}\partial_\mu\bigl(\sqrt g\,\sigma_{16}\gamma^\mu\bigr)=\sqrt g\,\sigma_{16}\gamma^\mu\Gamma_\mu$
and $\partial_4\bigl(\sqrt g\,\sigma_{16}\gamma^4\bigr)=0$, the derivative part of $h$ is
$-\frac1H\sum_j\bigl(A_j\partial_j+\tfrac12\,\partial_jA_j\bigr)$ with $A_j=\sqrt g\,\sigma_{16}\gamma^j$
real \textbf{antisymmetric}, and for such an operator $\langle f,hg\rangle-\langle hf,g\rangle$ is a
total derivative. The connection term is exactly the term that makes the one-particle Hamiltonian
Hermitian. \cfproved{VII \S28}
\begin{itemize}
  \item \cfL{HAMILTONIAN [THE RESULT]: the Hamiltonian density Pi\_Psi d\_4 Psi + d\_4 Psibar Pi\_Psibar - L of the symmetrized Lagrangian is Sqrt[g][-(1/(2H)) Sum\_\{k != 4\}(Psibar gamma\textasciicircum{}k d\_k Psi - d\_k Psibar gamma\textasciicircum{}k Psi) + V]: NO spin connection (fields of all eight coordinates)}
  \item \cfL{HAMILTONIAN [THE RESULT]: (Sqrt[g]/H) sigma16 (gamma\textasciicircum{}mu D\_mu Psi - H V' Psi) == i Gtilde d\_4 Psi - h Psi for admissible Psi(x0..x4): the field equation is i Gtilde d\_4 Psi = h Psi}
  \item \cfL{HAMILTONIAN [THE RESULT]: (1/2) Sum\_\{mu != 4\} d\_mu(Sqrt[g] sigma16 gamma\textasciicircum{}mu) == Sqrt[g] sigma16 Sum gamma\textasciicircum{}mu Gamma\_mu, and d\_4(Sqrt[g] sigma16 gamma\textasciicircum{}4) == 0: the connection term is the one that makes h Hermitian}
  \item \cfL{HAMILTONIAN [has content]: A\_0 = Sqrt[g] sigma16 gamma\textasciicircum{}0 is real antisymmetric, and for h\_0 = -(A\_0 d\_0 + (1/2) A\_0'):  <f, h\_0 g> - <h\_0 f, g> == -d\_0(f\textasciicircum{}dagger A\_0 g), a pure boundary term}
\end{itemize}

$J$ commutes with $h$: every matrix in $h$ ($\sigma_{16}$, $\sigma_{16}\gamma^j$ for $j=0..3$,
$\sigma_{16}\gamma^\mu\Gamma_\mu$) commutes with $J$, for every $a_4$, every $x_0$-profile and every
$V$. A hidden derivative term would anticommute (control). The $x_4$-dependence of $h$ enters only
through the observed frame factor $1/q$; it mixes positive and negative frequencies (a
\textbf{Bogoliubov} transformation, particle production by the $x_4$-dependence of $a_4$)
\textbf{within} the admissible sector, and since it commutes with $J$ that mixing is unitary on the
Fock space. \cfproved{VII \S28}
\begin{itemize}
  \item \cfL{HAMILTONIAN [THE RESULT]: [J, h] == 0 on the admissible sector, for every a4, every x0-profile and every V' -- including the spin-connection term}
  \item \cfL{HAMILTONIAN [control]: a hidden derivative term breaks it -- [J, h] != 0 once Psi depends on x5 (witnessed on Psi = x5 e1)}
  \item \cfL{HAMILTONIAN [has content]: the x4-dependence of h enters only through 1/q: d\_4 h Psi == (H a4'[H x4]) x (the observed-direction part of h), which commutes with J}
\end{itemize}

\subsection{The \texorpdfstring{$x_0$}{x0} direction and the boundary condition at the wall}
\label{sec:6.12}

For the rescaled field $\Psi'=\Psi/\sqrt{\sin 6Hx_0}$ the measure becomes flat in the proper distance
$z=-\ln\bigl(\cos 6Hx_0\bigr)/(6H)$ along $x_0$:
\[
  \sqrt g\,dx_0\,|\Psi|^2=|\Psi'|^2\,dz\qquad(\sec\cdot\sin=\tan=dz/dx_0).
\]
$z$ runs over $[0,\infty)$ as $6Hx_0$ runs over $[0,\pi/2)$; the $x_0$ part of the rescaled Dirac
operator is exactly $\Tsix{0}\,\partial_z$, and the frame factors are bounded,
$\sin(6Hx_0)=\sqrt{1-e^{-12Hz}}$. So the one-particle space is $L^2(dz)$. The operator is
\textbf{regular} at the wall $z=0$, so a boundary condition is needed there, and
\textbf{limit-point} at $z\to\infty$ (by the standard theorem that a one-dimensional Dirac system with
locally integrable coefficients is always limit-point at an infinite endpoint), so no condition is
needed or allowed there. The condition at the wall, in this real-Clifford convention, is the
bag-type
\[
  \Tsix{0}\,\Psi'=+\Psi'\quad\text{or}\quad\Tsix{0}\,\Psi'=-\Psi'\qquad\text{at } z=0
  \qquad(\text{no factor } i\text{, because } \Tsix{0}^2=+1).
\]
It kills the normal flux $\Psi'^{\dagger}\beta\,\Tsix{0}\,\Psi'$ (because $\{\beta,\Tsix{0}\}=0$) and
commutes with $J$. It is part of the canonical quantization. \cfproved{VII \S28} It also gives
$\bar\Psi\Psi=0$ at the wall [derived here: $\sigma_{16}=J\beta$ anticommutes with $\Tsix{0}$, because
$J$ commutes and $\beta$ anticommutes with it, so $P\sigma_{16}P=0$ on either eigenspace $P$ of
$\Tsix{0}$; design review QK-9].
\begin{itemize}
  \item \cfL{BOUNDARY [has content]: the x0 direction in proper distance z = -Log[Cos[6 H x0]]/(6 H): dz/dx0 == Tan[6Hx0] == Sec Sin (so Sqrt[g] dx0 |Psi|\textasciicircum{}2 = |Psi'|\textasciicircum{}2 dz), z(0) == 0 and z -> infinity at 6 H x0 -> Pi/2, gamma\textasciicircum{}0 d\_0 == T16[0] d\_z, and Sin[6Hx0] == Sqrt[1 - Exp[-12 H z]] (bounded frame factors)}
  \item \cfL{BOUNDARY [THE RESULT]: the bag condition T16[0] Psi' = +-Psi' at the wall -- T16[0]\textasciicircum{}2 == 1, \{beta, T16[0]\} == 0, the normal-flux form vanishes on both eigenspaces (P beta T16[0] P == 0), and [J, T16[0]] == 0}
\end{itemize}
(The full family of self-adjoint, Lorentz-covariant, $J$-sector-preserving wall conditions is
$P=\Tsix{0}\,Q$, with $Q$ any Hermitian involution in the 8-dimensional commutant of $\Tsix{0..4}$,
which contains $J$; the uniform choices $Q=\pm1$ are the two above. That refinement belongs to the
treatment of the states along the hidden coordinate in the coupled system; it was adopted by the
design review, DFT-8, and is not asserted in Part~VII.)

\subsection{Krein modes, and why the mode basis must consist of \texorpdfstring{$J$}{J}-eigenvectors}
\label{sec:6.13}

For an admissible momentum $k=(k_0,k_1,k_2,k_3)$ and mass $m$ (flat local analysis; on the canonical
frame $k$ is the physical momentum, the coordinate momentum along $x_1..x_3$ divided by $q$), put
$A=Gh=E_k$ and $w=\sqrt{k\cdot k+m^2}$. Then $A^2=w^2$, $[A,J]=0$, $A$ is Hermitian, and the four
Hermitian projectors
\[
  P_{s,\eta}=\tfrac14\,\bigl(\IDs+s\,A/w\bigr)\bigl(\IDs+\eta J\bigr),\qquad s,\eta=+1,-1,
\]
have rank 4 each (trace 4), are mutually orthogonal, sum to $\IDs$, and
$\sum\eta\,P_{s,\eta}=J=G^{-1}$: the \textbf{completeness relation}
$\sum_u\eta_u\,u\,u^{\dagger}=G^{-1}$ of any $J$-orthonormal mode basis. So each frequency $\pm w$ is
8-fold degenerate, and the $G$-form restricted to each frequency eigenspace has signature $(4,4)$.
For \textbf{every} mode $u$ in the image of $P_{s,\eta}$, normalized by $u^{\dagger}Gu=\eta$, the
projector identities $PXP=cP$ give, exactly and for arbitrary $k$ and $m$,
\begin{align*}
  \eta\,u^{\dagger}h\,u&=s\,w=\omega_u &&\text{(energy = frequency)},\\
  \eta\,u^{\dagger}\sigma_{16}\,u&=m/\omega_u &&\text{(scalar density)},\\
  \eta\,u^{\dagger}\,(\partial h/\partial k_j)\,u&=k_j/\omega_u=\partial\omega_u/\partial k_j
    &&\text{(velocity: Hellmann--Feynman)},
\end{align*}
with $\partial h/\partial k_j=-i\,\sigma_{16}\Tsix{j}$.
The identities are proved with $w$ a symbol and $w^2$ reduced to $k\cdot k+m^2$ (exact polynomial
remainder); nothing is numerical. \cfproved{VII \S28}
\begin{itemize}
  \item \cfL{KREIN MODES [has content]: A = G h satisfies A\textasciicircum{}2 == w\textasciicircum{}2 ID16 and [A, J] == 0, and A is Hermitian (G and h commute) for every admissible k and m}
  \item \cfL{KREIN MODES [THE RESULT]: the four P\_\{s,eta\} are projectors of trace (rank) 4, mutually orthogonal, summing to ID16, and Sum eta P\_\{s,eta\} == J == G\textasciicircum{}-1 (completeness)}
  \item \cfL{KREIN MODES [THE RESULT]: for every mode, eta u\textasciicircum{}dagger h u == s w (energy = frequency), eta u\textasciicircum{}dagger sigma16 u == m/omega\_u, eta u\textasciicircum{}dagger (dh/dk\_j) u == k\_j/omega\_u: P X P == c P with c = eta s w, eta s m/w, eta s k\_j/w}
  \item \cfL{KREIN MODES [has content]: dh/dk\_j == -i sigma16 T16[j], and d omega/d k\_j == k\_j/omega for omega = +-w}
  \item \cfL{KREIN MODES [THE RESULT]: every mode of P\_\{s,eta\} is a J-eigenvector: J P\_\{s,eta\} == eta P\_\{s,eta\}}
\end{itemize}

\paragraph{The mode basis must consist of $J$-eigenvectors.} The Krein prescription
$b_u^{\dagger}:=\eta_u\,b_u^{\ddagger}$ depends on the basis. It reproduces the $J$-involution
$\Psi^{\dagger}=\Psi^{\ddagger}J$ only if every mode is a $J$-eigenvector, $Ju=\eta_u u$, which is what
the $P_{s,\eta}$ deliver ($JP_{s,\eta}=\eta P_{s,\eta}$), and which is possible because
$[J,Gh_k]=0$ for momenta along $x_0..x_3$. A $G$-orthonormal basis mixed by a Krein boost inside one
frequency eigenspace still diagonalizes the evolution, but it induces the involution
$J_U=G\,U\,U^{\dagger}\neq J$, and under it $s$ is \textbf{not} Hermitian (control, exact, at $k=3e_1$,
$m=4$, with the boost $\cosh=5/4$, $\sinh=3/4$). A generic generalized-eigenvector routine does not
choose $J$-eigenvectors; the Fock construction of \secref{sec:6.14} does. \cfproved{VII \S28}
\cfL{KREIN MODES [control]: a Krein-boosted pair u1 = (5/4)u+ + (3/4)u-, u2 = (3/4)u+ + (5/4)u- in the omega = +5 eigenspace is still G-orthonormal and still consists of modes, but not of J-eigenvectors; the induced J\_U = G U U\textasciicircum{}dagger differs from J, and under it s is NOT Hermitian}

\subsection{A finite Fock space at one momentum (Jordan--Wigner): the Dirac sea, 8 particles and 8 antiparticles}
\label{sec:6.14}

At the momentum $k=3e_1$ with $m=4$ (so $\omega=5$ exactly) the sixteen modes of one momentum are
represented on $\mathbb C^{2^{16}}=\mathbb C^{65536}$ by Jordan--Wigner operators $f_1..f_{16}$
(\cfcode{SparseArray}, exact integers), which obey the positive CAR $\{f_i,f_j^{\dagger}\}=\delta_{ij}$,
$\{f_i,f_j\}=0$. Two realizations are checked.

\paragraph{Field level.} $\Psi_a:=f_a$ (the factor $H/\sqrt g$ set to 1), $\Psi^{\dagger}=f^{\dagger}$,
$\Psi^{\ddagger}:=\Psi^{\dagger}J$. Then:
\begin{itemize}
  \item $\{\Psi_a,\Psi^{\ddagger}_b\}=J_{ab}=(G^{-1})_{ab}$: exactly the Dirac-bracket anticommutator,
        realized on a positive Fock space;
  \item the free Hamiltonian $H_{\mathrm{free}}=\Psi^{\ddagger}h\Psi=\Psi^{\dagger}(Jh)\Psi$ and
        $s=\Psi^{\ddagger}\sigma_{16}\Psi=\Psi^{\dagger}\beta\Psi$ are Hermitian operators, and so is
        the interacting $H_{\mathrm{free}}+\tfrac37\,s^2$; $p=\Psi^{\ddagger}C_+\Psi$ is
        anti-Hermitian (control); the one-body matrix $Jh$ has eigenvalues $+5$ and $-5$, eight
        times each;
  \item the \textbf{Heisenberg equation} holds: $[H_{\mathrm{free}},\Psi_a]=-(G^{-1}h\Psi)_a$ for all
        16 components, i.e.\ $i\,\partial_4\Psi=i\,[H,\Psi]$ reproduces $iG\,\partial_4\Psi=h\Psi$;
  \item with the \textbf{opposite} sign of the anticommutator ($\Psi^{\ddagger}:=-\Psi^{\dagger}J$)
        the same construction gives $H'=-H_{\mathrm{free}}$ and the time-\textbf{reversed} equation
        (control). This is the convention-free confirmation of the sign derived with the Dirac
        bracket.
\end{itemize}

\paragraph{Mode level.} $b_u:=f_u$ for the sixteen $J$-eigenmodes of \secref{sec:6.13} ($u=1..8$:
$\omega=+5$; $u=9..16$: $\omega=-5$; $Ju=\eta_u u$ with $\eta=+1$ for four of each and $-1$ for the
other four):
\begin{itemize}
  \item the Krein conjugate $b_u^{\ddagger}:=\eta_u\,b_u^{\dagger}$ obeys
        $\{b_u,b_v^{\ddagger}\}=\eta_u\,\delta_{uv}$, the indefinite CAR that the field anticommutator
        demands through $\sum_u\eta_u\,u\,u^{\dagger}=G^{-1}$, and the $J$-involution
        $b^{\dagger}=\eta\,b^{\ddagger}$ restores $\{b,b^{\dagger}\}=1$;
  \item the naive Krein ``norm'' $\langle0|\,b_u b_u^{\ddagger}\,|0\rangle=\eta_u$ is
        \textbf{negative} for $\eta=-1$: that is the indefinite metric the $J$-involution removes (the
        Fock norm $\langle0|\,b_u b_u^{\dagger}\,|0\rangle=1$);
  \item $\mathrm{Ham}=\sum_u\omega_u\,\eta_u\,b_u^{\ddagger}b_u=\sum_u\omega_u\,b_u^{\dagger}b_u$, and
        $[\mathrm{Ham},b_u]=-\omega_u\,b_u$ ($b_u(x_4)=b_u\,e^{-i\omega_u x_4}$, the time dependence
        of the classical mode);
  \item the spectrum on the 65536 states: a \textbf{unique ground state} at $-8\omega=-40$, the
        \textbf{Dirac sea} (all eight negative-frequency modes filled); the normal-ordered
        ${:}\mathrm{Ham}{:}=\mathrm{Ham}+40\geq0$;
  \item the first excited level, $\omega=5$ above the sea (energy $-35$), is \textbf{16-fold}:
        \textbf{8 particle states} (one positive-frequency mode added, $N=+1$) and \textbf{8
        antiparticle states} (one negative-frequency mode emptied, a hole, $N=-1$), with
        $N=\sum_u b_u^{\dagger}b_u-8$;
  \item in every Fock basis state tested (the sea and the 16 first excited states) the expectation
        of $b_u^{\ddagger}b_v$ is $\eta_u\,n_u\,\delta_{uv}$: the rule
        $\langle\Psi^{\ddagger}M\Psi\rangle=\sum_{\text{occupied}}\eta_u\,u^{\dagger}Mu$ used in
        \secref{sec:8}.
\end{itemize}
\cfproved{VII \S28}
\begin{itemize}
  \item \cfL{FOCK [solver regression]: the Jordan-Wigner operators obey the positive CAR \{f\_i, f\_j\textasciicircum{}dagger\} == delta\_ij, \{f\_i, f\_j\} == 0 on C\textasciicircum{}65536}
  \item \cfL{FOCK (field level) [THE RESULT]: \{Psi\_a, Psi\textasciicircum{}ddag\_b\} == J\_ab == (G\textasciicircum{}-1)\_ab with Psi = f, Psi\textasciicircum{}ddag = f\textasciicircum{}dagger J -- the Dirac-bracket anticommutator, realized on a positive Fock space}
  \item \cfL{FOCK (field level) [THE RESULT]: J h is Hermitian with eigenvalues +5 (8) and -5 (8); H\_free and s are Hermitian operators, so is H\_free + (3/7) s\textasciicircum{}2; p = Psi\textasciicircum{}ddag C+ Psi is ANTI-Hermitian}
  \item \cfL{FOCK (field level) [THE RESULT]: Heisenberg -- [H\_free, Psi\_a] == -(G\textasciicircum{}-1 h Psi)\_a for all 16 components, so i d\_4 Psi = i [H, Psi] IS the field equation i G d\_4 Psi = h Psi}
  \item \cfL{FOCK (field level) [control]: with the opposite sign of the anticommutator the Hamiltonian is -H\_free and [H', Psi\_a] == +(G\textasciicircum{}-1 h Psi)\_a: the time-REVERSED equation}
  \item \cfL{FOCK (mode level) [THE RESULT]: \{b\_u, b\_v\textasciicircum{}ddag\} == eta\_u delta\_uv (the indefinite CAR demanded by Sum eta u u\textasciicircum{}dagger = G\textasciicircum{}-1), and the J-involution b\textasciicircum{}dagger = eta b\textasciicircum{}ddag restores \{b, b\textasciicircum{}dagger\} == 1}
  \item \cfL{FOCK (mode level) [has content]: the naive Krein 'norm' <0| b\_u b\_u\textasciicircum{}ddag |0> equals eta\_u -- NEGATIVE for the four eta = -1 modes of each frequency -- while the Fock norm <0| b\_u b\_u\textasciicircum{}dagger |0> == 1}
  \item \cfL{FOCK (mode level) [THE RESULT]: Sum omega\_u eta\_u b\_u\textasciicircum{}ddag b\_u == Sum omega\_u b\_u\textasciicircum{}dagger b\_u, and [Ham, b\_u] == -omega\_u b\_u for every u (b\_u(x4) = b\_u e\textasciicircum{}(-i omega\_u x4))}
  \item \cfL{FOCK (mode level) [THE RESULT]: the spectrum -- a UNIQUE ground state at -8 omega = -40, the Dirac sea (modes 9..16 filled, 1..8 empty); :Ham: = Ham + 40 >= 0}
  \item \cfL{FOCK (mode level) [THE RESULT]: the first excited level, omega = 5 above the sea, is 16-fold: 8 particle states (N = +1) and 8 antiparticle states (N = -1), with N = Sum b\textasciicircum{}dagger b - 8}
  \item \cfL{FOCK (mode level) [THE RESULT]: in every Fock basis state tested (the sea and the 16 first excited states) <n| b\_u\textasciicircum{}ddag b\_v |n> == eta\_u n\_u delta\_uv -- the expectation rule of Section 29}
\end{itemize}
The Jordan--Wigner CAR check on the 65536-dimensional space took 0.749\,s and the Heisenberg check
1.868\,s; the whole Fock cell (Input cell 191) took 9.858\,s
(\file{claude-fable/run_fermion_fable_part7.log}, lines \cfcode{[0.749 s]}, \cfcode{[1.868 s]},
\cfcode{CELL 191  t=9.858 s}).

\subsection{The particle number}
\label{sec:6.15}

The particle-number density is
\[
  n^4=-\frac{i}{H}\,\bar\Psi\gamma^4\Psi=\frac1H\,\Psi^{\ddagger}G\Psi=\frac1H\,\Psi^{\dagger}\Psi\;\geq0
\]
and its mode form $\sum_u\eta_u\,b_u^{\ddagger}b_u=\sum_u b_u^{\dagger}b_u$ is the particle number
(normal-ordered: particles minus antiparticles). \cfproved{VII \S28} \cfL{CHARGE [THE RESULT]: n\textasciicircum{}4 = -(i/H) Psibar gamma\textasciicircum{}4 Psi == (1/H) Psi\textasciicircum{}ddag G Psi == (1/H) Psi\textasciicircum{}dagger Psi, whose mode form Sum\_u eta\_u b\_u\textasciicircum{}ddag b\_u == Sum\_u b\_u\textasciicircum{}dagger b\_u is the particle number (normal-ordered: particles - antiparticles)}

\subsection{The dynamical-frame field \texorpdfstring{$\chi$}{chi}}
\label{sec:6.16}

On a frame whose volume factor depends on $x_4$ the field to quantize is
$\chi=(\sqrt g/H)^{1/2}\,\Psi$, with $\{\chi,\chi^{\ddagger}\}=G\,\delta^7$ (\secref{sec:6.4}). On the
canonical frame $\sqrt g=\sec 6Hx_0$ does not depend on $x_4$, and quantizing $\chi$ changes nothing
(\cfL{ANTICOMMUTATOR [has content]: the rescaled field Psi' = Psi/Sqrt[Sin[6 H x0]] has \{Psi', Psi'\textasciicircum{}ddag\} == H Cot[6 H x0] J delta\textasciicircum{}7, and Sqrt[g] is x4-independent (so quantizing chi = (Sqrt[g]/H)\textasciicircum{}(1/2) Psi changes nothing on the canonical frame)}). On the time-dependent warped frames of the interacting system
(Part~VIII), the rescaled $\chi=\Psi\big/\bigl(\sqrt{\sin 6Hx_0}\,(A^3B^3C)^{-1/2}\bigr)$ has
$\{\chi,\chi^{\ddagger}\}=H\cot(6Hx_0)\,G$, independent of $x_4$ and of the scale factors
(\cfL{RESCALING [THE RESULT]: \{Psi, Psi\textasciicircum{}ddag\} == (H/Sqrt[|det g|]) G on the warped frame (Section 28's cfAntiPsiPsiDd at c = Sqrt[|det g|]/H), and for chi = Psi/f, f = Sqrt[Sin] (A\textasciicircum{}3 B\textasciicircum{}3 C)\textasciicircum{}(-1/2): \{chi, chi\textasciicircum{}ddag\} == H Cot[6 H x0] G -- the Sqrt[det g] cancels, independent of x4 and of A, B, C}).

% ===============================================================================================
\section{The energy--momentum tensor operator}
\label{sec:7}

This is the first half of Section~29 of the notebook (Input cells 192--194, 16 assertions).

\subsection{The definition, and where it comes from}
\label{sec:7.1}

The energy--momentum tensor of fermion fable is
\[
  \boxed{\;\hat T_{\mu\nu}=-\frac{1}{4H}\Bigl[\bar\Psi\gamma_\mu D_\nu\Psi-(D_\nu\bar\Psi)\gamma_\mu\Psi+(\mu\leftrightarrow\nu)\Bigr]+g_{\mu\nu}\hat L,
  \qquad \gamma_\mu=g_{\mu\nu}\gamma^\nu\;}
\]
normal-ordered with respect to the $J$-vacuum of \secref{sec:6} when it is an operator. It is the
\textbf{Hilbert} (symmetric-tetrad) tensor of the symmetrized action,
$T^{\mu\nu}=\frac{2}{\sqrt g}\,\frac{\delta S}{\delta g_{\mu\nu}}$, and it equals the covariant
expression above \textbf{off shell}. The notebook checks this the direct way: it perturbs the
canonical frame by a symmetric first-order tetrad perturbation $h(x_0,x_4)$ in the $(\mu,\nu)$ slot,
recomputes to first order everything that depends on the frame (the metric, the Christoffel symbols,
the spin connection from the vielbein postulate, the curved gammas, $\sqrt g$), and takes the
Euler--Lagrange derivative with respect to $h$. Three facts come out:
\begin{enumerate}
  \item The result is the tensor above, checked on the diagonal $(x_1,x_1)$, $(x_4,x_4)$ and off the
        diagonal $(x_1,x_4)$, $(x_0,x_5)$, $(x_5,x_4)$.
  \item The variation of the spin connection contributes \textbf{exactly nothing}. Its totally
        antisymmetric part does not change under a symmetric tetrad variation, and the symmetrized
        Lagrangian sees the connection only through that part, $\tfrac12\,\omega_{[cab]}\,\gamma^{cab}$.
        The covariant-derivative terms of the tensor come instead from the frame dependence of
        $\{\gamma^\mu,\Gamma_\mu\}$, which vanishes on the background while its first variation does
        not.
  \item The Hilbert tensor is \textbf{not} the partial-derivative tensor $\Tpar$ (control, at
        $(x_1,x_4)$): one never varies a partial-derivative Lagrangian with respect to the frame.
\end{enumerate}
\cfproved{VII \S29}
\begin{itemize}
  \item \cfL{EMT [THE RESULT]: the Hilbert (symmetric-tetrad) tensor of the symmetrized action IS That, OFF shell -- components (x1,x1), (x4,x4), (x1,x4), (x0,x5), (x5,x4)} (the five components took 0.693\,s)
  \item \cfL{EMT [has content]: the spin-connection variation contributes NOTHING -- dropping it leaves the Hilbert tensor unchanged at (x1,x4), (x0,x5), (x5,x4)}
  \item \cfL{EMT [control]: the Hilbert tensor is NOT the partial-derivative tensor: at (x1,x4) they differ -- witnessed on constant spinors (Psibar = e1, Psi = e11: (T16[4] T16[1] T16[0])\_\{1,11\} != 0)}
\end{itemize}

\subsection{Symmetric and Hermitian}
\label{sec:7.2}

$\hat T$ is symmetric, and it is Hermitian under the classical conjugation: real for $\Psi=u+iv$,
$\bar\Psi=(u-iv)^{T}\sigma_{16}$, in all 64 components. \cfproved{VII \S29}
\begin{itemize}
  \item \cfL{EMT [definition]: That is symmetric}
  \item \cfL{EMT [THE RESULT]: That is Hermitian under the classical conjugation: real for Psi = u + i v, Psibar = (u - i v)\textasciicircum{}T sigma16, all 64 components}
\end{itemize}

\subsection{Conservation on shell, and why the commuting proxy proves it for the fermion}
\label{sec:7.3}

Solve the two field equations for the $x_4$-derivatives:
\begin{align*}
  \partial_4\Psi&=F=-\gamma^4\bigl(H\,V'(s)\,\Psi-\gamma^0\partial_0\Psi-(\gamma^\mu\Gamma_\mu)\,\Psi\bigr),\\
  \partial_4\bar\Psi&=\bar F=\bigl(H\,V'(s)\,\bar\Psi+\partial_0\bar\Psi\,\gamma^0-\bar\Psi\,(\Gamma_\mu\gamma^\mu)\bigr)\gamma^4,
\end{align*}
and let \cfcode{cfOnShellC} replace every $x_4$-derivative of either field (of any order) by the
corresponding derivative of $F$ or $\bar F$, to a fixed point. Then
$\nabla^\mu\hat T_{\mu\nu}=0$ in all eight components, for generic $\Psi(x_0,x_4)$,
$\bar\Psi(x_0,x_4)$ and generic $V$.

\paragraph{Why a commuting computation proves the Grassmann statement.} After the substitution,
every term of the divergence is a bilinear with $\bar\Psi$ (or a derivative of it) on the left and
$\Psi$ (or a derivative) on the right, multiplied by functions of the \textbf{even} composites $s$ and
$\partial s$ (through $V'$, $V''$ and their derivatives). No two odd factors are ever exchanged, so
the same algebra holds verbatim for anticommuting components. (The argument is not valid for
arbitrary quartic identities; it is valid here because of that structure.) A configuration that is
not a solution has a non-zero divergence (control).

\paragraph{The connection does not drop out of the tensor.} $\hat T$ differs from the
partial-derivative tensor by
\[
  \hat T_{\mu\nu}-T^{\mathrm{par}}_{\mu\nu}=-\frac{1}{4H}\,\bar\Psi\bigl(\{\gamma_\mu,\Gamma_\nu\}+\{\gamma_\nu,\Gamma_\mu\}\bigr)\Psi ,
\]
which vanishes on the diagonal (so $\rho$ and every pressure are the same in both) but not off it,
and only $\hat T$ is conserved. (The symbolic divergence of $\Tpar$ is not attempted: it takes
minutes, and the off-diagonal difference already shows the two tensors are different.)
\cfproved{VII \S29}
\begin{itemize}
  \item \cfL{EMT CONSERVATION [definition]: F and Fbar solve the field equation and the adjoint equation for the x4-derivatives}
  \item \cfL{EMT CONSERVATION [THE RESULT]: nabla\textasciicircum{}mu That\_\{mu nu\} == 0 ON SHELL in all eight components, for generic Psi(x0, x4), Psibar(x0, x4) and generic V} (the divergence took 0.767\,s and its on-shell
        reduction 2.599\,s)
  \item \cfL{EMT CONSERVATION [control]: OFF shell the divergence is not zero -- witnessed on Psi = (x4 + x0\textasciicircum{}2) e1 + e13, Psibar = e1 + x4 e5 with V = s\textasciicircum{}2, which is not a solution}
  \item \cfL{EMT [THE RESULT]: That - T\_par == -(1/(4H)) Psibar (\{gamma\_mu, Gamma\_nu\} + \{gamma\_nu, Gamma\_mu\}) Psi: equal on the diagonal (every \{gamma\_mu, Gamma\_mu\} vanishes), different off it (the (x1, x4) witness above)}
\end{itemize}
The statement is proved for fields of $(x_0,x_4)$. That is the family in which $\rho$, $P$ and the
pre-universe results of \secref{sec:8} are evaluated.

\subsection{On shell: the Lagrangian, the density, the pressures and the trace}
\label{sec:7.4}

With \cfcode{cfOnShellC}, for generic $\Psi(x_0,x_4)$, $\bar\Psi(x_0,x_4)$ and generic $V$:
\begin{align*}
  \frac{1}{2H}\Bigl[\bar\Psi\gamma^\mu D_\mu\Psi-(D_\mu\bar\Psi)\gamma^\mu\Psi\Bigr]&=s\,V'(s)
    \qquad\text{(the kinetic bilinear; NOT zero)},\\
  \hat L&=s\,V'(s)-V(s),\\
  \rho=\hat T_{44}&=V(s)+K_h,\qquad K_h=-\frac{1}{2H}\bigl(\bar\Psi\gamma^0\partial_0\Psi-\partial_0\bar\Psi\,\gamma^0\Psi\bigr),\\
  T^i{}_i=T^h{}_h&=s\,V'(s)-V(s)\\
    &\qquad(i=1,2,3\text{ observed; } h=5,6,7\text{ hidden; no sum}),\\
  T^0{}_0&=s\,V'(s)-V(s)+K_h,\\
  T^\mu{}_\mu&=7\,s\,V'(s)-8\,V(s)\qquad(\text{off shell: } 7L_{\mathrm{kin}}-8V).
\end{align*}
The observer is $u=\partial/\partial x_4$, so $\rho=T_{44}$ (the array element \cfcode{[[5,5]]}) and
the pressure along a direction $j$ is $P_j=T^j{}_j$ (no sum). $P_i=P_h$ holds because
$\{\gamma^\mu,\Gamma_\mu\}=0$ for \textbf{each} direction separately: the kinetic parts of $T^i{}_i$
and $T^h{}_h$ vanish for fields of $(x_0,x_4)$, and both reduce to $\hat L$. The trace is consistent
with the list: $T^4{}_4=-\rho$, so $-(V+K_h)+(sV'-V+K_h)+6\,(sV'-V)=7sV'-8V$ \cfderived. In the real
commuting limit $K_h$ is Part~VI's \cfcode{cfKh} (\secref{sec:2.12.4}), and for the rescaled solutions
$\Psi=\sqrt{\sin 6Hx_0}\;\Psi'(x_4)$ (and the same for $\bar\Psi$) $K_h=0$ and $\rho=V(s)$: Part~VI's
formulas, now for the complex field. \cfproved{VII \S29}
\begin{itemize}
  \item \cfL{EMT ON SHELL [THE RESULT]: the kinetic bilinear == s V'(s) (not zero) and Lhat == s V'(s) - V(s)}
  \item \cfL{EMT ON SHELL [THE RESULT]: rho = That\_44 == V(s) + K\_h}
  \item \cfL{EMT ON SHELL [THE RESULT]: T\textasciicircum{}i\_i == T\textasciicircum{}h\_h == s V'(s) - V(s) for the three observed and the three hidden directions, and T\textasciicircum{}0\_0 == s V'(s) - V(s) + K\_h}
  \item \cfL{EMT [has content]: the trace T\textasciicircum{}mu\_mu == 7 Lkin - 8 V off shell, == 7 s V'(s) - 8 V(s) on shell}
  \item \cfL{EMT [fidelity]: K\_h in the real commuting limit IS Part VI's cfKh, and K\_h == 0 on the rescaled solutions Sqrt[Sin] Psi'(x4), Sqrt[Sin] Psi'bar(x4), where rho == V(s)}
\end{itemize}

\subsection{Which components are Hermitian operators: the \texorpdfstring{$J$}{J}-parity table}
\label{sec:7.5}

For fields independent of $x_5,x_6,x_7$ (the admissible sector, here $\Psi(x_0,\dots,x_4)$,
$\bar\Psi(x_0,\dots,x_4)$), replace $\Psi\to J\Psi$ and $\bar\Psi\to\bar\Psi J$ ($J^2=1$, $s$ is
unchanged). A component that goes into itself is built from $J$-even matrices and is a
\textbf{Hermitian} operator; one that goes into minus itself is \textbf{anti-Hermitian}
(\secref{sec:6.8}). The notebook computes the parity of all 64 components (3.561\,s) and asserts the
table:

{\small
\begin{longtable}{@{}p{0.34\linewidth}p{0.16\linewidth}p{0.44\linewidth}@{}}
\toprule
\raggedright component & \raggedright $J$-parity & \raggedright operator \tabularnewline
\midrule
\endhead
\raggedright $\hat T_{\mu\nu}$, $\mu,\nu$ in $\{0..4\}$ & \raggedright even & \raggedright Hermitian \tabularnewline
\raggedright $\hat T_{\mu h}$, $\mu$ in $\{0..4\}$, $h$ in $\{5,6,7\}$ & \raggedright odd & \raggedright anti-Hermitian: purely imaginary expectation values, zero in every state invariant under the hidden rotations ($\hat T_{\mu h}$ is a vector under them) \tabularnewline
\raggedright $\hat T_{hh'}$, $h\neq h'$ & \raggedright zero identically & \raggedright --- \tabularnewline
\raggedright $\hat T_{hh}$ & \raggedright even & \raggedright Hermitian; $\hat T_{hh}=g_{hh}\hat L$ exactly, because $\{\gamma_h,\Gamma_h\}=0$ \tabularnewline
\bottomrule
\end{longtable}
}

The hidden momenta $P_h=\int\sqrt g\,\hat T^4{}_h$ are $J$-odd operators with zero expectation in the
$J$-eigenmode Fock states of \secref{sec:6} \cfprose{VII \S29}. The anti-Hermitian components are
consistent with $G_{\mu h}=0$ on the canonical frame. \cfproved{VII \S29}
\begin{itemize}
  \item \cfL{EMT HERMITICITY [THE RESULT]: J-parity table -- (mu, nu) in \{0..4\}: EVEN (Hermitian); (mu in 0..4, h in 5..7): ODD (anti-Hermitian); (h, h'), h != h': ZERO identically; (h, h): even}
  \item \cfL{EMT HERMITICITY [has content]: for x5..x7-independent fields That\_\{hh\} == g\_hh Lhat exactly, because \{gamma\_h, Gamma\_h\} == 0}
\end{itemize}

\subsection{The real commuting limit}
\label{sec:7.6}

Set $\Psi\to\psi$ (real, commuting) and $\bar\Psi\to\psi^{T}\sigma_{16}$. Then $\hat T$ becomes
Part~VI's covariant tensor $\Tcov$ of \secref{sec:2.12.3} exactly, all 64 components:
\cfL{FIDELITY [4] [fidelity]: the energy-momentum tensor That in the real limit IS Part VI's covariant tensor cfTCov, all 64 components}. The classical tensor of 2026-09-16 is therefore the commuting shadow of the
operator derived here.

% ===============================================================================================
\section{Density, pressure and equations of state}
\label{sec:8}

This is the second half of Section~29 of the notebook (Input cells 195--197, 23 assertions), with the
numerically stable forms used by the solver and the vacuum-energy caveat added from the design
review.

\subsection{The operators}
\label{sec:8.1}

The observer is $u=\partial/\partial x_4$. The energy-density and pressure \textbf{operators} are the
normal-ordered components of \secref{sec:7}:
\begin{align*}
  \rho_{\mathrm{op}}&={:}\hat T_{44}{:} &&\text{(the energy density measured by the observer)},\\
  P_{j,\mathrm{op}}&={:}\hat T^j{}_j{:}\quad(\text{no sum}) &&j=1,2,3\text{: the observed sheet; } j=0\text{: the hidden spacelike direction;}\\
  & &&h=5,6,7\text{: the hidden timelike sheet},\\
  w&=\langle\Pobs\rangle/\langle\rho\rangle &&\Pobs=\text{the pressure along the observed sheet}.
\end{align*}
On shell, for fields of $(x_0,x_4)$: $\rho=V(s)+K_h$, $P_i=P_h=s\,V'(s)-V(s)$,
$P_0=s\,V'(s)-V(s)+K_h$ (\secref{sec:7.4}). These are operator identities in the sense of
\secref{sec:5.9} (ordering by the Weyl rule; exact for linear $V$).

\subsection{Expectation values: the mode sums, and one particle}
\label{sec:8.2}

Use the \textbf{canonical normalization} $\chi=\Psi/\sqrt H$. Then $V(s)=V(H\sigma)=:W(\sigma)$ with
$\sigma=\langle\bar\chi\chi\rangle$, and $m=W'(\sigma)$ is the mass in the mode equation. On a
flat-frame plane wave $\chi=u\,e^{i(k\cdot x-\omega x_4)}$, $\bar\chi=\bar u\,e^{-i(k\cdot x-\omega x_4)}$
($\bar u=u^{\ddagger}\sigma_{16}$), the kinetic parts of the tensor are
\begin{align*}
  T_{44}+\hat L&=\omega\,\bar u\,(-i\,\Tsix{4})\,u=\omega\,u^{\ddagger}G\,u,\\
  T^j{}_j-\hat L&=k_j\,\bar u\,(-i\,\Tsix{j})\,u=k_j\,u^{\ddagger}(\partial h/\partial k_j)\,u &&(j=0..3,\text{ flat frame}).
\end{align*}
By the expectation rule
$\langle\Psi^{\ddagger}N\Psi\rangle=\sum_{\text{occupied}}\eta_u\,u^{\dagger}Nu$ (\secref{sec:6.14}) and
the exact mode identities of \secref{sec:6.13}, each occupied positive-frequency mode contributes
$\omega$ to the energy density, $k_j^2/\omega$ to the pressure along $j$, and $m/\omega$ to $\sigma$;
each hole in the sea contributes the same with $\omega\to|\omega|$. So for a state of $N$ modes in a
coordinate volume $\mathrm{Vol}$:
\[
  \rho=\sum\frac{\omega}{\mathrm{Vol}}+\bigl(W-\sigma W'\bigr),\qquad
  P_j=\sum\frac{k_j^2}{\omega\,\mathrm{Vol}}+\bigl(\sigma W'-W\bigr),\qquad
  \sigma=\sum\frac{m}{\omega\,\mathrm{Vol}} .
\]
For \textbf{one particle}, $w=(k^2/3)/\omega^2$ (isotropic average) on top of the background
$W-\sigma W'$. \cfproved{VII \S29}
\begin{itemize}
  \item \cfL{MODE SUMS [THE RESULT]: on a flat-frame plane wave (canonical normalization), T\_44 + Lhat == omega ubar (-i T16[4]) u and T\textasciicircum{}j\_j - Lhat == k\_j ubar (-i T16[j]) u (j = 0..3): energy and pressure are the frequency and k\_j dh/dk\_j bilinears}
  \item \cfL{MODE SUMS [has content]: with ubar = u\textasciicircum{}ddag sigma16 these are omega u\textasciicircum{}ddag G u and k\_j u\textasciicircum{}ddag (dh/dk\_j) u, so by Section 28 each occupied mode contributes omega, k\_j\textasciicircum{}2/omega and (to sigma) m/omega}
\end{itemize}

\subsection{The vacuum: normal ordering, the sea energy, and the cosmological-constant problem}
\label{sec:8.3}

$\langle0|{:}\hat T_{\mu\nu}{:}|0\rangle=0$ by normal ordering with respect to the $J$-vacuum. That
statement is clean only for a \textbf{static} background ($a_4=\text{const}$) or in the adiabatic
limit. On an $x_4$-dependent background normal ordering is ambiguous: with respect to the
instantaneous vacuum it breaks covariant conservation, and with respect to a fixed in-vacuum it leaves
a divergence. A conserved renormalized $\langle T\rangle$ needs adiabatic subtraction or
point-splitting \cfprose{VII \S29}.

The \textbf{unrenormalized sea energy} per unit volume, $-g\int^{\Lambda}\frac{d^3k}{(2\pi)^3}\,\omega$,
is minus the Fermi-sea energy at $\kF=\Lambda$:
\[
  -\epsKS(\Lambda)=-\frac{g}{16\pi^2}\Bigl[2\Lambda^4+2m^2\Lambda^2+\frac{m^4}{4}-m^4\ln\frac{2\Lambda}{m}\Bigr]+o(1)
  \qquad(\Lambda\to\infty).
\]
Its finite, mass-dependent part contains $-\frac{g}{32\pi^2}\,m^4\ln m^2$. It is absorbed into $V$
(the renormalized, ``no-sea'' functional). \textbf{The cosmological-constant problem is not solved
here.} \cfproved{VII \S29} \cfL{VACUUM [THE RESULT]: the unrenormalized sea energy -eps\_KS(Lambda) == -(g/(16 pi\textasciicircum{}2))(2 Lambda\textasciicircum{}4 + 2 m\textasciicircum{}2 Lambda\textasciicircum{}2 + m\textasciicircum{}4/4 - m\textasciicircum{}4 Log[2 Lambda/m]) + o(1) as Lambda -> infinity; its m-dependent finite part contains -(g/(32 pi\textasciicircum{}2)) m\textasciicircum{}4 Log[m\textasciicircum{}2]}

\paragraph{For a mass that varies (the coupled system), the sea energy cannot be absorbed silently}
(design review DFT-4, adopted). The renormalized one-loop Dirac-sea energy that the no-sea functional
drops (relativistic Hartree approximation, Chin 1977), per 3-volume, with reference mass $M$ and
counterterms through $m^4$, is
\[
  \Delta E_{\mathrm{vac}}(m)=-\frac{g}{16\pi^2}\Bigl[m^4\ln\frac mM+M^3(M-m)-\frac72M^2(M-m)^2+\frac{13}{3}M(M-m)^3-\frac{25}{12}(M-m)^4\Bigr]
\]
[file: \file{fable-cosmology/rust/fable_fermion/src/kohn_sham.rs}, function \cfcode{vacuum_energy};
its unit test \file{vacuum_energy_vanishes_to_fifth_order_at_the_reference_mass} checks that it
vanishes like $(m-M)^5$ at $m=M$ and is even in $m$]. Since \texttt{2bc936d} the solver evaluates it
as $\Delta E_{\mathrm{vac}}=-\frac{g}{16\pi^2}M^4F(u)$ with $u=(m-M)/M$, summing the power series
$F(u)=u^5/5-u^6/30+u^7/105-\dots$ for $|u|\le0.75$ and the closed form (with \cfcode{log1p}) otherwise,
because the closed form above cancels to rounding noise near $m=M$; the unit test
\file{vacuum_energy_matches_high_precision_references_near_and_far_from_the_reference_mass} compares it
with 120-digit references, and the commit records agreement to $1.8\times10^{-15}$ [file: the same
function, \cfcode{vac_bracket} and their tests; commit \texttt{2bc936d}]. It is neither zero nor absorbable into $W$
without changing the gap equation; the solver reports it (column \cfcode{vac_over_rhoc}) and does not
include it in the dynamics. Reading $W$ as the fully renormalized effective potential absorbs it
formally, but that is a fine-tuning. Its size \cfderived: at $m=0$ the bracket is
$M^4\bigl(1-\frac72+\frac{13}{3}-\frac{25}{12}\bigr)=-M^4/4$, so
$\Delta E_{\mathrm{vac}}(0)-\Delta E_{\mathrm{vac}}(M)=gM^4/(64\pi^2)$. With $g=8$ and the unit
$E_c=\rho_{c0}^{1/4}=2.46261\times10^{-3}$\,eV (asserted in Part~VIII:
\cfL{fable4d RUNS [definition]: the units computed from CODATA 2018 -- E\_c = rho\_c0\textasciicircum{}(1/4) == 2.46261e-3 eV, Omega\_r0 == 9.2096e-5, Omega\_b0 = 0.02237/h\textasciicircum{}2 == 0.049243 (to the quoted digits)}), a mass variation of order $m$ changes the vacuum energy by
$\frac{g\,m^4}{64\pi^2}\big/\rho_{c0}=3.44$, $3.44\times10^{4}$, $3.44\times10^{8}$,
$3.44\times10^{12}$ for $m=0.01$, $0.1$, $1$, $10$\,eV (it scales as $m^4$). So for
$m\gtrsim10$\,meV it exceeds the critical density. (The design review quoted $3.44\times10^{-4}$ and
$3.44$ at $0.01$ and $0.1$\,eV; those two figures are off by $10^4$ and are corrected here.)

\subsection{The homogeneous Fermi sea (the Kohn--Sham ground state) in closed form}
\label{sec:8.4}

Fill the positive-frequency modes with $|k|<\kF$ along the observed sheet (zero modes along $x_0$ ---
the $\sqrt{\sin}$ profile --- and along $x_5,x_6,x_7$). The degeneracy is $\boldsymbol{g=8}$ = 4
flavours $\times$ 2 spins (sections \ref{sec:3.6} and~\ref{sec:6}: eight positive-frequency modes per
momentum). With $\wF=\sqrt{\kF^2+m^2}$ and $\mathcal L=\ln\bigl((\kF+\wF)/m\bigr)$ ($m>0$; everything
depends on $m^2$ except $\sigma$, which has the sign of $m$):
\begin{align*}
  n&=\frac{g\,\kF^3}{6\pi^2} &&=g\int\frac{d^3k}{(2\pi)^3}\,1,\\
  \sigKS&=\frac{g\,m}{4\pi^2}\bigl[\kF\wF-m^2\mathcal L\bigr] &&=g\int\frac{d^3k}{(2\pi)^3}\,\frac m\omega,\\
  \epsKS&=\frac{g}{16\pi^2}\bigl[\kF\wF\,(2\kF^2+m^2)-m^4\mathcal L\bigr] &&=g\int\frac{d^3k}{(2\pi)^3}\,\omega,\\
  \PKS&=\frac{g}{48\pi^2}\bigl[\kF\wF\,(2\kF^2-3m^2)+3m^4\mathcal L\bigr] &&=g\int\frac{d^3k}{(2\pi)^3}\,\frac{k^2}{3\omega}.
\end{align*}
Each closed form is proved equal to its integral: its $\kF$-derivative is the integrand times
$g\kF^2/(2\pi^2)$ and it vanishes at $\kF=0$; independently, Mathematica's \cfcode{Integrate}
returns the same function (with $\operatorname{artanh}(\kF/\wF)$, which equals $\mathcal L$); and a
numerical quadrature at $\kF=3/2$, $m=1$, $g=8$ agrees to $10^{-12}$ (a cross-check; the proof is the
symbolic one). \cfproved{VII \S29}
\begin{itemize}
  \item \cfL{KOHN-SHAM [definition]: the degeneracy g = 8 -- per momentum the positive-frequency eigenspace has rank 8 (Section 28's P\_\{+,+\} + P\_\{+,-\}), four flavours times two spins}
  \item \cfL{KOHN-SHAM [THE RESULT]: each closed form has k\_F-derivative (g/(2 pi\textasciicircum{}2)) k\_F\textasciicircum{}2 f(k\_F) with f = 1, m/omega, omega, k\textasciicircum{}2/(3 omega), and vanishes at k\_F = 0 -- so it IS the integral g Int d\textasciicircum{}3k/(2 pi)\textasciicircum{}3 f}
  \item \cfL{KOHN-SHAM [has content]: Integrate returns the same functions, with ArcTanh[k\_F/w\_F] in place of L -- and ArcTanh[k\_F/w\_F] == L (equal derivatives, both 0 at k\_F = 0)}
  \item \cfL{KOHN-SHAM [control]: numerical quadrature agrees at k\_F = 3/2, m = 1, g = 8 to 1e-12 (a cross-check; the proof is the symbolic one above)}
\end{itemize}

The exported closed forms at $g=8$ ($k_F=\kF$;
\file{provenance-latex/generated/fermion_fable_eos.tex}, unchanged):

{\footnotesize
% fermion fable: the Kohn-Sham Fermi sea in closed form, g = 8 (canonical normalization chi = Psi/Sqrt[H], m = W'(sigma) > 0)
% generated by claude-fable/export_fermion_fable_tex.wls from the notebook's Part VII objects, 2026-09-24
% notation: x_0..x_7 are the coordinates (x_4 = the observer's time), a_4 = the free function a4 (never given a value),
% q = e^{-a_4(H x_4)}/sin^{1/6}(6 H x_0), p = e^{+a_4(H x_4)}/sin^{1/6}(6 H x_0), s = bar(psi) psi.
\begin{align*}
n &= \frac{4 k_{F}^3}{3 \pi ^2} \\
\sigma_{KS} &= \frac{2 m \left(k_{F} \sqrt{k_{F}^2+m^2}-m^2 \log \left(\frac{\sqrt{k_{F}^2+m^2}+k_{F}}{m}\right)\right)}{\pi ^2} \\
\varepsilon_{KS} &= \frac{k_{F} \sqrt{k_{F}^2+m^2} \left(2 k_{F}^2+m^2\right)-m^4 \log \left(\frac{\sqrt{k_{F}^2+m^2}+k_{F}}{m}\right)}{2 \pi ^2} \\
P_{KS} &= \frac{k_{F} \sqrt{k_{F}^2+m^2} \left(2 k_{F}^2-3 m^2\right)+3 m^4 \log \left(\frac{\sqrt{k_{F}^2+m^2}+k_{F}}{m}\right)}{6 \pi ^2} \\
\rho &= \varepsilon_{KS} + W - \sigma W', \quad P_{obs} = P_{KS} + \sigma W' - W, \quad P_{hid} = \sigma W' - W, \quad \rho + P_{obs} = w_F\, n, \quad w_F = \sqrt{k_F^2+m^2} \\
w &= P_{KS}/\varepsilon_{KS} \in [0, 1/3] \text{ for the mass term } W = m\sigma \quad (w \to k_F^2/(5m^2) \text{ as } k_F/m\to 0,\ w\to 1/3 \text{ as } k_F/m\to\infty)
\end{align*}

}

\paragraph{The identities that organize them:}
\begin{align*}
  \epsKS+\PKS&=\wF\,n &&\text{(the zero-temperature Gibbs relation,}\\
  & &&\ \text{chemical potential } \mu=\wF),\\
  \epsKS-3\PKS&=m\,\sigKS &&\text{(the trace identity)},\\
  \partial\epsKS/\partial m&=\sigKS,\qquad \partial\epsKS/\partial\kF=\wF\,dn/d\kF .
\end{align*}
\cfproved{VII \S29} \cfL{KOHN-SHAM [THE RESULT]: eps + P == w\_F n (Gibbs), eps - 3 P == m sigma (trace), d eps/dm == sigma, d eps/dk\_F == w\_F dn/dk\_F}

\subsection{The numerically stable forms}
\label{sec:8.5}

The closed forms cancel catastrophically for $x=\kF/|m|\ll1$: $\varepsilon\sim mn$, $\sigma\sim n$ and
$P\sim n\kF^2/(5m)$ are small differences of large terms. The solver therefore uses three regimes, all
written with the integral representations
\[
  S(x)=2\int_0^x\frac{t^2}{\sqrt{1+t^2}}\,dt,\qquad E(x)=8\int_0^x t^2\sqrt{1+t^2}\,dt,\qquad
  Q(x)=8\int_0^x\frac{t^4}{\sqrt{1+t^2}}\,dt,
\]
with $\sigma=g\,m^3S/(4\pi^2)$, $\varepsilon=g\,m^4E/(16\pi^2)$, $P=g\,m^4Q/(48\pi^2)$ [file:
\file{fable-cosmology/rust/fable_fermion/src/kohn_sham.rs}, module documentation and constants
\cfcode{X_LO = 0.6}, \cfcode{X_HI = 100.0}]:

{\small
\begin{longtable}{@{}p{0.25\linewidth}p{0.69\linewidth}@{}}
\toprule
\raggedright regime & \raggedright form used \tabularnewline
\midrule
\endhead
\raggedright $x<X_{\mathrm{LO}}=0.6$ & \raggedright the convergent binomial series (radius 1) in $x^2$, e.g.\ $\varepsilon=|m|\,n\,\bigl[1+\frac{3}{10}x^2-\frac{3}{56}x^4+\dots\bigr]$, $\sigma=\sgnm(m)\,n\,\bigl[1-\frac{3}{10}x^2+\dots\bigr]$ \tabularnewline
\raggedright $0.6\leq x\leq X_{\mathrm{HI}}=100$ & \raggedright the closed forms (with $\asinh x=\mathcal L$), losing at most about 1.3 digits at $x=0.6$ \tabularnewline
\raggedright $x>100$ & \raggedright the ultra-relativistic series in $y=|m|/\kF$, with the logarithm kept exact \tabularnewline
\raggedright $m=0$ exactly & \raggedright $\sigma=0$, $\varepsilon=3P=g\,\kF^4/(8\pi^2)$ \tabularnewline
\bottomrule
\end{longtable}
}

The series coefficients are $\binom{-1/2}{j}$ for $S$ and $Q$ and $\binom{1/2}{j}$ for $E$:
\[
  S=\sum_j 2\binom{-1/2}{j}\frac{x^{2j+3}}{2j+3},\qquad
  E=\sum_j 8\binom{1/2}{j}\frac{x^{2j+3}}{2j+3},\qquad
  Q=\sum_j 8\binom{-1/2}{j}\frac{x^{2j+5}}{2j+5}
\]
(the \cfcode{series_nr} function of the same file). The commit that finished the solver
(\texttt{c5892bd}) records that these agree with quadrature to $10^{-12}$ from $x=10^{-8}$ to $10^{6}$,
and that \cfcode{cargo test --release} passes 21 + 8 tests, among them
\file{kohn_sham::tests::fermi_integrals_match_quadrature} and
\file{kohn_sham::tests::series_and_closed_forms_are_continuous_at_the_switch_points}. After the defect
fixes of \texttt{2bc936d} the final solver passes 34 of 34 tests (23 unit, 11 integration), these two
among them [commit \texttt{2bc936d}; the \texttt{\#[test]} functions of \file{src/*.rs} and
\file{tests/cosmology.rs}]. (The design
review had proposed a switch at $x=0.25$; the implemented switch is $0.6$.)

\subsection{The mean field}
\label{sec:8.6}

With $W(\sigma):=V(H\sigma)$, the gap equation $m=W'(\sigma)$ and self-consistency
$\sigma=\sigKS(\kF,m)$:
\begin{align*}
  \rho&=\epsKS+W-\sigma W',\\
  \Pobs&=\PKS+\sigma W'-W,\\
  \Phid=P_{x_0}&=\sigma W'-W &&\text{(along } x_0 \text{ and along } x_5,x_6,x_7\text{: zero modes}\\
  & &&\ \text{carry no kinetic pressure)},\\
  \rho+\Pobs&=\wF\,n &&\text{(exact)},\\
  d\rho/dn&=\wF+(\sigKS-\sigma)\,W''\,\sigma'(n), &&\text{i.e. } =\wF \text{ at the self-consistent point}.
\end{align*}
\cfproved{VII \S29} \cfL{MEAN FIELD [THE RESULT]: rho = eps + W - sigma W', P\_obs = P\_KS + sigma W' - W give rho + P\_obs == w\_F n exactly; and d rho/dn == w\_F + (sigma\_KS - sigma) W'' sigma'(n), i.e. == w\_F at the self-consistent point sigma = sigma\_KS (m = W'(sigma) throughout)}

The split $\rho=\epsKS+U$ with $U:=W-\sigma W'$ separates a \textbf{quasiparticle} part ($\epsKS$,
$\PKS$, with $0\leq\PKS/\epsKS\leq1/3$) from a \textbf{mean-field condensate} $U$, which has $P=-U$ in
the observed and in the hidden directions alike: an 8-dimensional vacuum-like energy that varies with
$\sigma$ [derived here from the formulas above].

\subsection{The classical limit is the non-relativistic limit, and the sign rule}
\label{sec:8.7}

At the self-consistent point the trace identity gives, \textbf{exactly},
\[
  \rho=W(\sigma)+3\PKS,\qquad \Pobs=\bigl(\sigma W'(\sigma)-W(\sigma)\bigr)+\PKS ,
\]
that is, Part~VI's $\rho=V(s)$, $P=sV'-V$ (with $s=H\sigma$) plus the degeneracy pressure and three
times it. Part~VI is recovered exactly in the limit $\PKS/\epsKS\to0$, which is $\kF/|m|\to0$: the
cold, \textbf{non-relativistic} limit of the Fermi sea. (It is not ``$\kF\to0$ at fixed $n$'', which
is impossible, and not a coherent zero mode, which the Pauli principle forbids.) \cfproved{VII \S29}
\cfL{CLASSICAL LIMIT [THE RESULT]: at the self-consistent point rho == W + 3 P\_KS and P\_obs == (sigma W' - W) + P\_KS exactly (sigma = sigma\_KS, m = W'); Part VI's rho = V(s), P = s V' - V is the limit P\_KS/eps\_KS -> 0, i.e. k\_F/m -> 0}

Expanded in $x=\kF/|m|$ with the series of \secref{sec:8.5} and $W'=m$ \cfderived:
\begin{align*}
  \sigma&=\sgnm(m)\,n\,\bigl[1-\tfrac{3}{10}x^2+\dots\bigr],\\
  \rho&=W\bigl(\sgnm(m)\,n\bigr)+\tfrac{3}{10}\,n\,\kF^2/|m|+\dots,\\
  \Pobs&=\bigl[\sigma W'-W\bigr]+n\,\kF^2/(5|m|)+\dots,\\
  \Phid&=\sigma W'-W
\end{align*}
($\rho=|m|\,n\,(1+\frac{3}{10}x^2)+W(\sigma)-|m|\,n\,(1-\frac{3}{10}x^2)$ with
$W(\sigma)=W(\sgnm(m)\,n)-\frac{3}{10}|m|\,n\,x^2+\dots$). \textbf{The sign rule:} self-consistency
forces $\sgnf(\sigma)=\sgnf(m)=\sgnf(W')$. This removes the Part~VI branch $Ms>0$. For the author's
mass term $W=-2M\sigma$ ($V=-(2M/H)\,s$, $m=-2M$), $\rho\to2|M|\,n>0$ in the classical limit.

\subsection{The author's mass term: \texorpdfstring{$w$}{w} falls from 1/3 to 0}
\label{sec:8.8}

For $W=m\sigma$ ($V=-(2M/H)\,s$ gives $m=-2M$), $W-\sigma W'=0$, so
\[
  \rho=\epsKS>0,\qquad \Pobs=\PKS\geq0,\qquad 0\leq w=\PKS/\epsKS\leq\tfrac13 ,
\]
and $w$ depends only on $x=\kF/|m|$: $w=I_2/(3I_1)$ with the $m=1$ integrals
\begin{gather*}
  I_1=\tfrac18\bigl(x\,w_x(2x^2+1)-\ln(x+w_x)\bigr),\qquad
  I_2=\tfrac18\bigl(x\,w_x(2x^2-3)+3\ln(x+w_x)\bigr),\\
  I_3=\tfrac12\bigl(x\,w_x-\ln(x+w_x)\bigr),\qquad w_x=\sqrt{x^2+1}.
\end{gather*}
$w$ rises \textbf{monotonically}, since $dw/dx=x^2\Delta/(3\,w_x I_1^2)$ with
$\Delta=x^2I_1-w_x^2I_2$, $\Delta(0)=0$ and $\Delta'=2xI_3>0$, from $w\to0$ ($x\to0$,
$w=x^2/5+O(x^4)$: dust) to $w\to1/3$ ($x\to\infty$: radiation). As a gas of fixed particle number
cools, $w$ falls from 1/3 to 0.

\paragraph{The classical sign problem disappears.} Part~VI needed $Ms<0$ for $\rho=-(2M/H)\,s>0$.
Quantum mechanically $\sigKS=g\int m/\omega$ is \textbf{odd} in $m$, and
$m\,\sigKS=\frac{g\,m^4}{2\pi^2}\,I_3(\kF/|m|)\geq0$ for either sign of $m$. With $m=-2M$,
$M\sigKS\leq0$: exactly the branch $Ms<0$ that Part~VI had to assume. \cfproved{VII \S29}
\begin{itemize}
  \item \cfL{AUTHOR'S MASS TERM [THE RESULT]: W = m sigma gives W - sigma W' == 0, so rho == eps\_KS, P == P\_KS, and w = P\_KS/eps\_KS depends only on x = k\_F/m: w == I2/(3 I1), I\_n the m = 1 integrals}
  \item \cfL{AUTHOR'S MASS TERM [THE RESULT]: 0 <= w <= 1/3 -- P\_KS >= 0 (I2' = x\textasciicircum{}4/w\_x > 0, I2(0) = 0) and eps - 3P = m sigma = (g m\textasciicircum{}4/(2 pi\textasciicircum{}2)) I3 >= 0 (I3' = x\textasciicircum{}2/w\_x > 0, I3(0) = 0)}
  \item \cfL{AUTHOR'S MASS TERM [THE RESULT]: w rises MONOTONICALLY -- dw/dx == x\textasciicircum{}2 Delta/(3 w\_x I1\textasciicircum{}2) with Delta = x\textasciicircum{}2 I1 - w\_x\textasciicircum{}2 I2, Delta(0) = 0 and Delta' = 2 x I3 > 0 -- from w -> 0 (x -> 0, w = x\textasciicircum{}2/5 + O(x\textasciicircum{}4): dust) to w -> 1/3 (x -> infinity: radiation)}
  \item \cfL{AUTHOR'S MASS TERM [has content]: the classical sign problem disappears -- sigma\_KS = g Int m/omega is ODD in m and m sigma\_KS = (g m\textasciicircum{}4/(2 pi\textasciicircum{}2)) I3(k\_F/|m|) >= 0 for either sign, so with m = -2M: M sigma\_KS <= 0, the branch M s < 0 that Part VI had to assume}
\end{itemize}

\subsection{\texorpdfstring{$w\geq-1$}{w >= -1} always: no phantom crossing for the quantized fable}
\label{sec:8.9}

From $\rho+\Pobs=\wF\,n$ (\secref{sec:8.6}), which is exact, and $\wF\,n\geq0$:
\[
  w_f+1=\frac{\rho+\Pobs}{\rho}=\frac{\wF\,n}{\rho}\geq0\qquad\text{wherever }\rho>0 .
\]
($w_f$ is the fable's equation of state $w=\langle\Pobs\rangle/\langle\rho\rangle$ of
\secref{sec:8.1}; $\wF=\sqrt{\kF^2+m^2}$ is the Fermi energy.)
The quantized fermion fable \textbf{cannot cross the phantom divide}, for any potential $W$ [derived
here from \cfL{MEAN FIELD [THE RESULT]: rho = eps + W - sigma W', P\_obs = P\_KS + sigma W' - W give rho + P\_obs == w\_F n exactly; and d rho/dn == w\_F + (sigma\_KS - sigma) W'' sigma'(n), i.e. == w\_F at the self-consistent point sigma = sigma\_KS (m = W'(sigma) throughout)}]. The classical Part~VI crossing at $V'(s_1)=0$
(\secref{sec:2.12.7}) is unreachable. In the quantum theory $\rho+\Pobs=\wF\,n>0$ at every non-zero
density, whatever the value of $m$. And for the \cfcode{expdamp} potential the gap equation keeps
$0<\sigma<\min(n,s_1)$ strictly, so $m=W'(\sigma)$ never reaches the classical crossing point at
finite density: $\sigma\to s_1$ and $m\to0$ only as $n\to\infty$ (design review DFT-6; solver tests
\file{kohn_sham::tests::expdamp_pins_sigma_below_s1_and_mass_goes_to_zero_at_high_density} and
\file{kohn_sham::tests::no_phantom_crossing_rho_plus_p_obs_is_nonnegative}). The classical phantom
crossing of 2026-09-16 therefore does \textbf{not} survive quantization. (What an observer who fits a
cold-dark-matter-plus-dark-energy model to the coupled system can infer --- an apparent crossing, when
the dust is subtracted at today's dark-matter density and the effective mass grows --- is a property
of that fit, not of the fluid; it is part of the coupled system, Effort~B, and is stated in
\secref{sec:1}.)

\subsection{On the pre-universe: the quantum gas cools, the classical field does not}
\label{sec:8.10}

Four facts, each proved on the canonical frame with $a_4$ free.

\paragraph{(1) Charge.} The U(1) current is conserved on shell for generic $\Psi(x_0,x_4)$,
$\bar\Psi(x_0,x_4)$. For the $x_0$-profile $\Psi=\sqrt{\sin 6Hx_0}\;\Psi'$ with $\Psi'$ independent of
$x_0$, $\sqrt g\,j^0$ is independent of $x_0$ ($\sec\cdot\sin\cdot\cot=1$), so
$\partial_4\bigl(\sqrt g\,j^4\bigr)=0$: the charge per coordinate volume does not depend on $x_4$, and
since $\sqrt g$ does not either, neither does the proper 8-density. The 7-volume is constant because
the observed sheet grows exactly as the hidden sheet shrinks ($q^3p^3$ is $x_4$-independent).
\cfproved{VII \S29}
\begin{itemize}
  \item \cfL{PRE-UNIVERSE (1) [THE RESULT]: the U(1) current is conserved on shell, (1/Sqrt[g]) d\_mu(Sqrt[g] j\textasciicircum{}mu) == 0, for generic Psi(x0, x4), Psibar(x0, x4)}
  \item \cfL{PRE-UNIVERSE (1) [THE RESULT]: for Psi = Sqrt[Sin] Psi'(x4), Sqrt[g] j\textasciicircum{}0 is x0-independent (Sec Sin Cot == 1), so d\_4(Sqrt[g] j\textasciicircum{}4) == 0 on shell; and Sqrt[g] and q\textasciicircum{}3 p\textasciicircum{}3 are x4-independent}
\end{itemize}

\paragraph{(2) Which $x_4$-only fields solve the equations:} $V''s'=0$ (\secref{sec:5.8}).

\paragraph{(3) Redshift.} A mode $\Psi=\sqrt{\sin}\;u(x_4)\,e^{ikx_1}$ (coordinate momentum $k$ along
the observed sheet) obeys $\Tsix{4}\,u'+i\,(k/q)\,\Tsix{1}\,u=HV'u$ exactly (mass term, $V'$
constant): its \textbf{physical} momentum is $k/q$, and $d\ln(k/q)/dx_4=H\,a_4'(Hx_4)=-\Hobs$. The
observed-sheet momenta redshift as $e^{a_4}\sim1/a$, and the instantaneous frequency is
$\sqrt{k^2/q^2+m^2}$ ($E^2=\omega^2$ at $k_{\mathrm{phys}}=k/q$, \secref{sec:6.10}). \cfproved{VII \S29}
\begin{itemize}
  \item \cfL{PRE-UNIVERSE (3) [THE RESULT]: for Psi = Sqrt[Sin] u(x4) Exp[i k x1] the field equation (mass term, V' = constant) is Sqrt[Sin] Exp[i k x1] (T16[4] u' + i (k/q) T16[1] u - H V' u): the physical momentum is k/q}
  \item \cfL{PRE-UNIVERSE (3) [THE RESULT]: d Log[k/q]/dx4 == H a4'[H x4] == -H\_obs (observed momenta redshift as Exp[a4]), and E\textasciicircum{}2 == (k\textasciicircum{}2/q\textasciicircum{}2 + m\textasciicircum{}2) ID16 at the physical momentum}
\end{itemize}

\paragraph{(4) Cooling.} The occupied coordinate momenta are fixed (the mode labels are conserved;
the $x_4$-dependence of $q$ only mixes $\pm\omega$, a Bogoliubov effect that vanishes adiabatically).
So $x=k_{F,\mathrm{phys}}/m$ falls as $e^{a_4}$: $d\ln x/dx_4=H\,a_4'$, and since $w(x)$ is increasing,
$\sgnf(dw/dx_4)=\sgnf(a_4')$. In the regime $a_4'<0$ (the observed sheet expanding, the sign
hypothesis of Part~VI, \secref{sec:2.11}) the quantum fable gas \textbf{cools}, $w:1/3\to0$, even
though its 8-density does not dilute. The classical Part~VI fable, by contrast, has $ds/dx_4=0$ on the
rescaled solutions and $w=sV'/V-1$ \textbf{frozen}. \cfproved{VII \S29}
\begin{itemize}
  \item \cfL{PRE-UNIVERSE (4) [THE RESULT]: x = k\_F,phys/m = x\_c/(q m) obeys d Log[x]/dx4 == H a4'[H x4]; since dw/dx > 0, dw/dx4 = w'(x) x H a4' < 0 exactly when a4' < 0 (the observed sheet expanding): the quantum gas COOLS}
  \item \cfL{PRE-UNIVERSE (4) [control]: the classical Part VI fable does NOT cool -- ds/dx4 == 0 on its rescaled solutions (Part VI's result, re-asserted), and its w = s V'/V - 1 is a function of s alone: frozen}
\end{itemize}

{\small
\begin{longtable}{@{}p{0.17\linewidth}p{0.33\linewidth}p{0.44\linewidth}@{}}
\toprule
\raggedright quantity & \raggedright classical fable (Part~VI) & \raggedright fermion fable (Part~VII, Kohn--Sham, mass term) \tabularnewline
\midrule
\endhead
\raggedright $\rho$ & \raggedright $V(s)$ & \raggedright $\epsKS(\kF,m)>0$ \tabularnewline
\raggedright $P$ (observed) & \raggedright $sV'-V$ ($=0$: dust) & \raggedright $\PKS$ in $[0,\rho/3]$ \tabularnewline
\raggedright $w$ & \raggedright $sV'/V-1$, frozen ($ds/dx_4=0$) & \raggedright $\PKS/\epsKS$: $1/3\to0$ as the observed sheet expands \tabularnewline
\raggedright 8-density & \raggedright constant & \raggedright constant (charge per coordinate volume conserved) \tabularnewline
\bottomrule
\end{longtable}
}
\noindent(This table is the notebook's displayed summary at the end of Section~29 \cfdisplayed{VII \S29}.)

Finally, \textbf{$a_4$ is still undefined} after Part~VII: \cfL{PART VII [control]: a4 is STILL UNDEFINED -- nothing in Part VII gave it a value}

% ===============================================================================================
% Section 9: the shared canonical-spin-connection section (docs-shared/canonical-spin-connection.tex),
% with the Part VII placeholders filled from claude-fable/run_fermion_fable_part7.log.
\begingroup
\setlength{\emergencystretch}{4em}
\tolerance=2000
\section{The canonical spin connection: a complete discussion}
\label{csc:top}\label{sec:9}

This section covers the canonical spin connection of the 4+4 pre-universe in full: the frame it
comes from, how it is computed, every non-zero component, its $16\times16$ spinor form, the
gauge-covariant derivative of the 16-component spinor, how it relates to the three bridges of
Part~IV and to the fable-5.1 (Weitzenb\"ock) bridge of Part~V, and what it does and does not do for
the complex fermion fable. It stands on its own. Every object is defined here and every result
is restated here, with the place where it is proved.

Throughout, ``Section n'' (capital S) means a section of the notebook. The numbered parts of this
discussion are called ``subsections'' 9.1--9.13.

% --------------------------------------------------------------------------------------------
\subsection{Sources, and how the status of each statement is marked}
\label{csc:sources}\label{sec:9.1}

\paragraph{Where the proofs live.} The proofs are in the Mathematica notebook
\cfcode{claude-fable/claude-fable_Einstein-Rosen-2-Planes.nb}. That notebook is generated from cell
manifests and is never edited by hand:

{\small
\begin{longtable}{@{}p{0.40\linewidth}p{0.07\linewidth}p{0.12\linewidth}p{0.31\linewidth}@{}}
\toprule
\raggedright manifest & \raggedright Part & \raggedright Sections & \raggedright subject \tabularnewline
\midrule
\endhead
\raggedright \cfcode{claude-fable/cells_part1.wl} & \raggedright I & \raggedright 0--4 & \raggedright session, coordinates, \cfcode{eta4488}, \cfcode{tau}, \cfcode{sigma} \tabularnewline
\raggedright \cfcode{claude-fable/cells_part2.wl} & \raggedright I & \raggedright 5--9 & \raggedright \cfcode{T16}, \cfcode{sigma16}, \cfcode{SAB}, matrix bases, Cartan triality \tabularnewline
\raggedright \cfcode{claude-fable/cells_part3.wl} & \raggedright II & \raggedright 10--14 & \raggedright the author's physics reproduced (\cfcode{Psi16}, \cfcode{La}, \cfcode{eLa}) \tabularnewline
\raggedright \cfcode{claude-fable/cells_part4.wl} & \raggedright III, IV & \raggedright 15--17, 18--20 & \raggedright frame, Christoffel symbols, canonical spin connection, spinor covariant derivative; the three bridges \tabularnewline
\raggedright \cfcode{claude-fable/cells_part5.wl} & \raggedright V & \raggedright 21--22 & \raggedright the fable-5.1 (Weitzenb\"ock) bridge; the master comparison \tabularnewline
\raggedright \cfcode{claude-fable/cells_part6.wl} & \raggedright VI & \raggedright 23--25 & \raggedright fableScalar and the real (classical) fable spinor \tabularnewline
\raggedright \cfcode{claude-fable/cells_part7.wl} & \raggedright VII & \raggedright 26--29 & \raggedright the complex fermion fable (new): refinement, Lagrangian and field equations, canonical quantization, energy--momentum tensor \tabularnewline
\raggedright \cfcode{claude-fable/cells_part8.wl} & \raggedright VIII & \raggedright 30--33 & \raggedright fable as a source of the Einstein equations of the primordial field (the coupled system; its spin connection is Section~32) \tabularnewline
\bottomrule
\end{longtable}
}

The run of Parts I--VI, \cfcode{claude-fable/run_fable51_part6.log}, evaluates 172 Input cells straight
out of the \cfcode{.nb}. It reports \cfcode{Assertions run: 358   passed: 358   failed: 0} and
\cfcode{cells w/ msgs   : 0}, with \cfcode{identities accepted on numerical evidence alone : 0} and
\cfcode{non-vanishing witnesses : 27}. The Part~VII run is recorded in \cfcode{claude-fable/run_fermion_fable_part7.log}. It evaluates 198 Input
cells and reports \cfcode{Assertions run: 528   passed: 528   failed: 0} and \cfcode{cells w/ msgs   : 0}, with
\cfcode{identities accepted on numerical evidence alone : 0} and \cfcode{non-vanishing witnesses : 38}
(Part~VII adds 170 assertions and 11 witnesses to Parts I--VI). The acceptance run of Part~VIII,
\cfcode{claude-fable/run_fermion_fable_part8.log}, reported \cfcode{Assertions run: 642   passed: 642   failed: 0};
its two comparisons of the Rust solver with the Mathematica reference runs were skipped, because the
solver's CSVs did not exist yet. The final run of the whole notebook,
\cfcode{claude-fable/run_fermion_fable_final.log}, evaluates 217 Input cells and reports
\cfcode{Assertions run: 644   passed: 644   failed: 0} and \cfcode{cells w/ msgs   : 0}, with
\cfcode{identities accepted on numerical evidence alone : 0} and \cfcode{non-vanishing witnesses : 46}.

\paragraph{How an identity is certified.} Each identity is certified in up to three stages:
\begin{enumerate}
  \item exact structural equality (the entry is literally \cfcode{0});
  \item \cfcode{Simplify} under the geometric assumptions, with a 5-second budget;
  \item a numerical check at three rational probe points to 30 significant digits.
\end{enumerate}
Stage~3 has two uses, and the notebook counts them separately. The first is to accept an
identity. The notebook reports \emph{zero} of those, and its last assertion fails if that count is
ever non-zero, so every identity quoted below is closed symbolically. The second is to witness
that something is \emph{not} zero, through \cfcode{cfNonZeroWitnessQ}. A witness succeeds only if
every entry evaluates to an actual number at some probe point and at least one of them exceeds
$10^{-10}$ in magnitude. That is a complete proof of non-vanishing.

The probe functions used in stage~3 are test inputs only, never definitions:
\cfcode{a4 -> Function[u, 3/7 + u/5 + Sin[u]/11]}, a probe for the boost rapidity of Bridge~1, and
probe values of the Bridge~3 parameter \cfcode{lambdaOct}.

\paragraph{Status marks used below.}

{\small
\begin{longtable}{@{}p{0.34\linewidth}p{0.60\linewidth}@{}}
\toprule
\raggedright mark & \raggedright meaning \tabularnewline
\midrule
\endhead
\raggedright \cfproved{P \S n} ``label'' or label & \raggedright a \cfcode{cfAssert} of Part P, Section n, that prints \cfcode{PASS} in \cfcode{run_fable51_part6.log} (Parts I--VI), \cfcode{run_fermion_fable_part7.log} (Part~VII) or \cfcode{run_fermion_fable_final.log} (Part~VIII, the final run); the label is quoted verbatim \tabularnewline
\raggedright \cfdisplayed{P \S n} & \raggedright the printed output of a cell of Part P, Section n (a table, a count or a closed form); an output, not an assertion \tabularnewline
\raggedright \cfprose{P \S n} & \raggedright stated only in prose in the notebook (a Text cell, or a comment inside an Input cell), not asserted \tabularnewline
\raggedright \cfderived & \raggedright a short derivation carried out in this section from results marked above, with every step shown \tabularnewline
\bottomrule
\end{longtable}
}

The notebook grades its own assertion labels. \cfcode{[definition]} means true by construction.
\cfcode{[structural]} means true for any input of the right shape. \cfcode{[solver regression]} is a
check of the code, not of the geometry. \cfcode{[has content]} could have failed.
\cfcode{[THE RESULT]} marks a headline result. \cfcode{[control]} shows that the wrong alternative
fails.

% --------------------------------------------------------------------------------------------
\subsection{Conventions}
\label{csc:conventions}\label{sec:9.2}

\begin{itemize}
  \item Coordinates \cfcode{X = {x0, x1, ..., x7}}. The Wolfram array position is the coordinate
    index plus one, so \cfcode{x4} is entry~5. Greek indices are curved and Latin indices are flat;
    both run $0..7$. $x_4$ is the observer's time ($g_{44}=-1$), $x_1,x_2,x_3$ are the observed
    3-space, $x_0$ is the hidden spacelike direction, and $x_5,x_6,x_7$ are the hidden timelike
    sheet. The notebook's own words for $x_5..x_7$ are ``the superluminal deflating directions''.
  \item Flat metric: \cfcode{eta4488 = DiagonalMatrix[{1,1,1,1,-1,-1,-1,-1}]}. \cfproved{I \S2}
    \cflab{signature of eta4488 is (4,4)}, \cflab{eta4488 . eta4488 == ID8}.
  \item The $8\times8$ generators have \cfcode{tau[0] = ID8}. The conjugates are
    \cfcode{taubar[A] = sigma . Transpose[sigma . tau[A]]}, with
    \cfcode{sigma = ArrayFlatten[{{0,ID4},{ID4,0}}]} (\cfproved{I \S4} \cflab{sigma . sigma == ID8},
    \cflab{sigma is symmetric}). It follows that \cfcode{taubar[0] = ID8}.
  \item The $16\times16$ Dirac matrices are E.~A.~Lord's reduced Brauer--Weyl doubling,
    \cfcode{T16[A] = ArrayFlatten[{{0, taubar[A]}, {tau[A], 0}}]} for $A=0..7$. They are real, and
    \cfcode{T16[0] = ArrayFlatten[{{0,ID8},{ID8,0}}]}. \cfproved{I \S5}
    \cflab{{T16[A], T16[B]}/2 == eta4488[[A+1,B+1]] ID16}. Throughout, $\Tsix{a}$ \emph{is}
    $\gamma^a$, with the flat index \emph{up}. Since \cfcode{eta4488} is its own inverse,
    $\gamma_a=\eta_{ab}\Tsix{b}$ differs from $\Tsix{a}$ by the sign $\eta_{aa}$.
  \item Spinor metric: \cfcode{sigma16 = T16[0].T16[1].T16[2].T16[3]}. \cfproved{I \S5}
    \cflab{sigma16 == ArrayFlatten[{{-sigma,0},{0,sigma}}]},
    \cflab{sigma16 . T16[A] is antisymmetric for A = 0..7},
    \cflab{sigma16 is symmetric and sigma16 . sigma16 == ID16}.
  \item Chirality: \cfcode{T16[8] = T16[0].T16[1]...T16[7] = diag(-ID8, +ID8)}. \cfproved{I \S5}
    \cflab{PL projects onto the upper (type-1) components, PR onto the lower (type-2) components}.
    A 16-spinor is $\Psi=(\psi_1,\psi_2)$. The upper 8 components $\psi_1$ are a split-octonion
    type-1 spinor, and the lower 8 components $\psi_2$ are a type-2 spinor.
  \item so(4,4) generators: \cfcode{SAB[[A+1,B+1]] = (1/4) (T16[A].T16[B] - T16[B].T16[A])}.
    \cfproved{I \S5} \cflab{SAB is antisymmetric in its two flat indices},
    \cflab{sigma16 . SAB is antisymmetric},
    \cflab{T16[8] commutes with every SAB: the chiral halves are each Spin(4,4)-invariant},
    \cflab{[SAB, T16] reproduces the vector representation}.
  \item $H$, $K$ and $M$ are Protected constants of the original notebook.
\end{itemize}

% --------------------------------------------------------------------------------------------
\subsection{The canonical frame (vielbein), the metric, and the volume factor}
\label{csc:frame}\label{sec:9.3}

A vielbein here is simply an 8-dimensional vierbein, that is, an 8-dimensional frame field. The
canonical frame is the diagonal matrix that the original notebook records as the value of the
symbol it writes ``gtrye with indices alpha and (A)'' \cfprose{III \S15}:
\begin{gather}
  e=\mathrm{diag}\bigl(\tan 6Hx_0,\ q,\ q,\ q,\ 1,\ p,\ p,\ p\bigr),\\
  q=\frac{e^{-a_4(Hx_4)}}{\sin^{1/6}6Hx_0},\qquad
  p=\frac{e^{+a_4(Hx_4)}}{\sin^{1/6}6Hx_0}.
\end{gather}
In the notebook: \cfcode{frameCanonical = DiagonalMatrix[{Tan[6 H x0], q, q, q, 1, p, p, p}]},
\cfcode{q = cfQminus} (observed directions $x_1,x_2,x_3$), \cfcode{p = cfQplus} (hidden directions
$x_5,x_6,x_7$).

\paragraph{Rows curved, columns flat.} \cfcode{frameCanonical[[mu+1, a+1]]} $=e_\mu{}^a$. The
coframe is the matrix inverse, stored with the flat index on the row:
\cfcode{coframeCanonical = Inverse[frameCanonical]}, \cfcode{coframeCanonical[[a+1, mu+1]]}
$=e_a{}^\mu$. \cfproved{III \S15}
\cflab{CANONICAL [structural]: frame . coframe == ID8  (certifies only that the frame is invertible)}.

\paragraph{The free function \texorpdfstring{$a_4$}{a4}.} The original notebook never defines
\cfcode{a4}. It appears only inside the recorded frame. It is carried symbolically everywhere and
is \emph{never given a value}. The notebook prints a notice to this effect when Section~15 runs
\cfprose{III \S15}, and Part~VI ends with \cfproved{VI \S25}
\cflab{PART VI [control]: a4 is STILL UNDEFINED -- nothing in Part VI gave it a value; and the FLRW scale factor aF never entered the canonical frame}.
Every result in this section holds for an arbitrary differentiable $a_4$. Below, $a_4'$ means
\cfcode{Derivative[1][a4][H x4]}.

\paragraph{Reciprocity.} $qp=\sin^{-1/3}6Hx_0$ does not depend on $x_4$: the observed directions
contract exactly as fast as the hidden directions expand \cfprose{III \S15}.

\paragraph{The metric.} The bridge between the curved and flat metrics is
\begin{equation}
  \text{[d]}\qquad g_{\mu\nu}=e_\mu{}^a\,\eta_{ab}\,e_\nu{}^b ,
\end{equation}
that is, \cfcode{gCanonical = frameCanonical . eta4488 . Transpose[frameCanonical]}, and the line
element is
\begin{equation}
  ds^2=\tan^2(6Hx_0)\,dx_0^2+q^2\,(dx_1^2+dx_2^2+dx_3^2)-dx_4^2-p^2\,(dx_5^2+dx_6^2+dx_7^2).
\end{equation}
\begin{itemize}
  \item \cfproved{III \S15} \cflab{CANONICAL [definition]: g == frame . eta4488 . Transpose[frame]  (bridge [d]; this IS how g was defined, so it cannot fail)}.
  \item \cfproved{III \S15} \cflab{CANONICAL [has content]: g == diag(Tan[6 H x0]^2, q^2, q^2, q^2, -1, -p^2, -p^2, -p^2), the stated line element}.
  \item \cfproved{III \S15} \cflab{CANONICAL [has content]: Diagonal[g] == Diagonal[eta4488] * Diagonal[frame]^2 exactly} (Sylvester's law made explicit).
  \item \cfproved{III \S15} \cflab{CANONICAL [has content]: g has signature (4,4), GIVEN a4 real and 0 < 6 H x0 < Pi/2}.
  \item \cfproved{III \S15} \cflab{Inverse[g] == Transpose[coframe] . eta4488 . coframe}, \cflab{g is symmetric}, \cflab{the metric is non-degenerate}.
\end{itemize}

\paragraph{The volume factor.} Section~15 displays \cfcode{Det[g] = Sec[6 H x0]^2} and
\cfcode{Sqrt[Abs[Det[g]]] = Abs[Sec[6 H x0]]} \cfdisplayed{III \S15}. \cfproved{V \S21}
\cflab{FABLE-5.1: det g == Sec[6 H x0]^2  (positive, as signature 4+4 requires)}. On the domain
$0<6Hx_0<\pi/2$ the root is $\sec 6Hx_0$. From Part~V on, the notebook uses
\cfcode{sqrtDetgTele = Sec[6 H x0]} rather than \cfcode{Abs[Sec[...]]}, so that \cfcode{Abs} is
never differentiated \cfprose{V \S21}. The same result follows from the frame directly
\cfderived:
\begin{equation}
  \det e=\tan(6Hx_0)\,q^3p^3=\frac{\tan 6Hx_0}{\sin 6Hx_0}=\sec 6Hx_0=\sqrt{|\det g|}.
\end{equation}
The volume factor does not depend on $x_4$.

\paragraph{The curved Dirac matrices.} \cfcode{gammaCurvedCanonical[[mu+1]] = Sum_a coframe[[a,mu]] T16[a]}.
The frame is diagonal, so \cfderived
\begin{equation}
\begin{aligned}
  \gamma^0&=\cot(6Hx_0)\,\Tsix{0},\\
  \gamma^j&=\tfrac1q\,\Tsix{j}=e^{a_4}\sin^{1/6}(6Hx_0)\,\Tsix{j}\qquad(j=1,2,3),\\
  \gamma^4&=\Tsix{4},\\
  \gamma^k&=\tfrac1p\,\Tsix{k}=e^{-a_4}\sin^{1/6}(6Hx_0)\,\Tsix{k}\qquad(k=5,6,7).
\end{aligned}
\end{equation}
\begin{itemize}
  \item \cfproved{III \S17} \cflab{{gammaCurved[mu], gammaCurved[nu]} == 2 gInv[mu,nu] ID16}.
  \item \cfproved{VI \S24} \cflab{FABLE [has content]: (gamma^4)^2 == -ID16, so gamma^4 d_4 Psi' == H V'(s) Psi' is solved by d_4 Psi' == -H V'(s) gamma^4 Psi'}.
\end{itemize}

% --------------------------------------------------------------------------------------------
\subsection{The generalized Christoffel symbols}
\label{csc:christoffel}\label{sec:9.4}

These are the Levi-Civita symbols of $g$:
\begin{equation}
  \Gamma^\rho{}_{\mu\nu}=\tfrac12\,g^{\rho s}\bigl(\partial_\mu g_{s\nu}+\partial_\nu g_{s\mu}-\partial_s g_{\mu\nu}\bigr).
\end{equation}
\cfcode{cfChristoffel[g, ginv, coords]} computes the first derivatives
\cfcode{dg[[k,m,p]] = D[g[[m,p]], coords[[k]]]} once, reuses them in the triple loop, and simplifies
each entry under the notebook's assumptions. The result is \cfcode{GammaCanonical[[rho+1, mu+1, nu+1]]}.
Section~16 lists its 37 non-zero entries (out of 512) \cfdisplayed{III \S16}; here
\cfcode{Gamma[r;m,n]} stands for $\Gamma^r{}_{mn}$:

{\scriptsize
\begin{verbatim}
Gamma[0;0,0] = 12*H*Csc[12*H*x0]
Gamma[0;j,j] =  (H*Cos[6*H*x0]^3)/(E^(2*a4[H*x4])*Sin[6*H*x0]^(10/3))   j = 1,2,3
Gamma[0;k,k] = -(E^(2*a4[H*x4])*H*Cos[6*H*x0]^3)/Sin[6*H*x0]^(10/3)      k = 5,6,7
Gamma[j;0,j] = Gamma[j;j,0] = -H*Cot[6*H*x0]                             j = 1,2,3
Gamma[j;j,4] = Gamma[j;4,j] = -H*a4'[H*x4]                               j = 1,2,3
Gamma[4;j,j] = -(H*a4'[H*x4])/(E^(2*a4[H*x4])*Sin[6*H*x0]^(1/3))         j = 1,2,3
Gamma[4;k,k] = -(E^(2*a4[H*x4])*H*a4'[H*x4])/Sin[6*H*x0]^(1/3)           k = 5,6,7
Gamma[k;0,k] = Gamma[k;k,0] = -H*Cot[6*H*x0]                             k = 5,6,7
Gamma[k;4,k] = Gamma[k;k,4] =  H*a4'[H*x4]                               k = 5,6,7
\end{verbatim}
}
\noindent The count is $1+3+3+6+6+3+3+6+6=37$.
\begin{itemize}
  \item \cfproved{III \S16} \cflab{Gamma is symmetric in its two lower indices (torsion-free affine connection)}. This check is structural: swapping the two lower indices in the formula reproduces the formula whenever $g$ is symmetric.
  \item \cfproved{VI \S23} \cflab{HELPER [solver regression]: cfCovariantDivergence of the metric itself vanishes (metric compatibility of Section 16's Gamma; a wrong index slot in the helper fails this)}.
\end{itemize}

% --------------------------------------------------------------------------------------------
\subsection{The vielbein postulate, and how \texorpdfstring{\texttt{cfSpinConnection}}{cfSpinConnection} solves it}
\label{csc:postulate}\label{sec:9.5}

The spin connection is fixed by the zero-torsion (Levi-Civita) condition, written as the
\emph{vielbein postulate}. The total covariant derivative of the frame vanishes when $\Gamma$ acts
on the curved index and $\omega$ acts on the flat index \cfprose{III \S16}:
\begin{equation}
  \partial_\mu e_\nu{}^a-\Gamma^\rho{}_{\mu\nu}\,e_\rho{}^a+\omega_\mu{}^a{}_b\,e_\nu{}^b=0 .
\end{equation}
In the notebook: \cfcode{D[e[nu,a], x[mu]] - Gamma[rho,mu,nu] e[rho,a] + omegaMixed[mu,a,b] e[nu,b] == 0},
with \cfcode{omegaMixed[mu,a,b]} $=\omega_\mu{}^a{}_b$. Contracting with the coframe frees $\nu$:
\begin{equation}
  \omega_\mu{}^a{}_c=\sum_\nu\Bigl(\sum_\rho\Gamma^\rho{}_{\mu\nu}\,e_\rho{}^a-\partial_\mu e_\nu{}^a\Bigr)e_c{}^\nu .
\end{equation}
The notebook's solver is exactly this formula (Section~16):

{\scriptsize
\begin{verbatim}
cfSpinConnection[frame_, coords_, etaFlat_, gamma_] :=
 Module[{n = Length[coords], cofr, dE},
  cofr = cfSimp[Inverse[frame]];                        (* cofr[[a,nu]] = e_a^nu *)
  dE = Table[D[frame[[nu, a]], coords[[mu]]], {mu, n}, {nu, n}, {a, n}];
  Table[
    cfSimp[Sum[(Sum[gamma[[r, mu, nu]] frame[[r, a]], {r, n}]
                - dE[[mu, nu, a]]) cofr[[c, nu]], {nu, n}]],
    {mu, n}, {a, n}, {c, n}]];

cfLowerFirstFlat[om_, etaFlat_] := Module[{n = Length[etaFlat]},
  Table[Sum[etaFlat[[a, c]] om[[mu, c, b]], {c, n}], {mu, n}, {a, n}, {b, n}]];

(* omega_mu^a_b *)
omegaCanonicalMixed =
  cfSpinConnection[frameCanonical, X, eta4488, GammaCanonical];
(* omega_{mu a b} *)
omegaCanonical = cfLowerFirstFlat[omegaCanonicalMixed, eta4488];
\end{verbatim}
}

The same routine produces every connection of the notebook that comes from the vielbein
postulate: the canonical one, those of Bridges 1 and 2, the Levi-Civita part \cfcode{omegaTriLC} of
Bridge~3 (to which Section~20 then adds the contortion), and the fable-5.1 connection. The
Christoffel symbols are computed once and passed in. So the comparisons below compare results, not
different pieces of code \cfprose{III \S16, IV, V \S21}.

\paragraph{What each certificate is worth.} The notebook grades its certificates as follows
\cfprose{III \S16}:
\begin{itemize}
  \item \cfproved{III \S16} \cflab{CANONICAL [solver regression]: vielbein postulate residual is zero}.
    This is an algebraic identity of the solver. \cfcode{cfSpinConnection} contracts the postulate
    with \cfcode{Inverse[frame]} and \cfcode{cfVielbeinResidual} contracts it back with
    \cfcode{frame}, and the two cancel for any affine connection and any invertible frame.
  \item \cfproved{III \S16} \cflab{CANONICAL [structural]: torsion vanishes}. Once the residual is
    zero, this reduces to the symmetry of $\Gamma$ in its lower indices, which is itself
    structural. The torsion 2-form used is
    $T^a{}_{\mu\nu}=\partial_\mu e_\nu{}^a-\partial_\nu e_\mu{}^a+\omega_\mu{}^a{}_b e_\nu{}^b-\omega_\nu{}^a{}_b e_\mu{}^b$
    (\cfcode{cfTorsion}).
  \item \cfproved{III \S16} \cflab{CANONICAL [has content]: omega[mu,a,b] == -omega[mu,b,a]  (metric compatibility)}.
    This is not imposed anywhere. It holds only because the affine connection fed in is the
    Levi-Civita connection of the metric the frame builds, so a wrong Christoffel symbol would
    break it. The connection is therefore $\mathfrak{so}(4,4)$-valued.
  \item \cfproved{III \S16} \cflab{INDEPENDENT CHECK: the frame-only derivation reproduces omegaCanonical exactly}
    and \cflab{INDEPENDENT CHECK: and it is antisymmetric in its two raised flat indices}. The
    connection is re-derived from the frame alone by the anholonomy formula, which never mentions
    $\Gamma$ or $g$ (\cfcode{cfSpinConnectionFromFrame}, with $e^{a\nu}=(\eta^{-1}\cdot\text{coframe})^{a\nu}$):
    \begin{equation}
    \begin{aligned}
      \omega_\mu{}^{ab}={}&\tfrac12 e^{a\nu}\bigl(\partial_\mu e_\nu{}^b-\partial_\nu e_\mu{}^b\bigr)
       -\tfrac12 e^{b\nu}\bigl(\partial_\mu e_\nu{}^a-\partial_\nu e_\mu{}^a\bigr)\\
       &-\tfrac12 e^{a\rho}e^{b\sigma}\bigl(\partial_\rho e_\sigma{}^c-\partial_\sigma e_\rho{}^c\bigr)\eta_{cd}\,e_\mu{}^d .
    \end{aligned}
    \end{equation}
  \item The canonical frame is diagonal and so equals its own transpose. Section~16 therefore
    cannot detect a transposition of the frame's two slots. Section~18 closes that gap on a frame
    that is not symmetric (the canonical frame times a constant boost of rapidity $1/3$):
    \cfproved{IV \S18} \cflab{the slot-test frame is genuinely not symmetric},
    \cflab{INDEPENDENT CHECK on a non-symmetric frame: the two solvers agree, so the curved/flat slots are placed correctly in both},
    \cflab{and the result is the canonical connection conjugated by the constant boost, as the gauge law demands}.
\end{itemize}

% --------------------------------------------------------------------------------------------
\subsection{The canonical spin connection: all its non-zero components}
\label{csc:components}\label{sec:9.6}

\cfcode{omegaCanonical} has \textbf{24 non-zero components out of 512} \cfdisplayed{III \S16}.
They are listed by \cfcode{cfShowConnection[omegaCanonical, "omega"]}. In this subsection and the
next only, write $c=\cos 6Hx_0$, $s=\sin 6Hx_0$ and $A=a_4(Hx_4)$. For $j=1,2,3$ and $k=5,6,7$,
with all indices down,
$\omega_{\mu ab}=$ \cfcode{omega[mu; a, b] = omegaCanonical[[mu+1, a+1, b+1]]}:

\begin{center}
\begin{tabular}{@{}ll@{}}
\toprule
component (both flat indices down) & value \\
\midrule
$\omega_{j0j}=-\omega_{jj0}$ & $H c^2\big/\bigl(e^{A}s^{13/6}\bigr)$ \\
$\omega_{j4j}=-\omega_{jj4}$ & $H a_4'\big/\bigl(e^{A}s^{1/6}\bigr)$ \\
$\omega_{k0k}=-\omega_{kk0}$ & $-e^{A}H c^2\big/s^{13/6}$ \\
$\omega_{k4k}=-\omega_{kk4}$ & $e^{A}H a_4'\big/s^{1/6}$ \\
\bottomrule
\end{tabular}
\end{center}

\noindent That gives $6\text{ directions}\times2\text{ planes}\times2\text{ orderings}=24$
components. Every other component is zero.
Part~VII re-asserts the count and the pattern: \cfproved{VII \S27} \cfL{SPIN CONNECTION [has content]: metric compatible (omega\_mu\textasciicircum{}\{ab\} antisymmetric) and every non-zero component is omega\_mu\textasciicircum{}\{mu 0\} or omega\_mu\textasciicircum{}\{mu 4\} (up to antisymmetry): omega\_0 == omega\_4 == 0}.
That label writes the connection with both flat indices up, $\omega_\mu{}^{ab}$. Neither the
antisymmetry nor the pattern of non-zero components depends on that position: raising both flat
indices multiplies each component by $\eta_{aa}\eta_{bb}=\pm1$ \cfderived. The \textbf{sign} of a
listed component does depend on it: raising both indices flips the components in the planes with
exactly one timelike index --- $(4,j)$ and $(0,k)$ here --- and leaves the $(0,j)$ and $(4,k)$
components unchanged. That is why the table above is written with both flat indices down, as
\cfcode{omegaCanonical} is, and why Part~VIII, Section~32, lists the components of the connections of
its dynamical frames the same way (\secref{sec:10.3}).

\paragraph{Structure \cfderived, from the table.}
\begin{itemize}
  \item $\omega_0=\omega_4=0$: the connection vanishes along the hidden direction $x_0$ and along
    the observer's time $x_4$.
  \item In each remaining direction $\mu$ ($1,2,3,5,6,7$) only the two flat planes $(0,\mu)$ and
    $(4,\mu)$ are excited.
  \item Whether the generator of a plane is a rotation or a boost depends on the signs of
    \cfcode{eta4488} in that plane. For $j=1,2,3$, $(0,j)$ is a rotation (both directions
    spacelike) and $(4,j)$ is a boost. For $k=5,6,7$, $(0,k)$ is a boost and $(4,k)$ is a
    rotation (both directions timelike).
  \item The $x_0$-dependent components (planes $(0,\mu)$) do not involve $a_4'$. The
    $a_4$-dependent components (planes $(4,\mu)$) are proportional to $a_4'$, so they all vanish
    when $a_4$ is constant.
  \item No component depends on $x_1,x_2,x_3,x_5,x_6,x_7$. The metric, the frame and the
    connection are invariant under translations along the observed and hidden sheets.
  \item The \emph{mixed} components $\omega_\mu{}^a{}_b=$ \cfcode{omegaCanonicalMixed} follow by
    raising the first flat index with \cfcode{eta4488}. For example, $\omega_j{}^4{}_j=-\omega_{j4j}$
    because $\eta_{44}=-1$.
\end{itemize}

\paragraph{Curvature.} The curvature 2-form is
$R_{\mu\nu}{}^a{}_b=\partial_\mu\omega_\nu{}^a{}_b-\partial_\nu\omega_\mu{}^a{}_b+[\omega_\mu,\omega_\nu]^a{}_b$
(\cfcode{cfCurvature}). It has \textbf{156 non-zero components} \cfdisplayed{III \S16}.
\cfproved{III \S16} \cflab{CANONICAL: curvature is antisymmetric in the two spacetime indices}.
The Ricci scalar is \cfdisplayed{III \S16}
\begin{equation}
  R=-3H^2\Bigl(\bigl(55+7\cos 12Hx_0\bigr)\cot^2(6Hx_0)\csc^2(6Hx_0)-2\,a_4'^2\Bigr)
\end{equation}
In the notebook:

{\scriptsize
\begin{verbatim}
RicciScalarCanonical = -3*H^2*((55 + 7*Cos[12*H*x0])*Cot[6*H*x0]^2*Csc[6*H*x0]^2
                          - 2*a4'[H*x4]^2)
\end{verbatim}
}
\noindent This is equivalent to $-6H^2\bigl(24\cot^2 6Hx_0+31\cot^4 6Hx_0-a_4'^2\bigr)$ \cfderived, using
$\cos 12Hx_0=1-2\sin^2 6Hx_0$ and $\csc^2=1+\cot^2$. It depends on $a_4$ only through $a_4'^2$. It
is the frame-independent invariant that every bridge below is checked against.

\paragraph{Summary of the certified properties of the canonical connection.} It is metric
compatible (antisymmetric in its flat indices, a statement with content), it is torsion-free
(structural), it is reproduced exactly by an independent frame-only derivation, and it is not
flat.

% --------------------------------------------------------------------------------------------
\subsection{The spinor connection \texorpdfstring{$\Gamma_\mu$}{Gamma-mu} and its exact normalization}
\label{csc:spinmatrix}\label{sec:9.7}

The $16\times16$ spin-connection matrix is built by \cfcode{cfSpinMatrix} in Section~17:

{\scriptsize
\begin{verbatim}
cfSpinMatrix[om_, mu_Integer] :=
  (1/8) Sum[om[[mu, a, b]] (T16[a - 1] . T16[b - 1] - T16[b - 1] . T16[a - 1]),
            {a, 8}, {b, 8}];
GammaSpinCanonical =
  Table[cfSimpArray[cfSpinMatrix[omegaCanonical, mu]], {mu, 8}];
\end{verbatim}
}

It is fed \cfcode{omegaCanonical}, the connection with \emph{both flat indices down}. Since
$\Tsix{a}=\gamma^a$ has its flat index \emph{up}, the exact convention is
\begin{equation}
\begin{aligned}
  \Gamma_\mu&=\tfrac18\,\omega_{\mu ab}\,[\gamma^a,\gamma^b]
   =\tfrac14\sum_{a,b=0}^{7}\omega_{\mu ab}\,\Tsix{a}\Tsix{b}\\
  &=\tfrac12\sum_{a<b}\omega_{\mu ab}\,\Tsix{a}\Tsix{b}
   =\tfrac12\,\omega_{\mu ab}\,\texttt{SAB[[a+1,b+1]]} .
\end{aligned}
\end{equation}
The second line uses the antisymmetry of $\omega_{\mu ab}$ and $\Tsix{a}\Tsix{b}=-\Tsix{b}\Tsix{a}$
for $a\neq b$ \cfderived. The equivalent form with raised connection indices is
$\tfrac14\,\omega_\mu{}^{ab}\,\gamma_a\gamma_b$, with $\gamma_a=\eta_{ab}\Tsix{b}$.

\paragraph{Index warning.} The expression $\tfrac14\,\omega_\mu{}^{ab}\,\Tsix{a}\Tsix{b}$, with the
raised connection contracted against the unlowered \cfcode{T16}, is \emph{not} the notebook's
$\Gamma_\mu$. It would multiply each plane by $\eta_{aa}\eta_{bb}$, and so flip the sign of the
$(4,j)$ and $(0,k)$ terms \cfderived.

The coefficient and the sign are not conventions left open. They are fixed against the vielbein
postulate:
\begin{itemize}
  \item \cfproved{III \S17} \cflab{[definition] (1/8) omega [gamma^a, gamma^b] == (1/2) omega SAB  (SAB is defined as (1/4)[gamma,gamma], so this cannot fail)}.
  \item \cfproved{III \S17} \cflab{[has content] the spinor connection is compatible with omega: [Gamma^spin_mu, gamma^a] + omega_mu^a_b gamma^b == 0}.
    The spinor connection rotates the Dirac matrices exactly as $\omega$ rotates a vector. With
    $1/4$ in place of $1/8$, or with the opposite sign, this check fails.
  \item \cfproved{III \S17} \cflab{[control] with the opposite sign the compatibility check FAILS (witnessed): the sign is fixed by the vielbein postulate, not free}.
  \item \cfproved{VII \S27} the normalization used throughout Part~VII, \cfL{SPIN CONNECTION [definition]: Gamma\textasciicircum{}spin\_mu == Sum\_\{a<b\} omega\_mu,ab (1/2) T16[a] T16[b]}.
\end{itemize}

\paragraph{The explicit matrices \cfderived, from the table of
subsection~\ref{csc:components} and the formula above.} Only the planes $(0,\mu)$ and $(4,\mu)$
contribute. For $j=1,2,3$ the ordered pair $(j,4)$ gives
$\omega_{jj4}\Tsix{j}\Tsix{4}=\omega_{j4j}\Tsix{4}\Tsix{j}$, so
\begin{equation}
\begin{aligned}
  \Gamma_0&=\Gamma_4=0,\\
  \Gamma_j&=\tfrac12\Bigl[\frac{Hc^2}{e^{A}s^{13/6}}\,\Tsix{0}\Tsix{j}+\frac{Ha_4'}{e^{A}s^{1/6}}\,\Tsix{4}\Tsix{j}\Bigr]\qquad(j=1,2,3),\\
  \Gamma_k&=\tfrac12\Bigl[-\frac{e^{A}Hc^2}{s^{13/6}}\,\Tsix{0}\Tsix{k}+\frac{e^{A}Ha_4'}{s^{1/6}}\,\Tsix{4}\Tsix{k}\Bigr]\qquad(k=5,6,7).
\end{aligned}
\end{equation}
The provenance script \cfcode{claude-fable/prov03_covariant_derivative.wls} (log
\cfcode{claude-fable/prov03.log}) prints the non-zero directions of \cfcode{Gamma^spin} as
\cfcode{{1, 2, 3, 5, 6, 7}}, each with 32 non-zero entries. That matches the two terms above: each
$\Tsix{a}\Tsix{b}$ with $a\neq b$ is a signed permutation matrix with 16 entries, and the two terms
occupy different positions.

\paragraph{The direct-sum structure.} Every $\Tsix{a}$ is off-diagonal in the type-1/type-2 split,
so every product of two of them is block-diagonal:
\cfcode{T16[A].T16[B] = ArrayFlatten[{{taubar[A].tau[B], 0}, {0, tau[A].taubar[B]}}]}. The spin
connection therefore never mixes the two split-octonion spinors.
\begin{itemize}
  \item \cfproved{III \S17} \cflab{Gamma^spin[mu] is block diagonal: it never mixes type-1 with type-2}.
  \item \cfproved{III \S17} \cflab{the type-1 block is (1/8) omega (taubar^a tau^b - taubar^b tau^a)}.
  \item \cfproved{III \S17} \cflab{the type-2 block is (1/8) omega (tau^a taubar^b - tau^b taubar^a)}.
\end{itemize}
Below, the upper (type-1) block is written $\Gamma^{(1)}_\mu$ and the lower (type-2) block
$\Gamma^{(2)}_\mu$. In the notebook these are \cfcode{GammaSpinType1} and \cfcode{GammaSpinType2}.

% --------------------------------------------------------------------------------------------
\subsection{The gauge-covariant derivative of the 16-component spinor}
\label{csc:dirac}\label{sec:9.8}

\begin{equation}
\begin{aligned}
  D_\mu\Psi&=\partial_\mu\Psi+\Gamma_\mu\Psi, &
  D_\mu\psi_1&=\partial_\mu\psi_1+\Gamma^{(1)}_\mu\psi_1,\\
  D_\mu\psi_2&=\partial_\mu\psi_2+\Gamma^{(2)}_\mu\psi_2, &
  \mathrm{Dirac}[\Psi]&=\gamma^\mu D_\mu\Psi .
\end{aligned}
\end{equation}
In the notebook these are \cfcode{cfDcov16[psi, mu] = D[psi, X[[mu]]] + GammaSpinCanonical[[mu]] . psi},
\cfcode{cfDcovType1}, \cfcode{cfDcovType2} and
\cfcode{cfDiracOperator[psi] = Sum[gammaCurvedCanonical[[mu]] . cfDcov16[psi, mu], {mu, 8}]}.
\begin{itemize}
  \item \cfproved{III \S17} \cflab{Dcov16 on the direct sum is the direct sum of Dcov1 and Dcov2}
    (on a generic spinor of all eight coordinates).
\end{itemize}

\paragraph{Covariant constancy of $\gamma^\mu$.} The Christoffel symbols act on the curved index,
the spinor connection acts by commutator, and the frame connects the two. If any of the three had
the wrong sign or index order, the curved Dirac matrices would not be covariantly constant:
\begin{equation}
  D_\mu\gamma^\nu=\partial_\mu\gamma^\nu+\Gamma^\nu{}_{\mu\rho}\gamma^\rho+[\Gamma_\mu,\gamma^\nu]=0 .
\end{equation}
\begin{itemize}
  \item \cfproved{III \S17} \cflab{[has content] and therefore the curved Dirac matrices are covariantly constant: D_mu gammaCurved^nu == 0}.
\end{itemize}

\paragraph{The commutator identity.} Contract $\mu=\nu$ and use
$\Gamma^\mu{}_{\mu\rho}=\partial_\rho\ln\sqrt g$ \cfderived:
\begin{equation}\label{csc:eq:commutator}
  \frac{1}{\sqrt g}\,\partial_\mu\bigl(\sqrt g\,\gamma^\mu\bigr)=[\gamma^\mu,\Gamma_\mu] .
\end{equation}
\begin{itemize}
  \item \cfproved{VI \S24} \cflab{FABLE [has content]: in general (1/Sqrt[g]) d_mu(Sqrt[g] gamma^mu) == [gamma^mu, Gamma^spin_mu]  (covariant constancy of gamma^mu, Section 17)}.
\end{itemize}

\paragraph{The connection term of the Dirac operator, direction by direction \cfderived.} From the
explicit $\Gamma_\mu$ of subsection~\ref{csc:spinmatrix} and the curved gammas of
subsection~\ref{csc:frame}, using $\Tsix{j}^2=+\IDs$, $\Tsix{k}^2=-\IDs$ and the anticommutation of
distinct \cfcode{T16}:
\begin{align}
  \gamma^j\Gamma_j&=\tfrac{1}{2q}\,\Tsix{j}\bigl(w_0\Tsix{0}\Tsix{j}+w_4\Tsix{4}\Tsix{j}\bigr)\notag\\
    &=-\tfrac{1}{2q}\bigl(w_0\Tsix{0}+w_4\Tsix{4}\bigr),\\
  \gamma^k\Gamma_k&=\tfrac{1}{2p}\,\Tsix{k}\bigl(v_0\Tsix{0}\Tsix{k}+v_4\Tsix{4}\Tsix{k}\bigr)\notag\\
    &=+\tfrac{1}{2p}\bigl(v_0\Tsix{0}+v_4\Tsix{4}\bigr),
\end{align}
where $w_0,w_4$ and $v_0,v_4$ are the coefficients in $\Gamma_j$ and $\Gamma_k$. Since
$w_0/q=Hc^2/s^2=H\cot^2 6Hx_0$, $w_4/q=Ha_4'$, $v_0/p=-H\cot^2 6Hx_0$ and $v_4/p=Ha_4'$:

\begin{center}
\begin{tabular}{@{}lll@{}}
\toprule
$\mu$ & $\gamma^\mu\Gamma_\mu$ (no sum) & $\{\gamma^\mu,\Gamma_\mu\}$ (no sum) \\
\midrule
$0$ & $0$ & $0$ \\
$j=1,2,3$ & $-\tfrac12H\cot^2(6Hx_0)\,\Tsix{0}-\tfrac12Ha_4'\,\Tsix{4}$ & $0$ \\
$4$ & $0$ & $0$ \\
$k=5,6,7$ & $-\tfrac12H\cot^2(6Hx_0)\,\Tsix{0}+\tfrac12Ha_4'\,\Tsix{4}$ & $0$ \\
\bottomrule
\end{tabular}
\end{center}

\noindent The anticommutator column follows because $\Gamma_j\gamma^j$ equals
$+\tfrac{1}{2q}(w_0\Tsix{0}+w_4\Tsix{4})$, the negative of $\gamma^j\Gamma_j$, and similarly for
$k$.

\paragraph{The sum.} Summing the table \cfderived{} gives $3(-\tfrac12)+3(-\tfrac12)=-3$ for the
$\Tsix{0}$ terms and $3(-\tfrac12)+3(+\tfrac12)=0$ for the $\Tsix{4}$ terms. The observed and hidden
$a_4'$ terms cancel:
\begin{equation}\label{csc:eq:sum}
  \gamma^\mu\Gamma_\mu=-3H\cot^2(6Hx_0)\,\Tsix{0},\qquad \sum_\mu\{\gamma^\mu,\Gamma_\mu\}=0 .
\end{equation}
\begin{itemize}
  \item \cfproved{VI \S24} \cflab{FABLE [has content]: the HYPOTHESIS that makes it a product -- the anticommutator {gamma^mu, Gamma^spin_mu} == 0 on the canonical frame (no totally antisymmetric part of the connection)} (the sum over $\mu$).
  \item \cfproved{VI \S24} \cflab{FABLE [has content]: hence (1/2)(1/Sqrt[g]) d_mu(Sqrt[g] gamma^mu) == gamma^mu Gamma^spin_mu, and on this geometry == -3 H Cot[6 H x0]^2 T16[0]}.
  \item \cfproved{VII \S27} the per-direction form: every $\gamma^\mu\Gamma_\mu$ is
    $\alpha_\mu\Tsix{0}+\beta_\mu\Tsix{4}$, with $\sum\alpha_\mu=-3H\cot^2$ and $\sum\beta_\mu=0$,
    \cfL{SPIN CONNECTION [THE RESULT]: direction by direction gamma\textasciicircum{}mu Gamma\_mu == alpha\_mu T16[0] + beta\_mu T16[4]; Sum alpha\_mu == -3 H Cot[6 H x0]\textasciicircum{}2 and Sum beta\_mu == 0 (the observed and hidden a4 terms cancel)}; and the anticommutator vanishes for
    \emph{each} $\mu$ separately, \cfL{LAGRANGIAN (c) [THE RESULT]: Lsym - Lsym\_d == (1/(2H)) Sum\_mu Psibar \{gamma\textasciicircum{}mu, Gamma\_mu\} Psi identically, and \{gamma\textasciicircum{}mu, Gamma\_mu\} == 0 for EACH mu on the canonical frame: Lsym == Lsym\_d}.
\end{itemize}

\paragraph{Why the anticommutator vanishes \cfderived.} For $a\neq b$,
$\{\gamma^c,\gamma^a\gamma^b\}=2\gamma^a\gamma^b\gamma^c+2\eta^{ca}\gamma^b-2\eta^{cb}\gamma^a$.
This is $2\gamma^{cab}$ (the antisymmetrized product) when $c$ differs from both $a$ and $b$, and
$0$ when $c=a$ or $c=b$. Hence
\begin{equation}\label{csc:eq:antisym}
  \sum_\mu\{\gamma^\mu,\Gamma_\mu\}=\tfrac12\,\omega_{[cab]}\,\gamma^{cab},\qquad \omega_{cab}=e_c{}^\mu\omega_{\mu ab},
\end{equation}
and only the totally antisymmetric part of $\omega_{cab}$ survives. For the diagonal canonical
frame, $c$ equals the direction $\mu$, and every non-zero $\omega_{\mu ab}$ has $\mu$ equal to $a$
or $b$ (the planes $(0,\mu)$, $(4,\mu)$). So the totally antisymmetric part is zero, and it is zero
direction by direction.

\paragraph{Contrast with an expanding frame.} On the separate 4-dimensional FLRW reference frame of
Part~VI, the same computation leaves a Hubble term: \cfproved{VI \S24}
\cflab{FLRW [THE RESULT]: gamma^mu Gamma^spin_mu == (3/2) (a'/a) gamma^4 -- the Hubble term of the FLRW Dirac equation}.
On the canonical frame the analogous $\Tsix{4}$ terms cancel between the observed sheet and the
hidden sheet ($\sum\beta_\mu=0$), because $qp$ does not depend on $x_4$ \cfderived: the
per-direction coefficient $-\tfrac12Ha_4'$ of the table equals $\tfrac12\partial_4\ln q$, and
$+\tfrac12Ha_4'$ equals $\tfrac12\partial_4\ln p$, so the sum is $\tfrac32\partial_4\ln(qp)=0$.

\paragraph{On the model's own spinor \texttt{Psi16} of Section 11.}
\begin{itemize}
  \item \cfproved{III \S17} \cflab{[definition] Dcov[mu] Psi16 - D[Psi16, x_mu] == Gamma^spin[mu] . Psi16}.
  \item \cfproved{III \S17} \cflab{in the six directions on which Psi16 does not depend, Dcov Psi16 is purely the connection term}.
  \item \cfproved{III \S17} \cflab{the connection term does not vanish on Psi16: the connection matrices themselves are non-zero (witnessed)}.
  \item \cfproved{VI \S24} \cflab{FABLE [has content]: the connection term does not vanish on an x0-independent spinor -- witnessed on the constant spinor e1: a spinor homogeneous in everything but x4 is NOT a solution}.
\end{itemize}
Section~17 constructs the operator and applies it. It does not re-derive the flat field equations
of Part~II from a covariant Lagrangian \cfprose{III \S17}.

% --------------------------------------------------------------------------------------------
\subsection{The three bridges of Part IV}
\label{csc:bridges}\label{sec:9.9}

\paragraph{The gauge law \cfprose{IV}, verified per bridge.} Suppose a new frame is
\cfcode{frameNew = frameOld . L} with $L(x)$ invertible on the flat indices. Substituting into the
vielbein postulate, with \cfcode{M = Transpose[L]}, gives
\begin{equation}
  \omega^{\text{new}}_\mu=M\,\omega^{\text{old}}_\mu\,M^{-1}-(\partial_\mu M)\,M^{-1}\qquad\text{(mixed components)}.
\end{equation}
The second term is the inhomogeneous gauge term, and it vanishes exactly when $L$ is constant.
The law is not assumed anywhere. Each new connection is derived independently by
\cfcode{cfSpinConnection}, and the law is then verified against the result. A frame
\cfcode{frameCanonical . L} with $L$ preserving the flat metric describes the \emph{same geometry},
and its Levi-Civita connection is the canonical one in a different gauge, \emph{not} a new
connection \cfprose{IV}.

\subsubsection{Bridge 1: a locally boosted frame, the same connection in a local gauge (Section 18)}
\label{sec:9.9.1}

A new object, not in the original notebook, is introduced here: the rapidity
$\theta=$ \cfcode{cfBoostRapidity[x0, x4]}, an arbitrary function. It generates a boost in the flat
$(0,4)$ plane:
\begin{equation}
  \Lambda[\theta]=\mathrm{ID8}+(\cosh\theta-1)(E_{00}+E_{44})+\sinh\theta\,(E_{04}+E_{40})=e^{\theta K},
\end{equation}
where $K=E_{04}+E_{40}$ is the $(0,4)$ boost generator on flat indices,
and \cfcode{frameBoost = frameCanonical . Lambda}.
\begin{itemize}
  \item \cfproved{IV \S18} \cflab{Lambda is in O(4,4): Lambda . eta . Transpose[Lambda] == eta}, \cflab{Lambda is NOT in O(8): it is a genuine boost, not a rotation}, \cflab{Lambda == MatrixExp[theta K] with K the (0,4) boost generator}.
  \item \cfproved{IV \S18} \cflab{BRIDGE 1 reproduces the SAME curved metric g}, \cflab{BRIDGE 1 frame really is different from the canonical frame}.
  \item \cfproved{IV \S18} \cflab{BRIDGE 1: vielbein postulate holds identically}, \cflab{BRIDGE 1: omega[mu,a,b] == -omega[mu,b,a]}, \cflab{BRIDGE 1: torsion still vanishes}.
  \item \cfproved{IV \S18} \cflab{COMPARISON 1: omegaBoost == M omegaCanonical Inverse[M] - D[M,x_mu] Inverse[M]}, \cflab{COMPARISON 1: the WHOLE difference is the inhomogeneous gauge term}, \cflab{COMPARISON 1: the difference equals -D[theta, x_mu] * (boost generator)}.
  \item \cfproved{IV \S18} \cflab{COMPARISON 1: curvature transforms covariantly, R' == M R Inverse[M]}, \cflab{COMPARISON 1: BRIDGE 1 has the SAME Ricci scalar as the canonical frame}.
  \item \cfdisplayed{IV \S18} \cfcode{omegaBoost} has 28 non-zero components. The difference
    \cfcode{deltaBoost} has 4: $-\partial_\mu\theta\,K$ for $\mu=0$ and $\mu=4$, in the two orderings of the plane $(0,4)$.
\end{itemize}
At the spinor level, the spin lift is \cfcode{S = MatrixExp[theta SAB[[1,5]]]}, with no extra factor
$1/2$, because \cfcode{SAB} already carries $1/4$ \cfprose{IV \S18}:
\begin{itemize}
  \item \cfproved{IV \S18} \cflab{BRIDGE 1 spin lift: Inverse[S] . S == ID16}, \cflab{BRIDGE 1 spin lift: S preserves the spinor metric, Transpose[S] . sigma16 . S == sigma16}, \cflab{BRIDGE 1 spin lift [has content]: Inverse[S] . gamma^a . S == Lambda^a_b gamma^b  (S is the spin lift of Lambda)}.
  \item \cfproved{IV \S18} \cflab{COMPARISON 1 at the spinor level [THE GAUGE LAW]: Gamma'_mu == S Gamma_mu Inverse[S] - D[S,x_mu] Inverse[S]}, \cflab{COMPARISON 1 at the spinor level: the inhomogeneous term is -D[theta,x_mu] SAB[[1,5]], the spin image of -D[theta,x_mu] K}, \cflab{COMPARISON 1 at the spinor level: and that is exactly cfSpinMatrix applied to the vector difference deltaBoost}.
  \item \cfproved{IV \S18} \cflab{BRIDGE 1 [has content]: D'_mu (S psi) == S D_mu psi for a generic spinor -- one covariant derivative, two gauges}, \cflab{BRIDGE 1: the curved Dirac matrices of the boosted frame are S gamma^mu Inverse[S]}, \cflab{BRIDGE 1: so the Dirac operator is one operator in two gauges, Dirac'[S psi] == S Dirac[psi]}.
\end{itemize}
A solution of the Dirac equation in one gauge is therefore a solution in the other, after
$\psi\to S\psi$. Part~V identifies the inhomogeneous term as the fable-5.1 connection of the
boosted frame (subsection~\ref{csc:fable51}).

\subsubsection{Bridge 2: the null (light-cone) frame, the same connection in a constant gauge (Section 19)}
\label{sec:9.9.2}

\begin{equation}
  U=\tfrac{1}{\sqrt2}\begin{pmatrix}\mathrm{ID4}&\mathrm{ID4}\\ \mathrm{ID4}&-\mathrm{ID4}\end{pmatrix},\qquad U=U^{T},\qquad U\cdot U=\mathrm{ID8},
\end{equation}
with \cfcode{etaNull = Transpose[U] . eta4488 . U} and \cfcode{frameNull = frameCanonical . U}.
The eight flat directions are paired as $(0,4)$, $(1,5)$, $(2,6)$ and $(3,7)$, one spacelike and
one timelike in each pair, and each pair is replaced by its two null combinations \cfprose{IV \S19}.
\begin{itemize}
  \item \cfproved{IV \S19} \cflab{U is an involution: U . U == ID8}, \cflab{U is symmetric}, \cflab{THE NULL TANGENT METRIC IS EXACTLY THE SPINOR METRIC sigma}, \cflab{etaNull is symmetric and squares to ID8}, \cflab{all eight frame directions are null: etaNull has zero diagonal}.
    The identity \cfcode{etaNull == sigma} is a statement about the flat metric, not about the connection.
  \item \cfproved{IV \S19} \cflab{BRIDGE 2 reproduces the SAME curved metric g}, \cflab{BRIDGE 2: vielbein postulate holds identically}, \cflab{BRIDGE 2: omega[mu,a,b] == -omega[mu,b,a] with indices lowered by etaNull}, \cflab{BRIDGE 2: torsion vanishes}.
  \item \cfproved{IV \S19} \cflab{COMPARISON 2: omegaNull == U . omegaCanonical . U exactly}, \cflab{COMPARISON 2: the inhomogeneous gauge term is ZERO because U is constant}, \cflab{COMPARISON 2: R_null == U . R_canonical . U}, \cflab{COMPARISON 2: BRIDGE 2 has the SAME Ricci scalar as the canonical frame}.
  \item \cfdisplayed{IV \S19} \cfcode{omegaNull} has 48 non-zero components. A constant change of
    basis spreads the canonical 24 over more entries, and the difference from $U\omega U$ is zero.
\end{itemize}

\subsubsection{Bridge 3: triality frame with totally antisymmetric split-octonion torsion, a different connection (Section 20)}
\label{sec:9.9.3}

The flat index is a \emph{type-1 split-octonion spinor index}, reached through the constant
triality bridge \cfcode{triVecToSpin} of Section~9. \cfproved{I \S9}
\cflab{THE TRIALITY METRIC IS EXACTLY THE SPINOR METRIC:  etaTri == sigma}.
The frame is \cfcode{frameTri = frameCanonical . triVecToSpin}, with flat tangent metric $\sigma$, and
the connection is
\begin{equation}
  \omega^{\text{Oct}}_{\mu ab}=\omega^{\text{TriLC}}_{\mu ab}+\lambda\,m^{\text{Skew}}_{abc}\,(e_{\text{Tri}})_\mu{}^{c},
  \qquad (e_{\text{Tri}})_\mu{}^{c}=\texttt{frameTri[[mu,c]]} .
\end{equation}
Here \cfcode{omegaTriLC} is the Levi-Civita connection in the triality frame.
$m^{\text{Skew}}=$ \cfcode{mSkewSpin} is the totally antisymmetrized split-octonion structure
constant \cfcode{mLower}, transported to the spinor basis. $\lambda=$ \cfcode{lambdaOct} is a free
real torsion-strength parameter introduced in Section~20, not in the original notebook.
\begin{itemize}
  \item \cfproved{IV \S20} \cflab{the triality bridge is Euclidean-orthogonal: Transpose[P] == Inverse[P]}, \cflab{but it is NOT in O(4,4): P . eta4488 . Transpose[P] != eta4488 (it carries eta4488 to sigma)}, \cflab{BRIDGE 3 reproduces the SAME curved metric g}, \cflab{the triality tangent metric is sigma}.
  \item \cfproved{IV \S20} \cflab{the antisymmetrized structure constants are totally antisymmetric}, \cflab{they agree with mLower on the seven imaginary directions}, \cflab{mSkewSpin is totally antisymmetric} (48 non-zero components \cfdisplayed{IV \S20}).
  \item \cfproved{IV \S20} \cflab{BRIDGE 3 (lambda=0): vielbein postulate holds identically}, \cflab{BRIDGE 3 (lambda=0): omega[mu,a,b] == -omega[mu,b,a]}, \cflab{BRIDGE 3 (lambda=0) is the constant triality conjugate of the canonical connection}.
  \item \cfproved{IV \S20} \cflab{BRIDGE 3: omegaOct[mu,a,b] == -omegaOct[mu,b,a]  (still metric compatible)}, \cflab{BRIDGE 3: at lambda = 0 it reduces to the Levi-Civita connection}.
  \item \cfproved{IV \S20} \cflab{BRIDGE 3: torsion vanishes at lambda = 0}, \cflab{BRIDGE 3: torsion is NOT zero for lambda != 0}, \cflab{BRIDGE 3: the flat torsion T[a,b,c] is TOTALLY ANTISYMMETRIC}, \cflab{BRIDGE 3: T[a,b,c] == -2 lambda mSkewSpin[a,b,c]}.
  \item \cfproved{IV \S20} \cflab{COMPARISON 3: omegaOct - (triality conjugate of omegaCanonical) == contortion}, \cflab{COMPARISON 3: the difference is linear in lambda and vanishes at lambda = 0}. The difference is a \emph{tensor} (a contortion), so this is a different connection.
\end{itemize}

\paragraph{Its curvature is quadratic in $\lambda$.} With $K$ the contortion,
$R(\omega^{\text{Oct}})=R(\omega^{\text{TriLC}})+\lambda\bigl(dK+[\omega^{\text{TriLC}},K]\bigr)+\lambda^2[K,K]$:
\begin{itemize}
  \item \cfproved{IV \S20} \cflab{BRIDGE 3: at lambda = 0 the curvature is the triality conjugate of the canonical curvature}, \cflab{BRIDGE 3 [has content]: the curvature DOES depend on lambda (witnessed)}, \cflab{BRIDGE 3: the lambda-dependence is a polynomial of degree exactly two}.
  \item \cfproved{IV \S20} \cflab{BRIDGE 3: the lambda^1 term is the Levi-Civita covariant exterior derivative of the contortion, dK + [omegaTriLC, K]}. Written out, it is $\partial_\mu K_\nu-\partial_\nu K_\mu+[\omega^{\text{TriLC}}_\mu,K_\nu]-[\omega^{\text{TriLC}}_\nu,K_\mu]$.
  \item \cfproved{IV \S20} \cflab{BRIDGE 3: the lambda^2 term is [K_mu, K_nu], the contortion commuted with itself}.
\end{itemize}

\paragraph{The Ricci shift.} $-42\lambda^2=-\tfrac14T_{abc}T^{abc}$:
\begin{itemize}
  \item \cfproved{IV \S20} \cflab{BRIDGE 3: at lambda = 0 the Ricci scalar is the canonical one}.
  \item \cfproved{IV \S20} \cflab{BRIDGE 3 [THE RESULT]: R(omegaOct) - R(canonical) == -(1/4) T_{abc} T^{abc}, a constant shift set by the torsion alone}.
  \item \cfproved{IV \S20} \cflab{BRIDGE 3: equivalently -lambda^2 (mSkew . mSkew), with mSkew . mSkew == 42, so the shift is exactly -42 lambda^2}.
    Check \cfderived: $T=-2\lambda m^{\text{Skew}}$ gives $T\cdot T=4\lambda^2\cdot42=168\lambda^2$, and $\tfrac14\cdot168=42$.
\end{itemize}
This connection resembles the Cartan--Schouten connections of a Lie group, but it differs in two
ways that must not be assumed away \cfprose{IV \S20}. It is not flat. And the octonion structure
constants violate the Jacobi identity, so they cannot be the anholonomy of any frame.

\paragraph{Spinor level.} \cfcode{cfSpinMatrix} is correct only for \cfcode{eta4488} vector indices.
Bridge~3's flat index is a triality-spinor index, so the matching Dirac matrices are
\cfcode{T16Tri[a] = Sum_A triSpinToVec[[a,A]] T16[A-1]}, and they are contracted by
\cfcode{cfSpinMatrixTri} \cfprose{IV \S20}:
\begin{itemize}
  \item \cfproved{IV \S20} \cflab{THE TRIALITY DIRAC MATRICES OBEY THE sigma CLIFFORD RELATION: {T16Tri[a],T16Tri[b]} == 2 sigma[[a,b]] ID16}, \cflab{they are NOT the vector-frame Dirac matrices, so the distinction is real}.
  \item \cfproved{IV \S20} \cflab{at lambda = 0 the Bridge 3 spinor connection equals the canonical one}, \cflab{the extra spinor coupling vanishes at lambda = 0}, \cflab{the extra spinor coupling is non-zero for lambda != 0}, \cflab{the extra spinor coupling is still block diagonal (type-1 + type-2 preserved)}.
  \item The extra coupling
    \cfcode{(lambda/8) mSkew[a,b,c] frameTri[[mu,c]] [T16Tri[a], T16Tri[b]]} is \emph{quadratic} in
    the gammas. The third index is contracted with the frame, not with a third gamma. The cubic
    Einstein--Cartan axial term would appear only after contracting with $\gamma^\mu$ in the Dirac
    operator, and the notebook does \emph{not} carry that contraction out \cfprose{IV \S20}. The
    extra coupling has 160 non-zero entries \cfdisplayed{IV \S20}.
\end{itemize}

% --------------------------------------------------------------------------------------------
\subsection{The fable-5.1 bridge of Part V: the Weitzenb\"ock (teleparallel) connection}
\label{csc:fable51}\label{sec:9.10}

The fable-5.1 bridge keeps the metric, the frame, \cfcode{eta4488} and the bridge [d] unchanged,
and changes only the rule for parallel transport. It declares the canonical frame itself to be
parallel \cfprose{V \S21}:
\begin{equation}
  (\Gamma^W)^\rho{}_{\mu\nu}=e_a{}^\rho\,\partial_\mu e_\nu{}^a ,
\end{equation}
in the notebook \cfcode{GammaWeitzenboeck[[rho, mu, nu]] = Sum_a coframe[[a, rho]] D[frame[[nu, a]], x[mu]]}.
\begin{itemize}
  \item \cfproved{V \S21} \cflab{FABLE-5.1 [definition]: the frame is parallel  --  nabla^W e == 0 for all eight frame vectors}, \cflab{FABLE-5.1 [has content]: Gamma^W is metric compatible, nabla^W g == 0}, \cflab{FABLE-5.1 [has content]: Gamma^W is NOT symmetric in its lower indices -- it has torsion}, \cflab{FABLE-5.1 [has content]: Gamma^W is NOT the Levi-Civita connection}.
  \item $\Gamma^W$ has 13 non-zero components \cfdisplayed{V \S22}.
\end{itemize}

\paragraph{The spin connection is zero, and so is the curvature.} $\Gamma^W$ is fed to the same
solver. Because the first two terms of the vielbein postulate then cancel by construction, the
solver returns $\omega=0$:
\begin{itemize}
  \item \cfproved{V \S21} \cflab{FABLE-5.1 [solver regression]: vielbein postulate residual is zero}, \cflab{FABLE-5.1 [THE RESULT]: the spin connection in the canonical frame is IDENTICALLY ZERO}, \cflab{FABLE-5.1: trivially antisymmetric, so trivially metric compatible}.
  \item \cfproved{V \S21} \cflab{FABLE-5.1 [has content]: the torsion is NOT zero}, \cflab{FABLE-5.1: T^rho_{mu nu} == Gamma^W[rho,mu,nu] - Gamma^W[rho,nu,mu]  (torsion is the antisymmetric part of Gamma^W)}. With $\omega=0$ the torsion is $T^a{}_{\mu\nu}=\partial_\mu e_\nu{}^a-\partial_\nu e_\mu{}^a$.
  \item \cfproved{V \S21} \cflab{FABLE-5.1 [THE RESULT]: the curvature is IDENTICALLY ZERO -- a flat connection on a curved metric}.
\end{itemize}

\paragraph{The canonical connection is minus the contortion of this torsion.} The contortion is
$K=\Gamma^W-\Gamma^{LC}$, and in terms of the torsion
$K_{\rho\mu\nu}=\tfrac12\bigl(T_{\rho\mu\nu}+T_{\mu\rho\nu}+T_{\nu\rho\mu}\bigr)$ (all indices down, first
index then raised). Substituting $\Gamma^{LC}=\Gamma^W-K$ into the vielbein postulate leaves
\cfcode{omegaCanonical[mu,a,b] = -K[rho,mu,nu] e[rho,a] coframe[[b,nu]]} \cfprose{V \S21}.
\begin{itemize}
  \item \cfproved{V \S21} \cflab{COMPARISON F [has content]: Gamma^W - Gamma^LC == (1/2)(T_{rmn} + T_{mrn} + T_{nrm}) with the index raised, the contortion formula}.
  \item \cfproved{V \S21} \cflab{COMPARISON F [THE RESULT]: omegaCanonical == - contortion(fable-5.1 torsion), pulled back to flat indices, exactly}.
  \item \cfproved{V \S21} \cflab{COMPARISON F: so omegaCanonical - omegaFable51 is a TENSOR, not a gauge term (omegaFable51 == 0)}.
  \item \cfcode{contortionFable51} has 24 non-zero components, the same 24 as
    \cfcode{omegaCanonical}. \cfcode{torsionFable51Flat} has 24, and \cfcode{RiemannFable51} has 0
    \cfdisplayed{V \S22}.
\end{itemize}

\paragraph{The kind of torsion.} A torsion tensor splits into a vector part, an axial part and a
tensor part:
\begin{itemize}
  \item \cfproved{V \S21} \cflab{FABLE-5.1 torsion vector T[mu] = T[nu,nu,mu] is a GRADIENT: T[mu] == D[Log[Sin[6 H x0]], x[mu]]}.
  \item \cfproved{V \S21} \cflab{FABLE-5.1 torsion vector has exactly one non-zero component, 6 H Cot[6 H x0] along x0}.
  \item \cfproved{V \S21} \cflab{FABLE-5.1: the AXIAL (totally antisymmetric) part of the torsion vanishes -- the complement of Bridge 3}, \cflab{FABLE-5.1: the torsion itself does not vanish, so the tensor part carries the rest}.
\end{itemize}
For any diagonal frame the torsion vector is, component by component (no sum) \cfderived,
\begin{equation}\label{csc:eq:torsionvector}
  T_\mu=\sum_\nu e_a{}^\nu\bigl(\partial_\nu e_\mu{}^a-\partial_\mu e_\nu{}^a\bigr)=\partial_\mu\ln\frac{e_\mu{}^\mu}{\det e}.
\end{equation}
On the canonical frame this is $\partial_0\ln(\tan 6Hx_0/\sec 6Hx_0)=\partial_0\ln\sin 6Hx_0=6H\cot 6Hx_0$
for $\mu=0$. It is zero for $\mu=4$, because $e_4{}^4=1$ and $\det e=\sec 6Hx_0$ do not depend on
$x_4$. It is zero for $\mu=1,2,3,5,6,7$, because nothing depends on those coordinates. This agrees
with the assertion.

\paragraph{$R=-T+B$ in closed form.} With the torsion scalar
$T=\tfrac14T^{\rho\mu\nu}T_{\rho\mu\nu}+\tfrac12T^{\rho\mu\nu}T_{\nu\mu\rho}-T^\mu T_\mu$ and the boundary term
$B=\frac{2}{\sqrt{\det g}}\,\partial_\mu\bigl(\sqrt{\det g}\,T^\mu\bigr)$, where
$\sqrt{\det g}=\sec 6Hx_0$:
\begin{align}
  R&=-3H^2\Bigl(\bigl(55+7\cos 12Hx_0\bigr)\cot^2(6Hx_0)\csc^2(6Hx_0)-2a_4'^2\Bigr),\\
  T&=-6H^2\bigl(5\cot^4(6Hx_0)+a_4'^2\bigr),\\
  B&=-36H^2\bigl(5+\cos 12Hx_0\bigr)\cot^2(6Hx_0)\csc^2(6Hx_0).
\end{align}
These closed forms are \cfdisplayed{V \S21} (\cfcode{RicciScalarCanonical},
\cfcode{torsionScalarFable51}, \cfcode{boundaryTermFable51}). The following are \cfproved{V \S21}:
\cflab{COMPARISON F [THE RESULT]: R == -T + B  --  the teleparallel equivalent of GR, in closed form for this geometry},
\cflab{FABLE-5.1: T is not zero and B is not zero -- neither side is trivial}. The Einstein--Hilbert
and teleparallel actions therefore differ by a boundary term on this geometry.

\paragraph{The spinor: one vector term, removed by a rescaling.} In the fable-5.1 gauge the spinor
connection is zero, so the Dirac operator is $\gamma^\mu\partial_\mu$. The canonical Dirac operator
carries in addition $\gamma^\mu\Gamma_\mu$, and this whole term is one vector term:
\begin{align}
  \gamma^\mu\Gamma_\mu&=-\tfrac12\,T_\mu\gamma^\mu=-\tfrac12\,\partial_\mu\bigl(\ln\sin 6Hx_0\bigr)\gamma^\mu
     \notag\\
     &=-3H\cot^2(6Hx_0)\,\Tsix{0},\label{csc:eq:vectorterm}\\
  \mathrm{Dirac}_{LC}\bigl[\sqrt{\sin 6Hx_0}\,\Psi'\bigr]&=\sqrt{\sin 6Hx_0}\;\gamma^\mu\partial_\mu\Psi'\qquad\text{for every }\Psi'.
\end{align}
The first line checks against subsection~\ref{csc:dirac} \cfderived:
$-\tfrac12(6H\cot)(\cot\,\Tsix{0})=-3H\cot^2\Tsix{0}$, because $\gamma^0=\cot\,\Tsix{0}$. The matrix
$\gamma^\mu\Gamma_\mu$ is $-3H\cot^2$ times the permutation matrix $\Tsix{0}$, so it has 16
non-zero entries.
\begin{itemize}
  \item \cfproved{V \S21} \cflab{FABLE-5.1 [THE RESULT]: gamma^mu Gamma^spin[mu] == -(1/2) T[mu] gamma^mu, one vector term}, \cflab{FABLE-5.1: equivalently -(1/2) D[Log[Sin[6 H x0]], x_mu] gamma^mu}.
  \item \cfproved{V \S21} \cflab{FABLE-5.1 [THE RESULT]: DiracCanonical[ Sqrt[Sin[6Hx0]] Psi' ] == Sqrt[Sin[6Hx0]] DiracFable51[Psi'], for a generic Psi'}, \cflab{FABLE-5.1: the rescaling is NOT a no-op -- the two Dirac operators differ on the constant spinor}.
  \item \cfproved{VI \S24} with a potential: \cflab{FABLE [THE RESULT]: (gamma^mu D_mu - H V'(s)) [Sqrt[Sin[6 H x0]] Psi'] == Sqrt[Sin[6 H x0]] (gamma^mu d_mu Psi' - H V'(s) Psi') for a GENERIC Psi'(x0, x4): the rescaling removes the connection term exactly}.
\end{itemize}
The rescaling works because the torsion vector is a gradient. A vector-torsion term that is a
gradient is removed by multiplying the spinor by
$\exp\bigl(\tfrac12\ln\sin 6Hx_0\bigr)=\sqrt{\sin 6Hx_0}$ \cfderived.

\paragraph{For the real spinor of Parts II--VI the connection drops out of the Lagrangian.}
\cfcode{sigma16 . T16[a]} is antisymmetric, so $\Psi^T\sigma_{16}\Tsix{a}\Psi=0$ for real commuting
$\Psi$. Since $\gamma^\mu\Gamma_\mu$ is a multiple of $\gamma^\mu$, the connection contributes
nothing to that Lagrangian:
\begin{itemize}
  \item \cfproved{V \S21} \cflab{FABLE-5.1 [THE RESULT]: Psi^T sigma16 gamma^mu Gamma^spin[mu] Psi == 0 for EVERY Psi -- the spin connection drops out of the Lagrangian}, \cflab{FABLE-5.1: the reason is Section 11's identity, Psi^T sigma16 T16[a] Psi == 0 (sigma16 T16[a] antisymmetric)}, \cflab{FABLE-5.1: but the connection term does NOT drop out of the Dirac OPERATOR (it is -(1/2) T[mu] gamma^mu, non-zero)}.
  \item \cfproved{V \S21} \cflab{FABLE-5.1: the covariant Lagrangian with frame -> ID8 and volume factor -> 1 IS the author's La of Section 11, term for term}, \cflab{FABLE-5.1: and the covariant Lagrangian has NO connection term to omit -- its kinetic term is already the fable-5.1 one}. The author's flat \cfcode{La} omits only the volume factor $\sec 6Hx_0$ and the coframe factors inside $\gamma^\mu$ \cfprose{V \S21}.
\end{itemize}
This mechanism is specific to a \emph{real} field. subsection~\ref{csc:complex} shows that it does
not carry over to the complex fermion fable.

\paragraph{Frame covariance: Bridge 1 explained.} In Bridge~1's boosted frame, the same $\Gamma^W$
fed to the same solver gives a fable-5.1 spin connection that is pure gauge:
\begin{itemize}
  \item \cfproved{V \S21} \cflab{FABLE-5.1 in the boosted frame: vielbein postulate residual is zero}, \cflab{FABLE-5.1 in the boosted frame [THE RESULT]: omega == -D[M,x_mu] Inverse[M], pure gauge and nothing else}, \cflab{FABLE-5.1 in the boosted frame: the lemma's prediction itself satisfies the vielbein postulate with Gamma^W}, \cflab{FABLE-5.1 in the boosted frame: the curvature is STILL identically zero}.
  \item \cfproved{V \S21} \cflab{and Bridge 1's inhomogeneous term of Section 18 IS this object: deltaBoost == omegaFable51Boost}, \cflab{so omegaBoost == (canonical, conjugated) + (fable-5.1 in the boosted frame)}.
  \item \cfcode{omegaFable51BoostMixed} has 4 non-zero components, the same 4 as \cfcode{deltaBoost} \cfdisplayed{V \S21}.
\end{itemize}

% --------------------------------------------------------------------------------------------
\subsection{Master comparison of the five connections (Section 22)}
\label{csc:master}\label{sec:9.11}

Section~22 prints two tables: the summary \cfcode{cfConnectionSummary}, and a grid of the
comparisons side by side \cfdisplayed{V \S22}. The table below is \cfcode{cfConnectionSummary},
with three entries of the side-by-side grid added to the column ``difference from canonical'':
Bridge~1's spinor-level law, and Bridge~3's curvature and Ricci shift. Part~VII restates the bridge
results it relies on, so that it can be read alone: \cfproved{VII \S27} \cfL{SPIN CONNECTION [has content]: Parts IV-V restated -- Bridge 1's difference is the pure gauge term, Bridge 2 is a constant conjugate, Bridge 3 has totally antisymmetric torsion that is not zero}
and \cfL{SPIN CONNECTION [has content]: Part V restated -- omega\textasciicircum{}W == 0 in the canonical frame, omega\textasciicircum{}LC == -contortion, T\_mu == d\_mu Log[Sin[6 H x0]], R == -T + B, and gamma\textasciicircum{}mu Gamma\_mu == -(1/2) T\_mu gamma\textasciicircum{}mu}. For all five connections
the curved metric $g$ is the same (``yes (reference)'' for the canonical one), and the connection
is metric compatible; those two columns of the notebook's table read ``yes'' throughout and are
not repeated below.

{\footnotesize
\setlength{\tabcolsep}{4pt}
\begin{longtable}{@{}p{0.29\linewidth}p{0.19\linewidth}p{0.29\linewidth}p{0.14\linewidth}@{}}
\toprule
\raggedright connection; frame field; flat metric & \raggedright torsion zero? curvature zero? & \raggedright difference from canonical & \raggedright verdict \tabularnewline
\midrule
\endhead
\raggedright \cfcode{omegaCanonical}; \cfcode{frameCanonical} (diagonal); \cfcode{eta4488} & \raggedright torsion: yes; curvature: no & \raggedright --- (reference) & \raggedright Levi-Civita \tabularnewline
\raggedright \cfcode{omegaBoost} (Bridge 1); \cfcode{frameCanonical . Lambda(x)}; \cfcode{eta4488} & \raggedright torsion: yes; curvature: no & \raggedright gauge term $-\partial_\mu\theta\,K$, which is the fable-5.1 connection of the boosted frame; at spinor level $\Gamma'=S\Gamma S^{-1}-(\partial S)S^{-1}$ with the explicit spin lift $S$ & \raggedright SAME connection, local gauge \tabularnewline
\raggedright \cfcode{omegaNull} (Bridge 2); \cfcode{frameCanonical . U}; \cfcode{etaNull == sigma} & \raggedright torsion: yes; curvature: no & \raggedright constant similarity $U\omega U$ & \raggedright SAME connection, constant gauge \tabularnewline
\raggedright \cfcode{omegaOct} (Bridge 3); \cfcode{frameCanonical . triVecToSpin}; \cfcode{etaTri == sigma} & \raggedright torsion: \textbf{no}, axial torsion $-2\lambda\,m^{\text{Skew}}$; curvature: no & \raggedright contortion $\lambda\,m^{\text{Skew}}_{abc}\,(e_{\text{Tri}})_\mu{}^{c}$ (with $e_{\text{Tri}}$ = \cfcode{frameTri}), a tensor; curvature quadratic in $\lambda$; Ricci shift $-42\lambda^2$ & \raggedright \textbf{DIF\-FER\-ENT} connection \tabularnewline
\raggedright \cfcode{omegaFable51} (fable-5.1); \cfcode{frameCanonical} (parallel); \cfcode{eta4488} & \raggedright torsion: \textbf{no}, vector + tensor torsion, axial part zero; curvature: \textbf{yes}, flat & \raggedright $\omega^{\text{canonical}}$ $=-$contortion($T_{\text{fable51}}$), a tensor; $\omega^{\text{fable51}}=0$ & \raggedright \textbf{DIF\-FER\-ENT} connection \tabularnewline
\bottomrule
\end{longtable}
}


Non-zero component counts (Section~22 fingerprint table and the counts printed in Sections
16--21, \textbf{[displayed]}):

\begin{center}
{\small
\begin{tabular}{@{}ll@{}}
\toprule
object & non-zero components \\
\midrule
\cfcode{GammaCanonical} (Christoffel) & 37 \\
\cfcode{GammaWeitzenboeck} & 13 \\
\cfcode{omegaCanonical} & 24 \\
\cfcode{omegaBoost} / its difference \cfcode{deltaBoost} & 28 / 4 \\
\cfcode{omegaNull} / its difference & 48 / 0 \\
\cfcode{omegaTriLC} (Bridge 3, $\lambda=0$) & 84 \\
\cfcode{contortionOct} (Bridge 3) & 78 \\
\cfcode{torsionOctFlat} (Bridge 3) & 48 \\
\cfcode{omegaFable51} (canonical frame) & 0 \\
\cfcode{contortionFable51} / \cfcode{torsionFable51Flat} & 24 / 24 \\
\cfcode{RiemannCanonical} & 156 \\
\cfcode{RiemannFable51} & 0 \\
\bottomrule
\end{tabular}
}
\end{center}

% --------------------------------------------------------------------------------------------
\subsection{What the canonical connection does for the complex fermion fable}
\label{csc:complex}\label{sec:9.12}

Everything in this subsection concerns the \emph{complex} 16-component fermion fable of Part~VII.
Two kinds of statement appear. Some are consequences of results already proved in Parts I--VI,
and the derivation is shown here. The others are assertions of the new Part~VII, marked
\cfproved{VII \S n}, with the label quoted verbatim from \cfcode{claude-fable/run_fermion_fable_part7.log}.

\subsubsection{The field and the symmetrized Lagrangian}
\label{sec:9.12.1}

Fermion fable is a complex 16-component Grassmann spinor $\Psi$. It is the minimal field that uses
the author's $\sigma_{16}$ conjugation, propagates, admits a potential $V(s)$, and reduces to
Part~VI's real field as its real commuting shadow. It is not the unique field with dynamics. This is
the conclusion of Section~26 as a whole (\secref{sec:3.8}); no single assertion states it. It rests on
\cfproved{VII \S26} \cfL{REFINEMENT (d) [THE RESULT]: s = Psibar Psi == i theta1\textasciicircum{}T sigma16 theta2: NON-ZERO (16 monomials) and HERMITIAN; p = Psibar T16[8] Psi is non-zero and Hermitian too}, \cfL{REFINEMENT (b) [THE RESULT]: Majorana with C+: the kinetic term is genuine -- Euler-Lagrange expression 2 C+ T16[a] phi with C+ T16[a] invertible -- but every invariant bilinear vanishes (massless, no V(s))},
\cfL{REFINEMENT (c) [THE RESULT]: Weyl and Majorana-Weyl: PL C T16[a] PL == PR C T16[a] PR == 0 for C = C-, C+ and every a (C chirality-diagonal, T16[a] chirality-odd) -- no kinetic term} and \cfL{FIDELITY [1] [fidelity]: Psi -> psi, Psibar -> psi\textasciicircum{}T sigma16 turns the symmetrized covariant Lhat into Part VI's cfLhatFable EXACTLY (canonical frame, generic V)}, among the others quoted in \secref{sec:3}. Its
Dirac conjugate and scalar bilinear are
\begin{equation}
  \bar\Psi=\Psi^{\ddagger}\sigma_{16},\qquad
  s=\bar\Psi\Psi=-\psi_1^{\ddagger}\sigma\psi_1+\psi_2^{\ddagger}\sigma\psi_2,
\end{equation}
where ${}^\ddagger$ is the classical conjugate transpose. The second form follows from
$\sigma_{16}=\mathrm{diag}(-\sigma,\sigma)$. In the quantized theory the Hilbert adjoint is
$\Psi^\dagger:=\Psi^\ddagger J$ (subsection~\ref{sec:9.12.6}).

The only Lagrangian used is the \emph{symmetrized covariant} one:
\begin{gather}
  L=\sqrt g\,\hat L,\qquad
  \hat L=\frac{1}{2H}\Bigl[\bar\Psi\gamma^\mu D_\mu\Psi-(D_\mu\bar\Psi)\gamma^\mu\Psi\Bigr]-V(s),\label{csc:eq:Lsym}\\
  D_\mu\Psi=\partial_\mu\Psi+\Gamma_\mu\Psi,\qquad
  D_\mu\bar\Psi=\partial_\mu\bar\Psi-\bar\Psi\Gamma_\mu .\notag
\end{gather}
$D_\mu\bar\Psi$ is the Dirac conjugate of $D_\mu\Psi$, and $\hat L$ is real, because
$\sigma_{16}\gamma^a$ and $\sigma_{16}\Gamma_\mu$ are \emph{real antisymmetric} matrices. That rests
on \cfproved{I \S5} \cflab{sigma16 . T16[A] is antisymmetric for A = 0..7} and
\cflab{sigma16 . SAB is antisymmetric}, together with
$\Gamma_\mu=\tfrac12\omega_{\mu ab}\,\texttt{SAB}$ with real $\omega$
(subsection~\ref{csc:spinmatrix}). \cfproved{VII \S27} \cfL{LAGRANGIAN (a) [definition]: every Gamma\textasciicircum{}spin\_mu is real and sigma16 Gamma\_mu is antisymmetric (Gamma\_mu\textasciicircum{}T sigma16 == -sigma16 Gamma\_mu): D\_mu Psibar is the Dirac conjugate of D\_mu Psi}.

\subsubsection{It drops out of the symmetrized Lagrangian because the matrix \texorpdfstring{$\{\gamma^\mu,\Gamma_\mu\}$}{anticommutator} vanishes}
\label{sec:9.12.2}

Expanding $D$ in \eqref{csc:eq:Lsym} gives \cfderived
\begin{equation}
  \hat L(D)=\hat L(\partial)+\frac{1}{2H}\,\bar\Psi\Bigl(\sum_\mu\{\gamma^\mu,\Gamma_\mu\}\Bigr)\Psi .
\end{equation}
The connection terms are $\bar\Psi\gamma^\mu\Gamma_\mu\Psi$ from the first product and
$+\bar\Psi\Gamma_\mu\gamma^\mu\Psi$ from $-(D_\mu\bar\Psi)\gamma^\mu\Psi$. On the canonical frame
$\{\gamma^\mu,\Gamma_\mu\}=0$ \emph{for each} $\mu$ (subsection~\ref{csc:dirac}), so
$\hat L(D)=\hat L(\partial)$ identically. That is, ``covariant $=$ partial'' holds for the
symmetrized complex Lagrangian on this frame. It holds on any frame where the totally antisymmetric
part $\omega_{[cab]}$ vanishes \eqref{csc:eq:antisym}, and fails in general.
\begin{itemize}
  \item \cfproved{VII \S27} \cfL{LAGRANGIAN (c) [THE RESULT]: Lsym - Lsym\_d == (1/(2H)) Sum\_mu Psibar \{gamma\textasciicircum{}mu, Gamma\_mu\} Psi identically, and \{gamma\textasciicircum{}mu, Gamma\_mu\} == 0 for EACH mu on the canonical frame: Lsym == Lsym\_d}.
\end{itemize}
This is \emph{not} Part~VI's mechanism. For the real commuting field,
$\Psi^T\sigma_{16}\gamma^\mu\Gamma_\mu\Psi$ vanishes because $\sigma_{16}\Tsix{0}$ is antisymmetric
(subsection~\ref{csc:fable51}). For a complex field,
$\bar\Psi\Tsix{0}\Psi=\Psi^\ddagger(\sigma_{16}\Tsix{0})\Psi$ with $\sigma_{16}\Tsix{0}$ real
antisymmetric is $i$ times the Hermitian form of $-i\sigma_{16}\Tsix{0}$, which is not zero in
general. So the drop-out for the complex field is a property of the \emph{matrix}
$\{\gamma^\mu,\Gamma_\mu\}$, not of a bilinear identity \cfderived. The other Lagrangians behave
differently:
\begin{itemize}
  \item The \emph{unsymmetrized covariant} Lagrangian $\frac1H\bar\Psi\gamma^\mu D_\mu\Psi-V$ equals
    the symmetrized one plus $\frac{1}{2H}\frac{1}{\sqrt g}\partial_\mu\bigl(\sqrt g\,\bar\Psi\gamma^\mu\Psi\bigr)$,
    a total divergence. Its connection term is
    $\frac1H\bar\Psi\gamma^\mu\Gamma_\mu\Psi=-3\cot^2(6Hx_0)\,\bar\Psi\Tsix{0}\Psi$, which is not
    zero. \cfproved{VII \S27}
    \cfL{LAGRANGIAN (d) [has content]: Lun - Lsym == (1/(2H)) (1/Sqrt[g]) d\_mu( Sqrt[g] Psibar gamma\textasciicircum{}mu Psi ), a total divergence},
    \cfL{LAGRANGIAN (e) [has content]: Lun - Lun\_d == (1/H) Psibar gamma\textasciicircum{}mu Gamma\_mu Psi == -3 Cot[6 H x0]\textasciicircum{}2 Psibar T16[0] Psi}, \cfL{LAGRANGIAN (e) [control]: that term is NOT zero for a complex field (witness u = e1, v = e13), and IS zero in the real commuting limit (the scope of Parts V-VI)}.
  \item The \emph{unsymmetrized partial} Lagrangian $\frac1H\bar\Psi\gamma^\mu\partial_\mu\Psi-V$ is
    \emph{inadmissible}. Its $\Psi$ and $\bar\Psi$ equations are not Dirac conjugates of each other:
    they differ by
    $\bar\Psi\sum_\mu[\gamma^\mu,\Gamma_\mu]=2\bar\Psi\gamma^\mu\Gamma_\mu=-6H\cot^2(6Hx_0)\,\bar\Psi\Tsix{0}$.
    \cfproved{VII \S27} \cfL{LAGRANGIAN (f) [THE RESULT]: Lun\_d varied with respect to Psibar gives (Sqrt[g]/H)(gamma\textasciicircum{}mu d\_mu Psi - H V' Psi): the connection term is MISSING}, \cfL{LAGRANGIAN (f) [THE RESULT]: varied with respect to Psi it gives -(Sqrt[g]/H)((d\_mu Psibar) gamma\textasciicircum{}mu + Psibar [gamma\textasciicircum{}mu, Gamma\_mu] + H V' Psibar), with Psibar [gamma\textasciicircum{}mu, Gamma\_mu] == 2 Psibar gamma\textasciicircum{}mu Gamma\_mu: the connection term is PRESENT}, \cfL{LAGRANGIAN (f) [control]: so Lun\_d is INADMISSIBLE -- its two field equations are not Dirac conjugates: they differ by 2 Psibar gamma\textasciicircum{}mu Gamma\_mu = -6 H Cot\textasciicircum{}2 Psibar T16[0], non-zero (witness u = e1, v = 0, V = s\textasciicircum{}2)}.
\end{itemize}

\subsubsection{It enters the field equation as \texorpdfstring{$-3H\cot^2\,\mathrm{T16}[0]$}{-3 H Cot2 T16[0]}, and the rescaling removes it}
\label{sec:9.12.3}

Varying $\bar\Psi$ and $\Psi$ independently in the symmetrized action gives the field equation
and its conjugate:
\begin{equation}
  \gamma^\mu D_\mu\Psi=H\,V'(s)\,\Psi,\qquad (D_\mu\bar\Psi)\gamma^\mu=-H\,V'(s)\,\bar\Psi .
\end{equation}
\cfproved{VII \S27} \cfL{FIELD EQUATION [THE RESULT]: EL\_Psibar == (Sqrt[g]/H)(gamma\textasciicircum{}mu D\_mu Psi - H V'(s) Psi) for generic Psi, Psibar of all eight coordinates: the field equation is gamma\textasciicircum{}mu D\_mu Psi == H V'(s) Psi}, \cfL{FIELD EQUATION [THE RESULT]: EL\_Psi == -(Sqrt[g]/H)((D\_mu Psibar) gamma\textasciicircum{}mu + H V'(s) Psibar): the adjoint equation is (D\_mu Psibar) gamma\textasciicircum{}mu == -H V'(s) Psibar}.
For the real field of Part~VI the same equation is \cfproved{VI \S24}
\cflab{FABLE [THE RESULT]: EL == Sqrt[g] (2/H) sigma16 . ( gamma^mu D_mu Psi - H V'(s) Psi ): the field equation of fable is EXACTLY the Levi-Civita covariant Dirac equation  gamma^mu D_mu Psi == H V'(s) Psi}.

On the canonical frame the connection enters as follows \cfderived, from \eqref{csc:eq:sum}. In
the equation it is $\gamma^\mu\Gamma_\mu\Psi=-3H\cot^2(6Hx_0)\,\Tsix{0}\Psi$. In the conjugate it is
$-\bar\Psi\Gamma_\mu\gamma^\mu=+\bar\Psi\gamma^\mu\Gamma_\mu=-3H\cot^2(6Hx_0)\,\bar\Psi\Tsix{0}$,
using the per-$\mu$ anticommutator. Because
$\Tsix{0}=\bigl(\begin{smallmatrix}0&\mathrm{ID8}\\ \mathrm{ID8}&0\end{smallmatrix}\bigr)$, in the 8+8
split-octonion form $\Psi=(\psi_1,\psi_2)$:
\begin{align}
  e_a{}^\mu\,\bar\tau_a\bigl(\partial_\mu+\Gamma^{(2)}_\mu\bigr)\psi_2&=H\,V'(s)\,\psi_1
    &&\text{(upper 8 rows)},\\
  e_a{}^\mu\,\tau_a\bigl(\partial_\mu+\Gamma^{(1)}_\mu\bigr)\psi_1&=H\,V'(s)\,\psi_2
    &&\text{(lower 8 rows)},
\end{align}
where the connection term is $-3H\cot^2(6Hx_0)\,\psi_2$ in the upper rows and
$-3H\cot^2(6Hx_0)\,\psi_1$ in the lower rows, and $\tau_a=$ \cfcode{tau[a]}, $\bar\tau_a=$ \cfcode{taubar[a]},
$s=-\psi_1^\ddagger\sigma\psi_1+\psi_2^\ddagger\sigma\psi_2$, $\bar\psi_1=-\psi_1^\ddagger\sigma$
and $\bar\psi_2=\psi_2^\ddagger\sigma$. The conjugate equations are
$(\partial\bar\psi_2-\bar\psi_2\Gamma^{(2)})\tau_a e_a{}^\mu=-HV'\bar\psi_1$ and
$(\partial\bar\psi_1-\bar\psi_1\Gamma^{(1)})\bar\tau_a e_a{}^\mu=-HV'\bar\psi_2$.
\cfproved{VII \S27} \cfL{WRITTEN OUT (ii) [THE RESULT]: the split-octonion 8+8 form -- the type-1 rows are the taubar equations for psi2, the type-2 rows the tau equations for psi1 (tau[0] == taubar[0] == ID8)}, \cfL{WRITTEN OUT (ii) [THE RESULT]: the covariant 8+8 form -- rows 1..8 are e\_a\textasciicircum{}mu taubar[a] (d\_mu + Gamma\textasciicircum{}(2)\_mu) psi2 - H V' psi1, rows 9..16 are e\_a\textasciicircum{}mu tau[a] (d\_mu + Gamma\textasciicircum{}(1)\_mu) psi1 - H V' psi2; and the conjugate rows}, \cfL{WRITTEN OUT (ii) [has content]: with Psibar = Psi\textasciicircum{}ddag sigma16 and sigma16 = diag(-sigma, sigma): s == -psi1\textasciicircum{}ddag sigma psi1 + psi2\textasciicircum{}ddag sigma psi2 (psibar1 = -psi1\textasciicircum{}ddag sigma, psibar2 = psi2\textasciicircum{}ddag sigma)}.

\paragraph{The rescaling.} The rescaling of Part~V applies unchanged, because it is a statement
about the operator. $\Psi=\sqrt{\sin 6Hx_0}\,\Psi'$ turns the equation into
$\sqrt{\sin 6Hx_0}\,\bigl(\gamma^\mu\partial_\mu\Psi'-HV'(s)\Psi'\bigr)=0$ with
$s=\sin(6Hx_0)\,\bar\Psi'\Psi'$. The connection term is gone. \cfproved{VII \S27}
\cfL{WRITTEN OUT (iv) [THE RESULT]: Psi = Sqrt[Sin[6Hx0]] Psi' turns the field equation into Sqrt[Sin] (Cot T16[0] d\_0 Psi' + (1/q) T16[i] d\_i Psi' + T16[4] d\_4 Psi' + (1/p) T16[h] d\_h Psi' - H V'(Sin s') Psi'): NO connection term}, \cfL{WRITTEN OUT (iv) [THE RESULT]: and the adjoint equation into Sqrt[Sin] (Cot d\_0 Psi'bar T16[0] + ... + H V'(Sin s') Psi'bar), with no connection term either}. The $x_4$-only family
$\Psi=\sqrt{\sin 6Hx_0}\,\Psi'(x_4)$ solves the equation if and only if $V''s'=0$. That holds for
linear $V$ (the author's mass term), or for null data $s'=0$. For the complex field $s'$ is
conserved on the $x_4$-only equation, so null data give $x_4$-only solutions for any $V$.
\cfproved{VII \S29} \cfL{PRE-UNIVERSE (2) [THE RESULT]: Psi = Sqrt[Sin] Psi'(x4) satisfies the field equation iff d\_x0 V'(Sin[6Hx0] s') == 0, and d\_x0 V'(Sin s') == 6 H Cos[6 H x0] s' V''(Sin s'): V linear OR s' == 0}, \cfL{PRE-UNIVERSE (2) [has content]: s' = Psi'bar Psi' is conserved by the x4-only equations (d\_4 Psi' = -H V' T16[4] Psi', d\_4 Psi'bar = H V' Psi'bar T16[4]) for ANY V; so null data s' = 0 are solutions for any V, with rho == V(0); and for V = s\textasciicircum{}2 with s' != 0 the x0-dependence is real (witness)}. The real-field version is
\cfproved{VI \S24}
\cflab{FABLE [has content]: for the mass term V = -(2M/H) s, V' is constant and Psi'(x4) IS an exact solution}.

\paragraph{The frame-independent form.} Covariant constancy of $\gamma^\mu$
\eqref{csc:eq:commutator} together with a vanishing $\sum_\mu\{\gamma^\mu,\Gamma_\mu\}$ gives, for
any frame on which that anticommutator vanishes,
\begin{equation}
  \gamma^\mu\Gamma_\mu=\frac12\,\frac{1}{\sqrt g}\,\partial_\mu\bigl(\sqrt g\,\gamma^\mu\bigr)=-\frac12\,T_\mu\gamma^\mu,
  \qquad T_\mu=\partial_\mu\ln\frac{e_\mu{}^\mu}{\det e}\quad\text{(no sum)},
\end{equation}
where $T_\mu$ is the Weitzenb\"ock torsion vector of a diagonal frame \eqref{csc:eq:torsionvector}.
The second equality is a one-line computation for a diagonal frame, since $\sqrt g=\det e$ and
$\gamma^\mu=\Tsix{\mu}/e_\mu{}^\mu$ \cfderived. This links the canonical connection directly to the
fable-5.1 bridge: the connection term of the Dirac operator is the Weitzenb\"ock torsion vector. On
the canonical frame $T_\mu=\partial_\mu\ln\sin 6Hx_0$ (subsection~\ref{csc:fable51}), and the
rescaling by $\exp(\tfrac12\ln\sin)=\sqrt{\sin 6Hx_0}$ removes it. \cfproved{VII \S27}
\cfL{SPIN CONNECTION [has content]: Part V restated -- omega\textasciicircum{}W == 0 in the canonical frame, omega\textasciicircum{}LC == -contortion, T\_mu == d\_mu Log[Sin[6 H x0]], R == -T + B, and gamma\textasciicircum{}mu Gamma\_mu == -(1/2) T\_mu gamma\textasciicircum{}mu} (on the canonical frame); and on the time-dependent warped frames of the coupled system, Part~VIII, Section~32: \cfL{WEITZENBOECK [THE RESULT]: the torsion vector of the warped frame is a GRADIENT, T\_mu == d\_mu Log[Sin[6 H x0]/(A\textasciicircum{}3 B\textasciicircum{}3 C)]; on the canonical member it is Part V's torsionVectorFable51}, \cfL{WEITZENBOECK [THE RESULT]: gamma\textasciicircum{}mu Gamma\_mu == -(1/2) T\_mu gamma\textasciicircum{}mu on the warped frame -- one vector term, a gradient}.

\subsubsection{It enters neither the canonical anticommutator nor the symmetrized Hamiltonian density}
\label{sec:9.12.4}

The admissible theory is a dimensional reduction $\Psi=\Psi(x_0,x_1,x_2,x_3,x_4)$. It uses a finite
hidden coordinate volume $V_{\text{hid}}=\int dx_5\,dx_6\,dx_7$, which is an assumption: compact
timelike directions contain closed timelike curves. Time is $x_4$. The canonical anticommutator on
a slice $x_4=y_4$ is
\begin{equation}
  \{\Psi_a(x),\Psi^\ddagger_b(y)\}=\frac{H}{V_{\text{hid}}\sqrt{|g|}}\,G_{ab}\,\delta^4(x_{0..3}-y_{0..3}),\qquad G=-i\,\sigma_{16}\gamma^4 .
\end{equation}
\cfproved{VII \S28} \cfL{DIRAC BRACKET [THE RESULT]: \{Theta\_i, Theta\_j\} = (Omega\textasciicircum{}-1)\_ij gives \{Psi\_a, Psi\textasciicircum{}ddag\_b\} == i (K\textasciicircum{}-1)\_ab == (Gtilde\textasciicircum{}-1)\_ab and \{Psi\_a, Psi\_b\} == 0}, \cfL{ANTICOMMUTATOR [THE RESULT]: on the canonical frame (c = Sqrt[g]/H) \{Psi, Psi\textasciicircum{}ddag\} == (H/Sqrt[g]) J == H Cos[6 H x0] J  (J\textasciicircum{}-1 == J), which IS the design's (H/Sqrt[g])(-i sigma16 gamma\textasciicircum{}4)}. No spin connection appears in it. It is
fixed by the coefficient of $\partial_4\Psi$ in the Lagrangian, and $\Gamma_4=0$ in any case. On
the canonical frame $\gamma^4=\Tsix{4}$, so
$G=-i\,\Tsix{0}\Tsix{1}\Tsix{2}\Tsix{3}\Tsix{4}=J$ \cfderived. $J$ is the fundamental symmetry
of subsection~\ref{sec:9.12.6}.

The Hamiltonian density of the symmetrized Lagrangian is
\begin{equation}
  \mathcal H=\sqrt g\,\Bigl[-\frac{1}{2H}\sum_{k\neq4}\bigl(\bar\Psi\gamma^k\partial_k\Psi-\partial_k\bar\Psi\gamma^k\Psi\bigr)+V(s)\Bigr],
\end{equation}
and it contains \emph{no spin connection} \cfproved{VII \S28}
\cfL{HAMILTONIAN [THE RESULT]: the Hamiltonian density Pi\_Psi d\_4 Psi + d\_4 Psibar Pi\_Psibar - L of the symmetrized Lagrangian is Sqrt[g][-(1/(2H)) Sum\_\{k != 4\}(Psibar gamma\textasciicircum{}k d\_k Psi - d\_k Psibar gamma\textasciicircum{}k Psi) + V]: NO spin connection (fields of all eight coordinates)}. This follows from
subsection~\ref{sec:9.12.2}: the symmetrized $\hat L$ does not depend on the connection. The connection re-emerges
in the equation of motion. Integrating the symmetrized spatial derivative terms by parts produces
$\gamma^\mu\Gamma_\mu=\tfrac12\frac{1}{\sqrt g}\partial_\mu(\sqrt g\,\gamma^\mu)$, because the volume
factor and the curved gammas depend on $x_0$. \cfproved{VII \S28}
\cfL{HAMILTONIAN [THE RESULT]: (Sqrt[g]/H) sigma16 (gamma\textasciicircum{}mu D\_mu Psi - H V' Psi) == i Gtilde d\_4 Psi - h Psi for admissible Psi(x0..x4): the field equation is i Gtilde d\_4 Psi = h Psi}, \cfL{HAMILTONIAN [THE RESULT]: (1/2) Sum\_\{mu != 4\} d\_mu(Sqrt[g] sigma16 gamma\textasciicircum{}mu) == Sqrt[g] sigma16 Sum gamma\textasciicircum{}mu Gamma\_mu, and d\_4(Sqrt[g] sigma16 gamma\textasciicircum{}4) == 0: the connection term is the one that makes h Hermitian}.

The same rescaling organizes the one-particle space. $\sqrt g\,dx_0\,|\Psi|^2=|\Psi'|^2\,dz$ for
$\Psi'=\Psi/\sqrt{\sin 6Hx_0}$, where $z=-\ln\cos(6Hx_0)/(6H)$ is proper distance along $x_0$. So the
one-particle space of $\Psi'$ is $L^2(dz)$. The $x_0$-operator is regular at the wall $z=0$ and
needs a boundary condition there. The Lorentz-covariant, $J$-compatible ones are
$\Tsix{0}\Psi'=\pm\Psi'$ at $z=0$. It is limit-point at $z\to\infty$, where no condition is needed.
\cfproved{VII \S28} \cfL{BOUNDARY [has content]: the x0 direction in proper distance z = -Log[Cos[6 H x0]]/(6 H): dz/dx0 == Tan[6Hx0] == Sec Sin (so Sqrt[g] dx0 |Psi|\textasciicircum{}2 = |Psi'|\textasciicircum{}2 dz), z(0) == 0 and z -> infinity at 6 H x0 -> Pi/2, gamma\textasciicircum{}0 d\_0 == T16[0] d\_z, and Sin[6Hx0] == Sqrt[1 - Exp[-12 H z]] (bounded frame factors)}, \cfL{BOUNDARY [THE RESULT]: the bag condition T16[0] Psi' = +-Psi' at the wall -- T16[0]\textasciicircum{}2 == 1, \{beta, T16[0]\} == 0, the normal-flux form vanishes on both eigenspaces (P beta T16[0] P == 0), and [J, T16[0]] == 0}.

\subsubsection{It does not drop out of the energy--momentum tensor}
\label{sec:9.12.5}

The energy--momentum tensor operator is
\begin{equation}
  \hat T_{\mu\nu}=-\frac{1}{4H}\Bigl[\bar\Psi\gamma_\mu D_\nu\Psi-(D_\nu\bar\Psi)\gamma_\mu\Psi+(\mu\leftrightarrow\nu)\Bigr]+g_{\mu\nu}\hat L,
  \qquad \gamma_\mu=g_{\mu\nu}\gamma^\nu .
\end{equation}
This is the Hilbert (symmetric-tetrad) tensor of the symmetrized action, and it equals the
covariant tensor $T_{\text{cov}}$ off shell. It is \emph{built with $D$, not $\partial$}. The
difference from the tensor $T_{\text{par}}$ built with partial derivatives is \cfderived
\begin{equation}
  T^{\text{cov}}_{\mu\nu}-T^{\text{par}}_{\mu\nu}=-\frac{1}{4H}\,\bar\Psi\bigl(\{\gamma_\mu,\Gamma_\nu\}+\{\gamma_\nu,\Gamma_\mu\}\bigr)\Psi .
\end{equation}
The $g_{\mu\nu}\hat L$ terms agree by subsection~\ref{sec:9.12.2}. For $\mu=\nu$ the difference is
$-\frac{1}{2H}\bar\Psi\{\gamma_\mu,\Gamma_\mu\}\Psi=0$ by the per-$\mu$ anticommutator, because
$\gamma_\mu=g_{\mu\mu}\gamma^\mu$ on the diagonal frame. So the energy density and all the pressures
are the same with $D$ or $\partial$. Not every off-diagonal component differs. The $(0,4)$
component coincides, because $\Gamma_0=\Gamma_4=0$. The $(5,6)$ component coincides too, because
there the two anticommutators cancel \cfderived: $\gamma_5=-p\,\Tsix{5}$, $\gamma_6=-p\,\Tsix{6}$,
the coefficients of $\Gamma_5$ and $\Gamma_6$ are equal, and
$\Tsix{5}\Tsix{a}\Tsix{6}=-\Tsix{6}\Tsix{a}\Tsix{5}$ for $a=0,4$. The pair at which the
notebook shows the two tensors to differ is given in the third item below.
\begin{itemize}
  \item \cfproved{VII \S29} Hilbert $=T_{\text{cov}}$, via a first-order symmetric perturbation of the
    frame, \cfL{EMT [THE RESULT]: the Hilbert (symmetric-tetrad) tensor of the symmetrized action IS That, OFF shell -- components (x1,x1), (x4,x4), (x1,x4), (x0,x5), (x5,x4)}.
  \item \cfproved{VII \S29} the piece coming from the variation of the spin connection is zero,
    \cfL{EMT [has content]: the spin-connection variation contributes NOTHING -- dropping it leaves the Hilbert tensor unchanged at (x1,x4), (x0,x5), (x5,x4)}. The Lagrangian depends on the connection
    only through $\sum_\mu\{\gamma^\mu,\Gamma_\mu\}=\tfrac12\omega_{[cab]}\gamma^{cab}$
    \eqref{csc:eq:antisym}, and under a symmetric variation of the tetrad
    $\delta\omega_{[cab]}=0$. So the spin-connection variation contributes \emph{exactly nothing}.
    The covariant terms of $\hat T$ come from the dependence of $\{\gamma^\mu,\Gamma_\mu\}$ on the
    frame. A partial-derivative Lagrangian must never be varied with respect to the frame.
    (Part~VI's Text cell described this contribution as ``a term antisymmetric in mu, nu that
    drops from the symmetrised tensor''. The statement above is the sharper one.)
  \item \cfproved{VII \S29} Hilbert $\neq T_{\text{par}}$, witnessed on the non-vacuous pair $(x_1,x_4)$ (the
    pairs $(0,4)$ and $(5,6)$ coincide on the canonical frame, as shown above),
    \cfL{EMT [control]: the Hilbert tensor is NOT the partial-derivative tensor: at (x1,x4) they differ -- witnessed on constant spinors (Psibar = e1, Psi = e11: (T16[4] T16[1] T16[0])\_\{1,11\} != 0)}.
  \item \cfproved{VII \S29} $T_{\text{cov}}$ is conserved on shell, \cfL{EMT CONSERVATION [THE RESULT]: nabla\textasciicircum{}mu That\_\{mu nu\} == 0 ON SHELL in all eight components, for generic Psi(x0, x4), Psibar(x0, x4) and generic V}.
\end{itemize}
The real-field analogues are already proved. \cfproved{VI \S24}
\cflab{FABLE [definition]: T_cov is symmetric by construction}. \cfproved{VI \S24}
\cflab{FABLE [has content]: the diagonal of T_cov equals the diagonal of the partial-derivative tensor, for the generic Psi(x0, x4): rho and every pressure are the same in both}.
\cfproved{VI \S24}
\cflab{FABLE [has content]: but the two tensors DIFFER off the diagonal -- witnessed on the (x3, x4) component T[[4,5]], a momentum density along the observed sheet, for the constant spinor e1 + e13 and V = s^2}.
\cfproved{VI \S24}
\cflab{FABLE [THE RESULT]: nabla^mu T_{mu nu} == 0 ON SHELL in all eight components, for a generic Psi(x0, x4) -- the covariant tensor is the conserved one}.
That only the covariant tensor is conserved, and $T_{\text{par}}$ is not, is \cfprose{VI \S24}: the
divergence of $T_{\text{par}}$ is not computed symbolically.

\subsubsection{The connection and the fundamental symmetry \texorpdfstring{$J$}{J}}
\label{sec:9.12.6}

\begin{equation}
  J=-i\,\Tsix{0}\Tsix{1}\Tsix{2}\Tsix{3}\Tsix{4},\qquad J^2=\IDs .
\end{equation}
$J$ is the fundamental symmetry that gives the positive quantization of the admissible sector.
With $\beta=-i\,\Tsix{4}$, $\sigma_{16}=J\beta$ holds, and the Hilbert adjoint is
$\Psi^\dagger=\Psi^\ddagger J$. \cfproved{VII \S28} \cfL{J [THE RESULT]: sigma16 == J beta with beta = -i T16[4] Hermitian, beta\textasciicircum{}2 == 1; T16[0..3] Hermitian with beta T16[j] beta == -T16[j]: beta is the observer's Dirac adjoint}, \cfL{J [THE RESULT]: with Psi\textasciicircum{}dagger := Psi\textasciicircum{}ddag J:  Psibar = Psi\textasciicircum{}dagger beta (J sigma16 == beta),  Psi\textasciicircum{}ddag G = Psi\textasciicircum{}dagger (J G == ID16),  \{Psi, Psi\textasciicircum{}dagger\} == (H/Sqrt[g]) ID16 (G\textasciicircum{}-1 J == ID16): POSITIVE}.

$J^2=\IDs$ follows directly \cfderived. The product $P=\Tsix{0}\cdots\Tsix{4}$ of five mutually
anticommuting matrices satisfies
$P^2=(-1)^{5\cdot4/2}\,\Tsix{0}^2\cdots\Tsix{4}^2=(+1)(1)(1)(1)(1)(-1)=-\IDs$, so
$(-iP)^2=+\IDs$.

The commutation rules also follow by counting sign changes \cfderived. A single $\Tsix{a}$
anticommutes with the four factors of $P$ other than itself and commutes with itself, if
$a\in\{0..4\}$. It anticommutes with all five if $a\in\{5,6,7\}$. Hence $J$ \emph{commutes} with
$\Tsix{0..4}$ and \emph{anticommutes} with $\Tsix{5},\Tsix{6},\Tsix{7}$. Applying this to the
explicit $\Gamma_\mu$ of subsection~\ref{csc:spinmatrix}:

\begin{center}
{\small
\begin{tabular}{@{}p{0.28\linewidth}p{0.19\linewidth}p{0.43\linewidth}@{}}
\toprule
\raggedright object & \raggedright relation to $J$ & \raggedright reason \tabularnewline
\midrule
\raggedright $\Gamma_0=\Gamma_4=0$ & \raggedright commute & \raggedright zero \tabularnewline
\raggedright $\Gamma_j$ ($j=1,2,3$) & \raggedright \textbf{commute} & \raggedright made of $\Tsix{0}\Tsix{j}$, $\Tsix{4}\Tsix{j}$, all factors in $\{0..4\}$ \tabularnewline
\raggedright $\Gamma_k$ ($k=5,6,7$) & \raggedright \textbf{anticommute} & \raggedright made of $\Tsix{0}\Tsix{k}$, $\Tsix{4}\Tsix{k}$: one factor in $\{0..4\}$, one in $\{5,6,7\}$ \tabularnewline
\raggedright each $\gamma^\mu\Gamma_\mu$ (no sum), and their sum & \raggedright \textbf{commute} & \raggedright each is $\alpha_\mu\Tsix{0}+\beta_\mu\Tsix{4}$ (subsection~\ref{csc:dirac}) \tabularnewline
\bottomrule
\end{tabular}
}
\end{center}

The hidden components $\Gamma_5$, $\Gamma_6$ and $\Gamma_7$ are $J$-odd. They enter the Dirac
operator only through $\gamma^k\Gamma_k$, a product of two $J$-odd matrices, which is $J$-even. So
every term of the admissible-sector Hamiltonian, the connection term included, is $J$-even. The
$J$-odd objects, namely the $(0,h)$, $(i,h)$ and $(4,h)$ components of $T_{\text{cov}}$ and the
hidden momenta $P_h$ ($h=5,6,7$), are nonzero operators with zero expectation in $J$-eigenmode Fock
states. \cfproved{VII \S28} \cfL{J TABLE [THE RESULT]: J commutes with sigma16, T16[0..4], sigma16 gamma\textasciicircum{}mu (mu = 0..4), every gamma\textasciicircum{}mu Gamma\_mu and Gamma\_1..Gamma\_3; it ANTIcommutes with T16[5..7], sigma16 gamma\textasciicircum{}h (h = 5..7), Gamma\_5..Gamma\_7, T16[8] and C+}. The same $J$ table holds on the dynamical
(warped and Bianchi-I) frames of the coupled system. Part~VIII, Section~32, proves it:
\cfL{J [THE RESULT]: [J, gamma\textasciicircum{}mu Gamma\_mu] == 0 for EACH mu separately (h = 5,6,7 included), on the warped and on the Bianchi-I frame}, \cfL{J [THE RESULT]: [J, Gamma\_mu] == 0 for mu = 0..4 and \{J, Gamma\_h\} == 0 for h = 5, 6, 7, on the warped and on the Bianchi-I frame} and
\cfL{J [has content]: the anticommuting Gamma\_h are not zero (witnessed on A = 1 + x4/2, B = 1 + x4/3, C = 1 + x4\textasciicircum{}2/3), so \{J, Gamma\_h\} == 0 is not vacuous}.

% --------------------------------------------------------------------------------------------
\subsection{Status ledger}
\label{csc:ledger}\label{sec:9.13}

{\footnotesize
\begin{longtable}{@{}p{0.47\linewidth}p{0.17\linewidth}p{0.28\linewidth}@{}}
\toprule
\raggedright statement & \raggedright status & \raggedright where (notebook Part and Section; assertion tag) \tabularnewline
\midrule
\endhead
\raggedright frame $e$, $g=e\,\eta\,e^T$, stated line element, signature (4,4) & \raggedright proved & \raggedright III \S15 \tabularnewline
\raggedright $a_4$ is free and never given a value & \raggedright proved & \raggedright VI \S25, VII \S29 and VIII (controls); notice printed in III \S15 \tabularnewline
\raggedright $\det g=\sec^2 6Hx_0$; $\sqrt{|\det g|}=\sec 6Hx_0$ on $0<6Hx_0<\pi/2$ & \raggedright displayed; proved & \raggedright displayed III \S15; proved V \S21 \tabularnewline
\raggedright 37 Christoffel symbols; symmetric in the lower indices; metric compatible & \raggedright displayed; proved & \raggedright III \S16; VI \S23 \tabularnewline
\raggedright the vielbein postulate is solved by \cfcode{cfSpinConnection} (residual zero) & \raggedright proved (solver regression) & \raggedright III \S16 \tabularnewline
\raggedright the 24 non-zero components of \cfcode{omegaCanonical} & \raggedright displayed; independent frame-only derivation proved; count 24 proved & \raggedright III \S16; VII \S27 (\cfcode{SPIN CONNECTION [has content]}) \tabularnewline
\raggedright antisymmetry in the flat indices (metric compatibility) & \raggedright proved (has content) & \raggedright III \S16; VII \S27 \tabularnewline
\raggedright torsion-free & \raggedright proved (structural) & \raggedright III \S16 \tabularnewline
\raggedright curvature: 156 components; Ricci scalar & \raggedright displayed & \raggedright III \S16 \tabularnewline
\raggedright $\Gamma_\mu=\tfrac18\omega_{\mu ab}[\Tsix{a},\Tsix{b}]$; sign and coefficient fixed by compatibility & \raggedright proved (with control) & \raggedright III \S17; VII \S27 (\cfcode{SPIN CONNECTION [definition]}) \tabularnewline
\raggedright block-diagonal; type-1 and type-2 blocks; \cfcode{Dcov16} = direct sum & \raggedright proved & \raggedright III \S17 \tabularnewline
\raggedright covariant constancy of $\gamma^\mu$ & \raggedright proved & \raggedright III \S17 \tabularnewline
\raggedright $\frac{1}{\sqrt g}\partial_\mu(\sqrt g\gamma^\mu)=[\gamma^\mu,\Gamma_\mu]$ & \raggedright proved & \raggedright VI \S24; VII \S27 (\cfcode{FIELD EQUATION [has content]}) \tabularnewline
\raggedright $\sum_\mu\{\gamma^\mu,\Gamma_\mu\}=0$; $\gamma^\mu\Gamma_\mu=-3H\cot^2(6Hx_0)\,\Tsix{0}$ & \raggedright proved & \raggedright VI \S24; VII \S27 \tabularnewline
\raggedright per-direction $\gamma^\mu\Gamma_\mu=\alpha_\mu\Tsix{0}+\beta_\mu\Tsix{4}$, $\sum\alpha_\mu=-3H\cot^2(6Hx_0)$, $\sum\beta_\mu=0$; anticommutator zero for each $\mu$ & \raggedright derived here; proved & \raggedright VII \S27 (\cfcode{SPIN CONNECTION [THE RESULT]}, \cfcode{LAGRANGIAN (c) [THE RESULT]}) \tabularnewline
\raggedright Bridges 1 and 2: the same connection in a local / constant gauge, at vector and spinor level & \raggedright proved & \raggedright IV \S18, IV \S19; restated VII \S27 \tabularnewline
\raggedright Bridge 3: a different connection; torsion $-2\lambda m^{\text{Skew}}$; curvature of degree 2 in $\lambda$; Ricci shift $-42\lambda^2=-\tfrac14T\cdot T$ & \raggedright proved & \raggedright IV \S20; restated VII \S27 \tabularnewline
\raggedright fable-5.1: $\omega=0$, zero curvature, $\omega^{LC}=-\text{contortion}(T)$, $T_\mu=\partial_\mu\ln\sin 6Hx_0$, axial part zero, $R=-T+B$ & \raggedright proved & \raggedright V \S21; restated VII \S27 \tabularnewline
\raggedright $\gamma^\mu\Gamma_\mu=-\tfrac12T_\mu\gamma^\mu$; the rescaling $\Psi=\sqrt{\sin 6Hx_0}\,\Psi'$ & \raggedright proved & \raggedright V \S21; VI \S24; VII \S27 (\cfcode{SPIN CONNECTION [has content]}, \cfcode{WRITTEN OUT (iv)}); VIII \S32 on the warped frames (\cfcode{WEITZENBOECK}, \cfcode{RESCALING}) \tabularnewline
\raggedright Bridge 1's gauge term is the fable-5.1 connection of the boosted frame & \raggedright proved & \raggedright V \S21 \tabularnewline
\raggedright master table of the five connections & \raggedright displayed & \raggedright V \S22 \tabularnewline
\raggedright complex fable: the connection drops out of the symmetrized $L$ because the matrix $\{\gamma^\mu,\Gamma_\mu\}$ vanishes & \raggedright proved & \raggedright VII \S27 (\cfcode{LAGRANGIAN (c) [THE RESULT]}) \tabularnewline
\raggedright unsymmetrized covariant $L$ = symmetrized + divergence; unsymmetrized partial $L$ inadmissible & \raggedright proved & \raggedright VII \S27 (\cfcode{LAGRANGIAN (d)}, \cfcode{(e)}, \cfcode{(f)}) \tabularnewline
\raggedright field equations; connection term $-3H\cot^2(6Hx_0)\,\Tsix{0}$; 8+8 form; rescaling; frame-independent form & \raggedright proved & \raggedright VII \S27 (\cfcode{FIELD EQUATION}, \cfcode{WRITTEN OUT (i)}--\cfcode{(iv)}, \cfcode{SPIN CONNECTION}); VIII \S32 \tabularnewline
\raggedright no connection in the canonical anticommutator or in the symmetrized Hamiltonian density; the connection makes $h$ Hermitian & \raggedright proved & \raggedright VII \S28 (\cfcode{QUANTIZATION [has content]: Gamma_4 == 0}, \cfcode{ANTICOMMUTATOR}, \cfcode{HAMILTONIAN}) \tabularnewline
\raggedright EMT: Hilbert $=T_{\text{cov}}$; spin-connection variation zero; Hilbert $\neq T_{\text{par}}$; $T_{\text{cov}}$ conserved & \raggedright proved & \raggedright VII \S29 (\cfcode{EMT}, \cfcode{EMT CONSERVATION}) \tabularnewline
\raggedright $J$ commutes with each $\gamma^\mu\Gamma_\mu$ and with $\Gamma_0..\Gamma_4$, and anticommutes with $\Gamma_5..\Gamma_7$ & \raggedright derived here; proved & \raggedright VII \S28 (\cfcode{J TABLE [THE RESULT]}); VIII \S32 on the warped and Bianchi-I frames (\cfcode{J [THE RESULT]}) \tabularnewline
\bottomrule
\end{longtable}
}
\endgroup

% ===============================================================================================
\section{How it was verified}
\label{sec:10}

\subsection{Part VII of the notebook: sections, cells, assertions}
\label{sec:10.1}

Part~VII is the manifest \file{claude-fable/cells_part7.wl} (1777 lines; 170 \cfcode{cfAssert} calls).
Built into the notebook by \file{claude-fable/build_tools.py} (\secref{sec:11.1}), it adds 26 Input
cells, 173--198, to the 172 of Parts I--VI. The section-to-cell map is printed by
\file{claude-fable/nb_section_map.wls} (\file{claude-fable/nb_section_map.log}, \secref{sec:11.4}):

{\small
\begin{longtable}{@{}p{0.09\linewidth}p{0.35\linewidth}p{0.09\linewidth}p{0.09\linewidth}p{0.12\linewidth}p{0.1\linewidth}@{}}
\toprule
\raggedright notebook Section & \raggedright subject & \raggedright Input cells & \raggedright assertions & \raggedright seconds (sum of the cells' times) & \raggedright this page \tabularnewline
\midrule
\endhead
\raggedright 26 & \raggedright the refinement: why fermion fable is a complex 16-spinor & \raggedright 173--179 & \raggedright 38 & \raggedright 1.957 & \raggedright \ref{sec:3} \tabularnewline
\raggedright 27 & \raggedright the Lagrangian, the field equations in the primordial gravitational field, and the canonical spin connection & \raggedright 180--184 & \raggedright 36 & \raggedright 2.291 & \raggedright \ref{sec:4}, \ref{sec:5}, \ref{sec:9} \tabularnewline
\raggedright 28 & \raggedright canonical quantization in 4+4 dimensions: the Dirac bracket, the Krein structure $J$, and the admissible sector & \raggedright 185--191 & \raggedright 56 & \raggedright 12.311 & \raggedright \ref{sec:6} \tabularnewline
\raggedright 29 & \raggedright the energy--momentum tensor operator: definition, conservation, and $\rho$, $P$ and $w$ & \raggedright 192--198 & \raggedright 40 & \raggedright 17.874 & \raggedright \ref{sec:7}, \ref{sec:8} \tabularnewline
\midrule
\raggedright \textbf{Part VII} & & \raggedright \textbf{26 cells} & \raggedright \textbf{170} & \raggedright \textbf{34.433} & \tabularnewline
\bottomrule
\end{longtable}
}

The counts and times per cell, read from \file{claude-fable/run_fermion_fable_part7.log} (each cell's
\cfcode{PASS} lines are printed before its \cfcode{CELL n  t=...} line):

\begin{center}
{\footnotesize
\setlength{\tabcolsep}{3pt}
\begin{tabular}{@{}lccccccccccccc@{}}
\toprule
cell & 173 & 174 & 175 & 176 & 177 & 178 & 179 & 180 & 181 & 182 & 183 & 184 & 185 \\
\midrule
assertions & 4 & 2 & 9 & 7 & 7 & 4 & 5 & 10 & 5 & 3 & 5 & 13 & 5 \\
seconds & 0.016 & 0.541 & 0.117 & 0.592 & 0.053 & 0.064 & 0.574 & 1.005 & 0.403 & 0.395 & 0.209 & 0.279 & 0.062 \\
\bottomrule
\end{tabular}

\medskip
\begin{tabular}{@{}lccccccccccccc@{}}
\toprule
cell & 186 & 187 & 188 & 189 & 190 & 191 & 192 & 193 & 194 & 195 & 196 & 197 & 198 \\
\midrule
assertions & 5 & 9 & 10 & 9 & 6 & 12 & 5 & 4 & 7 & 3 & 11 & 9 & 1 \\
seconds & 0.017 & 0.344 & 0.129 & 0.221 & 1.680 & 9.858 & 2.248 & 4.431 & 4.080 & 2.098 & 4.945 & 0.072 & 0.000 \\
\bottomrule
\end{tabular}
}
\end{center}

The 170 assertions by tag (the text of a label before its first \texttt{[}), counted from the log:

{\small
\begin{longtable}{@{}p{0.33\linewidth}p{0.1\linewidth}p{0.4\linewidth}p{0.09\linewidth}@{}}
\toprule
\raggedright tag & \raggedright count & \raggedright tag & \raggedright count \tabularnewline
\midrule
\endhead
\raggedright \cfcode{REFINEMENT} & \raggedright 6 & \raggedright \cfcode{ADMISSIBILITY} (all items) & \raggedright 10 \tabularnewline
\raggedright \cfcode{GRASSMANN} & \raggedright 9 & \raggedright \cfcode{HAMILTONIAN} & \raggedright 7 \tabularnewline
\raggedright \cfcode{REFINEMENT (a)}, \cfcode{(b)}, \cfcode{(c)}, \cfcode{(d)} & \raggedright 3, 3, 1, 7 & \raggedright \cfcode{BOUNDARY} & \raggedright 2 \tabularnewline
\raggedright \cfcode{4D} & \raggedright 4 & \raggedright \cfcode{KREIN MODES} & \raggedright 6 \tabularnewline
\raggedright \cfcode{FIDELITY} & \raggedright 5 & \raggedright \cfcode{FOCK}, \cfcode{FOCK (field level)}, \cfcode{FOCK (mode level)} & \raggedright 1, 4, 6 \tabularnewline
\raggedright \cfcode{LAGRANGIAN (a)}--\cfcode{(f)} & \raggedright 1, 2, 1, 1, 2, 3 & \raggedright \cfcode{CHARGE} & \raggedright 1 \tabularnewline
\raggedright \cfcode{FIELD EQUATION} & \raggedright 5 & \raggedright \cfcode{EMT} & \raggedright 8 \tabularnewline
\raggedright \cfcode{CURRENT} & \raggedright 3 & \raggedright \cfcode{EMT CONSERVATION}, \cfcode{EMT ON SHELL}, \cfcode{EMT HERMITICITY} & \raggedright 3, 3, 2 \tabularnewline
\raggedright \cfcode{SPIN CONNECTION} & \raggedright 5 & \raggedright \cfcode{MODE SUMS}, \cfcode{VACUUM} & \raggedright 2, 1 \tabularnewline
\raggedright \cfcode{WRITTEN OUT (i)}--\cfcode{(iv)} & \raggedright 2, 3, 3, 3 & \raggedright \cfcode{KOHN-SHAM}, \cfcode{MEAN FIELD} & \raggedright 5, 1 \tabularnewline
\raggedright \cfcode{SYMMETRY} & \raggedright 2 & \raggedright \cfcode{AUTHOR'S MASS TERM}, \cfcode{CLASSICAL LIMIT} & \raggedright 4, 1 \tabularnewline
\raggedright \cfcode{QUANTIZATION} & \raggedright 5 & \raggedright \cfcode{PRE-UNIVERSE (1)}--\cfcode{(4)} & \raggedright 2, 2, 2, 2 \tabularnewline
\raggedright \cfcode{DIRAC BRACKET}, \cfcode{ANTICOMMUTATOR} & \raggedright 3, 2 & \raggedright \cfcode{PART VII} & \raggedright 1 \tabularnewline
\raggedright \cfcode{J}, \cfcode{J TABLE} & \raggedright 8, 1 & \raggedright ``no identity in this notebook was accepted on numerical evidence alone'' & \raggedright 1 \tabularnewline
\bottomrule
\end{longtable}
}

By grade: the labels carry \cfcode{[THE RESULT]} (a headline result), \cfcode{[has content]} (a
statement that could have failed), \cfcode{[definition]}, \cfcode{[structural]},
\cfcode{[solver regression]} (a check of the code), \cfcode{[control]} (the wrong alternative is
shown to fail) and \cfcode{[fidelity]} (agreement with the author's or an earlier Part's objects).

\paragraph{The final tally and the run summary,} verbatim from the end of
\file{claude-fable/run_fermion_fable_part7.log}:

\begin{Verbatim}[fontsize=\scriptsize]
identities accepted on numerical evidence alone : 0   (stage 3 returned True)
non-vanishing witnesses                         : 38   (cfNonZeroWitnessQ found a probe point at which every entry is a
                                                     number and one of them is non-zero: a complete proof of non-vanishing)
  PASS  no identity in this notebook was accepted on numerical evidence alone
Assertions run: 528   passed: 528   failed: 0
CELL 198  t=0.000 s
==================== RUN SUMMARY ====================
cells evaluated : 198
total seconds   : 232.146
cells w/ msgs   : 0
---- slowest cells ----
  cell 80  89.913 s
  cell 88  27.790 s
  cell 86  11.167 s
  cell 87  10.303 s
  cell 191  9.858 s
  cell 137  8.740 s
  cell 160  6.445 s
  cell 121  5.401 s
  cell 196  4.945 s
  cell 193  4.431 s
  cell 107  4.158 s
  cell 194  4.080 s
==================== ASSERTIONS ====================
assertions run  : 528
passed          : 528
FAILED          : 0
====================================================
RUN-DONE

\end{Verbatim}

So through Part~VII: 198 Input cells, \textbf{528 assertions, 528 passed, 0 failed}, \textbf{0 cells
with messages}, \textbf{0 identities accepted on numerical evidence alone} (every identity is closed
symbolically; the notebook's last assertion is \cfL{no identity in this notebook was accepted on numerical evidence alone}), and \textbf{38
non-vanishing witnesses} (19 from Parts I--V, 8 more from Part~VI, 11 more from Part~VII: the Part~V
and Part~VI tallies in the same log print 19 and 27). Parts I--VI account for 358 of the assertions;
Part~VII adds 170. The whole run took 232.146\,s; the slowest cells are Part~II's (cell 80, the
bilinears of the author's closed-form solution, 89.913\,s). The last line, \cfcode{RUN-DONE}, is not
printed by \file{run_from_nb.wls}: it was appended by the shell wrapper that ran it.

\subsection{Every Part VII assertion, verbatim}
\label{sec:10.2}

This is the complete output of Input cells 173--198 in \file{claude-fable/run_fermion_fable_part7.log}
(lines 634--844), from the line that closes cell 172 to the line that closes cell 198. Every label is
printed after \cfcode{PASS}; a bracketed \cfcode{[t s]} line is a timing printed by \cfcode{cfTimed}.
Lines longer than the page are broken, and the break is marked with a hook arrow.

\begin{Verbatim}[fontsize=\scriptsize]
CELL 172  t=0.000 s
  PASS  REFINEMENT [definition]: T16[8]^2 == sigma16^2 == C+^2 == ID16, C+ is symmetric, and C+ . T16[8] == sigma16 -- the four candidate mass matrices are only two
  PASS  REFINEMENT [has content]: C- T16[a] C-^-1 == -T16[a]^T and C+ T16[a] C+^-1 == +T16[a]^T for a = 0..7
  PASS  REFINEMENT [THE RESULT]: the symmetry of C Gamma depends only on the rank r of Gamma:  C-: S A A S S A A S S,  C+: S S A A S S A A S  (r = 0..8)
  PASS  REFINEMENT [has content]: for each C, 136 of the 256 products are symmetric and 120 antisymmetric
CELL 173  t=0.016 s
  PASS  REFINEMENT [THE RESULT]: the Spin(4,4)-invariant bilinear forms on the 16-spinor form a 2-dimensional space, spanned by C- = sigma16 and C+ = sigma16 T16[8]
  PASS  REFINEMENT [has content]: every invariant form is SYMMETRIC and chirality-diagonal, and the two chiral pieces PL sigma16 PL and PR sigma16 PR are invariant separately
CELL 174  t=0.541 s
  PASS  GRASSMANN [solver regression]: generators anticommute and square to zero (all pairs of 1..6)
  PASS  GRASSMANN [solver regression]: the product is associative, (x y) z == x (y z), on elements of degree 1, 2, 3 (and the products tested are not zero)
  PASS  GRASSMANN [solver regression]: graded commutativity  x y == (-1)^(|x| |y|) y x  (degrees 1, 2, 3), and x^2 == 0 for x odd
  PASS  GRASSMANN [solver regression]: the conjugation is an involutive anti-automorphism,  (x y)^ddag == y^ddag x^ddag,  (x^ddag)^ddag == x
  PASS  GRASSMANN [solver regression]: grD is an even derivation, d(x y) == d(x) y + x d(y);  grDL obeys d_i(x y) == d_i(x) y + (-1)^|x| x d_i(y)
  PASS  GRASSMANN [THE RESULT]: theta^T N theta == theta^T N_A theta and theta^T N_S theta == 0 for a general 16x16 N -- an anticommuting bilinear sees only the ANTISYMMETRIC part
  PASS  GRASSMANN [THE RESULT]: theta^T N phi - (1/2) d(theta^T N theta) == theta^T N_S phi for a general N
  PASS  GRASSMANN [THE RESULT]: the Euler-Lagrange expression of L = theta^T N phi (left derivatives) is (N + N^T) phi: EMPTY field equations exactly when N is antisymmetric
  PASS  GRASSMANN [control]: for COMMUTING components the roles are reversed -- theta^T N theta sees only N_S, and the Euler-Lagrange expression of theta^T N phi is 2 N_A phi
CELL 175  t=0.117 s
  PASS  REFINEMENT (a) [THE RESULT]: real Grassmann + sigma16: both invariant mass terms vanish, and theta^T sigma16 T16[a] phi == (1/2) d(theta^T sigma16 T16[a] theta) for every a (a total derivative)
  PASS  REFINEMENT (a) [has content]: its Euler-Lagrange expression is identically empty in every direction -- while the kinetic term itself is NOT zero
  PASS  REFINEMENT (a) [THE RESULT]: FRAME-INDEPENDENT -- sigma16 {T16[a], S^bc} is SYMMETRIC for all a, b, c (so theta^T sigma16 {gamma^mu, Gamma_mu} theta == 0 on every frame), and d_mu(Sqrt[g] sigma16 gamma^mu) == Sqrt[g] sigma16 [gamma^mu, Gamma_mu] on the canonical frame
  PASS  REFINEMENT (b) [THE RESULT]: Majorana with C+: the kinetic term is genuine -- Euler-Lagrange expression 2 C+ T16[a] phi with C+ T16[a] invertible -- but every invariant bilinear vanishes (massless, no V(s))
  PASS  REFINEMENT (b) [has content]: the scope -- the quartic Lorentz scalars v.v, B.B, C.C of the Majorana spinor are NOT zero and are all proportional to ONE quartic (same monomials, one constant ratio each)
  PASS  REFINEMENT (b) [has content]: and the commuting shadow of the C+ Majorana spinor has NO kinetic term: for commuting components theta^T C+ T16[a] phi is the total derivative (1/2) d(theta^T C+ T16[a] theta)
  PASS  REFINEMENT (c) [THE RESULT]: Weyl and Majorana-Weyl: PL C T16[a] PL == PR C T16[a] PR == 0 for C = C-, C+ and every a (C chirality-diagonal, T16[a] chirality-odd) -- no kinetic term
CELL 176  t=0.592 s
  PASS  REFINEMENT (d) [definition]: Psi^ddag == (theta1 - i theta2)^T/Sqrt[2] in the algebra
  PASS  REFINEMENT (d) [THE RESULT]: s = Psibar Psi == i theta1^T sigma16 theta2: NON-ZERO (16 monomials) and HERMITIAN; p = Psibar T16[8] Psi is non-zero and Hermitian too
  PASS  REFINEMENT (d) [has content]: s is a sum of 16 commuting nilpotent even elements x_a (x_a x_b == x_b x_a, x_a^2 == 0), so s^17 == 0 and V(s) is a finite polynomial
  PASS  REFINEMENT (d) [definition]: the 32x32 matrix A read off the algebra reproduces Psi^ddag K phi exactly
  PASS  REFINEMENT (d) [THE RESULT]: its symmetric part is A_S == (i/2) Omega, Omega = [[0, K], [-K, 0]] real SYMMETRIC and invertible: a genuine kinetic term
  PASS  REFINEMENT (d) [has content]: the symmetrized kinetic term is HERMITIAN and equals Psi^ddag K phi - (1/2) d(Psi^ddag K Psi); the unsymmetrized one is not Hermitian
  PASS  REFINEMENT (d) [has content]: the same holds in every direction a = 0..7: the symmetrized Psi^ddag sigma16 T16[a] phi is Hermitian, with symmetric part (i/2)[[0, sigma16 T16[a]], [-sigma16 T16[a], 0]] != 0
CELL 177  t=0.053 s
  PASS  4D [THE RESULT]: the 16 ordered products of T16[1..4] are linearly independent (the observer's algebra is M4(C)), and its commutant in M16(C) is 16-dimensional
  PASS  4D [has content]: Omega4^2 == -ID16, the hidden gammas h_A = Omega4 T16[A] (A = 0,5,6,7) commute with T16[1..4] and obey {h_A, h_B} == -2 eta_AB
  PASS  4D [THE RESULT]: the 16 products of the hidden gammas span the commutant, and the two algebras meet only in the multiples of ID16 (rank of the union 16 + 16 - 1 = 31): C^16 = C^4 (x) C^4, four 4D Dirac fermions
  PASS  4D [has content]: the observer's Lorentz generators T16[i] T16[j] (i < j in 1..4) commute with every hidden gamma
CELL 178  t=0.064 s
  [0.024 s]  Euler-Lagrange expressions of the complex fable with respect to Psibar, generic (x0, x4)
  [0.023 s]  Euler-Lagrange expressions of the complex fable with respect to Psi, generic (x0, x4)
  PASS  FIDELITY [1] [fidelity]: Psi -> psi, Psibar -> psi^T sigma16 turns the symmetrized covariant Lhat into Part VI's cfLhatFable EXACTLY (canonical frame, generic V)
  PASS  FIDELITY [2] [fidelity]: on the flat frame with V = -(2M/H) s and no volume factor, the same limit on Psi16 IS the author's La[] of Section 11
  PASS  FIDELITY [3] [fidelity]: sigma16 . EL_Psibar + EL_Psi in the real limit == Part VI's explicit variation cfELFable, component for component
  PASS  FIDELITY [3] [fidelity]: and on the flat frame with V = -(2M/H) s the same combination IS the author's eLa of Section 12
  PASS  FIDELITY [4] [fidelity]: the energy-momentum tensor That in the real limit IS Part VI's covariant tensor cfTCov, all 64 components
CELL 179  t=0.574 s
  PASS  LAGRANGIAN (a) [definition]: every Gamma^spin_mu is real and sigma16 Gamma_mu is antisymmetric (Gamma_mu^T sigma16 == -sigma16 Gamma_mu): D_mu Psibar is the Dirac conjugate of D_mu Psi
  PASS  LAGRANGIAN (b) [THE RESULT]: Lsym is HERMITIAN -- real for Psi = u + i v, Psibar = (u - i v)^T sigma16, generic u, v of all eight coordinates
  PASS  LAGRANGIAN (b) [control]: Lun is NOT Hermitian -- its imaginary part is not zero, witnessed on the constant spinors u = e1, v = e13 (with V = s^2)
  PASS  LAGRANGIAN (c) [THE RESULT]: Lsym - Lsym_d == (1/(2H)) Sum_mu Psibar {gamma^mu, Gamma_mu} Psi identically, and {gamma^mu, Gamma_mu} == 0 for EACH mu on the canonical frame: Lsym == Lsym_d
  PASS  LAGRANGIAN (d) [has content]: Lun - Lsym == (1/(2H)) (1/Sqrt[g]) d_mu( Sqrt[g] Psibar gamma^mu Psi ), a total divergence
  PASS  LAGRANGIAN (e) [has content]: Lun - Lun_d == (1/H) Psibar gamma^mu Gamma_mu Psi == -3 Cot[6 H x0]^2 Psibar T16[0] Psi
  PASS  LAGRANGIAN (e) [control]: that term is NOT zero for a complex field (witness u = e1, v = e13), and IS zero in the real commuting limit (the scope of Parts V-VI)
  PASS  LAGRANGIAN (f) [THE RESULT]: Lun_d varied with respect to Psibar gives (Sqrt[g]/H)(gamma^mu d_mu Psi - H V' Psi): the connection term is MISSING
  PASS  LAGRANGIAN (f) [THE RESULT]: varied with respect to Psi it gives -(Sqrt[g]/H)((d_mu Psibar) gamma^mu + Psibar [gamma^mu, Gamma_mu] + H V' Psibar), with Psibar [gamma^mu, Gamma_mu] == 2 Psibar gamma^mu Gamma_mu: the connection term is PRESENT
  PASS  LAGRANGIAN (f) [control]: so Lun_d is INADMISSIBLE -- its two field equations are not Dirac conjugates: they differ by 2 Psibar gamma^mu Gamma_mu = -6 H Cot^2 Psibar T16[0], non-zero (witness u = e1, v = 0, V = s^2)
CELL 180  t=1.005 s
  [0.134 s]  Euler-Lagrange expressions with respect to Psibar, generic fields of all eight coordinates
  [0.118 s]  Euler-Lagrange expressions with respect to Psi, generic fields of all eight coordinates
  PASS  FIELD EQUATION [has content]: the identity the derivation uses -- (1/Sqrt[g]) d_mu(Sqrt[g] gamma^mu) == [gamma^mu, Gamma_mu] (covariant constancy, Section 17)
  PASS  FIELD EQUATION [THE RESULT]: EL_Psibar == (Sqrt[g]/H)(gamma^mu D_mu Psi - H V'(s) Psi) for generic Psi, Psibar of all eight coordinates: the field equation is gamma^mu D_mu Psi == H V'(s) Psi
  PASS  FIELD EQUATION [THE RESULT]: EL_Psi == -(Sqrt[g]/H)((D_mu Psibar) gamma^mu + H V'(s) Psibar): the adjoint equation is (D_mu Psibar) gamma^mu == -H V'(s) Psibar
  PASS  FIELD EQUATION [has content]: the adjoint equation IS the Dirac conjugate of the field equation: with Psi = u + i v, Psibar = (u - i v)^T sigma16, adjoint residual == -(field residual)^* . sigma16
  PASS  FIELD EQUATION [control]: the opposite sign, (D Psibar) gamma == +H V' Psibar, is NOT the conjugate -- witnessed on u = e1 + e5, v = 0 with V = s^2 (s = -2 there)
CELL 181  t=0.403 s
  PASS  CURRENT [THE RESULT]: the Noether current of the U(1) phase is (Sqrt[g]/H) i Psibar gamma^mu Psi
  PASS  CURRENT [has content]: i sigma16 gamma^mu is Hermitian for every mu, so j^mu is real
  PASS  CURRENT [THE RESULT]: OFF SHELL (1/Sqrt[g]) d_mu(Sqrt[g] j^mu) == i (Ebar Psi + Psibar E) for generic fields of all eight coordinates -- conserved on shell
CELL 182  t=0.395 s
  PASS  SPIN CONNECTION [has content]: metric compatible (omega_mu^{ab} antisymmetric) and every non-zero component is omega_mu^{mu 0} or omega_mu^{mu 4} (up to antisymmetry): omega_0 == omega_4 == 0
  PASS  SPIN CONNECTION [definition]: Gamma^spin_mu == Sum_{a<b} omega_mu,ab (1/2) T16[a] T16[b]
  PASS  SPIN CONNECTION [THE RESULT]: direction by direction gamma^mu Gamma_mu == alpha_mu T16[0] + beta_mu T16[4]; Sum alpha_mu == -3 H Cot[6 H x0]^2 and Sum beta_mu == 0 (the observed and hidden a4 terms cancel)
  PASS  SPIN CONNECTION [has content]: Parts IV-V restated -- Bridge 1's difference is the pure gauge term, Bridge 2 is a constant conjugate, Bridge 3 has totally antisymmetric torsion that is not zero
  PASS  SPIN CONNECTION [has content]: Part V restated -- omega^W == 0 in the canonical frame, omega^LC == -contortion, T_mu == d_mu Log[Sin[6 H x0]], R == -T + B, and gamma^mu Gamma_mu == -(1/2) T_mu gamma^mu
CELL 183  t=0.209 s
  PASS  WRITTEN OUT (i) [THE RESULT]: the field equation on the canonical frame is Cot T16[0] d_0 Psi + (1/q) T16[i] d_i Psi + T16[4] d_4 Psi + (1/p) T16[h] d_h Psi - 3 H Cot^2 T16[0] Psi == H V'(s) Psi
  PASS  WRITTEN OUT (i) [THE RESULT]: the adjoint equation is Cot d_0 Psibar T16[0] + ... - 3 H Cot^2 Psibar T16[0] == -H V'(s) Psibar (the connection enters with the sign fixed by {gamma^mu, Gamma_mu} = 0)
  PASS  WRITTEN OUT (ii) [THE RESULT]: the split-octonion 8+8 form -- the type-1 rows are the taubar equations for psi2, the type-2 rows the tau equations for psi1 (tau[0] == taubar[0] == ID8)
  PASS  WRITTEN OUT (ii) [THE RESULT]: the covariant 8+8 form -- rows 1..8 are e_a^mu taubar[a] (d_mu + Gamma^(2)_mu) psi2 - H V' psi1, rows 9..16 are e_a^mu tau[a] (d_mu + Gamma^(1)_mu) psi1 - H V' psi2; and the conjugate rows
  PASS  WRITTEN OUT (ii) [has content]: with Psibar = Psi^ddag sigma16 and sigma16 = diag(-sigma, sigma): s == -psi1^ddag sigma psi1 + psi2^ddag sigma psi2 (psibar1 = -psi1^ddag sigma, psibar2 = psi2^ddag sigma)
  PASS  WRITTEN OUT (iii) [definition]: the term lists reproduce all sixteen field equations and all sixteen adjoint equations exactly
  PASS  WRITTEN OUT (iii) [has content]: every component equation has exactly one derivative per coordinate x0..x7, the connection term multiplies the component carrying the x0-derivative, and -H V' multiplies psi_r itself
  PASS  SYMMETRY [THE RESULT]: the only internal U(1) is the vector phase -- the commutant of all eight T16[a] in M16(C) is one-dimensional (multiples of ID16)
  PASS  SYMMETRY [has content]: the axial phase exp(i alpha T16[8]) leaves s invariant but not the kinetic term: exp(-i alpha T8) T16[a] exp(i alpha T8) == T16[a] exp(2 i alpha T8), so Psibar gamma^mu d_mu Psi changes at first order by 2 i alpha Psibar gamma^mu T16[8] d_mu Psi, which is not zero (witness)
  PASS  WRITTEN OUT (iii) [has content]: for fields of all eight coordinates the sixteen equations form ONE coupled block (the four blocks of four of Section 13 belong to the (x0, x4) reduction)
  PASS  WRITTEN OUT (iv) [THE RESULT]: Psi = Sqrt[Sin[6Hx0]] Psi' turns the field equation into Sqrt[Sin] (Cot T16[0] d_0 Psi' + (1/q) T16[i] d_i Psi' + T16[4] d_4 Psi' + (1/p) T16[h] d_h Psi' - H V'(Sin s') Psi'): NO connection term
  PASS  WRITTEN OUT (iv) [THE RESULT]: and the adjoint equation into Sqrt[Sin] (Cot d_0 Psi'bar T16[0] + ... + H V'(Sin s') Psi'bar), with no connection term either
  PASS  WRITTEN OUT (iv) [definition]: under the rescaling s == Sin[6 H x0] s',  s' = Psi'bar Psi'
CELL 184  t=0.279 s
  PASS  QUANTIZATION [definition]: lapse 1 and shift 0 (g_44 == -1, g_4mu == 0), and the slice x4 = const has signature (4,3) -- x0..x3 spacelike, x5..x7 timelike (given a4 real and 0 < 6 H x0 < Pi/2)
  PASS  QUANTIZATION [THE RESULT]: the momenta of Lsym are Pi_Psi == (Sqrt[g]/(2H)) Psibar gamma^4 and Pi_Psibar == -(Sqrt[g]/(2H)) gamma^4 Psi
  PASS  QUANTIZATION [has content]: theta_un - theta_sym == delta[(Sqrt[g]/(2H)) Psibar gamma^4 Psi] exactly: the same symplectic form, the same brackets
  PASS  QUANTIZATION [THE RESULT]: G = -i sigma16 gamma^4 == J = -i T16[0].T16[1].T16[2].T16[3].T16[4] on the canonical frame; J is Hermitian, J^2 == ID16, eigenvalues +1 (8) and -1 (8): signature (8,8)
  PASS  QUANTIZATION [has content]: Gamma_4 == 0 -- the time-derivative part of the Lagrangian, which fixes the anticommutator, contains no spin connection
CELL 185  t=0.062 s
  PASS  DIRAC BRACKET [THE RESULT]: in the Grassmann algebra Psi^ddag K phi == (i/2) Theta^T Omega_c Phi + (1/4) d(theta1^T K theta1 + theta2^T K theta2), for K = c sigma16 T16[4] with a free scalar c
  PASS  DIRAC BRACKET [has content]: Omega_c is real symmetric and invertible, Omega_c^-1 == [[0, -K^-1], [K^-1, 0]]
  PASS  DIRAC BRACKET [THE RESULT]: {Theta_i, Theta_j} = (Omega^-1)_ij gives {Psi_a, Psi^ddag_b} == i (K^-1)_ab == (Gtilde^-1)_ab and {Psi_a, Psi_b} == 0
  PASS  ANTICOMMUTATOR [THE RESULT]: on the canonical frame (c = Sqrt[g]/H) {Psi, Psi^ddag} == (H/Sqrt[g]) J == H Cos[6 H x0] J  (J^-1 == J), which IS the design's (H/Sqrt[g])(-i sigma16 gamma^4)
  PASS  ANTICOMMUTATOR [has content]: the rescaled field Psi' = Psi/Sqrt[Sin[6 H x0]] has {Psi', Psi'^ddag} == H Cot[6 H x0] J delta^7, and Sqrt[g] is x4-independent (so quantizing chi = (Sqrt[g]/H)^(1/2) Psi changes nothing on the canonical frame)
CELL 186  t=0.017 s
  PASS  J [THE RESULT]: J == -i Omega4 T16[0] == -i h_0 is a hidden flavour gamma: it lies in the commutant of T16[1..4] (in the span of the hidden-gamma products of Section 26)
  PASS  J [has content]: on the flavour factor J has signature (2,2): cfIdemObs is a rank-4 projector commuting with J, and J restricted to its image has eigenvalues +1, +1, -1, -1
  PASS  J [THE RESULT]: sigma16 == J beta with beta = -i T16[4] Hermitian, beta^2 == 1; T16[0..3] Hermitian with beta T16[j] beta == -T16[j]: beta is the observer's Dirac adjoint
  PASS  J [THE RESULT]: with Psi^dagger := Psi^ddag J:  Psibar = Psi^dagger beta (J sigma16 == beta),  Psi^ddag G = Psi^dagger (J G == ID16),  {Psi, Psi^dagger} == (H/Sqrt[g]) ID16 (G^-1 J == ID16): POSITIVE
  PASS  J TABLE [THE RESULT]: J commutes with sigma16, T16[0..4], sigma16 gamma^mu (mu = 0..4), every gamma^mu Gamma_mu and Gamma_1..Gamma_3; it ANTIcommutes with T16[5..7], sigma16 gamma^h (h = 5..7), Gamma_5..Gamma_7, T16[8] and C+
  PASS  J [has content]: the Hermiticity rule -- for classically Hermitian M, (J M)^dagger == J M iff [J, M] == 0 and == -J M iff {J, M} == 0: so s = Psi^ddag sigma16 Psi is Hermitian and p = Psi^ddag C+ Psi is ANTI-Hermitian
  PASS  J [THE RESULT]: 13 of the 28 Lorentz generators commute with J (a, b both in {0..4}: Spin(4,1), or both in {5,6,7}: Spin(3)); the other 15 anticommute -- 12 boosts x0..x3 <-> x5..x7 and 3 rotations x4 <-> x5..x7
  PASS  J [has content]: the Krein form G is invariant (S^T G + G S == 0) under exactly the 21 generators with a, b != 4 -- Spin(4,3), the group of the slice
  PASS  J [has content]: chirality -- {J, T16[8]} == 0, each chirality subspace is G-NULL (PL G PL == PR G PR == 0), and the two chiralities are orthogonal in the Hilbert product (PL, PR Hermitian, PL PR == 0)
CELL 187  t=0.344 s
  PASS  ADMISSIBILITY (1) [THE RESULT]: E_k^2 == (k_s^2 - k_h^2 + m^2) ID16 for every k: omega^2 = k_s^2 - k_h^2 + m^2
  PASS  ADMISSIBILITY (2) [THE RESULT]: for admissible k (k_h = 0), [J, E_k] == 0, E_k is J-Hermitian, and G J == ID16 is positive: J quantizes every admissible momentum positively
  PASS  ADMISSIBILITY (3) [THE RESULT]: below threshold (k = 5 e1 + 3 e5, m = 3) the boosted J' = S^-1 J S commutes with E_k and beta, J'^2 == 1, and G J' == S^2 is Hermitian positive definite
  PASS  ADMISSIBILITY (4) [THE RESULT]: AT threshold (k = 4 e1 + 5 e5, m = 3) E_k != 0 but E_k^2 == 0: nilpotent, not diagonalizable (rank 8: Jordan blocks of size 2)
  PASS  ADMISSIBILITY (5) [THE RESULT]: ABOVE threshold (k = (0, 3/4, 1, 0 | k5 = 2), m = 1) omega^2 == -23/16: every eigenvalue of E_k is +-i Sqrt[23]/4 = +-1.19896 i and every eigenvector is G-NEUTRAL (u^dagger G u == 0)
  PASS  ADMISSIBILITY [has content]: the general identity [E_k, beta] == 2 i T_k, T_k = Sum_a k_a T16[a], for every k (admissible and hidden components alike)
  PASS  ADMISSIBILITY (6) [THE RESULT]: traceless certificate -- every X commuting with beta, the four admissible E_k and one hidden-momentum E_k has Tr(G X) == 0, so no positive G J' exists; J itself is not in that commutant
  PASS  ADMISSIBILITY (6) [THE RESULT]: the analytic proof for a purely hidden momentum -- [E_k, beta] == 2 i k5 T16[5]; T16[5] anticommutes with G and T16[5]^dagger G T16[5] == -G (so a J' commuting with T16[5] has T5^dagger (G J') T5 == -(G J'))
  PASS  ADMISSIBILITY (7) [THE RESULT]: J is UNIQUE -- the joint commutant of beta and the admissible E_k is 8-dimensional, contains J, and every element of it commutes with J (so X^2 == 1 with J X > 0 forces X == J)
  PASS  ADMISSIBILITY [has content]: [P_4, P_h] == 0 -- the metric, the frame and the spin connection do not depend on x5, x6, x7
CELL 188  t=0.129 s
  PASS  HAMILTONIAN [THE RESULT]: the Hamiltonian density Pi_Psi d_4 Psi + d_4 Psibar Pi_Psibar - L of the symmetrized Lagrangian is Sqrt[g][-(1/(2H)) Sum_{k != 4}(Psibar gamma^k d_k Psi - d_k Psibar gamma^k Psi) + V]: NO spin connection (fields of all eight coordinates)
  PASS  HAMILTONIAN [THE RESULT]: (Sqrt[g]/H) sigma16 (gamma^mu D_mu Psi - H V' Psi) == i Gtilde d_4 Psi - h Psi for admissible Psi(x0..x4): the field equation is i Gtilde d_4 Psi = h Psi
  PASS  HAMILTONIAN [THE RESULT]: (1/2) Sum_{mu != 4} d_mu(Sqrt[g] sigma16 gamma^mu) == Sqrt[g] sigma16 Sum gamma^mu Gamma_mu, and d_4(Sqrt[g] sigma16 gamma^4) == 0: the connection term is the one that makes h Hermitian
  PASS  HAMILTONIAN [has content]: A_0 = Sqrt[g] sigma16 gamma^0 is real antisymmetric, and for h_0 = -(A_0 d_0 + (1/2) A_0'):  <f, h_0 g> - <h_0 f, g> == -d_0(f^dagger A_0 g), a pure boundary term
  PASS  BOUNDARY [has content]: the x0 direction in proper distance z = -Log[Cos[6 H x0]]/(6 H): dz/dx0 == Tan[6Hx0] == Sec Sin (so Sqrt[g] dx0 |Psi|^2 = |Psi'|^2 dz), z(0) == 0 and z -> infinity at 6 H x0 -> Pi/2, gamma^0 d_0 == T16[0] d_z, and Sin[6Hx0] == Sqrt[1 - Exp[-12 H z]] (bounded frame factors)
  PASS  BOUNDARY [THE RESULT]: the bag condition T16[0] Psi' = +-Psi' at the wall -- T16[0]^2 == 1, {beta, T16[0]} == 0, the normal-flux form vanishes on both eigenspaces (P beta T16[0] P == 0), and [J, T16[0]] == 0
  PASS  HAMILTONIAN [THE RESULT]: [J, h] == 0 on the admissible sector, for every a4, every x0-profile and every V' -- including the spin-connection term
  PASS  HAMILTONIAN [control]: a hidden derivative term breaks it -- [J, h] != 0 once Psi depends on x5 (witnessed on Psi = x5 e1)
  PASS  HAMILTONIAN [has content]: the x4-dependence of h enters only through 1/q: d_4 h Psi == (H a4'[H x4]) x (the observed-direction part of h), which commutes with J
CELL 189  t=0.221 s
  PASS  KREIN MODES [has content]: A = G h satisfies A^2 == w^2 ID16 and [A, J] == 0, and A is Hermitian (G and h commute) for every admissible k and m
  PASS  KREIN MODES [THE RESULT]: the four P_{s,eta} are projectors of trace (rank) 4, mutually orthogonal, summing to ID16, and Sum eta P_{s,eta} == J == G^-1 (completeness)
  PASS  KREIN MODES [THE RESULT]: for every mode, eta u^dagger h u == s w (energy = frequency), eta u^dagger sigma16 u == m/omega_u, eta u^dagger (dh/dk_j) u == k_j/omega_u: P X P == c P with c = eta s w, eta s m/w, eta s k_j/w
  PASS  KREIN MODES [has content]: dh/dk_j == -i sigma16 T16[j], and d omega/d k_j == k_j/omega for omega = +-w
  PASS  KREIN MODES [THE RESULT]: every mode of P_{s,eta} is a J-eigenvector: J P_{s,eta} == eta P_{s,eta}
  PASS  KREIN MODES [control]: a Krein-boosted pair u1 = (5/4)u+ + (3/4)u-, u2 = (3/4)u+ + (5/4)u- in the omega = +5 eigenspace is still G-orthonormal and still consists of modes, but not of J-eigenvectors; the induced J_U = G U U^dagger differs from J, and under it s is NOT Hermitian
CELL 190  t=1.680 s
  [0.749 s]  the CAR of 16 Jordan-Wigner operators
  PASS  FOCK [solver regression]: the Jordan-Wigner operators obey the positive CAR {f_i, f_j^dagger} == delta_ij, {f_i, f_j} == 0 on C^65536
  PASS  FOCK (field level) [THE RESULT]: {Psi_a, Psi^ddag_b} == J_ab == (G^-1)_ab with Psi = f, Psi^ddag = f^dagger J -- the Dirac-bracket anticommutator, realized on a positive Fock space
  PASS  FOCK (field level) [THE RESULT]: J h is Hermitian with eigenvalues +5 (8) and -5 (8); H_free and s are Hermitian operators, so is H_free + (3/7) s^2; p = Psi^ddag C+ Psi is ANTI-Hermitian
  [1.868 s]  the Heisenberg equation on the Fock space
  PASS  FOCK (field level) [THE RESULT]: Heisenberg -- [H_free, Psi_a] == -(G^-1 h Psi)_a for all 16 components, so i d_4 Psi = i [H, Psi] IS the field equation i G d_4 Psi = h Psi
  PASS  FOCK (field level) [control]: with the opposite sign of the anticommutator the Hamiltonian is -H_free and [H', Psi_a] == +(G^-1 h Psi)_a: the time-REVERSED equation
  PASS  FOCK (mode level) [THE RESULT]: {b_u, b_v^ddag} == eta_u delta_uv (the indefinite CAR demanded by Sum eta u u^dagger = G^-1), and the J-involution b^dagger = eta b^ddag restores {b, b^dagger} == 1
  PASS  FOCK (mode level) [has content]: the naive Krein 'norm' <0| b_u b_u^ddag |0> equals eta_u -- NEGATIVE for the four eta = -1 modes of each frequency -- while the Fock norm <0| b_u b_u^dagger |0> == 1
  PASS  FOCK (mode level) [THE RESULT]: Sum omega_u eta_u b_u^ddag b_u == Sum omega_u b_u^dagger b_u, and [Ham, b_u] == -omega_u b_u for every u (b_u(x4) = b_u e^(-i omega_u x4))
  PASS  FOCK (mode level) [THE RESULT]: the spectrum -- a UNIQUE ground state at -8 omega = -40, the Dirac sea (modes 9..16 filled, 1..8 empty); :Ham: = Ham + 40 >= 0
  PASS  FOCK (mode level) [THE RESULT]: the first excited level, omega = 5 above the sea, is 16-fold: 8 particle states (N = +1) and 8 antiparticle states (N = -1), with N = Sum b^dagger b - 8
  PASS  FOCK (mode level) [THE RESULT]: in every Fock basis state tested (the sea and the 16 first excited states) <n| b_u^ddag b_v |n> == eta_u n_u delta_uv -- the expectation rule of Section 29
  PASS  CHARGE [THE RESULT]: n^4 = -(i/H) Psibar gamma^4 Psi == (1/H) Psi^ddag G Psi == (1/H) Psi^dagger Psi, whose mode form Sum_u eta_u b_u^ddag b_u == Sum_u b_u^dagger b_u is the particle number (normal-ordered: particles - antiparticles)
CELL 191  t=9.858 s
  [0.693 s]  Hilbert tensor from first-order symmetric tetrad perturbations, five components
  PASS  EMT [THE RESULT]: the Hilbert (symmetric-tetrad) tensor of the symmetrized action IS That, OFF shell -- components (x1,x1), (x4,x4), (x1,x4), (x0,x5), (x5,x4)
  PASS  EMT [has content]: the spin-connection variation contributes NOTHING -- dropping it leaves the Hilbert tensor unchanged at (x1,x4), (x0,x5), (x5,x4)
  PASS  EMT [control]: the Hilbert tensor is NOT the partial-derivative tensor: at (x1,x4) they differ -- witnessed on constant spinors (Psibar = e1, Psi = e11: (T16[4] T16[1] T16[0])_{1,11} != 0)
  PASS  EMT [definition]: That is symmetric
  PASS  EMT [THE RESULT]: That is Hermitian under the classical conjugation: real for Psi = u + i v, Psibar = (u - i v)^T sigma16, all 64 components
CELL 192  t=2.248 s
  PASS  EMT CONSERVATION [definition]: F and Fbar solve the field equation and the adjoint equation for the x4-derivatives
  [0.767 s]  nabla^mu That_{mu nu}, generic Psi(x0, x4), Psibar(x0, x4)
  [2.599 s]  the on-shell divergence, eight components
  PASS  EMT CONSERVATION [THE RESULT]: nabla^mu That_{mu nu} == 0 ON SHELL in all eight components, for generic Psi(x0, x4), Psibar(x0, x4) and generic V
  PASS  EMT CONSERVATION [control]: OFF shell the divergence is not zero -- witnessed on Psi = (x4 + x0^2) e1 + e13, Psibar = e1 + x4 e5 with V = s^2, which is not a solution
  PASS  EMT [THE RESULT]: That - T_par == -(1/(4H)) Psibar ({gamma_mu, Gamma_nu} + {gamma_nu, Gamma_mu}) Psi: equal on the diagonal (every {gamma_mu, Gamma_mu} vanishes), different off it (the (x1, x4) witness above)
CELL 193  t=4.431 s
  PASS  EMT ON SHELL [THE RESULT]: the kinetic bilinear == s V'(s) (not zero) and Lhat == s V'(s) - V(s)
  PASS  EMT ON SHELL [THE RESULT]: rho = That_44 == V(s) + K_h
  PASS  EMT ON SHELL [THE RESULT]: T^i_i == T^h_h == s V'(s) - V(s) for the three observed and the three hidden directions, and T^0_0 == s V'(s) - V(s) + K_h
  PASS  EMT [has content]: the trace T^mu_mu == 7 Lkin - 8 V off shell, == 7 s V'(s) - 8 V(s) on shell
  PASS  EMT [fidelity]: K_h in the real commuting limit IS Part VI's cfKh, and K_h == 0 on the rescaled solutions Sqrt[Sin] Psi'(x4), Sqrt[Sin] Psi'bar(x4), where rho == V(s)
  [3.561 s]  the J-parity of all 64 components of That, admissible fields
  PASS  EMT HERMITICITY [THE RESULT]: J-parity table -- (mu, nu) in {0..4}: EVEN (Hermitian); (mu in 0..4, h in 5..7): ODD (anti-Hermitian); (h, h'), h != h': ZERO identically; (h, h): even
  PASS  EMT HERMITICITY [has content]: for x5..x7-independent fields That_{hh} == g_hh Lhat exactly, because {gamma_h, Gamma_h} == 0
CELL 194  t=4.080 s
  PASS  MODE SUMS [THE RESULT]: on a flat-frame plane wave (canonical normalization), T_44 + Lhat == omega ubar (-i T16[4]) u and T^j_j - Lhat == k_j ubar (-i T16[j]) u (j = 0..3): energy and pressure are the frequency and k_j dh/dk_j bilinears
  PASS  MODE SUMS [has content]: with ubar = u^ddag sigma16 these are omega u^ddag G u and k_j u^ddag (dh/dk_j) u, so by Section 28 each occupied mode contributes omega, k_j^2/omega and (to sigma) m/omega
  PASS  VACUUM [THE RESULT]: the unrenormalized sea energy -eps_KS(Lambda) == -(g/(16 pi^2))(2 Lambda^4 + 2 m^2 Lambda^2 + m^4/4 - m^4 Log[2 Lambda/m]) + o(1) as Lambda -> infinity; its m-dependent finite part contains -(g/(32 pi^2)) m^4 Log[m^2]
CELL 195  t=2.098 s
  PASS  KOHN-SHAM [definition]: the degeneracy g = 8 -- per momentum the positive-frequency eigenspace has rank 8 (Section 28's P_{+,+} + P_{+,-}), four flavours times two spins
  PASS  KOHN-SHAM [THE RESULT]: each closed form has k_F-derivative (g/(2 pi^2)) k_F^2 f(k_F) with f = 1, m/omega, omega, k^2/(3 omega), and vanishes at k_F = 0 -- so it IS the integral g Int d^3k/(2 pi)^3 f
  PASS  KOHN-SHAM [has content]: Integrate returns the same functions, with ArcTanh[k_F/w_F] in place of L -- and ArcTanh[k_F/w_F] == L (equal derivatives, both 0 at k_F = 0)
  PASS  KOHN-SHAM [control]: numerical quadrature agrees at k_F = 3/2, m = 1, g = 8 to 1e-12 (a cross-check; the proof is the symbolic one above)
  PASS  KOHN-SHAM [THE RESULT]: eps + P == w_F n (Gibbs), eps - 3 P == m sigma (trace), d eps/dm == sigma, d eps/dk_F == w_F dn/dk_F
  PASS  MEAN FIELD [THE RESULT]: rho = eps + W - sigma W', P_obs = P_KS + sigma W' - W give rho + P_obs == w_F n exactly; and d rho/dn == w_F + (sigma_KS - sigma) W'' sigma'(n), i.e. == w_F at the self-consistent point sigma = sigma_KS (m = W'(sigma) throughout)
  PASS  AUTHOR'S MASS TERM [THE RESULT]: W = m sigma gives W - sigma W' == 0, so rho == eps_KS, P == P_KS, and w = P_KS/eps_KS depends only on x = k_F/m: w == I2/(3 I1), I_n the m = 1 integrals
  PASS  AUTHOR'S MASS TERM [THE RESULT]: 0 <= w <= 1/3 -- P_KS >= 0 (I2' = x^4/w_x > 0, I2(0) = 0) and eps - 3P = m sigma = (g m^4/(2 pi^2)) I3 >= 0 (I3' = x^2/w_x > 0, I3(0) = 0)
  PASS  AUTHOR'S MASS TERM [THE RESULT]: w rises MONOTONICALLY -- dw/dx == x^2 Delta/(3 w_x I1^2) with Delta = x^2 I1 - w_x^2 I2, Delta(0) = 0 and Delta' = 2 x I3 > 0 -- from w -> 0 (x -> 0, w = x^2/5 + O(x^4): dust) to w -> 1/3 (x -> infinity: radiation)
  PASS  AUTHOR'S MASS TERM [has content]: the classical sign problem disappears -- sigma_KS = g Int m/omega is ODD in m and m sigma_KS = (g m^4/(2 pi^2)) I3(k_F/|m|) >= 0 for either sign, so with m = -2M: M sigma_KS <= 0, the branch M s < 0 that Part VI had to assume
  PASS  CLASSICAL LIMIT [THE RESULT]: at the self-consistent point rho == W + 3 P_KS and P_obs == (sigma W' - W) + P_KS exactly (sigma = sigma_KS, m = W'); Part VI's rho = V(s), P = s V' - V is the limit P_KS/eps_KS -> 0, i.e. k_F/m -> 0
CELL 196  t=4.945 s
  PASS  PRE-UNIVERSE (1) [THE RESULT]: the U(1) current is conserved on shell, (1/Sqrt[g]) d_mu(Sqrt[g] j^mu) == 0, for generic Psi(x0, x4), Psibar(x0, x4)
  PASS  PRE-UNIVERSE (1) [THE RESULT]: for Psi = Sqrt[Sin] Psi'(x4), Sqrt[g] j^0 is x0-independent (Sec Sin Cot == 1), so d_4(Sqrt[g] j^4) == 0 on shell; and Sqrt[g] and q^3 p^3 are x4-independent
  PASS  PRE-UNIVERSE (2) [THE RESULT]: Psi = Sqrt[Sin] Psi'(x4) satisfies the field equation iff d_x0 V'(Sin[6Hx0] s') == 0, and d_x0 V'(Sin s') == 6 H Cos[6 H x0] s' V''(Sin s'): V linear OR s' == 0
  PASS  PRE-UNIVERSE (2) [has content]: s' = Psi'bar Psi' is conserved by the x4-only equations (d_4 Psi' = -H V' T16[4] Psi', d_4 Psi'bar = H V' Psi'bar T16[4]) for ANY V; so null data s' = 0 are solutions for any V, with rho == V(0); and for V = s^2 with s' != 0 the x0-dependence is real (witness)
  PASS  PRE-UNIVERSE (3) [THE RESULT]: for Psi = Sqrt[Sin] u(x4) Exp[i k x1] the field equation (mass term, V' = constant) is Sqrt[Sin] Exp[i k x1] (T16[4] u' + i (k/q) T16[1] u - H V' u): the physical momentum is k/q
  PASS  PRE-UNIVERSE (3) [THE RESULT]: d Log[k/q]/dx4 == H a4'[H x4] == -H_obs (observed momenta redshift as Exp[a4]), and E^2 == (k^2/q^2 + m^2) ID16 at the physical momentum
  PASS  PRE-UNIVERSE (4) [THE RESULT]: x = k_F,phys/m = x_c/(q m) obeys d Log[x]/dx4 == H a4'[H x4]; since dw/dx > 0, dw/dx4 = w'(x) x H a4' < 0 exactly when a4' < 0 (the observed sheet expanding): the quantum gas COOLS
  PASS  PRE-UNIVERSE (4) [control]: the classical Part VI fable does NOT cool -- ds/dx4 == 0 on its rescaled solutions (Part VI's result, re-asserted), and its w = s V'/V - 1 is a function of s alone: frozen
  PASS  PART VII [control]: a4 is STILL UNDEFINED -- nothing in Part VII gave it a value
CELL 197  t=0.072 s
identities accepted on numerical evidence alone : 0   (stage 3 returned True)
non-vanishing witnesses                         : 38   (cfNonZeroWitnessQ found a probe point at which every entry is a
                                                     number and one of them is non-zero: a complete proof of non-vanishing)
  PASS  no identity in this notebook was accepted on numerical evidence alone
Assertions run: 528   passed: 528   failed: 0
CELL 198  t=0.000 s
\end{Verbatim}

\subsection{The same assertions in the runs of Parts I--VIII}
\label{sec:10.3}

After Part~VIII (the coupled system, notebook Sections 30--33) was added, the whole notebook was run
twice more. \textbf{Part~VII's record is its own run} (\secref{sec:10.1}, \secref{sec:10.2}); the two
later runs show that adding Part~VIII changed nothing in it. With the \cfcode{CELL} and
\cfcode{[t s]} timing lines removed, the lines printed by cells 173--198 are identical in all three
logs --- the Part~VII run, the acceptance run of Part~VIII and the final run --- 174 lines each, 170 of
them \cfcode{PASS} lines [checked for this page with \cfcode{diff}].

\paragraph{The acceptance run of Part~VIII} (\file{claude-fable/run_fermion_fable_part8.log}, committed
in \texttt{bc04ed1}; a historical record). The Part~VII cells took 29.706\,s in that run. Its summary:

\begin{Verbatim}[fontsize=\scriptsize]
==================== RUN SUMMARY ====================
cells evaluated : 217
total seconds   : 233.289
cells w/ msgs   : 0
---- slowest cells ----
  cell 80  80.343 s
  cell 88  23.494 s
  cell 204  10.364 s
  cell 86  9.821 s
  cell 137  9.306 s
  cell 191  9.225 s
  cell 87  8.387 s
  cell 121  6.069 s
  cell 160  4.250 s
  cell 196  4.185 s
  cell 107  3.836 s
  cell 193  3.140 s
==================== ASSERTIONS ====================
assertions run  : 642
passed          : 642
FAILED          : 0
====================================================
RUN-DONE
\end{Verbatim}

The last Part~VII line of that log is still \cfcode{Assertions run: 528   passed: 528   failed: 0}
(line 843), and the run ends with \cfcode{Assertions run: 642   passed: 642   failed: 0} and 46
non-vanishing witnesses. Two of Part~VIII's assertions did not run in it: its comparisons of the Rust
solver with the Mathematica reference runs are made only if the solver's CSVs exist, and they did not
exist yet (two \cfcode{NOTE} lines say \cfcode{Rust CSV not found; that comparison was skipped, not failed}).
And three of its labels, in Section~32, write the listed spin-connection components with both flat
indices up --- the notation that \texttt{3a5020d} corrected:
\begin{itemize}
  \item \cfL{SPIN CONNECTION [THE RESULT]: the complete list of non-zero omega\_mu\textasciicircum{}\{ab\} of the WARPED frame is the thirteen stated components (and their antisymmetric partners) -- nothing else}
  \item \cfL{SPIN CONNECTION [THE RESULT]: the complete list for the Bianchi-I frame is omega\_0\textasciicircum{}\{04\} = C', omega\_i\textasciicircum{}\{i4\} = A', omega\_h\textasciicircum{}\{4h\} = B' -- nothing else}
  \item \cfL{SPIN CONNECTION [fidelity]: on the canonical member (labelled substitution) the warped connection IS omegaCanonical of Section 16, and the new component omega\_0\textasciicircum{}\{04\} vanishes there}
\end{itemize}
The arrays these assertions compare (\cfcode{cfOmegaWarp}, \cfcode{cfOmegaBI}) carry both flat indices
\textbf{down}, so the tests themselves were right and are unchanged. Raising both flat indices
multiplies a component by $\eta_{aa}\eta_{bb}$, which is $-1$ exactly in the planes with one timelike
index and one spacelike index ($(0,4)$, $(i,4)$, $(0,h)$) and $+1$ in the $(0,i)$ and $(4,h)$ planes
\cfderived. So, read with the indices up as written, the second label states the wrong sign for
$\omega_0{}^{04}$ and $\omega_i{}^{i4}$, while $\omega_h{}^{4h}=B'$ holds in either position; and the
list of values in the old Section~32 Text cell had the wrong sign for its $(0,4)$, $(i,4)$ and $(0,h)$
components. (The corrected Text cell of \texttt{3a5020d} says that raising both indices ``would flip
the sign of every plane that contains x4 or a hidden timelike direction''; exactly, it flips the
planes with one timelike and one spacelike index, and leaves the $(4,h)$ planes, whose two indices are
both timelike, unchanged.)

\paragraph{The final run} (\file{claude-fable/run_fermion_fable_final.log}, first recorded in
\texttt{1671155} and re-recorded after the label correction in \texttt{3a5020d}; the notebook in its
committed state). The Part~VII cells took 26.118\,s and the Part~VIII cells 20.994\,s in that run. Its
summary:

\begin{Verbatim}[fontsize=\scriptsize]
==================== RUN SUMMARY ====================
cells evaluated : 217
total seconds   : 198.352
cells w/ msgs   : 0
---- slowest cells ----
  cell 80  68.159 s
  cell 88  22.525 s
  cell 204  9.360 s
  cell 86  8.072 s
  cell 191  7.887 s
  cell 87  7.193 s
  cell 137  6.985 s
  cell 121  4.579 s
  cell 160  3.822 s
  cell 107  3.505 s
  cell 196  3.334 s
  cell 193  2.970 s
==================== ASSERTIONS ====================
assertions run  : 644
passed          : 644
FAILED          : 0
====================================================
\end{Verbatim}

The last Part~VII line of that log is again \cfcode{Assertions run: 528   passed: 528   failed: 0}
(line 843), and the run ends with \cfcode{Assertions run: 644   passed: 644   failed: 0} (line 1005),
\cfcode{identities accepted on numerical evidence alone : 0} and 46 non-vanishing witnesses; no
\cfcode{RUN-DONE} line was appended to this log. Its \cfcode{PASS} lines differ from those of the
acceptance run in exactly five places [checked for this page with \cfcode{diff}]. The two comparisons
with the Rust solver now run and pass, each after a line
\cfcode{Rust <path>: 71 rows compared; max difference per column ...}:
\begin{itemize}
  \item \cfL{fable4d RUNS [fidelity]: the Rust solver's mass30eV run agrees with this cell to 1e-6 in every column, at every one of the rows its grid covers}
  \item \cfL{fable4d RUNS [fidelity]: the Rust solver's power run agrees with this cell to 1e-6 in every column, at every one of the rows its grid covers}
\end{itemize}
and the three Section~32 labels write the components with both flat indices down, $\omega_{\mu ab}$,
as the arrays they compare have them (the tests are unchanged):
\begin{itemize}
  \item \cfL{SPIN CONNECTION [THE RESULT]: the complete list of non-zero omega\_\{mu ab\} (both flat indices down) of the WARPED frame is the thirteen stated components (and their antisymmetric partners) -- nothing else}
  \item \cfL{SPIN CONNECTION [THE RESULT]: the complete list for the Bianchi-I frame is omega\_\{0 04\} = C', omega\_\{i i4\} = A', omega\_\{h 4h\} = B' (both flat indices down) -- nothing else}
  \item \cfL{SPIN CONNECTION [fidelity]: on the canonical member (labelled substitution) the warped connection IS omegaCanonical of Section 16, and the new component omega\_\{0 04\} vanishes there}
\end{itemize}

\subsection{The design review, and every correction it made to Effort A}
\label{sec:10.4}

The design (revision~1) was reviewed by four adversarial reviewers --- quantization (findings
\texttt{QK-*}); field equations, energy--momentum tensor and spin connection (\texttt{EMT-*}); Einstein
equations and cosmology (\texttt{E-*}); density-functional theory and numerics (\texttt{DFT-*}) ---
each followed by an independent skeptic who tried to refute each finding. \textbf{No finding was
refuted.} Several were confirmed only in part, and then the corrected fix was adopted. The review
records are working documents and are not part of the repository; every adopted correction that
concerns Effort A is restated here, with where it now lives.

{\footnotesize
\setlength{\tabcolsep}{3pt}
\begin{longtable}{@{}p{0.12\linewidth}p{0.11\linewidth}p{0.2\linewidth}p{0.34\linewidth}p{0.16\linewidth}@{}}
\toprule
\raggedright design item & \raggedright findings & \raggedright what revision~1 said, or lacked & \raggedright what was adopted & \raggedright where on this page (notebook) \tabularnewline
\midrule
\endhead
\raggedright refinement & \raggedright QK-7 & \raggedright ``fermion fable must be complex'' without scope; F2(a) only on the flat frame & \raggedright the invariant forms are exactly $\operatorname{span}\{\sigma_{16},\sigma_{16}\Tsix{8}\}$; Majorana has one quartic scalar and no commuting kinetic shadow; Weyl/Majorana--Weyl cannot propagate; the complex spinor is the \textbf{minimal}, not the unique, field; F2(a) holds on \textbf{every} frame & \raggedright \ref{sec:3.2}, \ref{sec:3.4}, \ref{sec:3.8} (\cfcode{REFINEMENT (a)}, \cfcode{(b)}) \tabularnewline
\raggedright Lagrangian & \raggedright QK-4 (corrected), EMT-5 & \raggedright several Lagrangians used interchangeably; Part~VI's drop-out mechanism assumed to carry over & \raggedright only the symmetrized covariant Lagrangian; the unsymmetrized partial one is inadmissible (its two equations are not conjugate); the drop-out is a property of the matrix $\{\gamma^\mu,\Gamma_\mu\}$, per $\mu$, on diagonal frames only & \raggedright \ref{sec:4} (\cfcode{LAGRANGIAN (a)}--\cfcode{(f)}) \tabularnewline
\raggedright field equations & \raggedright EMT-9, EMT-11 & \raggedright 8+8 form sketched; operator ordering ``normal-ordered or Weyl''; $x_4$-only solutions assumed & \raggedright the 8+8 form with the conjugate rows; the axial phase is not a symmetry; the ordering rule with its scalar and $\gamma^4$ contraction parts; the $x_4$-only family needs $V''s'=0$, and $s'$ is conserved & \raggedright \ref{sec:5.3}, \ref{sec:5.8}, \ref{sec:5.9}, \ref{sec:5.11} \tabularnewline
\raggedright quantization: truncation & \raggedright QK-2 & \raggedright ``physical states obey $P_5=P_6=P_7=0$'' (a superselection rule) and Tomonaga--Schwinger integrability & \raggedright a \textbf{truncation} with finite $\Vhid$ (closed timelike curves stated); anticommutator with $H/(\Vhid\sqrt{|g|})$ and $\delta^4$; only $[P_4,P_h]=0$ kept & \raggedright \ref{sec:6.4}, \ref{sec:6.10} \tabularnewline
\raggedright quantization: theorem & \raggedright QK-1 (corrected) & \raggedright admissibility asserted from one example & \raggedright the theorem: positive quantization iff the momentum span is spacelike-definite, with items (1)--(7) (threshold Jordan block, $G$-neutral modes above threshold, traceless certificate, uniqueness) & \raggedright \ref{sec:6.10} (\cfcode{ADMISSIBILITY}) \tabularnewline
\raggedright quantization: $J$ & \raggedright QK-5 (corrected), plus the missed basis point & \raggedright Krein prescription $b^{\dagger}:=\eta\,b^{\ddagger}$, basis unspecified; reflection positivity claimed & \raggedright $J=-i\,\Omega_4\Tsix{0}$, $\sigma_{16}=J\beta$, flavour metric $(2,2)$; the adjoint $\Psi^{\dagger}:=\Psi^{\ddagger}J$ is mode-independent and \textbf{requires a $J$-eigenmode basis} (control: a Krein boost breaks Hermiticity of $s$); why it fails for bosons; reflection positivity \textbf{dropped} & \raggedright \ref{sec:6.5}, \ref{sec:6.6}, \ref{sec:6.13} (\cfcode{J}, \cfcode{KREIN MODES [control]}) \tabularnewline
\raggedright quantization: other & \raggedright QK-6, QK-8, QK-9, QK-10, QK-11, EMT-2 & \raggedright symmetry count, dynamical frames, wall condition, current sign and $J$-connection relation incomplete or wrong & \raggedright 13 kept / 15 lost generators, $G$ invariant under Spin(4,3); $\chi=(\sqrt g/H)^{1/2}\Psi$; one-particle space $L^2(dz)$, regular at the wall, limit-point at infinity, bag condition $\Tsix{0}\Psi'=\pm\Psi'$; $n^\mu=-\frac iH\bar\Psi\gamma^\mu\Psi$; the $J$ table with $\Gamma_{5..7}$ anticommuting & \raggedright \ref{sec:6.9}, \ref{sec:6.12}, \ref{sec:6.15}, \ref{sec:6.16}, \ref{sec:6.8} \tabularnewline
\raggedright EMT & \raggedright EMT-1, EMT-3, EMT-4, EMT-8, QK-3 & \raggedright EMT stated without derivation; conservation ``for independent $\Psi$, $\bar\Psi$'' without scope; ``$\langle{:}T{:}\rangle=0$ in the vacuum'' & \raggedright Hilbert $=\Tcov$ off shell; spin-connection variation exactly zero; Hilbert $\neq\Tpar$ (witness at a non-vacuous pair); why the commuting proxy proves conservation; the $J$-parity table; the vacuum statement only for static backgrounds; do not run the symbolic divergence of $\Tpar$ & \raggedright \ref{sec:7} (\cfcode{EMT}, \cfcode{EMT CONSERVATION}, \cfcode{EMT HERMITICITY}), \ref{sec:8.3} \tabularnewline
\raggedright expectation values & \raggedright DFT-5, DFT-6, DFT-12, DFT-15, DFT-4 & \raggedright classical limit ``$\kF\to0$ at fixed $n$'' (impossible) or a coherent zero mode; unstable closed forms; phantom crossing not examined; sea energy only for constant $m$ & \raggedright the classical limit is the \textbf{non-relativistic} limit, with the sign rule; the numerically stable forms; $w_f\geq-1$ (no phantom crossing); the pair-excitation thresholds ($\wF+|m|$ at $q=\kF$, $2\wF$ at $q=0$; not asserted in Part~VII); $\Delta E_{\mathrm{vac}}$ for mass-varying $W$ & \raggedright \ref{sec:8.5}, \ref{sec:8.7}, \ref{sec:8.9}, \ref{sec:8.3} \tabularnewline
\raggedright spin connection & \raggedright EMT-6, EMT-10 & \raggedright ``the connection enters the Hamiltonian'' & \raggedright it enters \textbf{neither} the anticommutator \textbf{nor} the Hamiltonian density; it re-emerges in the equation of motion as the term that makes $h$ Hermitian; $\gamma^\mu\Gamma_\mu=-\tfrac12T_\mu\gamma^\mu$ (the Weitzenb\"ock torsion vector) & \raggedright \ref{sec:6.11}, \ref{sec:9.12.3}, \ref{sec:9.12.4} \tabularnewline
\bottomrule
\end{longtable}
}

Two of these adopted corrections are also stated in Part~VII's own Text cells: the
Hamiltonian-density correction (EMT-6; ``This corrects the design's `it enters the Hamiltonian'\,'',
Section~28), and the identification of the classical limit with the non-relativistic one (DFT-5;
Section~29). One numerical slip in the review's vacuum-energy figures is corrected on this page
(\secref{sec:8.3}).

\subsection{The Fock check}
\label{sec:10.5}

The finite Fock space of \secref{sec:6.14} is the independent check of the quantization: it realizes
the Dirac-bracket anticommutator on a positive Hilbert space of dimension $2^{16}=65536$, with exact
integer sparse matrices, and confirms the sign of the anticommutator through the Heisenberg equation
(the opposite sign gives the time-reversed equation). It shows the unique Dirac-sea ground state at
$-40$, the 16-fold first excited level (8 particles, 8 antiparticles), the negative Krein ``norms''
that the $J$-involution removes, and the expectation rule used for $\rho$, $P$ and $\sigma$. Twelve
assertions (\cfcode{FOCK}, \cfcode{CHARGE}), Input cell 191, 9.858\,s.

\subsection{The TeX export}
\label{sec:10.6}

\file{claude-fable/export_fermion_fable_tex.wls} evaluates the notebook's Input cells (output
suppressed), checks that the Part~VII objects exist, and writes fourteen files into
\file{provenance-latex/generated/}:

{\small
\begin{longtable}{@{}p{0.3\linewidth}p{0.36\linewidth}p{0.28\linewidth}@{}}
\toprule
\raggedright file (\file{.tex} and \file{.txt}) & \raggedright content & \raggedright taken from the notebook object \tabularnewline
\midrule
\endhead
\raggedright \file{fermion_fable_field_equations} & \raggedright the 16 component field equations on the canonical frame & \raggedright \cfcode{cfFieldEqTerms} \tabularnewline
\raggedright \file{fermion_fable_adjoint_equations} & \raggedright the 16 component adjoint equations & \raggedright \cfcode{cfAdjEqTerms} \tabularnewline
\raggedright \file{fermion_fable_field_equations_8plus8} & \raggedright the split-octonion 8+8 form, compact and by component & \raggedright \cfcode{cfFieldOp16}, \cfcode{cfUpper8}, \cfcode{cfLower8}, \cfcode{cfFieldEqTerms} \tabularnewline
\raggedright \file{fermion_fable_field_equations_rescaled} & \raggedright the 16 equations for $\Psi'=\Psi/\sqrt{\sin 6Hx_0}$ & \raggedright the rescaled operator of Section~27 (iv), asserted equal to \cfcode{cfE8} under the rescaling \tabularnewline
\raggedright \file{fermion_fable_anticommutator} & \raggedright $\{\Psi_a,\Psi^{\ddagger}_b\}=H\cos(6Hx_0)\,J_{ab}\,\delta^7$, with the matrix & \raggedright \cfcode{cfAntiPsiPsiDd} at $c=\sqrt g/H$ \tabularnewline
\raggedright \file{fermion_fable_krein_J} & \raggedright $J$, $\beta$, $\sigma_{16}=J\beta$ & \raggedright \cfcode{cfJK}, \cfcode{cfBeta} \tabularnewline
\raggedright \file{fermion_fable_eos} & \raggedright the Kohn--Sham closed forms $n$, $\sigma$, $\varepsilon$, $P$, $w$ at $g=8$ & \raggedright \cfcode{cfNKS}, \cfcode{cfSigKS}, \cfcode{cfEpsKS}, \cfcode{cfPKS} \tabularnewline
\bottomrule
\end{longtable}
}

Every expression comes from an object that Part~VII asserts; the script does no mathematics of its
own beyond formatting (its header says so, and its code only renders term lists and matrices). The
\file{.tex} files are \verb|\input| by this document in sections \ref{sec:5}, \ref{sec:6} and
\ref{sec:8}; the \file{.txt} files hold the same content in plain text. The export was run on
2026-09-24 at 20:41 (PDT) on the 198-cell notebook, and it printed:

\begin{Verbatim}[fontsize=\scriptsize]
evaluating 198 Input cells of C:\Users\nsh\Documents\8-dim\Pre-Universe_14SEP26-77\claude-fable\claude-fable_Einstein-Rosen-2-Planes.nb (output suppressed)
assertions: 528   failed: 0   cells with messages: 0
  wrote C:\Users\nsh\Documents\8-dim\Pre-Universe_14SEP26-77\provenance-latex\generated\fermion_fable_field_equations.tex / .txt
  wrote C:\Users\nsh\Documents\8-dim\Pre-Universe_14SEP26-77\provenance-latex\generated\fermion_fable_adjoint_equations.tex / .txt
  wrote C:\Users\nsh\Documents\8-dim\Pre-Universe_14SEP26-77\provenance-latex\generated\fermion_fable_field_equations_8plus8.tex / .txt
  wrote C:\Users\nsh\Documents\8-dim\Pre-Universe_14SEP26-77\provenance-latex\generated\fermion_fable_field_equations_rescaled.tex / .txt
  wrote C:\Users\nsh\Documents\8-dim\Pre-Universe_14SEP26-77\provenance-latex\generated\fermion_fable_anticommutator.tex / .txt
  wrote C:\Users\nsh\Documents\8-dim\Pre-Universe_14SEP26-77\provenance-latex\generated\fermion_fable_krein_J.tex / .txt
  wrote C:\Users\nsh\Documents\8-dim\Pre-Universe_14SEP26-77\provenance-latex\generated\fermion_fable_eos.tex / .txt
done: 14 files in C:\Users\nsh\Documents\8-dim\Pre-Universe_14SEP26-77\provenance-latex\generated
\end{Verbatim}

(That console output was captured outside the repository; the fourteen files it wrote were committed
in \texttt{4fdfc40}.) A test document that \verb|\input|s all seven \file{.tex} fragments (with
\cfcode{amsmath}) compiled with pdfLaTeX (MiKTeX) to 9 pages with no errors, warnings or overfull
boxes (a scratch test on 2026-09-24, outside the repository).

% ===============================================================================================
\section{Every command, in full}
\label{sec:11}

The commands run in Git Bash on Windows 11 (they are plain POSIX shell and run the same way on Linux
and macOS). The toolchain recorded in the logs: Wolfram Language 15.0.1 for Microsoft Windows
(64-bit) (July 2, 2026), called as \cfcode{wolframscript} (the fourth line of each of the three run logs); Python 3 for
the builder (Python 3.14.5 when \secref{sec:11.1}'s output was recorded); and MiKTeX (pdfTeX
3.141592653-2.6-1.40.29, MiKTeX 26.5; Latexmk 4.88) for the LaTeX twin, installed on this machine at
\file{C:/Program Files/MiKTeX/miktex/bin/x64}.

\subsection{Build the notebook from its manifests}
\label{sec:11.1}

The notebook is never edited by hand. \file{build_tools.py} reads every
\file{claude-fable/cells_part*.wl} in order and writes both the notebook and \file{run_all.wls} (the
same Input cells as a script).

\begin{Verbatim}
cd "$(git rev-parse --show-toplevel)/claude-fable"
python build_tools.py          # use python3 where python is absent
\end{Verbatim}

It prints (recorded on 2026-09-24 for this page, with the eight committed manifests, in a scratch copy
of the directory so that nothing in the repository was touched; it ran in 0.115\,s):

\begin{Verbatim}[fontsize=\scriptsize]
manifest files : ['cells_part1.wl', 'cells_part2.wl', 'cells_part3.wl', 'cells_part4.wl', 'cells_part5.wl', 'cells_part6.wl', 'cells_part7.wl', 'cells_part8.wl']
cells parsed   : 384 {'Title': 8, 'Subtitle': 1, 'Subsubtitle': 1, 'Text': 123, 'Section': 34, 'Input': 217}
run_all.wls    : 217 Input cells
notebook       : 384 cells -> <repo>\claude-fable\claude-fable_Einstein-Rosen-2-Planes.nb
\end{Verbatim}

(the last line prints the absolute path of the directory it ran in). The rebuilt
\file{claude-fable_Einstein-Rosen-2-Planes.nb} and \file{run_all.wls} were \textbf{byte-identical} to
the committed ones (\cfcode{cmp} reported no difference). The check was repeated on 2026-09-25 on the
manifests of \texttt{3a5020d}, whose \file{cells_part8.wl} carries the corrected Section~32 labels: the
same four lines, and the rebuilt notebook and \file{run_all.wls} again byte-identical to the committed
ones. With only the seven manifests of Parts
I--VII present --- the state in which Part~VII was run --- it prints
\begin{Verbatim}[fontsize=\scriptsize]
cells parsed   : 340 {'Title': 7, 'Subtitle': 1, 'Subsubtitle': 1, 'Text': 103, 'Section': 30, 'Input': 198}
\end{Verbatim}
and \cfcode{run_all.wls    : 198 Input cells}, which is the style tally that the checker reported for
the Part~VII notebook (\secref{sec:11.3}).

\subsection{Run the notebook end to end, straight out of the \texttt{.nb}}
\label{sec:11.2}

\begin{Verbatim}
cd "$(git rev-parse --show-toplevel)/claude-fable"
wolframscript -file run_from_nb.wls > run_fermion_fable_final.log 2>&1
git checkout -- claude-fable_Einstein-Rosen-2-Planes-eLa.mx claude-fable_Einstein-Rosen-2-Planes-eLazt.mx
\end{Verbatim}

\begin{itemize}
  \item \textbf{What it does.} \file{run_from_nb.wls} imports the \file{.nb}, takes its Input cells in
        order, and evaluates each one as a single unit (a multi-line cell becomes one
        \cfcode{CompoundExpression}, exactly as the front end evaluates it), through the harness
        \file{runner_header.wl}. The harness prints \cfcode{CELL n  t=... s} after each cell, with
        \cfcode{MESSAGES: ...} appended if the cell raised any message, and ends with the run summary
        and the assertion counts.
  \item \textbf{Why from inside \file{claude-fable/}.} Run headless, the notebook cannot ask for its own
        file name, so Section~10 looks for the author's helper packages
        \file{ConvertMapleToMathematicaV2.wl} and \file{EtoExp.wl} in the working directory (and in
        \file{Pre-Universe_14SEP26-77/} below it, and one level up). They sit in \file{claude-fable/}.
        The reference CSVs of Parts VI and VIII are found as
        \verb|claude-fable\..\fable-cosmology\reference\...| (the log prints these paths). The script
        itself is self-locating (\cfcode{$InputFileName}) and finds the notebook from any directory.
  \item \textbf{The \file{.mx} restore.} Section~12 of the notebook writes
        \file{claude-fable_Einstein-Rosen-2-Planes-eLa.mx} and \file{...-eLazt.mx} with
        \cfcode{DumpSave}, whose bytes are not reproducible from run to run. The two files are restored
        to the committed version after every run (commits \texttt{4fdfc40} and \texttt{bc04ed1}
        record that this was done).
  \item \textbf{How long.} 198.352\,s for the 217 cells of the final run; 232.146\,s for the 198 cells
        of the Part~VII run and 233.289\,s for the 217 cells of the acceptance run of Part~VIII (the
        \cfcode{total seconds} lines of the three logs).
  \item \textbf{What it prints.} The first line is \cfcode{evaluating 217 Input cells straight out of the .nb}
        (\cfcode{evaluating 198 ...} for the Part~VII state). Then the provenance banner of the original
        notebook and every cell's line, \cfcode{PASS} labels and notes. The Part~VII portion is quoted
        in full in \secref{sec:10.2}. The end of the run is quoted in sections \ref{sec:10.1} and
        \ref{sec:10.3}.
\end{itemize}

The final log was produced by this command on the committed notebook of \texttt{3a5020d}. The
Part~VII log was produced by the same command with the output redirected to
\file{run_fermion_fable_part7.log}, when the notebook had Parts I--VII, and the acceptance log of
Part~VIII with the output redirected to \file{run_fermion_fable_part8.log}, before the Rust solver's
CSVs existed. The last line of those two, \cfcode{RUN-DONE}, was appended by the wrapper that ran the
command; the final log has no such line.

\paragraph{Checking a log} (the numbers are the results on the three committed logs):

\begin{Verbatim}[fontsize=\small]
cd "$(git rev-parse --show-toplevel)/claude-fable"
grep -c '^  PASS' run_fermion_fable_part7.log         # 528
grep -c '^  PASS' run_fermion_fable_part8.log         # 642
grep -c '^  PASS' run_fermion_fable_final.log         # 644
grep -n '^  FAIL' run_fermion_fable_final.log         # no output: no assertion failed
grep -c 'MESSAGES: ' run_fermion_fable_final.log      # 0: no cell raised a message
grep -E 'identities accepted|non-vanishing witnesses' run_fermion_fable_final.log | tail -2
#   identities accepted on numerical evidence alone : 0   (stage 3 returned True)
#   non-vanishing witnesses                         : 46   (cfNonZeroWitnessQ found a probe point at which every entry is a
sed -n '/RUN SUMMARY/,$p' run_fermion_fable_final.log  # the summary quoted in section 10.3
diff <(grep '^  PASS' run_fermion_fable_part8.log) <(grep '^  PASS' run_fermion_fable_final.log)
#   five differences: the three Section 32 labels, and the two Rust comparisons added (section 10.3)
\end{Verbatim}

A plain \cfcode{grep FAIL} also finds the summary line \cfcode{FAILED          : 0} and the two
\cfcode{PASS} labels whose text contains the word ``FAILS'' (control assertions that show a wrong
sign failing, lines 298 and 921 of the final log, and of the Part~VIII log). The run prints four
\cfcode{EXPECTED MESSAGE(S)} lines (three in Part~II, one in Part~V: messages that the original
notebook's time budgets or its underdetermined \cfcode{Solve} calls may raise, announced in advance)
and one \cfcode{NOTE} line in Part~III (that $a_4$ is undefined); in Part~VIII it prints the two lines
\cfcode{Rust <path>: 71 rows compared; ...} of the comparisons with the Rust solver's CSVs
\file{fable-cosmology/results/nb06_fable4d_mass30eV.csv} and \file{nb06_fable4d_power.csv} (committed in
\texttt{aca5102}). In the acceptance run of Part~VIII those CSVs were not yet present, and two
\cfcode{NOTE} lines said so instead (\cfcode{Rust CSV not found; that comparison was skipped, not failed}).
None of these lines is a message raised by a cell.

\subsection{Check the generated notebook without evaluating it}
\label{sec:11.3}

\begin{Verbatim}
wolframscript -file "$(git rev-parse --show-toplevel)/claude-fable/verify_nb.wls"
\end{Verbatim}

It imports the \file{.nb}, counts its cells by style, extracts the Input cells, and checks that every
one of them parses. It takes about 4\,s (4.118\,s when timed for the draft of this page). For the
committed notebook it prints (\file{claude-fable/verify_nb_part8.log}, re-recorded in \texttt{3a5020d}
after the Section~32 labels were corrected; before that correction it printed
\cfcode{file bytes: 545408} and \cfcode{total input characters: 347187}, the rest unchanged):

\begin{Verbatim}[fontsize=\scriptsize]
file bytes: 545713
Head: Notebook
cells: 384
style tally: {{Title, 8}, {Subtitle, 1}, {Subsubtitle, 1}, {Text, 123}, {Section, 34}, {Input, 217}}
input cells with plain-string BoxData: 217
total input characters: 347237
input cells that FAIL to parse: {}
first cell: InputForm[(* --- provenance banner, as in the original notebook ------------------------------------]
last cell : InputForm[cfAssert["PART VIII [control]: a4 is STILL UNDEFINED -- nothing in Part VIII gave it a val]
\end{Verbatim}

For the Part~VII notebook it printed (\file{claude-fable/verify_nb_part7.log}):

\begin{Verbatim}[fontsize=\scriptsize]
file bytes: 440320
Head: Notebook
cells: 340
style tally: {{Title, 7}, {Subtitle, 1}, {Subsubtitle, 1}, {Text, 103}, {Section, 30}, {Input, 198}}
input cells with plain-string BoxData: 198
total input characters: 279955
input cells that FAIL to parse: {}
first cell: InputForm[(* --- provenance banner, as in the original notebook ------------------------------------]
last cell : InputForm[(* --- final tally of every assertion made in this notebook, Parts I-VII -----------------]
\end{Verbatim}

\subsection{Print the section-to-cell map}
\label{sec:11.4}

\begin{Verbatim}
wolframscript -file "$(git rev-parse --show-toplevel)/claude-fable/nb_section_map.wls" > "$(git rev-parse --show-toplevel)/claude-fable/nb_section_map.log"
\end{Verbatim}

It reads the \file{.nb} without evaluating anything and prints, in document order, each Part title
and each Section with the range of its Input cells (3.340\,s when timed for this page). Its output,
\file{claude-fable/nb_section_map.log}:

\begin{Verbatim}[fontsize=\scriptsize]
== claude-fable_Einstein-Rosen-2-Planes
0.  What this notebook does, and how it is organized  |  Input cells none
1.  Session setup  |  Input cells 1-6
2.  Coordinates and the flat 4+4 Minkowski metric eta  |  Input cells 7-12
3.  SO(4): self-dual and anti-self-dual antisymmetric 4x4 matrices  |  Input cells 13-15
4.  The real 8x8 Clifford generators tau, their conjugates, and the spinor metric sigma  |  Input cells 16-21
5.  The 16x16 Dirac matrices T16, the chirality operator, sigma16 and the so(4,4) generators  |  Input cells 22-29
6.  A complete basis of the 16x16 matrix algebra  |  Input cells 30-35
7.  A complete basis of the 8x8 matrix algebra, and the induced flat metric  |  Input cells 36-38
8.  The 4x4 Dirac matrices, and a complete basis of the 4x4 matrix algebra  |  Input cells 39-44
9.  Cartan triality: the constant bridge between the vector and spinor realizations,
    and the split-octonion structure constants  |  Input cells 45-52
== PART II  --  The wave function of the un-universe, its field equations, and the 3 generations
10.  Housekeeping: where this notebook lives, and the two helper packages  |  Input cells 53-55
11.  The wave function of the un-universe, and its Lagrangian  |  Input cells 56-60
12.  The Euler-Lagrange equations, the coupling pattern, and the (z,t) chart  |  Input cells 61-66
13.  Decoupling into four blocks of four, and the closed-form solutions  |  Input cells 67-78
14.  Bilinears, the two solution branches, and M6 = 3 generations of Einstein-Rosen 2-planes  |  Input cells 79-90
== PART III  --  The curved 4+4 spacetime, its frame field, and the canonical spin connection
15.  The local flat 4+4 Minkowski coordinate system at every spacetime point  |  Input cells 91-99
16.  The generalized Christoffel symbols, and the canonical spin connection  |  Input cells 100-110
17.  The gauge-covariant derivative of the 16-component spinor  |  Input cells 111-117
== PART IV  --  Three bridges between curved and flat indices: two gauge transformations, and one new connection
18.  Bridge 1 -- a locally boosted frame: the same connection in a local gauge  |  Input cells 118-125
19.  Bridge 2 -- the null (light-cone) frame: the same connection in a constant gauge, and the spinor metric sigma  |  Input cells 126-130
20.  Bridge 3 -- the triality frame with totally antisymmetric split-octonion torsion  |  Input cells 131-139
== PART V  --  The fable-5.1 bridge: the connection in which the frame itself is parallel
21.  The fable-5.1 bridge -- the Weitzenboeck (teleparallel) connection of the canonical frame  |  Input cells 140-148
22.  Master comparison of the five spin connections  |  Input cells 149-152
== PART VI  --  fableScalar and fable: two fields on the pre-universe, and the equation of state of dark energy
23.  fableScalar -- a scalar field on the pre-universe, its energy-momentum tensor, and wScalar  |  Input cells 153-160
24.  fable -- the 16-component spinor field, its energy-momentum tensor, and w  |  Input cells 161-169
25.  Reading the two mechanisms side by side  |  Input cells 170-172
== PART VII  --  fermion fable: the complex 16-spinor, its canonical quantization in 4+4 dimensions, its field equations in the primordial gravitational field, and its energy-momentum tensor operator
26.  The refinement: why fermion fable is a complex 16-spinor  |  Input cells 173-179
27.  The Lagrangian, the field equations in the primordial gravitational field, and the canonical spin connection  |  Input cells 180-184
28.  Canonical quantization in 4+4 dimensions: the Dirac bracket, the Krein structure J, and the admissible sector  |  Input cells 185-191
29.  The energy-momentum tensor operator of fermion fable: definition, conservation, and rho, P and w  |  Input cells 192-198
== PART VIII  --  fable as a source of the Einstein equations of the primordial gravitational field
30.  The Einstein tensor of the canonical metric, and what it demands of its source  |  Input cells 199-203
31.  The generalized warped frame, its asymptotic region, and the 8-dimensional Bianchi-I limit  |  Input cells 204-207
32.  The canonical spin connection of the interacting system  |  Input cells 208-210
33.  The coupled equations of fable and the primordial field, and their solution from the radiation era to today  |  Input cells 211-217
total Input cells: 217
\end{Verbatim}

\subsection{Export the field equations, the anticommutator, \texorpdfstring{$J$}{J} and the equations of state to TeX and text}
\label{sec:11.5}

\begin{Verbatim}
wolframscript -file "$(git rev-parse --show-toplevel)/claude-fable/export_fermion_fable_tex.wls"
\end{Verbatim}

It is self-locating and sets its own working directory to \file{claude-fable/}, so it runs from
anywhere. It evaluates every Input cell of the notebook with the output suppressed (so it takes about
as long as a run; its duration was not recorded), prints the assertion count, stops with exit code 1
if the Part~VII objects are missing, and writes the fourteen files listed in \secref{sec:10.6}. Its
recorded output is quoted there. On today's 217-cell notebook the first two lines would read
\cfcode{evaluating 217 Input cells ...} and \cfcode{assertions: 644 ...} (644 because the Rust solver's
CSVs are committed, so Part~VIII's two comparisons run, as in the final run of \secref{sec:10.3}); the
Part~VII objects it exports are the same. The \file{.tex} files carry the date of the run in their header comment, so a
re-run on another day changes that one line; the \file{.txt} files carry no date. Like every
evaluation of the notebook, it rewrites the two \file{.mx} files of Section~12, which are then
restored as in \secref{sec:11.2}.

\subsection{Build the LaTeX twin and its PDF}
\label{sec:11.6}

\begin{Verbatim}
bash "$(git rev-parse --show-toplevel)/provenance-latex/build_all.sh"
\end{Verbatim}

or, in PowerShell,

\begin{Verbatim}
powershell -ExecutionPolicy Bypass -File provenance-latex\build_all.ps1
\end{Verbatim}

For each of the three documents (\file{PROVENANCE-14-FERMION-FABLE-CANONICAL-QUANTIZATION},
\file{...-15-...}, \file{...-16-...}) the script runs
\cfcode{latexmk -pdf -interaction=nonstopmode -halt-on-error <doc>.tex}, with the output in
\file{provenance-latex/<doc>.latexmk.log}; it stops on the first failure, printing
\cfcode{FAILED; see provenance-latex/<doc>.latexmk.log and <doc>.log}; it warns if the \file{.log}
contains ``undefined'' or ``Rerun to get''; and it lists each PDF. So it prints \cfcode{== <doc>} and
an \cfcode{ls -la} line per document, and finally \cfcode{== all three PDFs built}. Every document
begins with \verb|% preamble.tex -- shared preamble of the LaTeX twins of the provenance pages
%   PROVENANCE-14-FERMION-FABLE-CANONICAL-QUANTIZATION
%   PROVENANCE-15-FERMION-FABLE-AND-THE-PRIMORDIAL-GRAVITATIONAL-FIELD
%   PROVENANCE-16-THE-COMPLETE-SOLUTION-AND-ITS-COMMANDS
% Each document starts with % preamble.tex -- shared preamble of the LaTeX twins of the provenance pages
%   PROVENANCE-14-FERMION-FABLE-CANONICAL-QUANTIZATION
%   PROVENANCE-15-FERMION-FABLE-AND-THE-PRIMORDIAL-GRAVITATIONAL-FIELD
%   PROVENANCE-16-THE-COMPLETE-SOLUTION-AND-ITS-COMMANDS
% Each document starts with % preamble.tex -- shared preamble of the LaTeX twins of the provenance pages
%   PROVENANCE-14-FERMION-FABLE-CANONICAL-QUANTIZATION
%   PROVENANCE-15-FERMION-FABLE-AND-THE-PRIMORDIAL-GRAVITATIONAL-FIELD
%   PROVENANCE-16-THE-COMPLETE-SOLUTION-AND-ITS-COMMANDS
% Each document starts with \input{preamble}.  Build all three from this directory with
%     bash build_all.sh            (Linux, macOS, Git Bash on Windows)
%     powershell -ExecutionPolicy Bypass -File build_all.ps1
% or one of them with
%     latexmk -pdf -interaction=nonstopmode -halt-on-error PROVENANCE-14-FERMION-FABLE-CANONICAL-QUANTIZATION.tex
% The style follows fable-cosmology/latex/fable_cosmology.tex (pdfLaTeX, Latin Modern, a4, 2.5 cm).
\documentclass[11pt,a4paper]{article}

\usepackage[T1]{fontenc}
\usepackage[utf8]{inputenc}
\usepackage{lmodern}
\usepackage[margin=2.5cm]{geometry}
\usepackage{microtype}
\usepackage{amsmath,amssymb,mathtools}
\usepackage{bm}
\usepackage{graphicx}
\usepackage{booktabs}
\usepackage{tabularx}
\usepackage{longtable}
\usepackage{array}
\usepackage{caption}
\usepackage{subcaption}
\usepackage{placeins}
\usepackage{enumitem}
\usepackage{fancyvrb}
\usepackage{upquote}
\usepackage{xcolor}
\usepackage[hyphens]{url}
\usepackage[hidelinks]{hyperref}

\graphicspath{{figures/}}
\captionsetup{font=small,labelfont=bf}
\setlist{itemsep=2pt,topsep=4pt}
\fvset{fontsize=\footnotesize,xleftmargin=1em}
\newcolumntype{Y}{>{\raggedright\arraybackslash}X}
\setcounter{tocdepth}{2}
\allowdisplaybreaks

% ---- the objects of the notebook, named as the notebook names them ----------------------------
\newcommand{\fable}{\textsf{fable}}
\newcommand{\fableScalar}{\textsf{fableScalar}}
\newcommand{\Lhat}{\hat{L}}
\newcommand{\sig}{\sigma_{16}}
\newcommand{\T}[1]{T_{16}[#1]}
\newcommand{\sgn}{\sqrt{|g|}}
\newcommand{\dd}{\mathrm{d}}
\newcommand{\diag}{\operatorname{diag}}
\newcommand{\sgnf}{\operatorname{sign}}
\newcommand{\Hobs}{H_{\mathrm{obs}}}
\newcommand{\Vhid}{V_{\mathrm{hid}}}
\newcommand{\Psibar}{\bar{\Psi}}
\newcommand{\cconj}{\ddagger}          % classical conjugation of a Grassmann field
\newcommand{\hadj}{\dagger}            % Hilbert-space adjoint of the J-quantization
\newcommand{\Jm}{\mathcal{J}}          % the Krein fundamental symmetry J = -i T0 T1 T2 T3 T4
\newcommand{\KS}{\mathrm{KS}}
\newcommand{\kF}{k_{\mathrm{F}}}
\newcommand{\wF}{\omega_{\mathrm{F}}}
\newcommand{\meff}{m_{\mathrm{eff}}}
\newcommand{\Tr}{\operatorname{Tr}}
\newcommand{\file}[1]{\path{#1}}
\newcommand{\assertion}[1]{\textsc{assert:}~\emph{#1}}
\pdfstringdefDisableCommands{\def\fableScalar{fableScalar}\def\fable{fable}}
.  Build all three from this directory with
%     bash build_all.sh            (Linux, macOS, Git Bash on Windows)
%     powershell -ExecutionPolicy Bypass -File build_all.ps1
% or one of them with
%     latexmk -pdf -interaction=nonstopmode -halt-on-error PROVENANCE-14-FERMION-FABLE-CANONICAL-QUANTIZATION.tex
% The style follows fable-cosmology/latex/fable_cosmology.tex (pdfLaTeX, Latin Modern, a4, 2.5 cm).
\documentclass[11pt,a4paper]{article}

\usepackage[T1]{fontenc}
\usepackage[utf8]{inputenc}
\usepackage{lmodern}
\usepackage[margin=2.5cm]{geometry}
\usepackage{microtype}
\usepackage{amsmath,amssymb,mathtools}
\usepackage{bm}
\usepackage{graphicx}
\usepackage{booktabs}
\usepackage{tabularx}
\usepackage{longtable}
\usepackage{array}
\usepackage{caption}
\usepackage{subcaption}
\usepackage{placeins}
\usepackage{enumitem}
\usepackage{fancyvrb}
\usepackage{upquote}
\usepackage{xcolor}
\usepackage[hyphens]{url}
\usepackage[hidelinks]{hyperref}

\graphicspath{{figures/}}
\captionsetup{font=small,labelfont=bf}
\setlist{itemsep=2pt,topsep=4pt}
\fvset{fontsize=\footnotesize,xleftmargin=1em}
\newcolumntype{Y}{>{\raggedright\arraybackslash}X}
\setcounter{tocdepth}{2}
\allowdisplaybreaks

% ---- the objects of the notebook, named as the notebook names them ----------------------------
\newcommand{\fable}{\textsf{fable}}
\newcommand{\fableScalar}{\textsf{fableScalar}}
\newcommand{\Lhat}{\hat{L}}
\newcommand{\sig}{\sigma_{16}}
\newcommand{\T}[1]{T_{16}[#1]}
\newcommand{\sgn}{\sqrt{|g|}}
\newcommand{\dd}{\mathrm{d}}
\newcommand{\diag}{\operatorname{diag}}
\newcommand{\sgnf}{\operatorname{sign}}
\newcommand{\Hobs}{H_{\mathrm{obs}}}
\newcommand{\Vhid}{V_{\mathrm{hid}}}
\newcommand{\Psibar}{\bar{\Psi}}
\newcommand{\cconj}{\ddagger}          % classical conjugation of a Grassmann field
\newcommand{\hadj}{\dagger}            % Hilbert-space adjoint of the J-quantization
\newcommand{\Jm}{\mathcal{J}}          % the Krein fundamental symmetry J = -i T0 T1 T2 T3 T4
\newcommand{\KS}{\mathrm{KS}}
\newcommand{\kF}{k_{\mathrm{F}}}
\newcommand{\wF}{\omega_{\mathrm{F}}}
\newcommand{\meff}{m_{\mathrm{eff}}}
\newcommand{\Tr}{\operatorname{Tr}}
\newcommand{\file}[1]{\path{#1}}
\newcommand{\assertion}[1]{\textsc{assert:}~\emph{#1}}
\pdfstringdefDisableCommands{\def\fableScalar{fableScalar}\def\fable{fable}}
.  Build all three from this directory with
%     bash build_all.sh            (Linux, macOS, Git Bash on Windows)
%     powershell -ExecutionPolicy Bypass -File build_all.ps1
% or one of them with
%     latexmk -pdf -interaction=nonstopmode -halt-on-error PROVENANCE-14-FERMION-FABLE-CANONICAL-QUANTIZATION.tex
% The style follows fable-cosmology/latex/fable_cosmology.tex (pdfLaTeX, Latin Modern, a4, 2.5 cm).
\documentclass[11pt,a4paper]{article}

\usepackage[T1]{fontenc}
\usepackage[utf8]{inputenc}
\usepackage{lmodern}
\usepackage[margin=2.5cm]{geometry}
\usepackage{microtype}
\usepackage{amsmath,amssymb,mathtools}
\usepackage{bm}
\usepackage{graphicx}
\usepackage{booktabs}
\usepackage{tabularx}
\usepackage{longtable}
\usepackage{array}
\usepackage{caption}
\usepackage{subcaption}
\usepackage{placeins}
\usepackage{enumitem}
\usepackage{fancyvrb}
\usepackage{upquote}
\usepackage{xcolor}
\usepackage[hyphens]{url}
\usepackage[hidelinks]{hyperref}

\graphicspath{{figures/}}
\captionsetup{font=small,labelfont=bf}
\setlist{itemsep=2pt,topsep=4pt}
\fvset{fontsize=\footnotesize,xleftmargin=1em}
\newcolumntype{Y}{>{\raggedright\arraybackslash}X}
\setcounter{tocdepth}{2}
\allowdisplaybreaks

% ---- the objects of the notebook, named as the notebook names them ----------------------------
\newcommand{\fable}{\textsf{fable}}
\newcommand{\fableScalar}{\textsf{fableScalar}}
\newcommand{\Lhat}{\hat{L}}
\newcommand{\sig}{\sigma_{16}}
\newcommand{\T}[1]{T_{16}[#1]}
\newcommand{\sgn}{\sqrt{|g|}}
\newcommand{\dd}{\mathrm{d}}
\newcommand{\diag}{\operatorname{diag}}
\newcommand{\sgnf}{\operatorname{sign}}
\newcommand{\Hobs}{H_{\mathrm{obs}}}
\newcommand{\Vhid}{V_{\mathrm{hid}}}
\newcommand{\Psibar}{\bar{\Psi}}
\newcommand{\cconj}{\ddagger}          % classical conjugation of a Grassmann field
\newcommand{\hadj}{\dagger}            % Hilbert-space adjoint of the J-quantization
\newcommand{\Jm}{\mathcal{J}}          % the Krein fundamental symmetry J = -i T0 T1 T2 T3 T4
\newcommand{\KS}{\mathrm{KS}}
\newcommand{\kF}{k_{\mathrm{F}}}
\newcommand{\wF}{\omega_{\mathrm{F}}}
\newcommand{\meff}{m_{\mathrm{eff}}}
\newcommand{\Tr}{\operatorname{Tr}}
\newcommand{\file}[1]{\path{#1}}
\newcommand{\assertion}[1]{\textsc{assert:}~\emph{#1}}
\pdfstringdefDisableCommands{\def\fableScalar{fableScalar}\def\fable{fable}}
| (\file{provenance-latex/preamble.tex}: pdfLaTeX, Latin Modern,
A4, 2.5\,cm margins, and macros for the notebook's objects, with separate symbols for the classical
conjugation \verb|\cconj| and the Hilbert adjoint \verb|\hadj|). MiKTeX's \cfcode{latexmk} needs
Perl on the \cfcode{PATH}; Git for Windows ships one (\verb|$env:PATH += ";C:\Program Files\Git\usr\bin"|
in PowerShell, as \file{build_all.ps1} notes).

The preamble was test-compiled when it was added: commit \texttt{0a789e2} records that a minimal test
document compiles with \cfcode{latexmk} with no warnings. The LaTeX twin of this page, in its final
state, was compiled on its own on 2026-09-25, inside \file{provenance-latex/}, with the command of its
header,
\cfcode{latexmk -pdf -interaction=nonstopmode -halt-on-error PROVENANCE-14-FERMION-FABLE-CANONICAL-QUANTIZATION.tex}:
it ran with no errors, no warnings and no undefined references, and wrote a PDF of 114 pages
whose only two overfull lines are 5.5\,pt and 0.6\,pt too wide. The run of \file{build_all.sh} over the three pages belongs to the fresh-clone
verification of \secref{sec:12}.

% ===============================================================================================
\section{Push and repository verification}
\label{sec:12}

\paragraph{The repository.} \texttt{origin} is \url{https://github.com/once-ere/Pre-Universe_with_Claude.git},
and the work is pushed to its \texttt{main} branch. (The remote named \texttt{upstream}, the author's
other repository, is never pushed to.) The repository carries a \file{.gitattributes} with
\cfcode{* -text}, so that git never rewrites line endings and every file checks out with the bytes that
were committed.

\paragraph{The commits of the work, 2026-09-24 and 2026-09-25} (\cfcode{git log --oneline -18}, times
from \cfcode{git log --date=iso}, PDT; the work ran past midnight, so the last four commits are dated
2026-09-25):

{\small
\begin{longtable}{@{}p{0.1\linewidth}p{0.1\linewidth}p{0.74\linewidth}@{}}
\toprule
\raggedright commit & \raggedright time & \raggedright what it contains \tabularnewline
\midrule
\endhead
\raggedright \texttt{e34d454} & \raggedright 18:52:36 & \raggedright the 2026-09-16 work, as it stood: Part~VI of the notebook (Sections 23--25; 172 Input cells, 358/358), the \file{fable-cosmology/} solver, notebooks, results and references; with its known gaps listed in the commit message \tabularnewline
\raggedright \texttt{028b379} & \raggedright 19:24:32 & \raggedright the reference generator \file{make_reference.wls} no longer prints harmless underflow messages (the values it writes are unchanged) \tabularnewline
\raggedright \texttt{2f6d322} & \raggedright 19:57:29 & \raggedright the 2026-09-16 gaps closed: the student instructions and the paper of the fable-cosmology work, its provenance page completed, seven corrections to its design; a fresh local clone ran its \file{setup.sh} and \file{run_all.sh} end to end \tabularnewline
\raggedright \texttt{8459c57} & \raggedright 20:02:04 & \raggedright a work-in-progress snapshot of the fermion-fable solver crate \file{fable-cosmology/rust/fable_fermion} \tabularnewline
\raggedright \texttt{e130337} & \raggedright 20:14:25 & \raggedright a second work-in-progress snapshot of that crate \tabularnewline
\raggedright \texttt{0a789e2} & \raggedright 20:17:18 & \raggedright the LaTeX build tree of the three new provenance pages: \file{provenance-latex/preamble.tex}, \file{build_all.sh}, \file{build_all.ps1} \tabularnewline
\raggedright \texttt{d0f2a99} & \raggedright 20:18:53 & \raggedright \file{build_all.sh} tracked without the executable bit, like the repository's other scripts \tabularnewline
\raggedright \texttt{c5892bd} & \raggedright 20:38:49 & \raggedright the finished fermion-fable solver: the Kohn--Sham fermion gas (stable integrals, per-7-volume mean field, gap equation), the stabilized and 8D models, \cfcode{cargo test} 21 + 8 passing, and the independent scipy cross-check \tabularnewline
\raggedright \textbf{\texttt{4fdfc40}} & \raggedright \textbf{20:47:28} & \raggedright \textbf{Part~VII (this effort):} \file{claude-fable/cells_part7.wl}, the rebuilt notebook (198 Input cells, 528/528, 0 cells with messages), \file{run_fermion_fable_part7.log}, \file{verify_nb_part7.log}, \file{export_fermion_fable_tex.wls} and the fourteen generated TeX/text files \tabularnewline
\raggedright \texttt{b7ea18d} & \raggedright 20:48:42 & \raggedright a wording correction, in an earlier provenance page, about the canonical connection's planes: for $k=5,6,7$ the $(0,k)$ plane is a \textbf{boost} and the $(4,k)$ plane a \textbf{rotation} (the fact is stated in \secref{sec:9.6} of this page) \tabularnewline
\raggedright \texttt{bc04ed1} & \raggedright 21:24:04 & \raggedright Part~VIII (the coupled system, Sections 30--33; 217 Input cells, 642/642, with its two comparisons of the Rust solver skipped because the solver's CSVs did not exist yet), its acceptance run log \file{run_fermion_fable_part8.log}, and the Mathematica fermion references \tabularnewline
\raggedright \texttt{18a2501} & \raggedright 22:13:42 & \raggedright the wall-state solver of the coupled system (\file{fable-cosmology/fermion/waveguide*}): its derivation from the notebook's own \cfcode{T16} (94/94 checks pass), its Python solver (validation 64/64), its independent Mathematica cross-check, and its report \tabularnewline
\raggedright \texttt{fee5603} & \raggedright 22:31:20 & \raggedright the draft of this page, its LaTeX twin and its PDF (110 pages), with section~1 still a placeholder \tabularnewline
\raggedright \texttt{2bc936d} & \raggedright 23:12:51 & \raggedright six defects of the solver \file{fable_fermion} fixed, among them the stable evaluation of the Dirac-sea energy \cfcode{vac_over_rhoc} (\secref{sec:8.3}); \cfcode{cargo test} 34/34 (23 unit + 11 integration, previously 21 + 8) \tabularnewline
\raggedright \texttt{aca5102} & \raggedright 2026-09-25 00:23:46 & \raggedright notebooks 05--07 of the fermion fable executed and committed with their results (all seven notebooks pass the checker, \cfcode{nbcheck} 7/7; \file{fable-cosmology/run_all.sh} passes end to end); the wall-state solver's input \file{fable-cosmology/fermion/waveguide_T16.json} committed (its derivation log still reports \cfcode{94 checks, 94 PASS, 0 FAIL}); \cfcode{waveguide.py validate} exits 1 whenever a check fails \tabularnewline
\raggedright \texttt{1671155} & \raggedright 00:28:54 & \raggedright the final full run of the notebook recorded, \file{claude-fable/run_fermion_fable_final.log} (217 Input cells, 644/644, 0 cells with messages, 208.715\,s, 46 witnesses): with the solver's CSVs now committed, Part~VIII's two comparisons of the Rust solver with the Mathematica reference runs run and pass ($642+2=644$) \tabularnewline
\raggedright \texttt{38701ef} & \raggedright 05:40:19 & \raggedright the drafts of the pages of Efforts B and C (markdown, LaTeX, PDF), the notebook figures their PDFs use (\file{provenance-latex/figures/}), and the repository's front page updated for Parts VII--VIII and the three pages \tabularnewline
\raggedright \texttt{3a5020d} & \raggedright 05:54:25 & \raggedright Part~VIII, Section~32: its Text cell and three assertion labels now write the listed spin-connection components with \textbf{both flat indices down} --- $\omega_{\mu ab}$, $\omega_{0\,04}$, $\omega_{i\,i4}$, $\omega_{h\,4h}$ --- as the arrays they compare have them (raising both indices flips the sign of every component in a plane with one timelike and one spacelike index; \secref{sec:10.3}); no computation changed. The final run re-recorded in \file{run_fermion_fable_final.log} (217 Input cells, 644/644, 0 cells with messages, 198.352\,s, 46 witnesses, 0 identities accepted on numerical evidence alone, both Rust-vs-Mathematica comparisons passing); \file{verify_nb_part8.log} re-recorded (\secref{sec:11.3}); a comment in the solver's \file{models.rs} \tabularnewline
\bottomrule
\end{longtable}
}

\paragraph{Checking that the remote has what was pushed:}

\begin{Verbatim}
cd "$(git rev-parse --show-toplevel)"
git fetch origin
git rev-parse HEAD origin/main              # both must print the same hash
git ls-remote origin refs/heads/main        # the hash GitHub serves for main
git status --short                          # what is not yet committed
\end{Verbatim}

When the draft of this page was written (2026-09-24, about 21:28 PDT),
\cfcode{git ls-remote origin refs/heads/main} printed

\begin{Verbatim}
bc04ed16096e134857111ae16df5220d4ff18435	refs/heads/main
\end{Verbatim}

which was the local \texttt{main} (\texttt{bc04ed1}, ``Add Part VIII: fable as a source of the Einstein
equations of the primordial field''). Everything of this effort that exists as code, log or generated
file --- \file{cells_part7.wl}, the rebuilt notebook, the Part~VII run log, the checker's logs, the
export script and the generated TeX/text files --- was in that pushed commit.

\paragraph{What was not yet committed when the draft was written, and where it went.} The draft of
this page, its LaTeX twin and its PDF were committed in \texttt{fee5603}. The working-tree changes of
the other efforts that were pending at that moment are all committed since: the changes to the
solver's \file{constants.rs} and to \file{fable-cosmology/fermion/crosscheck.py} in \texttt{2bc936d},
and the new notebook \file{05_fermion_fable_quantum_eos.ipynb}, the edits to the other notebooks and
their results in \texttt{aca5102}. None of the numbers on this page is taken from a file that is not
committed. The two \file{.mx} files of Section~12 were also showing as modified in the working tree
then, rewritten by a notebook evaluation running at the time; they are restored as in
\secref{sec:11.2} before any commit. This final version of the page, its LaTeX twin and its PDF are
committed together with the final versions of the pages of Efforts B and C.

\paragraph{The final push and the fresh-clone verification.} The final state is verified from a fresh
clone of the pushed \texttt{main}, with these commands (\texttt{<scratch>} is any empty directory
outside the repository):

\begin{Verbatim}
cd <scratch>
git clone https://github.com/once-ere/Pre-Universe_with_Claude.git fresh
cd fresh
git log --oneline -1                                   # the hash of the pushed main
cd claude-fable
cp claude-fable_Einstein-Rosen-2-Planes.nb committed.nb
python build_tools.py                                  # prints the four lines of section 11.1
cmp claude-fable_Einstein-Rosen-2-Planes.nb committed.nb && echo "notebook rebuilds byte for byte"
wolframscript -file verify_nb.wls                      # section 11.3
wolframscript -file run_from_nb.wls > fresh_run.log 2>&1
sed -n '/RUN SUMMARY/,$p' fresh_run.log                # 217 cells, 644/644, cells w/ msgs 0
diff <(grep '^  PASS' fresh_run.log) <(grep '^  PASS' run_fermion_fable_final.log) && echo "every label identical"
cd .. && bash provenance-latex/build_all.sh            # section 11.6
\end{Verbatim}

The Rust solver's two CSVs that Part~VIII compares with its Mathematica reference runs,
\file{fable-cosmology/results/nb06_fable4d_mass30eV.csv} and \file{nb06_fable4d_power.csv}, are
committed (since \texttt{aca5102}), so a fresh clone runs both comparisons and its summary is compared
with the final run, 644/644. Before the final one, the most recent completed fresh-clone verification
in the repository's history is the one recorded in \texttt{2f6d322} (for the fable-cosmology work).
The delivered copy of the notebook outside the repository,
\file{C:/Users/nsh/Documents/8-dim/claude-fable_Einstein-Rosen-2-Planes.nb}, is byte-identical to the
committed notebook of \texttt{3a5020d} (545713 bytes, the \cfcode{file bytes} of \secref{sec:11.3};
compared with \cfcode{cmp} when this page was brought to its final state on 2026-09-25).

The final commit hashes, the \cfcode{git ls-remote} line of the final push, and the results of the
fresh-clone verification of that pushed state (the rebuild, the run summary, the comparison of every
label, the LaTeX build) are recorded here: \texttt{\{\{FINAL-PUSH\}\}}

\end{document}
